\input{preamble}

% OK, start here.
%
\begin{document}

\title{Dualité pour les schémas}


\maketitle

\phantomsection
\label{section-phantom}

\tableofcontents

\section{Introduction}
\label{section-introduction}

\noindent
Ce chapitre étudie la dualité relative pour les morphismes de schémas et le complexe dualisant sur un schéma. On pourra consulter \cite{RD}.

\medskip\noindent
Les complexes dualisants sur les anneaux noethériens ont été définis et étudiés dans Complexes dualisants, section \ref{dualizing-section-dualizing} et suivantes. Nous poursuivons ici cette étude en considérant les complexes dualisants sur les schémas, voir la section \ref{section-dualizing-schemes}.

\medskip\noindent
L'essentiel de ce chapitre est consacré à l'étude de l'adjoint à droite de l'image directe dans le cadre des catégories dérivées de faisceaux de modules à faisceaux de cohomologie quasi-cohérents. Voir les sections \ref{section-twisted-inverse-image}, \ref{section-restriction-to-opens}, \ref{section-base-change-map}, \ref{section-base-change-II}, \ref{section-trace}, \ref{section-compare-with-pullback}, \ref{section-sections-with-exact-support}, \ref{section-duality-finite}, \ref{section-perfect-proper}, \ref{section-dualizing-Cartier} et \ref{section-examples}. Nous suivons ici les articles \cite{Neeman-Grothendieck}, \cite{LN}, \cite{Lipman-notes} et \cite{Neeman-improvement}.

\medskip\noindent
Nous présentons les importants foncteurs image inverse extraordinaire $f^!$ pour les morphismes séparés de type fini entre schémas noethériens dans les sections \ref{section-upper-shriek}, \ref{section-upper-shriek-properties} et \ref{section-base-change-shriek}, qui aboutissent à la synthèse de la section \ref{section-duality}.

\medskip\noindent
Dans la section \ref{section-glue}, nous exposons une autre théorie de la dualité et des complexes dualisants lorsque l'on travaille sur une base localement noethérienne fixée, munie d'un complexe dualisant (cette section correspond à une remarque du livre de Hartshorne).

\medskip\noindent
Les dernières sections présentent quelques applications.

\medskip\noindent
Ce chapitre se prolonge par celui consacré à la dualité sur les espaces algébriques, voir Dualité pour les espaces, section \ref{spaces-duality-section-introduction}.







\section{Complexes dualisants sur les schémas}
\label{section-dualizing-schemes}

\noindent
Nous appelons complexe dualisant sur un schéma localement noethérien un complexe qui, localement pour la topologie des ouverts affines, provient d'un complexe dualisant sur l'anneau correspondant. Cette définition n'est pas tout à fait usuelle, mais elle coïncide avec toutes celles de la littérature pour les schémas noethériens de dimension finie.

\begin{lemma}
\label{lemma-equivalent-definitions}
Soit $X$ un schéma localement noethérien. Soit $K$ un objet de $D(\mathcal{O}_X)$. Les conditions suivantes sont équivalentes :
\begin{enumerate}
\item Pour tout ouvert affine $U = \Spec(A) \subset X$, il existe un complexe dualisant $\omega_A^\bullet$ sur $A$ tel que $K|_U$ soit isomorphe à l'image de $\omega_A^\bullet$ par le foncteur $\widetilde{} : D(A) \to D(\mathcal{O}_U)$.
\item Il existe un recouvrement ouvert affine $X = \bigcup U_i$, $U_i = \Spec(A_i)$, tel que, pour tout $i$, il existe un complexe dualisant $\omega_i^\bullet$ sur $A_i$ tel que $K|_{U_i}$ soit isomorphe à l'image de $\omega_i^\bullet$ par le foncteur $\widetilde{} : D(A_i) \to D(\mathcal{O}_{U_i})$.
\end{enumerate}
\end{lemma}

\begin{proof}
Supposons (2) et soit $U = \Spec(A)$ un ouvert affine de $X$. Comme la condition (2) implique que $K$ appartient à $D_\QCoh(\mathcal{O}_X)$, il existe un objet $\omega_A^\bullet$ de $D(A)$ dont le complexe de faisceaux quasi-cohérents associé est isomorphe à $K|_U$, voir Catégories dérivées des schémas, lemme \ref{perfect-lemma-affine-compare-bounded}. Nous allons montrer que $\omega_A^\bullet$ est un complexe dualisant sur $A$, ce qui achèvera la démonstration.

\medskip\noindent
Puisque $X = \bigcup U_i$ est un recouvrement ouvert, on peut trouver un recouvrement par des ouverts principaux $U = D(f_1) \cup \ldots \cup D(f_m)$ tel que chaque $D(f_j)$ soit un ouvert principal de l'un des ouverts affines $U_i$, voir Schémas, lemme \ref{schemes-lemma-standard-open-two-affines}. Écrivons $D(f_j) = D(g_j)$, où $g_j \in A_{i_j}$. Alors $A_{f_j} \cong (A_{i_j})_{g_j}$ et l'on a
$$
(\omega_A^\bullet)_{f_j} \cong (\omega_i^\bullet)_{g_j}
$$
dans la catégorie dérivée, d'après Catégories dérivées des schémas, lemme \ref{perfect-lemma-affine-compare-bounded}. D'après Complexes dualisants, lemme \ref{dualizing-lemma-dualizing-localize}, le complexe $(\omega_A^\bullet)_{f_j}$ est dualisant sur $A_{f_j}$ pour $j = 1, \ldots, m$. Il en résulte que $\omega_A^\bullet$ est dualisant, d'après Complexes dualisants, lemme \ref{dualizing-lemma-dualizing-glue}.
\end{proof}

\begin{definition}
\label{definition-dualizing-scheme}
Soit $X$ un schéma localement noethérien. Un objet $K$ de $D(\mathcal{O}_X)$ est appelé {\it complexe dualisant} si $K$ vérifie les conditions équivalentes du lemme \ref{lemma-equivalent-definitions}.
\end{definition}

\noindent
Voir les remarques faites au début de cette section.

\begin{lemma}
\label{lemma-affine-duality}
Soient $A$ un anneau noethérien et $X = \Spec(A)$. Soient $K, L$ des objets de $D(A)$. Si $K \in D_{\textit{Coh}}(A)$ et si $L$ est de dimension injective finie, alors
$$
R\SheafHom_{\mathcal{O}_X}(\widetilde{K}, \widetilde{L})
=
\widetilde{R\Hom_A(K, L)}
$$
dans $D(\mathcal{O}_X)$.
\end{lemma}

\begin{proof}
On peut supposer que $L$ est représenté par un complexe fini $I^\bullet$ de $A$-modules injectifs. Par récurrence sur la longueur de $I^\bullet$ et par compatibilité des constructions avec les triangles distingués, on se ramène au cas où $L = I[0]$, avec $I$ un $A$-module injectif. Dans ce cas, Catégories dérivées des schémas, lemme \ref{perfect-lemma-quasi-coherence-internal-hom}, affirme que le faisceau de cohomologie de degré $n$ de $R\SheafHom_{\mathcal{O}_X}(\widetilde{K}, \widetilde{L})$ est le faisceau associé au préfaisceau
$$
D(f) \longmapsto \Ext^n_{A_f}(K \otimes_A A_f, I \otimes_A A_f)
$$
Puisque $A$ est noethérien, le $A_f$-module $I \otimes_A A_f$ est injectif (Complexes dualisants, lemme \ref{dualizing-lemma-localization-injective-modules}). On obtient donc
\begin{align*}
\Ext^n_{A_f}(K \otimes_A A_f, I \otimes_A A_f)
& =
\Hom_{A_f}(H^{-n}(K \otimes_A A_f), I \otimes_A A_f) \\
& =
\Hom_{A_f}(H^{-n}(K) \otimes_A A_f, I \otimes_A A_f) \\
& =
\Hom_A(H^{-n}(K), I) \otimes_A A_f
\end{align*}
La dernière égalité résulte du fait que $H^{-n}(K)$ est un $A$-module de type fini, voir Algèbre, lemme \ref{algebra-lemma-hom-from-finitely-presented}. Cela montre que l'application canonique
$$
\widetilde{R\Hom_A(K, L)}
\longrightarrow
R\SheafHom_{\mathcal{O}_X}(\widetilde{K}, \widetilde{L})
$$
est un quasi-isomorphisme dans ce cas, ce qui achève la démonstration.
\end{proof}

\begin{lemma}
\label{lemma-internal-hom-evaluate-isom}
Soit $X$ un schéma noethérien. Soient $K, L, M \in D_\QCoh(\mathcal{O}_X)$. Alors l'application
$$
R\SheafHom(L, M) \otimes_{\mathcal{O}_X}^\mathbf{L} K
\longrightarrow
R\SheafHom(R\SheafHom(K, L), M)
$$
de Cohomologie, lemme \ref{cohomology-lemma-internal-hom-evaluate}, est un isomorphisme dans les deux cas suivants :
\begin{enumerate}
\item $K \in D^-_{\textit{Coh}}(\mathcal{O}_X)$, $L \in D^+_{\textit{Coh}}(\mathcal{O}_X)$, et $M$ est localement sur les ouverts affines de dimension injective finie (voir la démonstration) ; ou
\item \begingroup\emergencystretch=12pt $K$ et $L$ appartiennent à $D_{\textit{Coh}}(\mathcal{O}_X)$, l'objet $R\SheafHom(L, M)$ est de Tor-dimension finie, et $L$ et $M$ sont localement sur les ouverts affines de dimension injective finie (en particulier, $L$ et $M$ sont bornés).\par\endgroup
\end{enumerate}
\end{lemma}

\begin{proof}
Démonstration de (1). On dit que $M$ est localement sur les ouverts affines de dimension injective finie si $X$ possède un recouvrement par des ouverts affines $U = \Spec(A)$ tels que l'objet de $D(A)$ correspondant à $M|_U$ (Catégories dérivées des schémas, lemme \ref{perfect-lemma-affine-compare-bounded}) soit de dimension injective finie\footnote{Cette condition ne dépend pas du choix du recouvrement ouvert affine du schéma noethérien $X$. Les détails sont omis.}. Pour démontrer le lemme, on peut remplacer $X$ par $U$, c'est-à-dire supposer que $X = \Spec(A)$ pour un anneau noethérien $A$. Observons que $R\SheafHom(K, L)$ appartient à $D^+_{\textit{Coh}}(\mathcal{O}_X)$ d'après Catégories dérivées des schémas, lemme \ref{perfect-lemma-coherent-internal-hom}. De plus, la formation des deux membres de la flèche commute au foncteur $D(A) \to D_\QCoh(\mathcal{O}_X)$ d'après le lemme \ref{lemma-affine-duality} et Catégories dérivées des schémas, lemme \ref{perfect-lemma-quasi-coherence-internal-hom} (on utilise bien ici les hypothèses sur $K$, $L$, $M$ et ce que l'on vient de montrer au sujet de $R\SheafHom(K, L)$). Enfin, la flèche est un isomorphisme d'après Compléments d'algèbre, lemme \ref{more-algebra-lemma-internal-hom-evaluate-isomorphism}, (2).

\medskip\noindent
Démonstration de (2). On raisonne comme ci-dessus. La seule différence est que l'on obtient ici $R\SheafHom(K, L)$ dans $D_{\textit{Coh}}(\mathcal{O}_X)$, car on peut, localement sur les ouverts affines (ce que permet le lemme \ref{lemma-affine-duality}), appliquer Complexes dualisants, lemme \ref{dualizing-lemma-finite-ext-into-bounded-injective}. On conclut alors par Compléments d'algèbre, lemme \ref{more-algebra-lemma-internal-hom-evaluate-isomorphism-technical}.
\end{proof}

\begin{lemma}
\label{lemma-dualizing-schemes}
Soit $K$ un complexe dualisant sur un schéma localement noethérien $X$. Alors $K$ est un objet de $D_{\textit{Coh}}(\mathcal{O}_X)$ et $D = R\SheafHom_{\mathcal{O}_X}(-, K)$ induit une anti-équivalence
$$
D :
D_{\textit{Coh}}(\mathcal{O}_X)
\longrightarrow
D_{\textit{Coh}}(\mathcal{O}_X)
$$
munie d'un isomorphisme canonique $\text{id} \to D \circ D$. Si $X$ est quasi-compact, alors $D$ échange $D^+_{\textit{Coh}}(\mathcal{O}_X)$ et $D^-_{\textit{Coh}}(\mathcal{O}_X)$ et induit une anti-équivalence $D^b_{\textit{Coh}}(\mathcal{O}_X) \to D^b_{\textit{Coh}}(\mathcal{O}_X)$.
\end{lemma}

\begin{proof}
Soit $U \subset X$ un ouvert affine. Écrivons $U = \Spec(A)$ et soit $\omega_A^\bullet$ un complexe dualisant sur $A$ correspondant à $K|_U$ comme au lemme \ref{lemma-equivalent-definitions}. D'après le lemme \ref{lemma-affine-duality}, le diagramme
$$
\xymatrix{
D_{\textit{Coh}}(A) \ar[r] \ar[d]_{R\Hom_A(-, \omega_A^\bullet)} &
D_{\textit{Coh}}(\mathcal{O}_U) \ar[d]^{R\SheafHom_{\mathcal{O}_X}(-, K|_U)} \\
D_{\textit{Coh}}(A) \ar[r] &
D(\mathcal{O}_U)
}
$$
est commutatif. On en conclut que $D$ envoie $D_{\textit{Coh}}(\mathcal{O}_X)$ dans $D_{\textit{Coh}}(\mathcal{O}_X)$. De plus, l'application canonique
$$
L
\longrightarrow
R\SheafHom_{\mathcal{O}_X}(K, K) \otimes_{\mathcal{O}_X}^\mathbf{L} L
\longrightarrow
R\SheafHom_{\mathcal{O}_X}(R\SheafHom_{\mathcal{O}_X}(L, K), K)
$$
(où l'on utilise Cohomologie, lemme \ref{cohomology-lemma-internal-hom-evaluate}, pour la seconde flèche) est un isomorphisme pour tout $L$, car cela est vrai sur les ouverts affines d'après Complexes dualisants, lemme \ref{dualizing-lemma-dualizing}\footnote{On peut aussi commencer par montrer que $R\SheafHom_{\mathcal{O}_X}(K, K) = \mathcal{O}_X$ en travaillant localement sur les ouverts affines, puis utiliser le lemme \ref{lemma-internal-hom-evaluate-isom}, (2), pour voir que l'application est un isomorphisme.}, et nous avons déjà vu que, sur les ouverts affines, on retrouve la construction algébrique. L'énoncé relatif aux propriétés de bornitude du foncteur $D$ dans le cas quasi-compact résulte également des énoncés correspondants de Complexes dualisants, lemme \ref{dualizing-lemma-dualizing}.
\end{proof}

\noindent
Soit $X$ un espace localement annelé. Rappelons qu'un objet $L$ de $D(\mathcal{O}_X)$ est {\it inversible} s'il est inversible pour la structure monoïdale symétrique sur $D(\mathcal{O}_X)$ définie par le produit tensoriel dérivé. Dans Cohomologie, lemme \ref{cohomology-lemma-invertible-derived}, nous avons vu que cela signifie que $L$ est parfait et qu'il existe un recouvrement ouvert $X = \bigcup U_i$ tel que $L|_{U_i} \cong \mathcal{O}_{U_i}[-n_i]$ pour certains entiers $n_i$. Dans ce cas, la fonction
$$
x \mapsto n_x,\quad
\text{où }n_x\text{ est l'unique entier tel que }
H^{n_x}(L_x) \not = 0
$$
est localement constante sur $X$. En particulier, on a $L = \bigoplus H^n(L)[-n]$, ce qui définit un complexe de $\mathcal{O}_X$-modules (à différentielles nulles) représentant $L$.

\begin{lemma}
\label{lemma-dualizing-unique-schemes}
Soit $X$ un schéma localement noethérien. Si $K$ et $K'$ sont des complexes dualisants sur $X$, alors $K'$ est isomorphe à $K \otimes_{\mathcal{O}_X}^\mathbf{L} L$ pour un objet inversible $L$ de $D(\mathcal{O}_X)$.
\end{lemma}

\begin{proof}
Posons
$$
L = R\SheafHom_{\mathcal{O}_X}(K, K')
$$
C'est un objet inversible de $D(\mathcal{O}_X)$, car il en est ainsi localement sur les ouverts affines, voir Complexes dualisants, lemme \ref{dualizing-lemma-dualizing-unique} et sa démonstration. L'application d'évaluation $L \otimes_{\mathcal{O}_X}^\mathbf{L} K \to K'$ est un isomorphisme pour la même raison.
\end{proof}

\begin{lemma}
\label{lemma-dimension-function-scheme}
Soit $X$ un schéma localement noethérien. Soit $\omega_X^\bullet$ un complexe dualisant sur $X$. Alors $X$ est universellement caténaire et la fonction $X \to \mathbf{Z}$ définie par
$$
x \longmapsto \delta(x)\text{ tel que }
\omega_{X, x}^\bullet[-\delta(x)]
\text{ soit un complexe dualisant normalisé sur }
\mathcal{O}_{X, x}
$$
est une fonction de dimension.
\end{lemma}

\begin{proof}
Cela résulte immédiatement du cas affine, Complexes dualisants, lemme \ref{dualizing-lemma-dimension-function}, et des définitions.
\end{proof}

\begin{lemma}
\label{lemma-sitting-in-degrees}
Soit $X$ un schéma localement noethérien. Soit $\omega_X^\bullet$ un complexe dualisant sur $X$, de fonction de dimension associée $\delta$. Soit $\mathcal{F}$ un $\mathcal{O}_X$-module cohérent. Posons $\mathcal{E}^i = \SheafExt^{-i}_{\mathcal{O}_X}(\mathcal{F}, \omega_X^\bullet)$. Alors $\mathcal{E}^i$ est un $\mathcal{O}_X$-module cohérent et, pour $x \in X$, on a :
\begin{enumerate}
\item $\mathcal{E}^i_x$ ne peut être non nul que pour $\delta(x) \leq i \leq \delta(x) + \dim(\text{Supp}(\mathcal{F}_x))$ ;
\item $\dim(\text{Supp}(\mathcal{E}^{i + \delta(x)}_x)) \leq i$ ;
\item $\text{depth}(\mathcal{F}_x)$ est le plus petit entier $i \geq 0$ tel que $\mathcal{E}_x^{i + \delta(x)} \not = 0$ ;
\item on a $x \in \text{Supp}(\bigoplus_{j \leq i} \mathcal{E}^j)
\Leftrightarrow
\text{depth}_{\mathcal{O}_{X, x}}(\mathcal{F}_x) + \delta(x) \leq i$.
\end{enumerate}
\end{lemma}

\begin{proof}
Le lemme \ref{lemma-dualizing-schemes} affirme que $\mathcal{E}^i$ est cohérent. En choisissant un voisinage affine de $x$ et en utilisant Catégories dérivées des schémas, lemme \ref{perfect-lemma-quasi-coherence-internal-hom}, ainsi que Compléments d'algèbre, lemme \ref{more-algebra-lemma-base-change-RHom}, (3), on obtient
$$
\mathcal{E}^i_x =
\SheafExt^{-i}_{\mathcal{O}_X}(\mathcal{F}, \omega_X^\bullet)_x =
\Ext^{-i}_{\mathcal{O}_{X, x}}(\mathcal{F}_x,
\omega_{X, x}^\bullet) =
\Ext^{\delta(x) - i}_{\mathcal{O}_{X, x}}(\mathcal{F}_x,
\omega_{X, x}^\bullet[-\delta(x)])
$$
La construction de $\delta$ au lemme \ref{lemma-dimension-function-scheme} ramène ainsi les assertions (1), (2) et (3) à Complexes dualisants, lemme \ref{dualizing-lemma-sitting-in-degrees}. L'assertion (4) est une conséquence formelle de (3) et (1).
\end{proof}




\section{Adjoint à droite de l'image directe}
\label{section-twisted-inverse-image}

\noindent
Pour cette section et les suivantes, on pourra consulter \cite{Neeman-Grothendieck}, \cite{LN}, \cite{Lipman-notes} et \cite{Neeman-improvement}.

\medskip\noindent
Soit $f : X \to Y$ un morphisme de schémas. Nous considérons dans cette section l'adjoint à droite du foncteur $Rf_* : D_\QCoh(\mathcal{O}_X) \to D_\QCoh(\mathcal{O}_Y)$. Dans la littérature, lorsqu'il existe, ce foncteur est parfois noté $f^{\times}$. Cette notation n'est pas universellement admise et nous ne l'utiliserons pas. Nous n'emploierons pas non plus la notation $f^!$ pour ce foncteur, car elle serait incompatible (pour des morphismes $f$ généraux) avec celle de \cite{RD}.

\begin{lemma}
\label{lemma-twisted-inverse-image}
\begin{reference}
Cet énoncé est presque identique à \cite[Example 4.2]{Neeman-Grothendieck}.
\end{reference}
Soit $f : X \to Y$ un morphisme entre schémas quasi-séparés et quasi-compacts. Le foncteur $Rf_* : D_\QCoh(\mathcal{O}_X) \to D_\QCoh(\mathcal{O}_Y)$ admet un adjoint à droite.
\end{lemma}

\begin{proof}
Nous allons montrer l'existence d'un adjoint à droite en vérifiant les hypothèses de Catégories dérivées, proposition \ref{derived-proposition-brown}. Tout d'abord, la catégorie $D_\QCoh(\mathcal{O}_X)$ admet des sommes directes, voir Catégories dérivées des schémas, lemme \ref{perfect-lemma-quasi-coherence-direct-sums}. La catégorie $D_\QCoh(\mathcal{O}_X)$ est engendrée par des objets compacts d'après Catégories dérivées des schémas, théorème \ref{perfect-theorem-bondal-van-den-Bergh}. Puisque $X$ et $Y$ sont quasi-compacts et quasi-séparés, il en est de même de $f$, voir Schémas, lemmes \ref{schemes-lemma-compose-after-separated} et \ref{schemes-lemma-quasi-compact-permanence}. Le foncteur $Rf_*$ commute donc aux sommes directes, voir Catégories dérivées des schémas, lemme \ref{perfect-lemma-quasi-coherence-pushforward-direct-sums}. Cela achève la démonstration.
\end{proof}

\begin{example}
\label{example-affine-twisted-inverse-image}
Soit $A \to B$ un homomorphisme d'anneaux. Soient $Y = \Spec(A)$ et $X = \Spec(B)$, et soit $f : X \to Y$ le morphisme correspondant à $A \to B$. Alors $Rf_* : D_\QCoh(\mathcal{O}_X) \to D_\QCoh(\mathcal{O}_Y)$ correspond à la restriction $D(B) \to D(A)$ par les équivalences $D(B) \to D_\QCoh(\mathcal{O}_X)$ et $D(A) \to D_\QCoh(\mathcal{O}_Y)$. L'adjoint à droite correspond donc au foncteur $K \longmapsto R\Hom(B, K)$ de Complexes dualisants, section \ref{dualizing-section-trivial}.
\end{example}

\begin{example}
\label{example-does-not-preserve-coherent}
Si $f : X \to Y$ est un morphisme séparé de type fini entre schémas noethériens, l'adjoint à droite de $Rf_* : D_\QCoh(\mathcal{O}_X) \to D_\QCoh(\mathcal{O}_Y)$ n'envoie pas nécessairement $D_{\textit{Coh}}(\mathcal{O}_Y)$ dans $D_{\textit{Coh}}(\mathcal{O}_X)$. En effet, soit $k$ un corps et considérons le morphisme $f : \mathbf{A}^1_k \to \Spec(k)$. D'après l'exemple \ref{example-affine-twisted-inverse-image}, cela revient à demander si $R\Hom(B, -)$ envoie $D_{\textit{Coh}}(A)$ dans $D_{\textit{Coh}}(B)$, où $A = k$ et $B = k[x]$. Ce n'est pas le cas, car
$$
R\Hom(k[x], k) = \left(\prod\nolimits_{n \geq 0} k\right)[0]
$$
n'est pas un $k[x]$-module de type fini. Ainsi, $a(\mathcal{O}_Y)$ n'a pas des faisceaux de cohomologie cohérents.
\end{example}

\begin{example}
\label{example-does-not-preserve-bounded-above}
Si $f : X \to Y$ est un morphisme propre, ou même fini, entre schémas noethériens, l'adjoint à droite de $Rf_* : D_\QCoh(\mathcal{O}_X) \to D_\QCoh(\mathcal{O}_Y)$ n'envoie pas nécessairement $D_\QCoh^-(\mathcal{O}_Y)$ dans $D_\QCoh^-(\mathcal{O}_X)$. En effet, soit $k$ un corps, soit $k[\epsilon]$ l'anneau des nombres duaux sur $k$, et soient $X = \Spec(k)$ et $Y = \Spec(k[\epsilon])$. Alors $\Ext^i_{k[\epsilon]}(k, k)$ est non nul pour tout $i \geq 0$. Ainsi, $a(\mathcal{O}_Y)$ n'est pas borné supérieurement, d'après l'exemple \ref{example-affine-twisted-inverse-image}.
\end{example}

\begin{lemma}
\label{lemma-twisted-inverse-image-bounded-below}
Soit $f : X \to Y$ un morphisme de schémas quasi-compacts et quasi-séparés. Soit $a : D_\QCoh(\mathcal{O}_Y) \to D_\QCoh(\mathcal{O}_X)$ l'adjoint à droite de $Rf_*$ du lemme \ref{lemma-twisted-inverse-image}. Alors $a$ envoie $D^+_\QCoh(\mathcal{O}_Y)$ dans $D^+_\QCoh(\mathcal{O}_X)$. Plus précisément, il existe un entier $N$ tel que $H^i(K) = 0$ pour $i \leq c$ implique $H^i(a(K)) = 0$ pour $i \leq c - N$.
\end{lemma}

\begin{proof}
D'après Catégories dérivées des schémas, lemme \ref{perfect-lemma-quasi-coherence-direct-image}, le foncteur $Rf_*$ est de dimension cohomologique finie. Autrement dit, il existe un entier $N$ tel que $H^i(Rf_*L) = 0$ pour $i \geq N + c$ si $H^i(L) = 0$ pour $i \geq c$. Supposons que $K \in D^+_\QCoh(\mathcal{O}_Y)$ vérifie $H^i(K) = 0$ pour $i \leq c$. Alors
$$
\Hom_{D(\mathcal{O}_X)}(\tau_{\leq c - N}a(K), a(K)) =
\Hom_{D(\mathcal{O}_Y)}(Rf_*\tau_{\leq c - N}a(K), K) = 0
$$
d'après ce qui précède. Cela implique clairement que $H^i(a(K)) = 0$ pour $i \leq c - N$.
\end{proof}

\noindent
Soit $f : X \to Y$ un morphisme de schémas quasi-séparés et quasi-compacts. Notons $a$ l'adjoint à droite de $Rf_* : D_\QCoh(\mathcal{O}_X) \to D_\QCoh(\mathcal{O}_Y)$. Pour tous $K \in D_\QCoh(\mathcal{O}_Y)$ et $L \in D_\QCoh(\mathcal{O}_X)$, on obtient une application canonique
\begin{equation}
\label{equation-sheafy-trace}
Rf_*R\SheafHom_{\mathcal{O}_X}(L, a(K))
\longrightarrow
R\SheafHom_{\mathcal{O}_Y}(Rf_*L, K)
\end{equation}
Cette application est définie comme la composée
$$
Rf_*R\SheafHom_{\mathcal{O}_X}(L, a(K)) \to
R\SheafHom_{\mathcal{O}_Y}(Rf_*L, Rf_*a(K)) \to
R\SheafHom_{\mathcal{O}_Y}(Rf_*L, K)
$$
où la première flèche est celle de Cohomologie, remarque \ref{cohomology-remark-projection-formula-for-internal-hom}, et la seconde est la co-unité $Rf_*a(K) \to K$ de l'adjonction.

\begin{lemma}
\label{lemma-iso-on-RSheafHom}
Soit $f : X \to Y$ un morphisme de schémas quasi-compacts et quasi-séparés. Soit $a$ l'adjoint à droite de $Rf_* : D_\QCoh(\mathcal{O}_X) \to D_\QCoh(\mathcal{O}_Y)$. Soient $L \in D_\QCoh(\mathcal{O}_X)$ et $K \in D_\QCoh(\mathcal{O}_Y)$. Alors l'application (\ref{equation-sheafy-trace})
$$
Rf_*R\SheafHom_{\mathcal{O}_X}(L, a(K))
\longrightarrow
R\SheafHom_{\mathcal{O}_Y}(Rf_*L, K)
$$
\begingroup\emergencystretch=16pt
devient un isomorphisme après application du foncteur
$DQ_Y : D(\mathcal{O}_Y) \to D_\QCoh(\mathcal{O}_Y)$
étudié dans Catégories dérivées des schémas, section \ref{perfect-section-better-coherator}.\par\endgroup
\end{lemma}

\begin{proof}
L'énoncé a un sens puisque $DQ_Y$ existe d'après Catégories dérivées des schémas, lemme \ref{perfect-lemma-better-coherator}. Comme $DQ_Y$ est l'adjoint à droite du foncteur d'inclusion $D_\QCoh(\mathcal{O}_Y) \to D(\mathcal{O}_Y)$, il suffit de montrer que, pour tout $M \in D_\QCoh(\mathcal{O}_Y)$, l'application (\ref{equation-sheafy-trace}) induit une bijection
$$
\Hom_Y(M, Rf_*R\SheafHom_{\mathcal{O}_X}(L, a(K)))
\longrightarrow
\Hom_Y(M, R\SheafHom_{\mathcal{O}_Y}(Rf_*L, K))
$$
Pour cela, on utilise la suite d'égalités suivante :
\begin{align*}
\Hom_Y(M, Rf_*R\SheafHom_{\mathcal{O}_X}(L, a(K)))
& =
\Hom_X(Lf^*M, R\SheafHom_{\mathcal{O}_X}(L, a(K))) \\
& =
\Hom_X(Lf^*M \otimes_{\mathcal{O}_X}^\mathbf{L} L, a(K)) \\
& =
\Hom_Y(Rf_*(Lf^*M \otimes_{\mathcal{O}_X}^\mathbf{L} L), K) \\
& =
\Hom_Y(M \otimes_{\mathcal{O}_Y}^\mathbf{L} Rf_*L, K) \\
& =
\Hom_Y(M, R\SheafHom_{\mathcal{O}_Y}(Rf_*L, K))
\end{align*}
La première égalité résulte de Cohomologie, lemme \ref{cohomology-lemma-adjoint}. La deuxième résulte de Cohomologie, lemme \ref{cohomology-lemma-internal-hom}. La troisième résulte de la construction de $a$. La quatrième résulte de Catégories dérivées des schémas, lemme \ref{perfect-lemma-cohomology-base-change} (c'est l'étape essentielle). La cinquième résulte de Cohomologie, lemme \ref{cohomology-lemma-internal-hom}.
\end{proof}

\begin{example}
\label{example-iso-on-RSheafHom}
L'énoncé du lemme \ref{lemma-iso-on-RSheafHom} est faux si l'on n'applique pas le « cohérateur » $DQ_Y$. Supposons en effet $Y = \Spec(R)$ et $X = \mathbf{A}^1_R$. Prenons $L = \mathcal{O}_X$ et $K = \mathcal{O}_Y$. Le membre de gauche de la flèche appartient à $D_\QCoh(\mathcal{O}_Y)$, tandis que celui de droite est isomorphe à $\prod_{n \geq 0} \mathcal{O}_Y$, qui n'est pas quasi-cohérent.
\end{example}

\begin{remark}
\label{remark-iso-on-RSheafHom}
Dans la situation du lemme \ref{lemma-iso-on-RSheafHom}, on a
$$
DQ_Y(Rf_*R\SheafHom_{\mathcal{O}_X}(L, a(K))) =
Rf_* DQ_X(R\SheafHom_{\mathcal{O}_X}(L, a(K)))
$$
\begingroup\emergencystretch=16pt
d'après Catégories dérivées des schémas, lemme \ref{perfect-lemma-pushforward-better-coherator}. Ainsi, si
\par\noindent\hfil
$R\SheafHom_{\mathcal{O}_X}(L, a(K)) \in D_\QCoh(\mathcal{O}_X)$,
\hfil\par\noindent
on peut « effacer » le $DQ_Y$ du membre de gauche de la flèche. D'autre part, si l'on sait que $R\SheafHom_{\mathcal{O}_Y}(Rf_*L, K) \in D_\QCoh(\mathcal{O}_Y)$, on peut « effacer » le $DQ_Y$ du membre de droite. Si les deux conditions sont satisfaites, on voit que (\ref{equation-sheafy-trace}) est un isomorphisme. En combinant cela avec Catégories dérivées des schémas, lemme \ref{perfect-lemma-quasi-coherence-internal-hom}, on voit que $Rf_*R\SheafHom_{\mathcal{O}_X}(L, a(K)) \to
R\SheafHom_{\mathcal{O}_Y}(Rf_*L, K)$ est un isomorphisme si\par\endgroup
\begin{enumerate}
\item $L$ et $Rf_*L$ sont parfaits ; ou
\item $K$ est borné inférieurement et $L$ et $Rf_*L$ sont pseudo-cohérents.
\end{enumerate}
Pour (2), on utilise le fait que $a(K)$ est borné inférieurement si $K$ l'est, voir le lemme \ref{lemma-twisted-inverse-image-bounded-below}.
\end{remark}

\begin{example}
\label{example-iso-on-RSheafHom-noetherian}
\begingroup\emergencystretch=16pt
Soit $f : X \to Y$ un morphisme propre de schémas noethériens, et soient $L \in D^-_{\textit{Coh}}(X)$ et $K \in D^+_{\QCoh}(\mathcal{O}_Y)$. Alors l'application $Rf_*R\SheafHom_{\mathcal{O}_X}(L, a(K)) \to
R\SheafHom_{\mathcal{O}_Y}(Rf_*L, K)$ est un isomorphisme. En effet, les complexes $L$ et $Rf_*L$ sont pseudo-cohérents d'après Catégories dérivées des schémas, lemmes \ref{perfect-lemma-identify-pseudo-coherent-noetherian} et \ref{perfect-lemma-direct-image-coherent}, et l'on peut appliquer la remarque \ref{remark-iso-on-RSheafHom}.\par\endgroup
\end{example}

\begin{lemma}
\label{lemma-iso-global-hom}
Soit $f : X \to Y$ un morphisme de schémas quasi-séparés et quasi-compacts. Pour tous $L \in D_\QCoh(\mathcal{O}_X)$ et $K \in D_\QCoh(\mathcal{O}_Y)$, (\ref{equation-sheafy-trace}) induit un isomorphisme $R\Hom_X(L, a(K)) \to R\Hom_Y(Rf_*L, K)$ entre les Hom dérivés globaux.
\end{lemma}

\begin{proof}
Par la construction de Cohomologie, section \ref{cohomology-section-global-RHom}, on a
$$
R\Hom_X(L, a(K)) =
R\Gamma(X, R\SheafHom_{\mathcal{O}_X}(L, a(K))) =
R\Gamma(Y, Rf_*R\SheafHom_{\mathcal{O}_X}(L, a(K)))
$$
et
$$
R\Hom_Y(Rf_*L, K) = R\Gamma(Y, R\SheafHom_{\mathcal{O}_Y}(Rf_*L, K))
$$
Le lemme est donc une conséquence du lemme \ref{lemma-iso-on-RSheafHom}. En effet, une application $E \to E'$ dans $D(\mathcal{O}_Y)$ qui induit un isomorphisme $DQ_Y(E) \to DQ_Y(E')$ induit un quasi-isomorphisme $R\Gamma(Y, E) \to R\Gamma(Y, E')$. On a en effet $H^i(Y, E) = \Ext^i_Y(\mathcal{O}_Y, E) = \Hom(\mathcal{O}_Y[-i], E) =
\Hom(\mathcal{O}_Y[-i], DQ_Y(E))$, car $\mathcal{O}_Y[-i]$ appartient à $D_\QCoh(\mathcal{O}_Y)$ et $DQ_Y$ est l'adjoint à droite du foncteur d'inclusion $D_\QCoh(\mathcal{O}_Y) \to D(\mathcal{O}_Y)$.
\end{proof}







\section{Adjoint à droite de l'image directe et restriction aux ouverts}
\sectionmark{Adjoint à droite et ouverts}
\label{section-restriction-to-opens}

\noindent
Nous étudions dans cette section dans quelle mesure l'adjoint à droite de l'image directe commute à la restriction aux sous-schémas ouverts. Il s'agit d'une question de changement de base ; commençons donc par l'aborder dans un cadre plus général.

\medskip\noindent
On souhaite souvent savoir si les adjoints à droite de l'image directe commutent au changement de base. Considérons donc un carré cartésien
\begin{equation}
\label{equation-base-change}
\vcenter{
\xymatrix{
X' \ar[r]_{g'} \ar[d]_{f'} & X \ar[d]^f \\
Y' \ar[r]^g & Y
}
}
\end{equation}
de schémas quasi-compacts et quasi-séparés. Notons
$$
a  : D_\QCoh(\mathcal{O}_Y) \to D_\QCoh(\mathcal{O}_X)
\quad\text{et}\quad
a'  : D_\QCoh(\mathcal{O}_{Y'}) \to D_\QCoh(\mathcal{O}_{X'})
$$
les adjoints à droite de $Rf_*$ et $Rf'_*$ (lemme \ref{lemma-twisted-inverse-image}). Considérons l'application de changement de base de Cohomologie, remarque \ref{cohomology-remark-base-change}. Elle induit une transformation de foncteurs
$$
Lg^* \circ Rf_* \longrightarrow Rf'_* \circ L(g')^*
$$
sur les catégories dérivées de faisceaux à cohomologie quasi-cohérente. On obtient donc une transformation entre les adjoints à droite, en sens opposé :
$$
a \circ Rg_* \longleftarrow Rg'_* \circ a'
$$

\begin{lemma}
\label{lemma-flat-precompose-pus}
Dans le diagramme (\ref{equation-base-change}), supposons que $g$ soit plat, ou plus généralement que $f$ et $g$ soient Tor-indépendants. Alors $a \circ Rg_* \leftarrow Rg'_* \circ a'$ est un isomorphisme.
\end{lemma}

\begin{proof}
Dans ce cas, l'application de changement de base
$Lg^* \circ Rf_* K \longrightarrow Rf'_* \circ L(g')^*K$
est un isomorphisme pour tout $K$ dans $D_\QCoh(\mathcal{O}_X)$ d'après Catégories dérivées des schémas, lemme \ref{perfect-lemma-compare-base-change}. La transformation correspondante entre foncteurs adjoints est donc également un isomorphisme.
\end{proof}

\noindent
Soit $f : X \to Y$ un morphisme de schémas quasi-compacts et quasi-séparés. Soit $V \subset Y$ un sous-schéma ouvert quasi-compact et posons $U = f^{-1}(V)$. On obtient un carré cartésien
$$
\xymatrix{
U \ar[r]_{j'} \ar[d]_{f|_U} & X \ar[d]^f \\
V \ar[r]^j & Y
}
$$
comme en (\ref{equation-base-change}). D'après le lemme \ref{lemma-flat-precompose-pus}, l'application $\xi : a \circ Rj_* \leftarrow Rj'_* \circ a'$ est un isomorphisme, où $a$ et $a'$ sont les adjoints à droite de $Rf_*$ et $R(f|_U)_*$. On obtient une transformation de foncteurs $D_\QCoh(\mathcal{O}_Y) \to D_\QCoh(\mathcal{O}_U)$
\begin{equation}
\label{equation-sheafy}
(j')^* \circ a \to
(j')^* \circ a \circ Rj_* \circ j^* \xrightarrow{\xi^{-1}}
(j')^* \circ Rj'_* \circ a' \circ j^* \to a' \circ j^*
\end{equation}
où la première flèche provient de $\text{id} \to Rj_* \circ j^*$ et la dernière de l'isomorphisme $(j')^* \circ Rj'_* \to \text{id}$. En particulier, (\ref{equation-sheafy}), évaluée en $K$, est un isomorphisme si et seulement si $a(K)|_U \to a(Rj_*(K|_V))|_U$ en est un.

\begin{example}
\label{example-not-supported-on-inverse-image}
Il existe un morphisme fini $f : X \to Y$ de schémas noethériens tel que (\ref{equation-sheafy}) ne soit pas un isomorphisme lorsqu'on l'évalue en un certain $K \in D_{\textit{Coh}}(\mathcal{O}_Y)$. En effet, prenons $X = \Spec(B) \to Y = \Spec(A)$ avec $A = k[x, \epsilon]$, où $k$ est un corps, et prenons $\epsilon^2 = 0$ et $B = k[x] = A/(\epsilon)$. Pour $n \in \mathbf{N}$, posons $M_n = A/(\epsilon, x^n)$. On observe que
$$
\Ext^i_A(B, M_n) = M_n,\quad i \geq 0
$$
car $B$ admet la résolution libre périodique $\ldots \to A \to A \to A$, dont les applications sont données par la multiplication par $\epsilon$. Considérons l'objet
$K = \bigoplus M_n[n] = \prod M_n[n]$
de $D_{\textit{Coh}}(A)$ (égalité dans $D(A)$ d'après Catégories dérivées, lemmes \ref{derived-lemma-direct-sums} et \ref{derived-lemma-products}). Alors $a(K)$ correspond à $R\Hom(B, K)$ d'après l'exemple \ref{example-affine-twisted-inverse-image}, et
$$
H^0(R\Hom(B, K)) = \Ext^0_A(B, K) =
\prod\nolimits_{n \geq 1} \Ext^n_A(B. M_n) = 
\prod\nolimits_{n \geq 1} M_n
$$
d'après ce qui précède. Mais ce module possède des éléments qui ne sont annulés par aucune puissance de $x$, tandis que chaque élément de la cohomologie du complexe $K$ est annulé par une puissance de $x$. Autrement dit, l'application (\ref{equation-sheafy}), avec $V = D(x)$ et $U = D(x)$, évaluée sur le complexe $K$, ne peut pas être un isomorphisme, car $(j')^*(a(K))$ est non nul et $a'(j^*K)$ est nul.
\end{example}

\begin{lemma}
\label{lemma-when-sheafy}
Soit $f : X \to Y$ un morphisme de schémas quasi-compacts et quasi-séparés. Soit $a$ l'adjoint à droite de $Rf_* : D_\QCoh(\mathcal{O}_X) \to D_\QCoh(\mathcal{O}_Y)$. Soit $V \subset Y$ un ouvert quasi-compact d'image inverse $U \subset X$.
\begin{enumerate}
\item Pour tout $Q \in D_\QCoh^+(\mathcal{O}_Y)$ à support dans $Y \setminus V$, l'image $a(Q)$ est à support dans $X \setminus U$ si et seulement si (\ref{equation-sheafy}) est un isomorphisme sur tous les $K$ de $D_\QCoh^+(\mathcal{O}_Y)$.
\item Pour tout $Q \in D_\QCoh(\mathcal{O}_Y)$ à support dans $Y \setminus V$, l'image $a(Q)$ est à support dans $X \setminus U$ si et seulement si (\ref{equation-sheafy}) est un isomorphisme sur tous les $K$ de $D_\QCoh(\mathcal{O}_Y)$.
\item Si $a$ commute aux sommes directes, les conditions équivalentes de (1) impliquent celles de (2).
\end{enumerate}
\end{lemma}

\begin{proof}
Démonstration de (1). Soit $K \in D_\QCoh^+(\mathcal{O}_Y)$. Choisissons un triangle distingué
$$
K \to Rj_*K|_V \to Q \to K[1]
$$
On observe que $Q$ appartient à $D_\QCoh^+(\mathcal{O}_Y)$ (Catégories dérivées des schémas, lemme \ref{perfect-lemma-quasi-coherence-direct-image}) et est à support dans $Y \setminus V$ (Catégories dérivées des schémas, définition \ref{perfect-definition-supported-on}). En appliquant $a$, on obtient un triangle distingué
$$
a(K) \to a(Rj_*K|_V) \to a(Q) \to a(K)[1]
$$
sur $X$. Si $a(Q)$ est à support dans $X \setminus U$, la restriction à $U$ de l'application $a(K)|_U \to a(Rj_*K|_V)|_U$ est un isomorphisme, c'est-à-dire que (\ref{equation-sheafy}) est un isomorphisme sur $K$. La réciproque est immédiate.

\medskip\noindent
La démonstration de (2) est exactement la même que celle de (1).

\medskip\noindent
Démonstration de (3). Supposons satisfaites les conditions équivalentes de (1). Posons $T = Y \setminus V$. Nous utiliserons les notations $D_{\QCoh, T}(\mathcal{O}_Y)$ et $D_{\QCoh, f^{-1}(T)}(\mathcal{O}_X)$ pour les complexes dont les faisceaux de cohomologie sont à support dans $T$ et $f^{-1}(T)$. Puisque $a$ commute aux sommes directes, la sous-catégorie triangulée, saturée et strictement pleine $\mathcal{D}$ dont les objets sont les
$$
\{Q \in D_{\QCoh, T}(\mathcal{O}_Y) \mid
a(Q) \in D_{\QCoh, f^{-1}(T)}(\mathcal{O}_X)\}
$$
est stable par sommes directes, donc par colimites dérivées. D'autre part, la catégorie $D_{\QCoh, T}(\mathcal{O}_Y)$ est engendrée par un objet parfait $E$ (voir Catégories dérivées des schémas, lemme \ref{perfect-lemma-generator-with-support}). Par hypothèse, on a $E \in \mathcal{D}$. D'après Catégories dérivées, lemme \ref{derived-lemma-write-as-colimit}, tout objet $Q$ de $D_{\QCoh, T}(\mathcal{O}_Y)$ est une colimite dérivée d'un système $Q_1 \to Q_2 \to Q_3 \to \ldots$ tel que les cônes des applications de transition soient des sommes directes de décalés de $E$. Par récurrence, on voit que $Q_n \in \mathcal{D}$ pour tout $n$, puis que $Q$ appartient à $\mathcal{D}$. Les conditions équivalentes de (2) sont donc satisfaites.
\end{proof}

\begin{lemma}
\label{lemma-proper-noetherian}
Soit $Y$ un schéma quasi-compact et quasi-séparé. Soit $f : X \to Y$ un morphisme propre. Si\footnote{Cette démonstration s'applique aux morphismes de schémas quasi-compacts et quasi-séparés tels que $Rf_*P$ soit pseudo-cohérent pour tout $P$ parfait sur $X$. Un théorème de Kiehl \cite{Kiehl} montre aisément qu'il en est ainsi si $f$ est propre et pseudo-cohérent. C'est le degré de généralité approprié pour ce lemme et certains autres résultats de ce chapitre.}
\begin{enumerate}
\item $f$ est plat et de présentation finie ; ou
\item $Y$ est noethérien,
\end{enumerate}
alors les conditions équivalentes du lemme \ref{lemma-when-sheafy}, (1), sont satisfaites pour tous les ouverts quasi-compacts $V$ de $Y$.
\end{lemma}

\begin{proof}
Soit $Q \in D^+_\QCoh(\mathcal{O}_Y)$ à support dans $Y \setminus V$. Pour obtenir une contradiction, supposons que $a(Q)$ ne soit pas à support dans $X \setminus U$. On peut alors trouver un complexe parfait $P_U$ sur $U$ et une application non nulle $P_U \to a(Q)|_U$ (cela résulte de Catégories dérivées des schémas, théorème \ref{perfect-theorem-bondal-van-den-Bergh}). En utilisant Catégories dérivées des schémas, lemme \ref{perfect-lemma-lift-map-from-perfect-complex-with-support}, on peut donc supposer qu'il existe un complexe parfait $P$ sur $X$ et une application $P \to a(Q)$ dont la restriction à $U$ est non nulle. Par définition de $a$, cette application est adjointe à une application $Rf_*P \to Q$.

\medskip\noindent
Le complexe $Rf_*P$ est pseudo-cohérent. Dans le cas (1), cela résulte de Catégories dérivées des schémas, lemme \ref{perfect-lemma-flat-proper-pseudo-coherent-direct-image-general}. Dans le cas (2), cela résulte de Catégories dérivées des schémas, lemmes \ref{perfect-lemma-direct-image-coherent} et \ref{perfect-lemma-identify-pseudo-coherent-noetherian}. On peut donc appliquer Catégories dérivées des schémas, lemme \ref{perfect-lemma-map-from-pseudo-coherent-to-complex-with-support}, pour obtenir une application $I \to \mathcal{O}_Y$ entre complexes parfaits, dont la restriction à $V$ est un isomorphisme et telle que la composée $I \otimes^\mathbf{L}_{\mathcal{O}_Y} Rf_*P \to Rf_*P \to Q$ soit nulle. D'après Catégories dérivées des schémas, lemme \ref{perfect-lemma-cohomology-base-change}, on a $I \otimes^\mathbf{L}_{\mathcal{O}_Y} Rf_*P =
Rf_*(Lf^*I \otimes^\mathbf{L}_{\mathcal{O}_X} P)$. On en conclut que la composée
$$
Lf^*I \otimes^\mathbf{L}_{\mathcal{O}_X} P \to P \to a(Q)
$$
est nulle. Cependant, sa restriction à $U$ est l'application $P|_U \to a(Q)|_U$, supposée non nulle. Cette contradiction achève la démonstration.
\end{proof}












\section{Adjoint à droite de l'image directe et changement de base, I}
\label{section-base-change-map}

\begingroup\emergencystretch=16pt
\noindent
L'application (\ref{equation-sheafy}) est un cas particulier d'une application de changement de base. Supposons en effet donné un diagramme cartésien
$$
\xymatrix{
X' \ar[r]_{g'} \ar[d]_{f'} & X \ar[d]^f \\
Y' \ar[r]^g & Y
}
$$
de schémas quasi-compacts et quasi-séparés, c'est-à-dire un diagramme comme en (\ref{equation-base-change}). Supposons $f$ et $g$ {\bf Tor-indépendants}. On peut alors considérer le morphisme de foncteurs
$D_\QCoh(\mathcal{O}_Y) \to D_\QCoh(\mathcal{O}_{X'})$
défini par la composée
\begin{equation}
\label{equation-base-change-map}
L(g')^* \circ a \to
L(g')^* \circ a \circ Rg_* \circ Lg^* \leftarrow
L(g')^* \circ Rg'_* \circ a' \circ Lg^* \to a' \circ Lg^*
\end{equation}
La première flèche provient de l'application d'adjonction $\text{id} \to Rg_* Lg^*$, et la dernière de l'application d'adjonction $L(g')^*Rg'_* \to \text{id}$. L'hypothèse de Tor-indépendance est nécessaire pour inverser la flèche du milieu, voir le lemme \ref{lemma-flat-precompose-pus}. On peut aussi voir (\ref{equation-base-change-map}), par adjonction entre $L(g')^*$ et $R(g')_*$, comme une transformation naturelle
$$
a \to a \circ Rg_* \circ Lg^* \leftarrow Rg'_* \circ a' \circ Lg^*
$$
où, là encore, la deuxième flèche est inversible. Si $M \in D_\QCoh(\mathcal{O}_X)$ et $K \in D_\QCoh(\mathcal{O}_Y)$, cette application est donnée sur les foncteurs de Yoneda par
\begin{align*}
\Hom_X(M, a(K))
& =
\Hom_Y(Rf_*M, K) \\
& \to
\Hom_Y(Rf_*M, Rg_* Lg^*K) \\
& =
\Hom_{Y'}(Lg^*Rf_*M, Lg^*K) \\
& \leftarrow
\Hom_{Y'}(Rf'_* L(g')^*M, Lg^*K) \\
& =
\Hom_{X'}(L(g')^*M, a'(Lg^*K)) \\
& =
\Hom_X(M, Rg'_*a'(Lg^*K))
\end{align*}
(où la flèche dirigée vers la gauche est inversible par le théorème de changement de base de Catégories dérivées des schémas, lemme \ref{perfect-lemma-compare-base-change}), ce qui explicite quelque peu la construction.\par\endgroup

\medskip\noindent
Nous montrons d'abord dans cette section que l'application de changement de base satisfait certaines compatibilités naturelles relatives à la juxtaposition de carrés, analogues à celles de Cohomologie, remarques \ref{cohomology-remark-compose-base-change} et \ref{cohomology-remark-compose-base-change-horizontal}, pour l'application usuelle de changement de base. On pourra omettre la suite de cette section en première lecture.

\begin{lemma}
\label{lemma-compose-base-change-maps}
Considérons un diagramme commutatif
$$
\xymatrix{
X' \ar[r]_k \ar[d]_{f'} & X \ar[d]^f \\
Y' \ar[r]^l \ar[d]_{g'} & Y \ar[d]^g \\
Z' \ar[r]^m & Z
}
$$
de schémas quasi-compacts et quasi-séparés dont les deux carrés sont cartésiens et où $f$ et $l$, ainsi que $g$ et $m$, sont Tor-indépendants. Alors la composée des applications (\ref{equation-base-change-map}) associées aux deux carrés est l'application de changement de base associée au rectangle extérieur (voir la démonstration pour un énoncé précis).
\end{lemma}

\begin{proof}
Les hypothèses impliquent que $g \circ f$ et $m$ sont Tor-indépendants (les détails sont omis), de sorte que l'énoncé a un sens. Dans cette démonstration, nous écrivons $k^*$ au lieu de $Lk^*$ et $f_*$ au lieu de $Rf_*$. Soient $a$, $b$ et $c$ les adjoints à droite du lemme \ref{lemma-twisted-inverse-image} pour $f$, $g$ et $g \circ f$, et de même pour les versions munies d'un prime. La flèche correspondant au carré supérieur est la composée
$$
\gamma_{top} :
k^* \circ a \to k^* \circ a \circ l_* \circ l^*
\xleftarrow{\xi_{top}} k^* \circ k_* \circ a' \circ l^* \to a' \circ l^*
$$
où $\xi_{top} : k_* \circ a' \to a \circ l_*$ est un isomorphisme (donc inversible), qui est la flèche « duale » de l'application de changement de base $l^* \circ f_* \to f'_* \circ k^*$. Les flèches extérieures proviennent des applications canoniques $1 \to l_* \circ l^*$ et $k^* \circ k_* \to 1$. De même, pour le deuxième carré, on a
$$
\gamma_{bot} :
l^* \circ b \to l^* \circ b \circ m_* \circ m^*
\xleftarrow{\xi_{bot}} l^* \circ l_* \circ b' \circ m^* \to b' \circ m^*
$$
Pour le rectangle extérieur, on obtient
$$
\gamma_{rect} :
k^* \circ c \to k^* \circ c \circ m_* \circ m^*
\xleftarrow{\xi_{rect}} k^* \circ k_* \circ c' \circ m^* \to c' \circ m^*
$$
On a $(g \circ f)_* = g_* \circ f_*$, donc $c = a \circ b$, et de même $c' = a' \circ b'$. Le lemme affirme que $\gamma_{rect}$ est égale à la composée
$$
k^* \circ c = k^* \circ a \circ b \xrightarrow{\gamma_{top}}
a' \circ l^* \circ b \xrightarrow{\gamma_{bot}}
a' \circ b' \circ m^* = c' \circ m^*
$$
Pour le voir, considérons le diagramme suivant :
$$
\xymatrix{
& & k^* \circ a \circ b \ar[d] \ar[lldd] \\
& & k^* \circ a \circ l_* \circ l^* \circ b \ar[ld] \\
k^* \circ a \circ b \circ m_* \circ m^* \ar[r] &
k^* \circ a \circ l_* \circ l^* \circ b \circ m_* \circ m^* &
k^* \circ k_* \circ a' \circ l^* \circ b \ar[u]_{\xi_{top}} \ar[d] \ar[ld] \\
& k^*\circ k_* \circ a' \circ l^* \circ b \circ m_* \circ m^*
\ar[u]_{\xi_{top}} \ar[rd] &
a' \circ l^* \circ b \ar[d] \\
k^* \circ k_* \circ a' \circ b' \circ m^* \ar[uu]_{\xi_{rect}} \ar[ddrr] &
k^*\circ k_* \circ a' \circ l^* \circ l_* \circ b' \circ m^*
\ar[u]_{\xi_{bot}} \ar[l] \ar[dr] &
a' \circ l^* \circ b \circ m_* \circ m^* \\
& & a' \circ l^* \circ l_* \circ b' \circ m^* \ar[u]_{\xi_{bot}} \ar[d] \\
& & a' \circ b' \circ m^*
}
$$
En descendant par le côté droit, on obtient la composée ; en descendant par le côté gauche, on obtient $\gamma_{rect}$. Tous les quadrilatères de la partie droite de ce diagramme sont commutatifs d'après Catégories, lemme \ref{categories-lemma-properties-2-cat-cats}, ou, plus simplement, d'après la discussion précédant Catégories, définition \ref{categories-definition-horizontal-composition}. Il suffit donc de montrer que le diagramme
$$
\xymatrix{
a \circ l_* \circ l^* \circ b \circ m_* &
a \circ b \circ m_* \ar[l] \\
k_* \circ a' \circ l^* \circ b \circ m_* \ar[u]_{\xi_{top}} & \\
k_* \circ a' \circ l^* \circ l_* \circ b' \ar[u]_{\xi_{bot}} \ar[r] &
k_* \circ a' \circ b' \ar[uu]_{\xi_{rect}}
}
$$
devient commutatif lorsqu'on inverse les flèches $\xi_{top}$, $\xi_{bot}$ et $\xi_{rect}$ (ce qui ne revient pas à demander que le diagramme soit commutatif). Or le diagramme
$$
\xymatrix{
& a \circ l_* \circ l^* \circ b \circ m_* \\
a \circ l_* \circ l^* \circ l_* \circ b'
\ar[ru]^{\xi_{bot}} & &
k_* \circ a' \circ l^* \circ b \circ m_* \ar[ul]_{\xi_{top}} \\
& k_* \circ a' \circ l^* \circ l_* \circ b'
\ar[ul]^{\xi_{top}} \ar[ur]_{\xi_{bot}}
}
$$
est commutatif d'après Catégories, lemme \ref{categories-lemma-properties-2-cat-cats}. Puisque les diagrammes
$$
\vcenter{
\xymatrix{
a \circ l_* \circ l^* \circ b \circ m_* & a \circ b \circ m \ar[l] \\
a \circ l_* \circ l^* \circ l_* \circ b' \ar[u] &
a \circ l_* \circ b' \ar[l] \ar[u]
}
}
\quad\text{et}\quad
\vcenter{
\xymatrix{
a \circ l_* \circ l^* \circ l_* \circ b' \ar[r] & a \circ l_* \circ b' \\
k_* \circ a' \circ l^* \circ l_* \circ b' \ar[u] \ar[r] &
k_* \circ a' \circ b' \ar[u]
}
}
$$
sont commutatifs (voir les références citées) et que la composée de $l_* \to l_* \circ l^* \circ l_* \to l_*$ est l'identité, il suffit de montrer que
$$
k \circ a' \circ b' \xrightarrow{\xi_{bot}} a \circ l_* \circ b
\xrightarrow{\xi_{top}} a \circ b \circ m_*
$$
est égale à $\xi_{rect}$ (via les identifications $a \circ b = c$ et $a' \circ b' = c'$). C'est l'énoncé dual de Cohomologie, remarque \ref{cohomology-remark-compose-base-change}, ce qui achève la démonstration.
\end{proof}

\begin{lemma}
\label{lemma-compose-base-change-maps-horizontal}
Considérons un diagramme commutatif
$$
\xymatrix{
X'' \ar[r]_{g'} \ar[d]_{f''} & X' \ar[r]_g \ar[d]_{f'} & X \ar[d]^f \\
Y'' \ar[r]^{h'} & Y' \ar[r]^h & Y
}
$$
de schémas quasi-compacts et quasi-séparés dont les deux carrés sont cartésiens et où $f$ et $h$, ainsi que $f'$ et $h'$, sont Tor-indépendants. Alors la composée des applications (\ref{equation-base-change-map}) associées aux deux carrés est l'application de changement de base associée au rectangle extérieur (voir la démonstration pour un énoncé précis).
\end{lemma}

\begin{proof}
Les hypothèses impliquent que $f$ et $h \circ h'$ sont Tor-indépendants (les détails sont omis), de sorte que l'énoncé a un sens. Dans cette démonstration, nous écrivons $g^*$ au lieu de $Lg^*$ et $f_*$ au lieu de $Rf_*$. Soient $a$, $a'$ et $a''$ les adjoints à droite du lemme \ref{lemma-twisted-inverse-image} pour $f$, $f'$ et $f''$. La flèche correspondant au carré de droite est la composée
$$
\gamma_{right} :
g^* \circ a \to g^* \circ a \circ h_* \circ h^*
\xleftarrow{\xi_{right}} g^* \circ g_* \circ a' \circ h^* \to a' \circ h^*
$$
où $\xi_{right} : g_* \circ a' \to a \circ h_*$ est un isomorphisme (donc inversible), qui est la flèche « duale » de l'application de changement de base $h^* \circ f_* \to f'_* \circ g^*$. Les flèches extérieures proviennent des applications canoniques $1 \to h_* \circ h^*$ et $g^* \circ g_* \to 1$. De même, pour le carré de gauche, on a
$$
\gamma_{left} :
(g')^* \circ a' \to (g')^* \circ a' \circ (h')_* \circ (h')^*
\xleftarrow{\xi_{left}}
(g')^* \circ (g')_* \circ a'' \circ (h')^* \to a'' \circ (h')^*
$$
Pour le rectangle extérieur, on obtient
$$
\gamma_{rect} :
k^* \circ a \to
k^* \circ a \circ m_* \circ m^* \xleftarrow{\xi_{rect}}
k^* \circ k_* \circ a'' \circ m^* \to
a'' \circ m^*
$$
où $k = g \circ g'$ et $m = h \circ h'$. On a $k^* = (g')^* \circ g^*$ et $m^* = (h')^* \circ h^*$. Le lemme affirme que $\gamma_{rect}$ est égale à la composée
$$
k^* \circ a =
(g')^* \circ g^* \circ a \xrightarrow{\gamma_{right}}
(g')^* \circ a' \circ h^* \xrightarrow{\gamma_{left}}
a'' \circ (h')^* \circ h^* = a'' \circ m^*
$$
Pour le voir, considérons le diagramme suivant :
$$
\xymatrix{
& (g')^* \circ g^* \circ a \ar[d] \ar[ddl] \\
& (g')^* \circ g^* \circ a \circ h_* \circ h^* \ar[ld] \\
(g')^* \circ g^* \circ a \circ h_* \circ (h')_* \circ (h')^* \circ h^* &
(g')^* \circ g^* \circ g_* \circ a' \circ h^*
\ar[u]_{\xi_{right}} \ar[d] \ar[ld] \\
(g')^* \circ g^* \circ g_* \circ a' \circ (h')_* \circ (h')^* \circ h^*
\ar[u]_{\xi_{right}} \ar[dr] &
(g')^* \circ a' \circ h^* \ar[d] \\
(g')^* \circ g^* \circ g_* \circ (g')_* \circ a'' \circ (h')^* \circ h^*
\ar[u]_{\xi_{left}} \ar[ddr] \ar[dr] &
(g')^* \circ a' \circ (h')_* \circ (h')^* \circ h^* \\
& (g')^*\circ (g')_* \circ a'' \circ (h')^* \circ h^*
\ar[u]_{\xi_{left}} \ar[d] \\
& a'' \circ (h')^* \circ h^*
}
$$
En descendant par le côté droit, on obtient la composée ; en descendant par le côté gauche, on obtient $\gamma_{rect}$. Tous les quadrilatères de la partie droite de ce diagramme sont commutatifs d'après Catégories, lemme \ref{categories-lemma-properties-2-cat-cats}, ou, plus simplement, d'après la discussion précédant Catégories, définition \ref{categories-definition-horizontal-composition}. Il suffit donc de montrer que
$$
g_* \circ (g')_* \circ a'' \xrightarrow{\xi_{left}}
g_* \circ a' \circ (h')_* \xrightarrow{\xi_{right}}
a \circ h_* \circ (h')_*
$$
est égale à $\xi_{rect}$. C'est l'énoncé dual de Cohomologie, remarque \ref{cohomology-remark-compose-base-change-horizontal}, ce qui achève la démonstration.
\end{proof}

\begin{remark}
\label{remark-going-around}
Considérons un diagramme commutatif
$$
\xymatrix{
X'' \ar[r]_{k'} \ar[d]_{f''} & X' \ar[r]_k \ar[d]_{f'} & X \ar[d]^f \\
Y'' \ar[r]^{l'} \ar[d]_{g''} & Y' \ar[r]^l \ar[d]_{g'} & Y \ar[d]^g \\
Z'' \ar[r]^{m'} & Z' \ar[r]^m & Z
}
$$
de schémas quasi-compacts et quasi-séparés dont tous les carrés sont cartésiens, et où $(f, l)$, $(g, m)$, $(f', l')$, $(g', m')$ sont des couples de morphismes Tor-indépendants. Soient $a$, $a'$, $a''$, $b$, $b'$, $b''$ les adjoints à droite du lemme \ref{lemma-twisted-inverse-image} pour $f$, $f'$, $f''$, $g$, $g'$, $g''$. Désignons les carrés du diagramme par $A$, $B$, $C$, $D$ comme suit :
$$
\begin{matrix}
A & B \\
C & D
\end{matrix}
$$
Les applications (\ref{equation-base-change-map}) associées aux carrés sont alors (en utilisant $k^* = Lk^*$, etc.)
$$
\begin{matrix}
\gamma_A : (k')^* \circ a' \to a'' \circ (l')^* &
\gamma_B : k^* \circ a \to a' \circ l^* \\
\gamma_C : (l')^* \circ b' \to b'' \circ (m')^* &
\gamma_D : l^* \circ b \to b' \circ m^*
\end{matrix}
$$
Pour les rectangles $2 \times 1$ et $1 \times 2$, on dispose de quatre autres applications de changement de base
$$
\begin{matrix}
\gamma_{A + B} : (k \circ k')^* \circ a \to a'' \circ (l \circ l')^* \\
\gamma_{C + D} : (l \circ l')^* \circ b \to b'' \circ (m \circ m')^* \\
\gamma_{A + C} : (k')^* \circ (a' \circ b') \to (a'' \circ b'') \circ (m')^* \\
\gamma_{B + D} : k^* \circ (a \circ b) \to (a' \circ b') \circ m^*
\end{matrix}
$$
D'après le lemme \ref{lemma-compose-base-change-maps-horizontal}, on a
$$
\gamma_{A + B} = \gamma_A \circ \gamma_B, \quad
\gamma_{C + D} = \gamma_C \circ \gamma_D
$$
et, d'après le lemme \ref{lemma-compose-base-change-maps}, on a
$$
\gamma_{A + C} = \gamma_C \circ \gamma_A, \quad
\gamma_{B + D} = \gamma_D \circ \gamma_B
$$
Il serait ici plus correct d'écrire $\gamma_{A + B} = (\gamma_A \star \text{id}_{l^*}) \circ
(\text{id}_{(k')^*} \star \gamma_B)$ avec les notations de Catégories, section \ref{categories-section-formal-cat-cat}, et de même pour les autres expressions. Nous conservons toutefois l'abus de notation des démonstrations des lemmes \ref{lemma-compose-base-change-maps} et \ref{lemma-compose-base-change-maps-horizontal}, en omettant les produits $\star$ avec les identités : on peut reconstituer ceux-ci dès lors que la source et le but de la transformation sont connus. Cela étant, on obtient a priori deux transformations
$$
(k')^* \circ k^* \circ a \circ b
\longrightarrow
a'' \circ b'' \circ (m')^* \circ m^*
$$
à savoir
$$
\gamma_C \circ \gamma_A \circ \gamma_D \circ \gamma_B =
\gamma_{A + C} \circ \gamma_{B + D}
$$
et
$$
\gamma_C \circ \gamma_D \circ \gamma_A \circ \gamma_B =
\gamma_{C + D} \circ \gamma_{A + B}
$$
Cette remarque a pour objet de signaler que ces transformations sont égales. Pour le voir, il suffit de montrer que
$$
\xymatrix{
(k')^* \circ a' \circ l^* \circ b \ar[r]_{\gamma_D} \ar[d]_{\gamma_A} &
(k')^* \circ a' \circ b' \circ m^* \ar[d]^{\gamma_A} \\
a'' \circ (l')^* \circ l^* \circ b \ar[r]^{\gamma_D} &
a'' \circ (l')^* \circ b' \circ m^*
}
$$
est commutatif. Cela résulte de Catégories, lemme \ref{categories-lemma-properties-2-cat-cats}, ou, plus simplement, de la discussion précédant Catégories, définition \ref{categories-definition-horizontal-composition}.
\end{remark}







\section{Adjoint à droite de l'image directe et changement de base, II}
\label{section-base-change-II}

\noindent
Nous montrons dans cette section que l'application de changement de base de la section \ref{section-base-change-map} est un isomorphisme dans certains cas. Observons d'abord qu'il suffit de le vérifier sur des ouverts affines, pourvu que la formation de l'adjoint à droite de l'image directe commute à la restriction aux ouverts.

\begin{remark}
\label{remark-check-over-affines}
Considérons un diagramme cartésien
$$
\xymatrix{
X' \ar[r]_{g'} \ar[d]_{f'} & X \ar[d]^f \\
Y' \ar[r]^g & Y
}
$$
de schémas quasi-compacts et quasi-séparés, avec $(g, f)$ Tor-indépendants. Soient $V \subset Y$ et $V' \subset Y'$ des ouverts affines tels que $g(V') \subset V$. Formons les diagrammes cartésiens
$$
\vcenter{
\xymatrix{
U \ar[r] \ar[d] & X \ar[d] \\
V \ar[r] & Y
}
}
\quad\text{et}\quad
\vcenter{
\xymatrix{
U' \ar[r] \ar[d] & X' \ar[d] \\
V' \ar[r] & Y'
}
}
$$
Supposons que (\ref{equation-sheafy}), pour $K$ et le premier diagramme, et (\ref{equation-sheafy}), pour $Lg^*K$ et le deuxième diagramme, soient des isomorphismes. Alors la restriction de l'application de changement de base (\ref{equation-base-change-map})
$$
L(g')^*a(K) \longrightarrow a'(Lg^*K)
$$
à $U'$ est isomorphe à l'application de changement de base (\ref{equation-base-change-map}) pour $K|_V$ et le diagramme cartésien
$$
\xymatrix{
U' \ar[r] \ar[d] & U \ar[d] \\
V' \ar[r] & V
}
$$
Cela résulte du fait que (\ref{equation-sheafy}) est un cas particulier de l'application de changement de base (\ref{equation-base-change-map}), et de la compatibilité de ces applications avec la juxtaposition horizontale de carrés, voir le lemme \ref{lemma-compose-base-change-maps-horizontal}. Pour vérifier que la restriction de l'application de changement de base à $U'$ est un isomorphisme, il suffit donc de travailler avec le dernier diagramme.
\end{remark}

\begin{lemma}
\label{lemma-more-base-change}
Dans le diagramme (\ref{equation-base-change}), supposons que
\begin{enumerate}
\item $g : Y' \to Y$ soit un morphisme de schémas affines ;
\item $f : X \to Y$ soit propre ;
\item $f$ et $g$ soient Tor-indépendants.
\end{enumerate}
Alors l'application de changement de base (\ref{equation-base-change-map}) induit un isomorphisme
$$
L(g')^*a(K) \longrightarrow a'(Lg^*K)
$$
dans les cas suivants :
\begin{enumerate}
\item pour tout $K \in D_\QCoh(\mathcal{O}_Y)$ si $f$ est plat de présentation finie ;
\item pour tout $K \in D_\QCoh(\mathcal{O}_Y)$ si $f$ est parfait et $Y$ noethérien ;
\item pour $K \in D_\QCoh^+(\mathcal{O}_Y)$ si $g$ est de Tor-dimension finie et $Y$ noethérien.
\end{enumerate}
\end{lemma}

\begin{proof}
Écrivons $Y = \Spec(A)$ et $Y' = \Spec(A')$. Étant obtenu par changement de base d'un morphisme affine, le morphisme $g'$ est affine. Soit $M$ un générateur parfait de $D_\QCoh(\mathcal{O}_X)$, voir Catégories dérivées des schémas, théorème \ref{perfect-theorem-bondal-van-den-Bergh}. Alors $L(g')^*M$ est un générateur de $D_\QCoh(\mathcal{O}_{X'})$, voir Catégories dérivées des schémas, remarque \ref{perfect-remark-pullback-generator}. Il suffit donc de montrer que (\ref{equation-base-change-map}) induit un isomorphisme
\begin{equation}
\label{equation-iso}
R\Hom_{X'}(L(g')^*M, L(g')^*a(K))
\longrightarrow
R\Hom_{X'}(L(g')^*M, a'(Lg^*K))
\end{equation}
entre les complexes Hom globaux, voir Cohomologie, section \ref{cohomology-section-global-RHom}, car cela impliquera que le cône de $L(g')^*a(K) \to a'(Lg^*K)$ est nul. La démonstration s'organise comme suit : nous montrerons d'abord que ces complexes Hom sont isomorphes, puis, dans la dernière partie, que l'isomorphisme est induit par (\ref{equation-iso}).

\medskip\noindent
Le membre de gauche. Puisque $M$ est parfait, l'application canonique
$$
R\Hom_X(M, a(K)) \otimes^\mathbf{L}_A A'
\longrightarrow
R\Hom_{X'}(L(g')^*M, L(g')^*a(K))
$$
est un isomorphisme d'après Catégories dérivées des schémas, lemme \ref{perfect-lemma-affine-morphism-and-hom-out-of-perfect}. En la combinant avec l'isomorphisme
$R\Hom_Y(Rf_*M, K) = R\Hom_X(M, a(K))$
du lemme \ref{lemma-iso-global-hom}, on obtient que le membre de gauche est égal à $R\Hom_Y(Rf_*M, K) \otimes^\mathbf{L}_A A'$.

\medskip\noindent
Le membre de droite. On utilise d'abord l'isomorphisme
$$
R\Hom_{X'}(L(g')^*M, a'(Lg^*K)) = R\Hom_{Y'}(Rf'_*L(g')^*M, Lg^*K)
$$
du lemme \ref{lemma-iso-global-hom}. On utilise ensuite le fait que l'application de changement de base $Lg^*Rf_*M \to Rf'_*L(g')^*M$ est un isomorphisme d'après Catégories dérivées des schémas, lemme \ref{perfect-lemma-compare-base-change}. On peut donc réécrire ce membre sous la forme
\par\noindent\hfil
$R\Hom_{Y'}(Lg^*Rf_*M, Lg^*K)$.
\hfil\par\noindent
Puisque $Y$, $Y'$ sont affines et que $K$, $Rf_*M$ appartiennent à $D_\QCoh(\mathcal{O}_Y)$ (Catégories dérivées des schémas, lemme \ref{perfect-lemma-quasi-coherence-direct-image}), on dispose d'une application canonique
$$
\beta :
R\Hom_Y(Rf_*M, K) \otimes^\mathbf{L}_A A'
\longrightarrow
R\Hom_{Y'}(Lg^*Rf_*M, Lg^*K)
$$
dans $D(A')$. C'est la flèche de Compléments d'algèbre, équation (\ref{more-algebra-equation-base-change-RHom}), où l'on utilise Catégories dérivées des schémas, lemmes \ref{perfect-lemma-affine-compare-bounded} et \ref{perfect-lemma-quasi-coherence-internal-hom}, pour passer au cadre algébrique et en revenir.
\begin{enumerate}
\item Si $f$ est plat de présentation finie, le complexe $Rf_*M$ est parfait sur $Y$ d'après Catégories dérivées des schémas, lemme \ref{perfect-lemma-flat-proper-perfect-direct-image-general}, et $\beta$ est un isomorphisme d'après Compléments d'algèbre, lemme \ref{more-algebra-lemma-base-change-RHom}, (1).
\item Si $f$ est parfait et $Y$ noethérien, le complexe $Rf_*M$ est parfait sur $Y$ d'après Compléments sur les morphismes, lemme \ref{more-morphisms-lemma-perfect-proper-perfect-direct-image}, et $\beta$ est un isomorphisme comme précédemment.
\item Si $g$ est de Tor-dimension finie et $Y$ est noethérien, le complexe $Rf_*M$ est pseudo-cohérent sur $Y$ (Catégories dérivées des schémas, lemmes \ref{perfect-lemma-direct-image-coherent} et \ref{perfect-lemma-identify-pseudo-coherent-noetherian}), et $\beta$ est un isomorphisme d'après Compléments d'algèbre, lemme \ref{more-algebra-lemma-base-change-RHom}, (4).
\end{enumerate}
On obtient donc le même résultat qu'au paragraphe précédent.

\medskip\noindent
\begingroup\emergencystretch=2em
Dans la suite de la démonstration, nous montrons que les identifications des membres de gauche et de droite de (\ref{equation-iso}) données aux deuxième et troisième paragraphes sont effectivement réalisées par (\ref{equation-iso}). Pour alléger les formules, nous utiliserons $(-, -)_X = R\Hom_X(-, -)$, écrirons $- \otimes A'$ au lieu de $- \otimes_A^\mathbf{L} A'$, et poserons les abréviations $g^* = Lg^*$ et $f_* = Rf_*$. Considérons le diagramme commutatif suivant :\par\endgroup
\begingroup\small
$$
\xymatrix{
((g')^*M, (g')^*a(K))_{X'} \ar[d] &
(M, a(K))_X \otimes A' \ar[l]^-\alpha \ar[d] &
(f_*M, K)_Y \otimes A' \ar@{=}[l] \ar[d] \\
((g')^*M, (g')^*a(g_*g^*K))_{X'} &
(M, a(g_*g^*K))_X \otimes A' \ar[l]^-\alpha &
(f_*M, g_*g^*K)_Y \otimes A' \ar@{=}[l] \ar@/_4pc/[dd]_{\mu'} \\
((g')^*M, (g')^*g'_*a'(g^*K))_{X'} \ar[u] \ar[d] &
(M, g'_*a'(g^*K))_X \otimes A' \ar[u] \ar[l]^-\alpha \ar[ld]^\mu &
(f_*M, K) \otimes A' \ar[d]^\beta \\
((g')^*M, a'(g^*K))_{X'} &
(f'_*(g')^*M, g^*K)_{Y'} \ar@{=}[l] \ar[r] &
(g^*f_*M, g^*K)_{Y'}
}
$$
\endgroup
Les flèches notées $\alpha$ sont les applications de Catégories dérivées des schémas, lemme \ref{perfect-lemma-affine-morphism-and-hom-out-of-perfect}, pour le diagramme de sommets $X', X, Y', Y$. La partie supérieure du diagramme est commutative, car les flèches horizontales sont fonctorielles en leurs arguments. Les flèches verticales centrales proviennent de la transformation inversible $g'_* \circ a' \to a \circ g_*$ du lemme \ref{lemma-flat-precompose-pus} ; le carré central est donc commutatif. La descente par le côté gauche est (\ref{equation-iso}). Les flèches horizontales supérieures fournissent les identifications utilisées au deuxième paragraphe de la démonstration. Les flèches horizontales inférieures, y compris $\beta$, fournissent celles du troisième paragraphe. Étant donnés $E \in D(A)$, $E' \in D(A')$ et $c : E \to E'$ dans $D(A)$, nous noterons $\mu_c : E \otimes A' \to E'$ l'application induite par $c$ et par l'adjonction entre restriction et changement de base ; lorsque $c$ est clair, nous écrirons $\mu = \mu_c$, en omettant donc $c$. L'application $\mu$ du diagramme est de cette forme, avec $c$ donné par l'identification
$(M, g'_*a(g^*K))_X = ((g')^*M, a'(g^*K))_{X'}$
; le triangle faisant intervenir $\mu$ est commutatif d'après Catégories dérivées des schémas, remarque \ref{perfect-remark-multiplication-map}.

\medskip\noindent
Observons que
$$
\xymatrix{
(M, a(g_*g^*K))_X &
(f_*M, g_* g^*K)_Y \ar@{=}[l] &
(g^*f_*M, g^*K)_{Y'} \ar@{=}[l] \\
(M, g'_* a'(g^*K))_X \ar[u] &
((g')^*M, a'(g^*K))_{X'} \ar@{=}[l] &
(f'_*(g')^*M, g^*K)_{Y'} \ar@{=}[l] \ar[u]
}
$$
est commutatif par définition même de la transformation $g'_* \circ a' \to a \circ g_*$. En prenant $\mu'$ comme ci-dessus, correspondant à l'identification $(f_*M, g_*g^*K)_X = (g^*f_*M, g^*K)_{Y'}$, l'hexagone est lui aussi commutatif. Il suffit donc de montrer que $\beta$ est égale à la composée de
$(f_*M, K)_Y \otimes A' \to (f_*M, g_*g^*K)_X \otimes A'$
et de $\mu'$. Pour cela, il suffit de prouver que les deux applications induites $(f_*M, K)_Y \to (g^*f_*M, g^*K)_{Y'}$ sont égales. Autrement dit, il suffit de montrer que le diagramme
$$
\xymatrix{
R\Hom_A(E, K) \ar[rr]_{\text{induite par }\beta} \ar[rd] & &
R\Hom_{A'}(E \otimes_A^\mathbf{L} A', K \otimes_A^\mathbf{L} A') \\
& R\Hom_A(E, K \otimes_A^\mathbf{L} A') \ar[ru]
}
$$
est commutatif pour tout $E, K \in D(A)$. Comme c'est ainsi que $\beta$ est construit dans Compléments d'algèbre, section \ref{more-algebra-section-base-change-RHom}, la démonstration est achevée.
\end{proof}








\section{Adjoint à droite de l'image directe et morphismes trace}
\label{section-trace}

\noindent
Soit $f : X \to Y$ un morphisme de schémas quasi-compacts et quasi-séparés. Soit $a : D_\QCoh(\mathcal{O}_Y) \to D_\QCoh(\mathcal{O}_X)$ l'adjoint à droite du lemme \ref{lemma-twisted-inverse-image}. D'après Catégories, section \ref{categories-section-adjoint}, on obtient une transformation de foncteurs
$$
\text{Tr}_f : Rf_* \circ a \longrightarrow \text{id}
$$
Le morphisme correspondant $\text{Tr}_{f, K} : Rf_*a(K) \longrightarrow K$, pour $K \in D_\QCoh(\mathcal{O}_Y)$, est parfois appelé {\it morphisme trace}. Il se caractérise par le fait que la bijection
$$
\Hom_X(L, a(K)) \longrightarrow \Hom_Y(Rf_*L, K)
$$
pour $L \in D_\QCoh(\mathcal{O}_X)$, qui caractérise l'adjoint à droite, est donnée par
$$
\varphi \longmapsto \text{Tr}_{f, K} \circ Rf_*\varphi
$$
L'application (\ref{equation-sheafy-trace})
$$
Rf_*R\SheafHom_{\mathcal{O}_X}(L, a(K))
\longrightarrow
R\SheafHom_{\mathcal{O}_Y}(Rf_*L, K)
$$
s'obtient par composition avec $\text{Tr}_{f, K}$. Tout morphisme trace considéré dans cette section sera un cas particulier de celui-ci. Avant d'aborder quelques cas particuliers, montrons que la formation de la trace commute au changement de base.

\begin{lemma}[Morphisme trace et changement de base]
\label{lemma-trace-map-and-base-change}
Supposons donné un diagramme (\ref{equation-base-change}) où $f$ et $g$ sont Tor-indépendants. Alors les applications $1 \star \text{Tr}_f : Lg^* \circ Rf_* \circ a \to Lg^*$ et $\text{Tr}_{f'} \star 1 : Rf'_* \circ a' \circ Lg^* \to Lg^*$ se correspondent par les applications de changement de base
$\beta : Lg^* \circ Rf_* \to Rf'_* \circ L(g')^*$
(Cohomologie, remarque \ref{cohomology-remark-base-change}) et $\alpha : L(g')^* \circ a \to a' \circ Lg^*$ (\ref{equation-base-change-map}). Plus précisément, le diagramme
$$
\xymatrix{
Lg^* \circ Rf_* \circ a
\ar[d]_{\beta \star 1} \ar[r]_-{1 \star \text{Tr}_f} &
Lg^* \\
Rf'_* \circ L(g')^* \circ a \ar[r]^{1 \star \alpha} &
Rf'_* \circ a' \circ Lg^* \ar[u]_{\text{Tr}_{f'} \star 1}
}
$$
de transformations de foncteurs est commutatif.
\end{lemma}

\begin{proof}
Dans cette démonstration, nous écrivons $f_*$ pour $Rf_*$ et $g^*$ pour $Lg^*$, et nous omettons les produits $\star$ avec les identités : on peut reconstituer ceux-ci dès lors que la source et le but de la transformation sont connus. Rappelons que $\beta : g^* \circ f_* \to f'_* \circ (g')^*$ est un isomorphisme et que $\alpha$ est défini à l'aide de l'isomorphisme $\beta^\vee : g'_* \circ a' \to a \circ g_*$, adjoint de $\beta$, voir le lemme \ref{lemma-flat-precompose-pus} et sa démonstration. Remarquons d'abord que la flèche horizontale supérieure du diagramme du lemme est égale à la composée
$$
g^* \circ f_* \circ a \to
g^* \circ f_* \circ a \circ g_* \circ g^* \to
g^* \circ g_* \circ g^* \to g^*
$$
où la première flèche est l'unité pour $(g^*, g_*)$, la deuxième est $\text{Tr}_f$ et la troisième est la co-unité pour $(g^*, g_*)$. Cela résulte simplement du fait que la composée $g^* \to g^* \circ g_* \circ g^* \to g^*$ de l'unité et de la co-unité est l'identité. Considérons le diagramme
\begingroup\small
$$
\xymatrix{
& g^* \circ f_* \circ a \ar[ld]_\beta \ar[d] \ar[r]_{\text{Tr}_f} & g^* \\
f'_* \circ (g')^* \circ a \ar[dr] &
g^* \circ f_* \circ a \circ g_* \circ g^* \ar[d]_\beta \ar[ru] &
g^* \circ f_* \circ g'_* \circ a' \circ g^* \ar[l]_{\beta^\vee} \ar[d]_\beta &
f'_* \circ a' \circ g^* \ar[lu]_{\text{Tr}_{f'}} \\
& f'_* \circ (g')^* \circ a \circ g_* \circ g^* &
f'_* \circ (g')^* \circ g'_* \circ a' \circ g^* \ar[ru] \ar[l]_{\beta^\vee}
}
$$
\endgroup
Les deux carrés y sont commutatifs d'après Catégories, lemme \ref{categories-lemma-properties-2-cat-cats}, ou, plus simplement, d'après la discussion précédant Catégories, définition \ref{categories-definition-horizontal-composition}. Le triangle est commutatif d'après ce qui précède. D'après Catégories, lemme \ref{categories-lemma-transformation-between-functors-and-adjoints}, le carré
$$
\xymatrix{
g^* \circ f_* \circ g'_* \circ a' \ar[d]_{\beta^\vee} \ar[r]_-\beta &
f'_* \circ (g')^* \circ g'_* \circ a' \ar[d] \\
g^* \circ f_* \circ a \circ g_* \ar[r] &
\text{id}
}
$$
est commutatif, ce qui implique la commutativité du pentagone dans le grand diagramme. Puisque $\beta$ et $\beta^\vee$ sont des isomorphismes et que le parcours extérieur du grand diagramme est, par définition, égal à $\text{Tr}_f \circ \alpha \circ \beta$, le lemme est démontré.
\end{proof}

\noindent
Soit $f : X \to Y$ un morphisme de schémas quasi-compacts et quasi-séparés. Soit $a : D_\QCoh(\mathcal{O}_Y) \to D_\QCoh(\mathcal{O}_X)$ l'adjoint à droite de $Rf_*$ du lemme \ref{lemma-twisted-inverse-image}. D'après Catégories, section \ref{categories-section-adjoint}, on obtient une transformation de foncteurs
$$
\eta_f : \text{id} \to  a \circ Rf_*
$$
appelée unité de l'adjonction.

\begin{lemma}
\label{lemma-unit-and-base-change}
\begingroup\emergencystretch=4em
Supposons donné un diagramme (\ref{equation-base-change}) où $f$ et $g$ sont Tor-indépendants. Alors les applications $1 \star \eta_f : L(g')^* \to L(g')^* \circ a \circ Rf_*$ et $\eta_{f'} \star 1 : L(g')^* \to a' \circ Rf'_* \circ L(g')^*$ se correspondent par les applications de changement de base
$\beta : Lg^* \circ Rf_* \to Rf'_* \circ L(g')^*$
(Cohomologie, remarque \ref{cohomology-remark-base-change}) et $\alpha : L(g')^* \circ a \to a' \circ Lg^*$ (\ref{equation-base-change-map}).\par\endgroup
Plus précisément, le diagramme
$$
\xymatrix{
L(g')^* \ar[r]_-{1 \star \eta_f} \ar[d]_{\eta_{f'} \star 1} &
L(g')^* \circ a \circ Rf_* \ar[d]^\alpha \\
a' \circ Rf'_* \circ L(g')^* &
a' \circ Lg^* \circ Rf_* \ar[l]_-\beta
}
$$
de transformations de foncteurs est commutatif.
\end{lemma}

\begin{proof}
Cette démonstration est duale de celle du lemme \ref{lemma-trace-map-and-base-change}. Nous écrivons $f_*$ pour $Rf_*$ et $g^*$ pour $Lg^*$, et omettons les produits $\star$ avec les identités : on peut reconstituer ceux-ci dès lors que la source et le but de la transformation sont connus. Rappelons que $\beta : g^* \circ f_* \to f'_* \circ (g')^*$ est un isomorphisme et que $\alpha$ est défini à l'aide de l'isomorphisme $\beta^\vee : g'_* \circ a' \to a \circ g_*$, adjoint de $\beta$, voir le lemme \ref{lemma-flat-precompose-pus} et sa démonstration. Remarquons d'abord que la flèche verticale de gauche du diagramme du lemme est égale à la composée
$$
(g')^* \to (g')^* \circ g'_* \circ (g')^* \to
(g')^* \circ g'_* \circ a' \circ f'_* \circ (g')^* \to
a' \circ f'_* \circ (g')^*
$$
où la première flèche est l'unité pour $((g')^*, g'_*)$, la deuxième est $\eta_{f'}$ et la troisième est la co-unité pour $((g')^*, g'_*)$. Cela résulte simplement du fait que la composée
$(g')^* \to (g')^* \circ (g')_* \circ (g')^* \to (g')^*$
de l'unité et de la co-unité est l'identité. Considérons le diagramme
$$
\xymatrix{
& (g')^* \circ a \circ f_* \ar[r] &
(g')^* \circ a \circ g_* \circ g^* \circ f_*
\ar[ld]_\beta \\
(g')^* \ar[ru]^{\eta_f} \ar[dd]_{\eta_{f'}} \ar[rd] &
(g')^* \circ a \circ g_* \circ f'_* \circ (g')^* &
(g')^* \circ g'_* \circ a' \circ g^* \circ f_*
\ar[u]_{\beta^\vee} \ar[ld]_\beta \ar[d] \\
& (g')^* \circ g'_* \circ a' \circ f'_* \circ (g')^*
\ar[ld] \ar[u]_{\beta^\vee} &
a' \circ g^* \circ f_* \ar[lld]^\beta \\
a' \circ f'_* \circ (g')^*
}
$$
Les deux carrés y sont commutatifs d'après Catégories, lemme \ref{categories-lemma-properties-2-cat-cats}, ou, plus simplement, d'après la discussion précédant Catégories, définition \ref{categories-definition-horizontal-composition}. Le triangle est commutatif d'après ce qui précède. Par l'énoncé dual de Catégories, lemme \ref{categories-lemma-transformation-between-functors-and-adjoints}, le carré
$$
\xymatrix{
\text{id} \ar[r] \ar[d] &
g'_* \circ a' \circ g^* \circ f_* \ar[d]^\beta \\
g'_* \circ a' \circ g^* \circ f_* \ar[r]^{\beta^\vee} &
a \circ g_* \circ f'_* \circ (g')^*
}
$$
est commutatif, ce qui implique la commutativité du pentagone dans le grand diagramme. Puisque $\beta$ et $\beta^\vee$ sont des isomorphismes et que le parcours extérieur du grand diagramme est, par définition, égal à $\beta \circ \alpha \circ \eta_f$, le lemme est démontré.
\end{proof}

\begin{example}
\label{example-trace-affine}
Soit $A \to B$ un homomorphisme d'anneaux. Soient $Y = \Spec(A)$ et $X = \Spec(B)$, et soit $f : X \to Y$ le morphisme correspondant à $A \to B$. Comme on l'a vu dans l'exemple \ref{example-affine-twisted-inverse-image}, l'adjoint à droite de
$Rf_* : D_\QCoh(\mathcal{O}_X) \to D_\QCoh(\mathcal{O}_Y)$
envoie un objet $K$ de $D(A) = D_\QCoh(\mathcal{O}_Y)$ sur $R\Hom(B, K)$ dans $D(B) = D_\QCoh(\mathcal{O}_X)$. Le morphisme trace est le morphisme
$$
\text{Tr}_{f, K} : R\Hom(B, K) \longrightarrow R\Hom(A, K) = K
$$
induite par l'application de $A$-modules $A \to B$.
\end{example}




\section{Adjoint à droite de l'image directe et image inverse}
\label{section-compare-with-pullback}

\noindent
Soit $f : X \to Y$ un morphisme de schémas quasi-compacts et quasi-séparés. Soit $a$ l'adjoint à droite de l'image directe du lemme \ref{lemma-twisted-inverse-image}. Pour $K, L \in D_\QCoh(\mathcal{O}_Y)$, on dispose d'une application canonique
$$
Lf^*K \otimes^\mathbf{L}_{\mathcal{O}_X} a(L)
\longrightarrow
a(K \otimes_{\mathcal{O}_Y}^\mathbf{L} L)
$$
Cette application est adjointe à une application
$$
Rf_*(Lf^*K \otimes^\mathbf{L}_{\mathcal{O}_X} a(L)) =
K \otimes^\mathbf{L}_{\mathcal{O}_Y} Rf_*(a(L))
\longrightarrow
K \otimes^\mathbf{L}_{\mathcal{O}_Y} L
$$
(égalité d'après Catégories dérivées des schémas, lemme \ref{perfect-lemma-cohomology-base-change}), pour laquelle on utilise le morphisme trace $Rf_*a(L) \to L$. Lorsque $L = \mathcal{O}_Y$, on obtient une application
\begin{equation}
\label{equation-compare-with-pullback}
Lf^*K \otimes^\mathbf{L}_{\mathcal{O}_X} a(\mathcal{O}_Y) \longrightarrow a(K)
\end{equation}
fonctorielle en $K$ et compatible avec les triangles distingués.

\begin{lemma}
\label{lemma-compare-with-pullback-perfect}
Soit $f : X \to Y$ un morphisme de schémas quasi-compacts et quasi-séparés. L'application
$Lf^*K \otimes^\mathbf{L}_{\mathcal{O}_X} a(L) \to
a(K \otimes_{\mathcal{O}_Y}^\mathbf{L} L)$
définie ci-dessus pour $K, L \in D_\QCoh(\mathcal{O}_Y)$ est un isomorphisme si $K$ est parfait. En particulier, (\ref{equation-compare-with-pullback}) est un isomorphisme si $K$ est parfait.
\end{lemma}

\begin{proof}
Soit $K^\vee$ le « dual » de $K$, voir Cohomologie, lemme \ref{cohomology-lemma-dual-perfect-complex}. Pour $M \in D_\QCoh(\mathcal{O}_X)$, on a
\begin{align*}
\Hom_{D(\mathcal{O}_Y)}(Rf_*M, K \otimes^\mathbf{L}_{\mathcal{O}_Y} L)
& =
\Hom_{D(\mathcal{O}_Y)}(
Rf_*M \otimes^\mathbf{L}_{\mathcal{O}_Y} K^\vee, L) \\
& =
\Hom_{D(\mathcal{O}_X)}(
M \otimes^\mathbf{L}_{\mathcal{O}_X} Lf^*K^\vee, a(L)) \\
& =
\Hom_{D(\mathcal{O}_X)}(M,
Lf^*K \otimes^\mathbf{L}_{\mathcal{O}_X} a(L))
\end{align*}
La deuxième égalité résulte de la définition de $a$ et de la formule de projection (Cohomologie, lemme \ref{cohomology-lemma-projection-formula-perfect}), ou, plus généralement, de Catégories dérivées des schémas, lemme \ref{perfect-lemma-cohomology-base-change}. Le résultat découle donc du lemme de Yoneda.
\end{proof}

\begin{lemma}
\label{lemma-restriction-compare-with-pullback}
\begingroup\emergencystretch=4em
Supposons donné un diagramme (\ref{equation-base-change}) où $f$ et $g$ sont Tor-indépendants. Soit $K \in D_\QCoh(\mathcal{O}_Y)$.\par\endgroup
Le diagramme
$$
\xymatrix{
L(g')^*(Lf^*K \otimes^\mathbf{L}_{\mathcal{O}_X} a(\mathcal{O}_Y))
\ar[r] \ar[d] & L(g')^*a(K) \ar[d] \\
L(f')^*Lg^*K \otimes_{\mathcal{O}_{X'}}^\mathbf{L} a'(\mathcal{O}_{Y'})
\ar[r] & a'(Lg^*K)
}
$$
est commutatif, où les flèches horizontales sont les applications (\ref{equation-compare-with-pullback}) pour $K$ et $Lg^*K$, et où les applications verticales sont construites à l'aide de Cohomologie, remarque \ref{cohomology-remark-base-change}, et de (\ref{equation-base-change-map}).
\end{lemma}

\begin{proof}
Dans cette démonstration, nous écrirons $f_*$ pour $Rf_*$ et $f^*$ pour $Lf^*$, etc., ainsi que $\otimes$ pour $\otimes^\mathbf{L}_{\mathcal{O}_X}$, etc. Écrivons (\ref{equation-compare-with-pullback}) comme la composée
\begin{align*}
f^*K \otimes a(\mathcal{O}_Y)
& \to
a(f_*(f^*K \otimes a(\mathcal{O}_Y))) \\
& \leftarrow
a(K \otimes f_*a(\mathcal{O}_K)) \\
& \to
a(K \otimes \mathcal{O}_Y) \\
& \to
a(K)
\end{align*}
La première flèche est l'unité $\eta_f$ ; la deuxième est $a$ appliqué à Cohomologie, équation (\ref{cohomology-equation-projection-formula-map}), qui est un isomorphisme d'après Catégories dérivées des schémas, lemme \ref{perfect-lemma-cohomology-base-change} ; la troisième est $a$ appliqué à $\text{id}_K \otimes \text{Tr}_f$ ; la quatrième est $a$ appliqué à l'isomorphisme $K \otimes \mathcal{O}_Y = K$. La démonstration du lemme consiste à montrer que chacune de ces applications donne un carré commutatif comme dans l'énoncé. Pour $\eta_f$ et $\text{Tr}_f$, cela résulte des lemmes \ref{lemma-unit-and-base-change} et \ref{lemma-trace-map-and-base-change}. Pour la flèche faisant intervenir Cohomologie, équation (\ref{cohomology-equation-projection-formula-map}), cela résulte de Cohomologie, remarque \ref{cohomology-remark-compatible-with-diagram}. Pour l'application de multiplication, c'est clair. Cela achève la démonstration.
\end{proof}

\begin{lemma}
\label{lemma-compare-on-open}
Soit $f : X \to Y$ un morphisme propre de schémas noethériens. Soit $V \subset Y$ un ouvert tel que $f^{-1}(V) \to V$ soit un isomorphisme. Alors, pour $K \in D_\QCoh^+(\mathcal{O}_Y)$, l'application (\ref{equation-compare-with-pullback}) se restreint à un isomorphisme au-dessus de $f^{-1}(V)$.
\end{lemma}

\begin{proof}
D'après le lemme \ref{lemma-proper-noetherian}, l'application (\ref{equation-sheafy}) est un isomorphisme pour les objets de $D_\QCoh^+(\mathcal{O}_Y)$. Le lemme \ref{lemma-restriction-compare-with-pullback} montre donc que la restriction de (\ref{equation-compare-with-pullback}), pour $K$, à $f^{-1}(V)$ est l'application (\ref{equation-compare-with-pullback}) pour $K|_V$ et $f^{-1}(V) \to V$. Il suffit donc de montrer que l'application est un isomorphisme lorsque $f$ est le morphisme identité, ce qui est clair.
\end{proof}

\begin{lemma}
\label{lemma-transitivity-compare-with-pullback}
Soient $f : X \to Y$ et $g : Y \to Z$ deux morphismes composables de schémas quasi-compacts et quasi-séparés, et posons $h = g \circ f$. Soient $a, b, c$ les adjoints du lemme \ref{lemma-twisted-inverse-image} pour $f, g, h$. Pour tout $K \in D_\QCoh(\mathcal{O}_Z)$, le diagramme
$$
\xymatrix{
Lf^*(Lg^*K \otimes_{\mathcal{O}_Y}^\mathbf{L}
b(\mathcal{O}_Z)) \otimes_{\mathcal{O}_X}^\mathbf{L} a(\mathcal{O}_Y)
\ar@{=}[d] \ar[r] &
a(Lg^*K \otimes_{\mathcal{O}_Y}^\mathbf{L} b(\mathcal{O}_Z)) \ar[r] &
a(b(K)) \ar@{=}[d] \\
Lh^*K \otimes_{\mathcal{O}_X}^\mathbf{L} Lf^*b(\mathcal{O}_Z)
\otimes_{\mathcal{O}_X}^\mathbf{L} a(\mathcal{O}_Y) \ar[r] &
Lh^*K \otimes_{\mathcal{O}_X}^\mathbf{L} c(\mathcal{O}_Z) \ar[r] &
c(K)
}
$$
est commutatif, où les flèches sont celles de (\ref{equation-compare-with-pullback}) et où l'on a utilisé $Lh^* = Lf^* \circ Lg^*$ et $c = a \circ b$.
\end{lemma}

\begin{proof}
Dans cette démonstration, nous écrirons $f_*$ pour $Rf_*$ et $f^*$ pour $Lf^*$, etc., ainsi que $\otimes$ pour $\otimes^\mathbf{L}_{\mathcal{O}_X}$, etc. La composée des flèches supérieures est adjointe à une application
$$
g_*f_*(f^*(g^*K \otimes b(\mathcal{O}_Z)) \otimes a(\mathcal{O}_Y)) \to K
$$
Le membre de gauche est égal à $K \otimes g_*f_*(f^*b(\mathcal{O}_Z) \otimes a(\mathcal{O}_Y))$ d'après Catégories dérivées des schémas, lemme \ref{perfect-lemma-cohomology-base-change}, et l'examen des définitions montre que l'application provient de l'application
$$
g_*f_*(f^*b(\mathcal{O}_Z) \otimes a(\mathcal{O}_Y))
\xleftarrow{g_*\epsilon}
g_*(b(\mathcal{O}_Z) \otimes f_*a(\mathcal{O}_Y)) \xrightarrow{g_*\alpha}
g_*(b(\mathcal{O}_Z)) \xrightarrow{\beta} \mathcal{O}_Z
$$
tensorisée par $\text{id}_K$. Ici, $\epsilon$ est l'isomorphisme de Catégories dérivées des schémas, lemme \ref{perfect-lemma-cohomology-base-change}, et $\beta$ provient de la co-unité $g_*b \to \text{id}$. De même, la composée des flèches horizontales inférieures est adjointe à $\text{id}_K$ tensorisé par la composée
$$
g_*f_*(f^*b(\mathcal{O}_Z) \otimes a(\mathcal{O}_Y)) \xrightarrow{g_*f_*\delta}
g_*f_*(ab(\mathcal{O}_Z)) \xrightarrow{g_*\gamma}
g_*(b(\mathcal{O}_Z)) \xrightarrow{\beta}
\mathcal{O}_Z
$$
où $\gamma$ provient de la co-unité $f_*a \to \text{id}$, et $\delta$ est l'application dont l'adjointe est la composée
$$
f_*(f^*b(\mathcal{O}_Z) \otimes a(\mathcal{O}_Y))
\xleftarrow{\epsilon}
b(\mathcal{O}_Z) \otimes f_*a(\mathcal{O}_Y) \xrightarrow{\alpha}
b(\mathcal{O}_Z)
$$
Les propriétés générales des foncteurs adjoints, des applications adjointes et des co-unités (voir Catégories, section \ref{categories-section-adjoint}) donnent $\gamma \circ f_*\delta = \alpha \circ \epsilon^{-1}$, comme voulu.
\end{proof}





\section{Adjoint à droite de l'image directe pour les immersions fermées}
\sectionmark{Adjoint à droite et immersions fermées}
\label{section-sections-with-exact-support}

\noindent
Soit $i : (Z, \mathcal{O}_Z) \to (X, \mathcal{O}_X)$ un morphisme d'espaces annelés tel que $i$ soit un homéomorphisme sur une partie fermée et que $i^\sharp : \mathcal{O}_X \to i_*\mathcal{O}_Z$ soit surjective (par exemple, une immersion fermée de schémas). Soit $\mathcal{I} = \Ker(i^\sharp)$. Pour un faisceau de $\mathcal{O}_X$-modules $\mathcal{F}$, le faisceau
$$
\SheafHom_{\mathcal{O}_X}(i_*\mathcal{O}_Z, \mathcal{F})
$$
est un faisceau de $\mathcal{O}_X$-modules annulé par $\mathcal{I}$. D'après Modules, lemme \ref{modules-lemma-i-star-equivalence}, il existe donc un faisceau de $\mathcal{O}_Z$-modules, que nous noterons $\SheafHom(\mathcal{O}_Z, \mathcal{F})$, tel que
$$
i_*\SheafHom(\mathcal{O}_Z, \mathcal{F}) =
\SheafHom_{\mathcal{O}_X}(i_*\mathcal{O}_Z, \mathcal{F})
$$
comme $\mathcal{O}_X$-modules. Explicitons cet énoncé.

\begin{lemma}
\label{lemma-compute-sheaf-with-exact-support}
Avec les notations précédentes, le foncteur $\SheafHom(\mathcal{O}_Z, -)$ est un adjoint à droite du foncteur $i_* : \textit{Mod}(\mathcal{O}_Z) \to \textit{Mod}(\mathcal{O}_X)$. Pour $V \subset Z$ ouvert, on a
$$
\Gamma(V, \SheafHom(\mathcal{O}_Z, \mathcal{F})) =
\{s \in \Gamma(U, \mathcal{F}) \mid \mathcal{I}s = 0\}
$$
où $U \subset X$ est un ouvert dont l'intersection avec $Z$ est $V$.
\end{lemma}

\begin{proof}
Soit $\mathcal{G}$ un faisceau de $\mathcal{O}_Z$-modules. Alors
$$
\Hom_{\mathcal{O}_X}(i_*\mathcal{G}, \mathcal{F}) =
\Hom_{i_*\mathcal{O}_Z}(i_*\mathcal{G},
\SheafHom_{\mathcal{O}_X}(i_*\mathcal{O}_Z, \mathcal{F})) =
\Hom_{\mathcal{O}_Z}(\mathcal{G}, \SheafHom(\mathcal{O}_Z, \mathcal{F}))
$$
La première égalité résulte de Modules, lemme \ref{modules-lemma-adjoint-tensor-restrict}, et la seconde de la pleine fidélité de $i_*$, voir Modules, lemme \ref{modules-lemma-i-star-equivalence}. La description des sections est laissée au lecteur.
\end{proof}

\noindent
Le foncteur
$$
\textit{Mod}(\mathcal{O}_X)
\longrightarrow
\textit{Mod}(\mathcal{O}_Z),
\quad
\mathcal{F} \longmapsto \SheafHom(\mathcal{O}_Z, \mathcal{F})
$$
est exact à gauche et admet un prolongement dérivé
$$
R\SheafHom(\mathcal{O}_Z, -) : D(\mathcal{O}_X) \to D(\mathcal{O}_Z).
$$

\begin{lemma}
\label{lemma-sheaf-with-exact-support-adjoint}
Avec les notations précédentes, le foncteur $R\SheafHom(\mathcal{O}_Z, -)$ est l'adjoint à droite du foncteur $Ri_* : D(\mathcal{O}_Z) \to D(\mathcal{O}_X)$.
\end{lemma}

\begin{proof}
Cela résulte du fait que $i_*$ et $\SheafHom(\mathcal{O}_Z, -)$ sont des foncteurs adjoints d'après le lemme \ref{lemma-compute-sheaf-with-exact-support}. Voir Catégories dérivées, lemme \ref{derived-lemma-derived-adjoint-functors}.
\end{proof}

\begin{lemma}
\label{lemma-sheaf-with-exact-support-ext}
Avec les notations précédentes, on a
$$
Ri_*R\SheafHom(\mathcal{O}_Z, K) =
R\SheafHom_{\mathcal{O}_X}(i_*\mathcal{O}_Z, K)
$$
dans $D(\mathcal{O}_X)$ pour tout $K$ de $D(\mathcal{O}_X)$.
\end{lemma}

\begin{proof}
Cela résulte immédiatement de la construction du foncteur
\par\noindent\hfil
$R\SheafHom(\mathcal{O}_Z, -)$.
\hfil\par
\end{proof}

\begin{lemma}
\label{lemma-sheaf-with-exact-support-internal-home}
Avec les notations précédentes, pour $M \in D(\mathcal{O}_Z)$, on a
$$
R\SheafHom_{\mathcal{O}_X}(Ri_*M, K) =
Ri_*R\SheafHom_{\mathcal{O}_Z}(M, R\SheafHom(\mathcal{O}_Z, K))
$$
dans $D(\mathcal{O}_Z)$ pour tout $K$ de $D(\mathcal{O}_X)$.
\end{lemma}

\begin{proof}
Cela résulte immédiatement de la construction du foncteur
\par\noindent\hfil
$R\SheafHom(\mathcal{O}_Z, -)$
\hfil\par\noindent
et du fait que, si $\mathcal{K}^\bullet$ est un complexe K-injectif de $\mathcal{O}_X$-modules, alors $\SheafHom(\mathcal{O}_Z, \mathcal{K}^\bullet)$ est un complexe K-injectif de $\mathcal{O}_Z$-modules, voir Catégories dérivées, lemme \ref{derived-lemma-adjoint-preserve-K-injectives}.
\end{proof}

\begin{lemma}
\label{lemma-sheaf-with-exact-support-quasi-coherent}
Soit $i : Z \to X$ une immersion fermée pseudo-cohérente de schémas (toute immersion fermée convient si $X$ est localement noethérien). Alors
\begin{enumerate}
\item $R\SheafHom(\mathcal{O}_Z, -)$ envoie $D^+_\QCoh(\mathcal{O}_X)$ dans $D^+_\QCoh(\mathcal{O}_Z)$ ;
\item si $X = \Spec(A)$ et $Z = \Spec(B)$, le diagramme
$$
\xymatrix{
D^+(B) \ar[r] & D_\QCoh^+(\mathcal{O}_Z) \\
D^+(A) \ar[r] \ar[u]^{R\Hom(B, -)} &
D_\QCoh^+(\mathcal{O}_X) \ar[u]_{R\SheafHom(\mathcal{O}_Z, -)}
}
$$
est commutatif.
\end{enumerate}
\end{lemma}

\begin{proof}
Pour justifier la remarque entre parenthèses, si $X$ est localement noethérien, alors $i$ est pseudo-cohérent d'après Compléments sur les morphismes, lemme \ref{more-morphisms-lemma-Noetherian-pseudo-coherent}.

\medskip\noindent
Soit $K$ un objet de $D^+_\QCoh(\mathcal{O}_X)$. Pour démontrer (1), il suffit, d'après Morphismes, lemme \ref{morphisms-lemma-i-star-equivalence}, de montrer que $i_*$ appliqué à $H^n(R\SheafHom(\mathcal{O}_Z, K))$ donne un module quasi-cohérent sur $X$. D'après le lemme \ref{lemma-sheaf-with-exact-support-ext}, cela revient à montrer que
\par\noindent\hfil
$R\SheafHom_{\mathcal{O}_X}(i_*\mathcal{O}_Z, K)$
\hfil\par\noindent
appartient à $D_\QCoh(\mathcal{O}_X)$. Puisque $i$ est pseudo-cohérent, le faisceau $\mathcal{O}_Z$ est un $\mathcal{O}_X$-module pseudo-cohérent. Le résultat découle donc de Catégories dérivées des schémas, lemme \ref{perfect-lemma-quasi-coherence-internal-hom}.

\medskip\noindent
Supposons $X = \Spec(A)$ et $Z = \Spec(B)$ comme en (2). Soit $I^\bullet$ un complexe borné inférieurement de $A$-modules injectifs représentant un objet $K$ de $D^+(A)$. On sait alors que $R\Hom(B, K) = \Hom_A(B, I^\bullet)$, considéré comme complexe de $B$-modules. Choisissons un quasi-isomorphisme
$$
\widetilde{I^\bullet} \longrightarrow \mathcal{I}^\bullet
$$
où $\mathcal{I}^\bullet$ est un complexe borné inférieurement de $\mathcal{O}_X$-modules injectifs. La description du foncteur $\SheafHom(\mathcal{O}_Z, -)$ au lemme \ref{lemma-compute-sheaf-with-exact-support} donne une application
$$
\Hom_A(B, I^\bullet)
\longrightarrow
\Gamma(Z, \SheafHom(\mathcal{O}_Z, \mathcal{I}^\bullet))
$$
Observons que $\SheafHom(\mathcal{O}_Z, \mathcal{I}^\bullet)$ représente $R\SheafHom(\mathcal{O}_Z, \widetilde{K})$. En appliquant la propriété universelle du foncteur $\widetilde{\ }$, on obtient une application
$$
\widetilde{\Hom_A(B, I^\bullet)}
\longrightarrow
R\SheafHom(\mathcal{O}_Z, \widetilde{K})
$$
dans $D(\mathcal{O}_Z)$. On peut vérifier que cette application est un isomorphisme dans $D(\mathcal{O}_Z)$ après application de $i_*$. Or, après application de $i_*$, on obtient l'isomorphisme de Catégories dérivées des schémas, lemme \ref{perfect-lemma-quasi-coherence-internal-hom}, via l'identification du lemme \ref{lemma-sheaf-with-exact-support-ext}.
\end{proof}

\begin{lemma}
\label{lemma-sheaf-with-exact-support-coherent}
Soit $i : Z \to X$ une immersion fermée de schémas. Supposons $X$ localement noethérien. Alors $R\SheafHom(\mathcal{O}_Z, -)$ envoie $D^+_{\textit{Coh}}(\mathcal{O}_X)$ dans $D^+_{\textit{Coh}}(\mathcal{O}_Z)$.
\end{lemma}

\begin{proof}
La question est locale sur $X$ ; on peut donc supposer $X$ affine. Écrivons $X = \Spec(A)$ et $Z = \Spec(B)$, avec $A$ noethérien et $A \to B$ surjectif. On peut alors appliquer le lemme \ref{lemma-sheaf-with-exact-support-quasi-coherent} pour traduire la question en algèbre. Le résultat algébrique correspondant est une conséquence de Complexes dualisants, lemme \ref{dualizing-lemma-exact-support-coherent}.
\end{proof}

\begin{lemma}
\label{lemma-twisted-inverse-image-closed}
Soit $X$ un schéma quasi-compact et quasi-séparé. Soit $i : Z \to X$ une immersion fermée pseudo-cohérente (si $X$ est noethérien, toute immersion fermée est pseudo-cohérente). Soit $a : D_\QCoh(\mathcal{O}_X) \to D_\QCoh(\mathcal{O}_Z)$ l'adjoint à droite de $Ri_*$. Il existe alors un isomorphisme fonctoriel
$$
a(K) = R\SheafHom(\mathcal{O}_Z, K)
$$
pour $K \in D_\QCoh^+(\mathcal{O}_X)$.
\end{lemma}

\begin{proof}
(L'énoncé entre parenthèses résulte de Compléments sur les morphismes, lemme \ref{more-morphisms-lemma-Noetherian-pseudo-coherent}.) D'après le lemme \ref{lemma-sheaf-with-exact-support-adjoint}, le foncteur $R\SheafHom(\mathcal{O}_Z, -)$ est un adjoint à droite de $Ri_* : D(\mathcal{O}_Z) \to D(\mathcal{O}_X)$. De plus, d'après les lemmes \ref{lemma-sheaf-with-exact-support-quasi-coherent} et \ref{lemma-twisted-inverse-image-bounded-below}, les foncteurs $R\SheafHom(\mathcal{O}_Z, -)$ et $a$ envoient tous deux $D_\QCoh^+(\mathcal{O}_X)$ dans $D_\QCoh^+(\mathcal{O}_Z)$. On obtient donc l'isomorphisme par unicité des foncteurs adjoints.
\end{proof}

\begin{example}
\label{example-trace-closed-immersion}
Si $i : Z \to X$ est une immersion fermée de schémas noethériens, le diagramme
$$
\xymatrix{
i_*a(K) \ar[rr]_-{\text{Tr}_{i, K}} \ar@{=}[d] & &
K \ar@{=}[d] \\
i_*R\SheafHom(\mathcal{O}_Z, K) \ar@{=}[r] &
R\SheafHom_{\mathcal{O}_X}(i_*\mathcal{O}_Z, K)
\ar[r] & K
}
$$
est commutatif pour $K \in D_\QCoh^+(\mathcal{O}_X)$. Le signe d'égalité horizontal correspond ici au lemme \ref{lemma-sheaf-with-exact-support-ext}, et la flèche horizontale inférieure est induite par l'application $\mathcal{O}_X \to i_*\mathcal{O}_Z$. La commutativité du diagramme est une conséquence du lemme \ref{lemma-twisted-inverse-image-closed}.
\end{example}




\section{Adjoint à droite de l'image directe pour les immersions fermées et changement de base}
\label{section-sections-with-exact-support-base-change}

\noindent
Considérons un diagramme cartésien de schémas
$$
\xymatrix{
Z' \ar[r]_{i'} \ar[d]_g & X' \ar[d]^f \\
Z \ar[r]^i & X
}
$$
où $i$ est une immersion fermée. Si $Z$ et $X'$ sont Tor-indépendants sur $X$, il existe une application canonique de changement de base
\begin{equation}
\label{equation-base-change-exact-support}
Lg^*R\SheafHom(\mathcal{O}_Z, K)
\longrightarrow
R\SheafHom(\mathcal{O}_{Z'}, Lf^*K)
\end{equation}
dans $D(\mathcal{O}_{Z'})$, fonctorielle en $K$ dans $D(\mathcal{O}_X)$. En effet, par l'adjonction du lemme \ref{lemma-sheaf-with-exact-support-adjoint}, une telle flèche revient à une application
$$
Ri'_*Lg^*R\SheafHom(\mathcal{O}_Z, K)
\longrightarrow
Lf^*K
$$
dans $D(\mathcal{O}_{X'})$. La Tor-indépendance donne $Ri'_* \circ Lg^* = Lf^* \circ Ri_*$ (voir Catégories dérivées des schémas, lemme \ref{perfect-lemma-compare-base-change-closed-immersion}). Cela revient donc à une application
$$
Lf^*Ri_*R\SheafHom(\mathcal{O}_Z, K)
\longrightarrow
Lf^*K
$$
Pour celle-ci, on peut prendre $Lf^*(can)$, où $can : Ri_* R\SheafHom(\mathcal{O}_Z, K) \to K$ est la co-unité de l'adjonction.

\begin{lemma}
\label{lemma-check-base-change-is-iso}
Dans la situation précédente, l'application (\ref{equation-base-change-exact-support}) est un isomorphisme si et seulement si l'application de changement de base
$$
Lf^*R\SheafHom_{\mathcal{O}_X}(\mathcal{O}_Z, K)
\longrightarrow
R\SheafHom_{\mathcal{O}_{X'}}(\mathcal{O}_{Z'}, Lf^*K)
$$
de Cohomologie, remarque \ref{cohomology-remark-prepare-fancy-base-change}, en est un.
\end{lemma}

\begin{proof}
L'énoncé a un sens, car $\mathcal{O}_{Z'} = Lf^*\mathcal{O}_Z$ par la Tor-indépendance supposée. Puisque $i'_*$ est exact et fidèle, il suffit de montrer que l'application (\ref{equation-base-change-exact-support}) est un isomorphisme après application de $Ri'_*$. Puisque $Ri'_* \circ Lg^* = Lf^* \circ Ri_*$ par la Tor-indépendance supposée et Catégories dérivées des schémas, lemme \ref{perfect-lemma-compare-base-change-closed-immersion}, on obtient une application
$$
Lf^*Ri_*R\SheafHom(\mathcal{O}_Z, K)
\longrightarrow
Ri'_*R\SheafHom(\mathcal{O}_{Z'}, Lf^*K)
$$
dont la source et le but sont ceux de l'énoncé, d'après le lemme \ref{lemma-sheaf-with-exact-support-ext}. Nous omettons de vérifier qu'il s'agit de la même application que celle construite dans Cohomologie, remarque \ref{cohomology-remark-prepare-fancy-base-change}.
\end{proof}

\begin{lemma}
\label{lemma-flat-bc-sheaf-with-exact-support}
Dans la situation précédente, supposons $f$ plat et $i$ pseudo-cohérente. Alors (\ref{equation-base-change-exact-support}) est un isomorphisme pour $K$ dans $D^+_\QCoh(\mathcal{O}_X)$.
\end{lemma}

\begin{proof}
Première démonstration. Pour montrer que cette application est un isomorphisme, on peut travailler localement. On peut donc supposer $X$, $X'$, $Z$, $Z'$ affines, correspondant respectivement aux anneaux $A$, $A'$, $B$, $B'$. Alors $B$ et $A'$ sont Tor-indépendants sur $A$. D'après le lemme \ref{lemma-check-base-change-is-iso}, il suffit de vérifier que
$$
R\Hom_A(B, K) \otimes_A^\mathbf{L} A' =
R\Hom_{A'}(B', K \otimes_A^\mathbf{L} A')
$$
dans $D(A')$ pour tout $K \in D^+(A)$. On utilise ici Catégories dérivées des schémas, lemme \ref{perfect-lemma-quasi-coherence-internal-hom}, et le fait que $B$, resp.\ $B'$, soit pseudo-cohérent comme $A$-module, resp.\ $A'$-module, pour comparer les Hom dérivés au niveau des anneaux et des schémas. L'égalité affichée résulte de Compléments d'algèbre, lemme \ref{more-algebra-lemma-internal-hom-evaluate-tensor-isomorphism}, (3). Voir aussi la discussion dans Complexes dualisants, section \ref{dualizing-section-base-change-trivial-duality}.

\medskip\noindent
Deuxième démonstration\footnote{Cette démonstration montre qu'il suffit de supposer que $K$ appartient à $D^+(\mathcal{O}_X)$.}. Soit $z' \in Z'$ d'image $z \in Z$. Montrons d'abord que (\ref{equation-base-change-exact-support}), sur les fibres en $z'$, induit l'application
$$
R\Hom(\mathcal{O}_{Z, z}, K_z)
\otimes_{\mathcal{O}_{Z, x}}^\mathbf{L} \mathcal{O}_{Z', z'}
\longrightarrow
R\Hom(\mathcal{O}_{Z', z'},
K_z \otimes_{\mathcal{O}_{X, z}}^\mathbf{L} \mathcal{O}_{X', z'})
$$
de Complexes dualisants, équation (\ref{dualizing-equation-base-change}). En effet, les constructions de ces applications sont identiques. On applique ensuite Complexes dualisants, lemme \ref{dualizing-lemma-flat-bc-surjection}.
\end{proof}

\begin{lemma}
\label{lemma-sheaf-with-exact-support-tensor}
Soit $i : Z \to X$ une immersion fermée pseudo-cohérente de schémas. Soit $M \in D_\QCoh(\mathcal{O}_X)$ localement de Tor-amplitude dans $[a, \infty)$. Soit $K \in D_\QCoh^+(\mathcal{O}_X)$. Il existe alors un isomorphisme canonique
$$
R\SheafHom(\mathcal{O}_Z, K) \otimes_{\mathcal{O}_Z}^\mathbf{L} Li^*M =
R\SheafHom(\mathcal{O}_Z, K \otimes_{\mathcal{O}_X}^\mathbf{L} M)
$$
dans $D(\mathcal{O}_Z)$.
\end{lemma}

\begin{proof}
Une application du membre de gauche vers celui de droite revient à une application
$$
Ri_*R\SheafHom(\mathcal{O}_Z, K) \otimes_{\mathcal{O}_X}^\mathbf{L} M
\longrightarrow
K \otimes_{\mathcal{O}_X}^\mathbf{L} M
$$
d'après les lemmes \ref{lemma-sheaf-with-exact-support-adjoint} et \ref{lemma-sheaf-with-exact-support-ext}. Pour cette application, on prend la co-unité
\par\noindent\hfil
$Ri_*R\SheafHom(\mathcal{O}_Z, K) \to K$
\hfil\par\noindent
tensorisée par $\text{id}_M$. Pour voir qu'elle est un isomorphisme sous les hypothèses données, on traduit l'énoncé en algèbre à l'aide du lemme \ref{lemma-sheaf-with-exact-support-quasi-coherent}, puis on utilise, par exemple, Compléments d'algèbre, lemme \ref{more-algebra-lemma-internal-hom-evaluate-tensor-isomorphism}, (3). Au lieu d'utiliser le lemme \ref{lemma-sheaf-with-exact-support-quasi-coherent}, on peut considérer les fibres, comme dans la deuxième démonstration du lemme \ref{lemma-flat-bc-sheaf-with-exact-support}.
\end{proof}



\section{Adjoint à droite de l'image directe pour les morphismes finis}
\sectionmark{Adjoint à droite et morphismes finis}
\label{section-duality-finite}

\noindent
Si $i : Z \to X$ est une immersion fermée de schémas, il existe un adjoint à droite $\SheafHom(\mathcal{O}_Z, -)$ du foncteur $i_* : \textit{Mod}(\mathcal{O}_Z) \to \textit{Mod}(\mathcal{O}_X)$ dont le prolongement dérivé $R\SheafHom(\mathcal{O}_Z, -)$ est l'adjoint à droite de $Ri_* : D(\mathcal{O}_Z) \to D(\mathcal{O}_X)$. Voir la section \ref{section-sections-with-exact-support}. Dans le cas d'un morphisme fini $f : Y \to X$, cette stratégie ne peut pas fonctionner, car le foncteur $f_* : \textit{Mod}(\mathcal{O}_Y) \to \textit{Mod}(\mathcal{O}_X)$ n'est pas exact en général et n'admet donc pas d'adjoint à droite. On peut le remplacer par le foncteur exact
$\textit{Mod}(f_*\mathcal{O}_Y) \to \textit{Mod}(\mathcal{O}_X)$
et considérer l'adjoint à droite correspondant ainsi que son prolongement dérivé.

\medskip\noindent
Soit $f : Y \to X$ un morphisme affine de schémas. Pour un faisceau de $\mathcal{O}_X$-modules $\mathcal{F}$, le faisceau
$$
\SheafHom_{\mathcal{O}_X}(f_*\mathcal{O}_Y, \mathcal{F})
$$
est un faisceau de $f_*\mathcal{O}_Y$-modules. On obtient un foncteur
$\textit{Mod}(\mathcal{O}_X) \to \textit{Mod}(f_*\mathcal{O}_Y)$
que nous noterons $\SheafHom(f_*\mathcal{O}_Y, -)$.

\begin{lemma}
\label{lemma-compute-sheafhom-affine}
Avec les notations précédentes, le foncteur $\SheafHom(f_*\mathcal{O}_Y, -)$ est un adjoint à droite du foncteur de restriction $\textit{Mod}(f_*\mathcal{O}_Y) \to \textit{Mod}(\mathcal{O}_X)$. Pour un ouvert affine $U \subset X$, on a
$$
\Gamma(U, \SheafHom(f_*\mathcal{O}_Y, \mathcal{F})) =
\Hom_A(B, \mathcal{F}(U))
$$
où $A = \mathcal{O}_X(U)$ et $B = \mathcal{O}_Y(f^{-1}(U))$.
\end{lemma}

\begin{proof}
L'adjonction résulte de Modules, lemme \ref{modules-lemma-adjoint-tensor-restrict}. Puisque $f$ est affine, $f_*\mathcal{O}_Y$ est le faisceau quasi-cohérent correspondant à $B$ considéré comme $A$-module. La description des sections sur $U$ découle donc de Schémas, lemme \ref{schemes-lemma-compare-constructions}.
\end{proof}

\noindent
Le foncteur $\SheafHom(f_*\mathcal{O}_Y, -)$ est exact à gauche. Soit
$$
R\SheafHom(f_*\mathcal{O}_Y, -) :
D(\mathcal{O}_X)
\longrightarrow
D(f_*\mathcal{O}_Y)
$$
son prolongement dérivé.

\begin{lemma}
\label{lemma-sheafhom-affine-adjoint}
Avec les notations précédentes, le foncteur $R\SheafHom(f_*\mathcal{O}_Y, -)$ est l'adjoint à droite du foncteur $D(f_*\mathcal{O}_Y) \to D(\mathcal{O}_X)$.
\end{lemma}

\begin{proof}
Cela résulte du lemme \ref{lemma-compute-sheafhom-affine} et de Catégories dérivées, lemme \ref{derived-lemma-derived-adjoint-functors}.
\end{proof}

\begin{lemma}
\label{lemma-sheafhom-affine-ext}
Avec les notations précédentes, la composée
$$
D(\mathcal{O}_X) \xrightarrow{R\SheafHom(f_*\mathcal{O}_Y, -)}
D(f_*\mathcal{O}_Y) \to D(\mathcal{O}_X)
$$
est le foncteur $K \mapsto R\SheafHom_{\mathcal{O}_X}(f_*\mathcal{O}_Y, K)$.
\end{lemma}

\begin{proof}
Cela résulte immédiatement de la construction.
\end{proof}

\begin{lemma}
\label{lemma-finite-twisted}
Soit $f : Y \to X$ un morphisme fini pseudo-cohérent de schémas (un morphisme fini de schémas noethériens est pseudo-cohérent). Le foncteur $R\SheafHom(f_*\mathcal{O}_Y, -)$ envoie $D_\QCoh^+(\mathcal{O}_X)$ dans $D_\QCoh^+(f_*\mathcal{O}_Y)$. Si $X$ est quasi-compact et quasi-séparé, le diagramme
$$
\xymatrix{
D_\QCoh^+(\mathcal{O}_X) \ar[rr]_a \ar[rd]_{R\SheafHom(f_*\mathcal{O}_Y, -)}
& & D_\QCoh^+(\mathcal{O}_Y) \ar[ld]^\Phi \\
& D_\QCoh^+(f_*\mathcal{O}_Y)
}
$$
est commutatif, où $a$ est l'adjoint à droite du lemme \ref{lemma-twisted-inverse-image} pour $f$ et où $\Phi$ est l'équivalence de Catégories dérivées des schémas, lemme \ref{perfect-lemma-affine-morphism-equivalence}.
\end{lemma}

\begin{proof}
\begingroup\emergencystretch=6em
(La remarque entre parenthèses résulte de Compléments sur les morphismes, lemme \ref{more-morphisms-lemma-Noetherian-pseudo-coherent}.) Puisque $f$ est pseudo-cohérent, le $\mathcal{O}_X$-module $f_*\mathcal{O}_Y$ est pseudo-cohérent, voir Compléments sur les morphismes, lemme \ref{more-morphisms-lemma-finite-pseudo-coherent}. Ainsi, $R\SheafHom(f_*\mathcal{O}_Y, -)$ envoie $D_\QCoh^+(\mathcal{O}_X)$ dans $D_\QCoh^+(f_*\mathcal{O}_Y)$, voir Catégories dérivées des schémas, lemme \ref{perfect-lemma-quasi-coherence-internal-hom}. Alors $\Phi \circ a$ et $R\SheafHom(f_*\mathcal{O}_Y, -)$ coïncident sur $D_\QCoh^+(\mathcal{O}_X)$, car ces deux foncteurs sont adjoints à droite du foncteur de restriction $D_\QCoh^+(f_*\mathcal{O}_Y) \to D_\QCoh^+(\mathcal{O}_X)$. Pour le voir, on utilise les lemmes \ref{lemma-twisted-inverse-image-bounded-below} et \ref{lemma-sheafhom-affine-adjoint}.\par\endgroup
\end{proof}

\begin{remark}
\label{remark-trace-map-finite}
Si $f : Y \to X$ est un morphisme fini de schémas noethériens, le diagramme
$$
\xymatrix{
Rf_*a(K) \ar[r]_-{\text{Tr}_{f, K}} \ar@{=}[d] & K \ar@{=}[d] \\
R\SheafHom_{\mathcal{O}_X}(f_*\mathcal{O}_Y, K) \ar[r] & K
}
$$
est commutatif pour $K \in D_\QCoh^+(\mathcal{O}_X)$. Cela résulte du lemme \ref{lemma-finite-twisted}. La flèche horizontale inférieure est induite par l'application $\mathcal{O}_X \to f_*\mathcal{O}_Y$, et la flèche horizontale supérieure est le morphisme trace de la section \ref{section-trace}.
\end{remark}






\section{Adjoint à droite de l'image directe pour les morphismes propres et plats}
\label{section-proper-flat}

\noindent
Pour les morphismes propres, plats et de présentation finie entre schémas quasi-compacts et quasi-séparés, l'adjoint à droite de l'image directe possède des propriétés remarquables.

\begin{lemma}
\label{lemma-proper-flat}
Soit $Y$ un schéma quasi-compact et quasi-séparé. Soit $f : X \to Y$ un morphisme de schémas propre, plat et de présentation finie. Soit $a$ l'adjoint à droite de $Rf_* : D_\QCoh(\mathcal{O}_X) \to D_\QCoh(\mathcal{O}_Y)$ du lemme \ref{lemma-twisted-inverse-image}. Alors $a$ commute aux sommes directes.
\end{lemma}

\begin{proof}
Soit $P$ un objet parfait de $D(\mathcal{O}_X)$. D'après Catégories dérivées des schémas, lemme \ref{perfect-lemma-flat-proper-perfect-direct-image-general}, le complexe $Rf_*P$ est parfait sur $Y$. Soit $K_i$ une famille d'objets de $D_\QCoh(\mathcal{O}_Y)$. Alors
\begin{align*}
\Hom_{D(\mathcal{O}_X)}(P, a(\bigoplus K_i))
& =
\Hom_{D(\mathcal{O}_Y)}(Rf_*P, \bigoplus K_i) \\
& =
\bigoplus \Hom_{D(\mathcal{O}_Y)}(Rf_*P, K_i) \\
& =
\bigoplus \Hom_{D(\mathcal{O}_X)}(P, a(K_i))
\end{align*}
car un objet parfait est compact (Catégories dérivées des schémas, proposition \ref{perfect-proposition-compact-is-perfect}). Puisque $D_\QCoh(\mathcal{O}_X)$ admet un générateur parfait (Catégories dérivées des schémas, théorème \ref{perfect-theorem-bondal-van-den-Bergh}), on en conclut que l'application $\bigoplus a(K_i) \to a(\bigoplus K_i)$ est un isomorphisme, c'est-à-dire que $a$ commute aux sommes directes.
\end{proof}

\begin{lemma}
\label{lemma-proper-flat-relative}
Soit $Y$ un schéma quasi-compact et quasi-séparé. Soit $f : X \to Y$ un morphisme de schémas propre, plat et de présentation finie. Soit $a$ l'adjoint à droite de $Rf_* : D_\QCoh(\mathcal{O}_X) \to D_\QCoh(\mathcal{O}_Y)$ du lemme \ref{lemma-twisted-inverse-image}. Alors
\begin{enumerate}
\item pour tout fermé $T \subset Y$, si $Q \in D_\QCoh(Y)$ est à support dans $T$, alors $a(Q)$ est à support dans $f^{-1}(T)$ ;
\item pour tout ouvert quasi-compact $V \subset Y$ et tout $K \in D_\QCoh(\mathcal{O}_Y)$, l'application (\ref{equation-sheafy}) est un isomorphisme.
\end{enumerate}
\end{lemma}

\begin{proof}
Cela résulte des lemmes \ref{lemma-when-sheafy}, \ref{lemma-proper-noetherian} et \ref{lemma-proper-flat}. À strictement parler, on n'obtient (1) que pour les $T$ dont le complémentaire est un ouvert quasi-compact ; le cas général s'en déduit néanmoins, car tout fermé est une intersection de tels fermés.
\end{proof}

\begin{lemma}
\label{lemma-compare-with-pullback-flat-proper}
Soit $Y$ un schéma quasi-compact et quasi-séparé. Soit $f : X \to Y$ un morphisme de schémas propre, plat et de présentation finie. L'application (\ref{equation-compare-with-pullback}) est un isomorphisme pour tout objet $K$ de $D_\QCoh(\mathcal{O}_Y)$.
\end{lemma}

\begin{proof}
D'après le lemme \ref{lemma-proper-flat}, $a$ commute aux sommes directes. La collection des objets de $D_\QCoh(\mathcal{O}_Y)$ pour lesquels (\ref{equation-compare-with-pullback}) est un isomorphisme est donc une sous-catégorie triangulée, saturée et strictement pleine de $D_\QCoh(\mathcal{O}_Y)$, qui est de plus stable par sommes directes. Puisque $D_\QCoh(\mathcal{O}_Y)$ est une catégorie de modules (Catégories dérivées des schémas, théorème \ref{perfect-theorem-DQCoh-is-Ddga}) engendrée par un unique objet parfait (Catégories dérivées des schémas, théorème \ref{perfect-theorem-bondal-van-den-Bergh}), on peut raisonner comme dans Compléments d'algèbre, remarque \ref{more-algebra-remark-P-resolution}, pour voir qu'il suffit de montrer que (\ref{equation-compare-with-pullback}) est un isomorphisme pour un unique objet parfait. Or le résultat vaut pour les objets parfaits, voir le lemme \ref{lemma-compare-with-pullback-perfect}.
\end{proof}

\noindent
Le lemme suivant montre que l'application de changement de base (\ref{equation-base-change-map}) est un isomorphisme pour les morphismes propres, plats et de présentation finie. Nous verrons dans l'exemple \ref{example-base-change-wrong} que cela ne reste pas vrai pour les morphismes parfaits et propres ; dans ce cas, une hypothèse de Tor-indépendance est nécessaire.

\begin{lemma}
\label{lemma-proper-flat-base-change}
Soit $g : Y' \to Y$ un morphisme de schémas quasi-compacts et quasi-séparés. Soit $f : X \to Y$ un morphisme propre et plat de présentation finie. Alors l'application de changement de base (\ref{equation-base-change-map}) est un isomorphisme pour tout $K \in D_\QCoh(\mathcal{O}_Y)$.
\end{lemma}

\begin{proof}
D'après le lemme \ref{lemma-proper-flat-relative}, la formation des foncteurs $a$ et $a'$ commute à la restriction aux ouverts de $Y$ et $Y'$. On peut donc supposer que $Y' \to Y$ est un morphisme de schémas affines, voir la remarque \ref{remark-check-over-affines}. Dans ce cas, l'énoncé résulte du lemme \ref{lemma-more-base-change}.
\end{proof}

\begin{remark}
\label{remark-relative-dualizing-complex}
Soit $Y$ un schéma quasi-compact et quasi-séparé. Soit $f : X \to Y$ un morphisme propre et plat de présentation finie. Soit $a$ l'adjoint du lemme \ref{lemma-twisted-inverse-image} pour $f$. Dans cette situation, $\omega_{X/Y}^\bullet = a(\mathcal{O}_Y)$ est parfois appelé {\it complexe dualisant relatif}. D'après le lemme \ref{lemma-compare-with-pullback-flat-proper}, il existe un isomorphisme fonctoriel
$a(K) = Lf^*K \otimes_{\mathcal{O}_X}^\mathbf{L} \omega_{X/Y}^\bullet$
pour $K \in D_\QCoh(\mathcal{O}_Y)$. De plus, le morphisme trace
$$
\text{Tr}_{f, \mathcal{O}_Y} : Rf_*\omega_{X/Y}^\bullet \to \mathcal{O}_Y
$$
de la section \ref{section-trace} induit le morphisme trace pour tout $K$ dans $D_\QCoh(\mathcal{O}_Y)$. Plus précisément, on a le diagramme
$$
\xymatrix{
Rf_*a(K) \ar[rrr]_{\text{Tr}_{f, K}} \ar@{=}[d] & & &
K \ar@{=}[d] \\
Rf_*(Lf^*K \otimes_{\mathcal{O}_X}^\mathbf{L} \omega_{X/Y}^\bullet)
\ar@{=}[r] &
K \otimes_{\mathcal{O}_Y}^\mathbf{L} Rf_*\omega_{X/Y}^\bullet
\ar[rr]^-{\text{id}_K \otimes \text{Tr}_{f, \mathcal{O}_Y}} & & K
}
$$
où l'égalité en bas à droite est celle de Catégories dérivées des schémas, lemme \ref{perfect-lemma-cohomology-base-change}. Si $g : Y' \to Y$ est un morphisme de schémas quasi-compacts et quasi-séparés et si $X' = Y' \times_Y X$, le lemme \ref{lemma-proper-flat-base-change} donne $\omega_{X'/Y'}^\bullet = L(g')^*\omega_{X/Y}^\bullet$, où $g' : X' \to X$ est la projection ; d'après le lemme \ref{lemma-trace-map-and-base-change}, le morphisme trace
$$
\text{Tr}_{f', \mathcal{O}_{Y'}} :
Rf'_*\omega_{X'/Y'}^\bullet \to \mathcal{O}_{Y'}
$$
pour $f' : X' \to Y'$ est obtenue par changement de base de $\text{Tr}_{f, \mathcal{O}_Y}$ via l'isomorphisme de changement de base.
\end{remark}

\begin{remark}
\label{remark-relative-dualizing-complex-relative-cup-product}
Soient $f : X \to Y$, $\omega^\bullet_{X/Y}$ et $\text{Tr}_{f, \mathcal{O}_Y}$ comme dans la remarque \ref{remark-relative-dualizing-complex}. Soient $K$ et $M$ dans $D_\QCoh(\mathcal{O}_X)$, avec $M$ pseudo-cohérent (par exemple parfait). Supposons donnée une application
$K \otimes_{\mathcal{O}_X}^\mathbf{L} M \to \omega^\bullet_{X/Y}$
correspondant à un isomorphisme
$K \to R\SheafHom_{\mathcal{O}_X}(M, \omega^\bullet_{X/Y})$
via Cohomologie, équation (\ref{cohomology-equation-internal-hom}). Alors le cup-produit relatif (Cohomologie, remarque \ref{cohomology-remark-cup-product})
$$
Rf_*K \otimes_{\mathcal{O}_Y}^\mathbf{L} Rf_*M
\to
Rf_*(K \otimes_{\mathcal{O}_X}^\mathbf{L} M)
\to
Rf_*\omega^\bullet_{X/Y}
\xrightarrow{\text{Tr}_{f, \mathcal{O}_Y}}
\mathcal{O}_Y
$$
détermine un isomorphisme $Rf_*K \to R\SheafHom_{\mathcal{O}_Y}(Rf_*M, \mathcal{O}_Y)$. En effet, puisque $\omega^\bullet_{X/Y} = a(\mathcal{O}_Y)$, l'application canonique (\ref{equation-sheafy-trace})
$$
Rf_*R\SheafHom_{\mathcal{O}_X}(M, \omega^\bullet_{X/Y}) \to
R\SheafHom_{\mathcal{O}_Y}(Rf_*M, \mathcal{O}_Y)
$$
est un isomorphisme d'après le lemme \ref{lemma-iso-on-RSheafHom} et la remarque \ref{remark-iso-on-RSheafHom}, et parce que $M$ et $Rf_*M$ sont pseudo-cohérents, voir Catégories dérivées des schémas, lemme \ref{perfect-lemma-flat-proper-pseudo-coherent-direct-image-general}. Pour voir que le cup-produit relatif induit cet isomorphisme, on utilise la commutativité du diagramme de Cohomologie, remarque \ref{cohomology-remark-relative-cup-and-composition}.
\end{remark}

\begin{lemma}
\label{lemma-properties-relative-dualizing}
Soit $Y$ un schéma quasi-compact et quasi-séparé. Soit $f : X \to Y$ un morphisme de schémas propre, plat et de présentation finie, de complexe dualisant relatif $\omega_{X/Y}^\bullet$ (remarque \ref{remark-relative-dualizing-complex}). Alors
\begin{enumerate}
\item $\omega_{X/Y}^\bullet$ est un objet $Y$-parfait de $D(\mathcal{O}_X)$ ;
\item les faisceaux de cohomologie de $Rf_*\omega_{X/Y}^\bullet$ sont nuls en degrés strictement positifs ;
\item $\mathcal{O}_X \to
R\SheafHom_{\mathcal{O}_X}(\omega_{X/Y}^\bullet, \omega_{X/Y}^\bullet)$ est un isomorphisme.
\end{enumerate}
\end{lemma}

\begin{proof}
Comme la formation de $\omega_{X/Y}^\bullet$ commute au changement de base (voir la remarque \ref{remark-relative-dualizing-complex}), on peut supposer, et l'on suppose, $Y$ affine. Pour un objet parfait $E$ de $D(\mathcal{O}_X)$, on a
\begin{align*}
Rf_*(E \otimes_{\mathcal{O}_X}^\mathbf{L} \omega_{X/Y}^\bullet)
& =
Rf_*R\SheafHom_{\mathcal{O}_X}(E^\vee, \omega_{X/Y}^\bullet) \\
& =
R\SheafHom_{\mathcal{O}_Y}(Rf_*E^\vee, \mathcal{O}_Y) \\
& =
(Rf_*E^\vee)^\vee
\end{align*}
Pour la première égalité, voir Cohomologie, lemme \ref{cohomology-lemma-dual-perfect-complex}. Pour la deuxième, voir le lemme \ref{lemma-iso-on-RSheafHom}, la remarque \ref{remark-iso-on-RSheafHom} et Catégories dérivées des schémas, lemme \ref{perfect-lemma-flat-proper-perfect-direct-image-general}. La troisième égalité est la définition du dual. En particulier, ces références montrent aussi que le résultat est un objet parfait de $D(\mathcal{O}_Y)$. On en conclut que $\omega_{X/Y}^\bullet$ est $Y$-parfait d'après Compléments sur les morphismes, lemme \ref{more-morphisms-lemma-characterize-relatively-perfect}. Cela démontre (1).

\medskip\noindent
Soit $M$ un objet de $D_\QCoh(\mathcal{O}_Y)$. Alors
\begin{align*}
\Hom_Y(M, Rf_*\omega_{X/Y}^\bullet) & =
\Hom_X(Lf^*M, \omega_{X/Y}^\bullet) \\
& =
\Hom_Y(Rf_*Lf^*M, \mathcal{O}_Y) \\
& =
\Hom_Y(M \otimes_{\mathcal{O}_Y}^\mathbf{L} Rf_*\mathcal{O}_X, \mathcal{O}_Y)
\end{align*}
La première égalité résulte de Cohomologie, lemme \ref{cohomology-lemma-adjoint}. La deuxième résulte de la construction de $a$. La troisième résulte de Catégories dérivées des schémas, lemme \ref{perfect-lemma-cohomology-base-change}. Rappelons que $Rf_*\mathcal{O}_X$ est parfait de Tor-amplitude dans $[0, N]$ pour un certain $N$, voir Catégories dérivées des schémas, lemme \ref{perfect-lemma-flat-proper-perfect-direct-image-general}. On peut donc représenter $Rf_*\mathcal{O}_X$ par un complexe de modules projectifs de type fini situés en degrés $[0, N]$ (en utilisant Compléments d'algèbre, lemme \ref{more-algebra-lemma-perfect}, et le fait que $Y$ est affine). Par conséquent, si $M = \mathcal{O}_Y[-i]$ pour un certain $i > 0$, le dernier groupe est nul. Puisque $Y$ est affine, on en conclut que $H^i(Rf_*\omega_{X/Y}^\bullet) = 0$ pour $i > 0$. Cela démontre (2).

\medskip\noindent
Soit $E$ un objet parfait de $D_\QCoh(\mathcal{O}_X)$. On a alors
\begin{align*}
\Hom_X(E, R\SheafHom_{\mathcal{O}_X}(\omega_{X/Y}^\bullet, \omega_{X/Y}^\bullet)
& =
\Hom_X(E \otimes_{\mathcal{O}_X}^\mathbf{L} \omega_{X/Y}^\bullet,
\omega_{X/Y}^\bullet) \\
& =
\Hom_Y(Rf_*(E \otimes_{\mathcal{O}_X}^\mathbf{L} \omega_{X/Y}^\bullet),
\mathcal{O}_Y) \\
& =
\Hom_Y(Rf_*(R\SheafHom_{\mathcal{O}_X}(E^\vee, \omega_{X/Y}^\bullet)),
\mathcal{O}_Y) \\
& =
\Hom_Y(R\SheafHom_{\mathcal{O}_Y}(Rf_*E^\vee, \mathcal{O}_Y),
\mathcal{O}_Y) \\
& =
R\Gamma(Y, Rf_*E^\vee) \\
& =
\Hom_X(E, \mathcal{O}_X)
\end{align*}
La première égalité résulte de Cohomologie, lemme \ref{cohomology-lemma-internal-hom}. La deuxième est la définition de $\omega_{X/Y}^\bullet$. La troisième provient de la construction du complexe parfait dual $E^\vee$, voir Cohomologie, lemme \ref{cohomology-lemma-dual-perfect-complex}. La quatrième résulte de l'égalité
\par\noindent\hfil
$Rf_*R\SheafHom_{\mathcal{O}_X}(E^\vee, \omega_{X/Y}^\bullet) =
R\SheafHom_{\mathcal{O}_Y}(Rf_*E^\vee, \mathcal{O}_Y)$
\hfil\par\noindent
montrée au premier paragraphe de la démonstration. La cinquième résulte de la bidualité des complexes parfaits (Cohomologie, lemme \ref{cohomology-lemma-dual-perfect-complex}) et du fait que $Rf_*E$ est parfait d'après Catégories dérivées des schémas, lemme \ref{perfect-lemma-flat-proper-perfect-direct-image-general}. La dernière égalité est celle de Leray pour $f$. Cette suite d'égalités démontre essentiellement (3) par le lemme de Yoneda. En effet, l'objet
$R\SheafHom(\omega_{X/Y}^\bullet, \omega_{X/Y}^\bullet)$
appartient à $D_\QCoh(\mathcal{O}_X)$ d'après Catégories dérivées des schémas, lemme \ref{perfect-lemma-quasi-coherence-internal-hom}. En prenant $E = \mathcal{O}_X$ ci-dessus, on obtient une application
$\alpha : \mathcal{O}_X \to
R\SheafHom_{\mathcal{O}_X}(\omega_{X/Y}^\bullet, \omega_{X/Y}^\bullet)$
correspondant à $\text{id}_{\mathcal{O}_X} \in \Hom_X(\mathcal{O}_X, \mathcal{O}_X)$. Comme tous les isomorphismes ci-dessus sont fonctoriels en $E$, le cône de $\alpha$ est un objet $C$ de $D_\QCoh(\mathcal{O}_X)$ tel que $\Hom(E, C) = 0$ pour tout $E$ parfait. Puisque les objets parfaits engendrent la catégorie (Catégories dérivées des schémas, théorème \ref{perfect-theorem-bondal-van-den-Bergh}), on en conclut que $\alpha$ est un isomorphisme.
\end{proof}

\begin{lemma}[Rigidité]
\label{lemma-van-den-bergh}
Soit $Y$ un schéma quasi-compact et quasi-séparé. Soit $f : X \to Y$ un morphisme propre et plat de présentation finie, de complexe dualisant relatif $\omega_{X/Y}^\bullet$ (remarque \ref{remark-relative-dualizing-complex}). Il existe un isomorphisme canonique
\begin{equation}
\label{equation-pre-rigid}
\mathcal{O}_X =
c(L\text{pr}_1^*\omega_{X/Y}^\bullet) =
c(L\text{pr}_2^*\omega_{X/Y}^\bullet)
\end{equation}
et un isomorphisme canonique
\begin{equation}
\label{equation-rigid}
\omega_{X/Y}^\bullet =
c\left(L\text{pr}_1^*\omega_{X/Y}^\bullet
\otimes_{\mathcal{O}_{X \times_Y X}}^\mathbf{L}
L\text{pr}_2^*\omega_{X/Y}^\bullet\right)
\end{equation}
où $c$ est l'adjoint à droite du lemme \ref{lemma-twisted-inverse-image} pour la diagonale $\Delta : X \to X \times_Y X$.
\end{lemma}

\begin{proof}
Soit $a$ l'adjoint à droite de $Rf_*$ du lemme \ref{lemma-twisted-inverse-image}. Considérons le carré cartésien
$$
\xymatrix{
X \times_Y X \ar[r]_q \ar[d]_p & X \ar[d]_f \\
X \ar[r]^f & Y
}
$$
Soit $b$ l'adjoint à droite pour $p$ du lemme \ref{lemma-twisted-inverse-image}. Alors
\begin{align*}
\omega_{X/Y}^\bullet
& =
c(b(\omega_{X/Y}^\bullet)) \\
& =
c(Lp^*\omega_{X/Y}^\bullet
\otimes_{\mathcal{O}_{X \times_Y X}}^\mathbf{L} b(\mathcal{O}_X)) \\
& =
c(Lp^*\omega_{X/Y}^\bullet
\otimes_{\mathcal{O}_{X \times_Y X}}^\mathbf{L}
Lq^*a(\mathcal{O}_Y)) \\
& =
c(Lp^*\omega_{X/Y}^\bullet
\otimes_{\mathcal{O}_{X \times_Y X}}^\mathbf{L}
Lq^*\omega_{X/Y}^\bullet)
\end{align*}
comme en (\ref{equation-rigid}). Les justifications sont les suivantes :
\begin{enumerate}
\item La première égalité résulte de $\text{id} = c \circ b$, car $\text{id}_X = p \circ \Delta$.
\item La deuxième résulte du lemme \ref{lemma-compare-with-pullback-flat-proper}.
\item La troisième résulte du lemme \ref{lemma-proper-flat-base-change} et du fait que $\mathcal{O}_X = Lf^*\mathcal{O}_Y$.
\item La quatrième résulte de $\omega_{X/Y}^\bullet = a(\mathcal{O}_Y)$.
\end{enumerate}
L'équation (\ref{equation-pre-rigid}) se démontre exactement de la même manière.
\end{proof}

\begin{remark}
\label{remark-van-den-bergh}
Le lemme \ref{lemma-van-den-bergh} signifie que notre complexe dualisant relatif est {\it rigide}, en un sens analogue à la notion introduite dans \cite{vdB-rigid}. En effet, le foncteur de droite dans (\ref{equation-rigid}) étant « quadratique » en $\omega_{X/Y}^\bullet$, tandis que celui de gauche dans (\ref{equation-rigid}) est « linéaire », cela détermine dans une certaine mesure le complexe $\omega_{X/Y}^\bullet$. Il existe une approche de la théorie de la dualité au moyen de complexes dualisants (relatifs) « rigides », voir par exemple \cite{Neeman-rigid}, \cite{Yekutieli-rigid} et \cite{Yekutieli-Zhang}. Nous y reviendrons à la section \ref{section-relative-dualizing-complexes}.
\end{remark}





\section{Adjoint à droite de l'image directe pour les morphismes parfaits et propres}
\label{section-perfect-proper}

\noindent
Le cadre général approprié à cette section serait celui des morphismes parfaits et propres entre schémas quasi-compacts et quasi-séparés, voir \cite{LN}.

\begin{lemma}
\label{lemma-proper-flat-noetherian}
Soit $f : X \to Y$ un morphisme parfait et propre de schémas noethériens. Soit $a$ l'adjoint à droite de $Rf_* : D_\QCoh(\mathcal{O}_X) \to D_\QCoh(\mathcal{O}_Y)$ du lemme \ref{lemma-twisted-inverse-image}. Alors $a$ commute aux sommes directes.
\end{lemma}

\begin{proof}
Soit $P$ un objet parfait de $D(\mathcal{O}_X)$. D'après Compléments sur les morphismes, lemme \ref{more-morphisms-lemma-perfect-proper-perfect-direct-image}, le complexe $Rf_*P$ est parfait sur $Y$. Soit $K_i$ une famille d'objets de $D_\QCoh(\mathcal{O}_Y)$. Alors
\begin{align*}
\Hom_{D(\mathcal{O}_X)}(P, a(\bigoplus K_i))
& =
\Hom_{D(\mathcal{O}_Y)}(Rf_*P, \bigoplus K_i) \\
& =
\bigoplus \Hom_{D(\mathcal{O}_Y)}(Rf_*P, K_i) \\
& =
\bigoplus \Hom_{D(\mathcal{O}_X)}(P, a(K_i))
\end{align*}
car un objet parfait est compact (Catégories dérivées des schémas, proposition \ref{perfect-proposition-compact-is-perfect}). Puisque $D_\QCoh(\mathcal{O}_X)$ admet un générateur parfait (Catégories dérivées des schémas, théorème \ref{perfect-theorem-bondal-van-den-Bergh}), on en conclut que l'application $\bigoplus a(K_i) \to a(\bigoplus K_i)$ est un isomorphisme, c'est-à-dire que $a$ commute aux sommes directes.
\end{proof}

\begin{lemma}
\label{lemma-proper-flat-noetherian-relative}
Soit $f : X \to Y$ un morphisme parfait et propre de schémas noethériens. Soit $a$ l'adjoint à droite de $Rf_* : D_\QCoh(\mathcal{O}_X) \to D_\QCoh(\mathcal{O}_Y)$ du lemme \ref{lemma-twisted-inverse-image}. Alors
\begin{enumerate}
\item pour tout fermé $T \subset Y$, si $Q \in D_\QCoh(Y)$ est à support dans $T$, alors $a(Q)$ est à support dans $f^{-1}(T)$ ;
\item pour tout ouvert $V \subset Y$ et tout $K \in D_\QCoh(\mathcal{O}_Y)$, l'application (\ref{equation-sheafy}) est un isomorphisme, et
\end{enumerate}
\end{lemma}

\begin{proof}
Cela résulte des lemmes \ref{lemma-when-sheafy}, \ref{lemma-proper-noetherian} et \ref{lemma-proper-flat-noetherian}.
\end{proof}

\begin{lemma}
\label{lemma-compare-with-pullback-flat-proper-noetherian}
Soit $f : X \to Y$ un morphisme parfait et propre de schémas noethériens. L'application (\ref{equation-compare-with-pullback}) est un isomorphisme pour tout objet $K$ de $D_\QCoh(\mathcal{O}_Y)$.
\end{lemma}

\begin{proof}
D'après le lemme \ref{lemma-proper-flat-noetherian}, $a$ commute aux sommes directes. La collection des objets de $D_\QCoh(\mathcal{O}_Y)$ pour lesquels (\ref{equation-compare-with-pullback}) est un isomorphisme est donc une sous-catégorie triangulée, saturée et strictement pleine de $D_\QCoh(\mathcal{O}_Y)$, qui est de plus stable par sommes directes. Puisque $D_\QCoh(\mathcal{O}_Y)$ est une catégorie de modules (Catégories dérivées des schémas, théorème \ref{perfect-theorem-DQCoh-is-Ddga}) engendrée par un unique objet parfait (Catégories dérivées des schémas, théorème \ref{perfect-theorem-bondal-van-den-Bergh}), on peut raisonner comme dans Compléments d'algèbre, remarque \ref{more-algebra-remark-P-resolution}, pour voir qu'il suffit de montrer que (\ref{equation-compare-with-pullback}) est un isomorphisme pour un unique objet parfait. Or le résultat vaut pour les objets parfaits, voir le lemme \ref{lemma-compare-with-pullback-perfect}.
\end{proof}

\begin{lemma}
\label{lemma-proper-perfect-base-change}
Soit $f : X \to Y$ un morphisme parfait et propre de schémas noethériens. Soit $g : Y' \to Y$ un morphisme, avec $Y'$ noethérien. Si $X$ et $Y'$ sont Tor-indépendants sur $Y$, l'application de changement de base (\ref{equation-base-change-map}) est un isomorphisme pour tout $K \in D_\QCoh(\mathcal{O}_Y)$.
\end{lemma}

\begin{proof}
D'après le lemme \ref{lemma-proper-flat-noetherian-relative}, la formation des foncteurs $a$ et $a'$ commute à la restriction aux ouverts de $Y$ et $Y'$. On peut donc supposer que $Y' \to Y$ est un morphisme de schémas affines, voir la remarque \ref{remark-check-over-affines}. Dans ce cas, l'énoncé résulte du lemme \ref{lemma-more-base-change}.
\end{proof}



\section{Adjoint à droite de l'image directe pour les diviseurs de Cartier effectifs}
\label{section-dualizing-Cartier}

\noindent
Soit $X$ un schéma et soit $i : D \to X$ l'inclusion d'un diviseur de Cartier effectif. Notons $\mathcal{N} = i^*\mathcal{O}_X(D)$ le faisceau normal de $i$, voir Morphismes, section \ref{morphisms-section-conormal-sheaf}, et Diviseurs, section \ref{divisors-section-effective-Cartier-divisors}. Rappelons que $R\SheafHom(\mathcal{O}_D, -)$ désigne l'adjoint à droite de $i_* : D(\mathcal{O}_D) \to D(\mathcal{O}_X)$ et vérifie $i_*R\SheafHom(\mathcal{O}_D, -) =
R\SheafHom_{\mathcal{O}_X}(i_*\mathcal{O}_D, -)$, voir la section \ref{section-sections-with-exact-support}.

\begin{lemma}
\label{lemma-compute-for-effective-Cartier}
Comme ci-dessus, soit $X$ un schéma et soit $D \subset X$ un diviseur de Cartier effectif. Il existe un isomorphisme canonique $R\SheafHom(\mathcal{O}_D, \mathcal{O}_X) = \mathcal{N}[-1]$ dans $D(\mathcal{O}_D)$.
\end{lemma}

\begin{proof}
Cela revient à dire que $R\SheafHom(\mathcal{O}_D, \mathcal{O}_X)$ possède un unique faisceau de cohomologie non nul, en degré $1$, et que ce faisceau est isomorphe à $\mathcal{N}$. Puisque $i_*$ est exact et pleinement fidèle, il suffit de montrer que $i_*R\SheafHom(\mathcal{O}_D, \mathcal{O}_X)$ est isomorphe à $i_*\mathcal{N}[-1]$. On a
$i_*R\SheafHom(\mathcal{O}_D, \mathcal{O}_X) =
R\SheafHom_{\mathcal{O}_X}(i_*\mathcal{O}_D, \mathcal{O}_X)$
d'après le lemme \ref{lemma-sheaf-with-exact-support-ext}. On dispose d'une résolution
$$
0 \to \mathcal{I} \to \mathcal{O}_X \to i_*\mathcal{O}_D \to 0
$$
où $\mathcal{I}$ est le faisceau d'idéaux de $D$, qui permet d'effectuer le calcul. Puisque
\par\noindent\hfil
$R\SheafHom_{\mathcal{O}_X}(\mathcal{O}_X, \mathcal{O}_X) = \mathcal{O}_X$
\hfil\par\noindent
et $R\SheafHom_{\mathcal{O}_X}(\mathcal{I}, \mathcal{O}_X) = \mathcal{O}_X(D)$, par un calcul local, on voit que
$$
R\SheafHom_{\mathcal{O}_X}(i_*\mathcal{O}_D, \mathcal{O}_X) =
(\mathcal{O}_X \to \mathcal{O}_X(D))
$$
où le membre de droite comporte $\mathcal{O}_X$ en degré $0$ et $\mathcal{O}_X(D)$ en degré $1$. Le résultat découle de la suite exacte courte
$$
0 \to \mathcal{O}_X \to \mathcal{O}_X(D) \to i_*\mathcal{N} \to 0
$$
qui provient du fait que $D$ est le schéma des zéros de la section canonique de $\mathcal{O}_X(D)$, et du fait que $\mathcal{N} = i^*\mathcal{O}_X(D)$.
\end{proof}

\noindent
Pour tout objet $K$ de $D(\mathcal{O}_X)$, il existe une application canonique
\begin{equation}
\label{equation-map-effective-Cartier}
Li^*K
\otimes_{\mathcal{O}_D}^\mathbf{L}
R\SheafHom(\mathcal{O}_D, \mathcal{O}_X)
\longrightarrow
R\SheafHom(\mathcal{O}_D, K)
\end{equation}
dans $D(\mathcal{O}_D)$, fonctorielle en $K$ et compatible avec les triangles distingués. Cette application est adjointe à une application
$$
i_*(Li^*K \otimes^\mathbf{L}_{\mathcal{O}_D}
R\SheafHom(\mathcal{O}_D, \mathcal{O}_X)) =
K \otimes^\mathbf{L}_{\mathcal{O}_X}
R\SheafHom_{\mathcal{O}_X}(i_*\mathcal{O}_D, \mathcal{O}_X)
\longrightarrow K
$$
où l'égalité est celle de Cohomologie, lemme \ref{cohomology-lemma-projection-formula-closed-immersion}, et où la flèche provient de l'application canonique
$R\SheafHom_{\mathcal{O}_X}(i_*\mathcal{O}_D, \mathcal{O}_X) \to \mathcal{O}_X$
induite par $\mathcal{O}_X \to i_*\mathcal{O}_D$.

\medskip\noindent
Si $K \in D_\QCoh(\mathcal{O}_X)$, alors (\ref{equation-map-effective-Cartier}) est égale à (\ref{equation-compare-with-pullback}) via l'identification $a(K) = R\SheafHom(\mathcal{O}_D, K)$ du lemme \ref{lemma-twisted-inverse-image-closed}. Si $K \in D_\QCoh(\mathcal{O}_X)$ et si $X$ est noethérien, le lemme suivant est un cas particulier du lemme \ref{lemma-compare-with-pullback-flat-proper-noetherian}.

\begin{lemma}
\label{lemma-sheaf-with-exact-support-effective-Cartier}
Comme ci-dessus, soit $X$ un schéma et soit $D \subset X$ un diviseur de Cartier effectif. Alors (\ref{equation-map-effective-Cartier}), combinée au lemme \ref{lemma-compute-for-effective-Cartier}, définit un isomorphisme
$$
Li^*K \otimes_{\mathcal{O}_D}^\mathbf{L} \mathcal{N}[-1]
\longrightarrow
R\SheafHom(\mathcal{O}_D, K)
$$
fonctoriel en $K$ dans $D(\mathcal{O}_X)$.
\end{lemma}

\begin{proof}
Puisque $i_*$ est exact et pleinement fidèle sur les modules, il suffit, pour montrer que l'application est un isomorphisme, de le vérifier après application de $i_*$. Nous utiliserons sans autre mention les suites exactes courtes
$0 \to \mathcal{I} \to \mathcal{O}_X \to i_*\mathcal{O}_D \to 0$
et
$0 \to \mathcal{O}_X \to \mathcal{O}_X(D) \to i_*\mathcal{N} \to 0$
de la démonstration du lemme \ref{lemma-compute-for-effective-Cartier}. D'après Cohomologie, lemme \ref{cohomology-lemma-projection-formula-closed-immersion}, utilisé pour définir l'application (\ref{equation-map-effective-Cartier}), le membre de gauche devient
$$
K \otimes_{\mathcal{O}_X}^\mathbf{L} i_*\mathcal{N}[-1] =
K \otimes_{\mathcal{O}_X}^\mathbf{L} (\mathcal{O}_X \to \mathcal{O}_X(D))
$$
Le membre de droite devient
\begin{align*}
R\SheafHom_{\mathcal{O}_X}(i_*\mathcal{O}_D, K) & =
R\SheafHom_{\mathcal{O}_X}((\mathcal{I} \to \mathcal{O}_X), K) \\
& =
R\SheafHom_{\mathcal{O}_X}((\mathcal{I} \to \mathcal{O}_X), \mathcal{O}_X)
\otimes_{\mathcal{O}_X}^\mathbf{L} K
\end{align*}
la dernière égalité résultant de Cohomologie, lemme \ref{cohomology-lemma-dual-perfect-complex}. Puisque l'application provient de l'isomorphisme
$$
R\SheafHom_{\mathcal{O}_X}((\mathcal{I} \to \mathcal{O}_X), \mathcal{O}_X)
= (\mathcal{O}_X \to \mathcal{O}_X(D))
$$
le lemme est clair.
\end{proof}





\section{Exemples d'adjoints à droite de l'image directe}
\label{section-examples}

\noindent
Dans cette section, nous calculons l'adjoint à droite de l'image directe dans quelques exemples. Les isomorphismes sont canoniques, mais seulement au sens le plus faible : nous ne démontrons ni n'affirmons leur compatibilité avec les diverses opérations telles que le changement de base et la composition des morphismes. Ces questions font l'objet d'une vaste littérature ; on pourra commencer par \cite{RD}, \cite{Conrad-GD} (ces références adoptent un autre point de départ pour la dualité, mais abordent la construction de représentants canoniques des complexes dualisants relatifs), puis consulter les travaux de Joseph Lipman et de ses collaborateurs.

\begin{lemma}
\label{lemma-upper-shriek-P1}
Soit $Y$ un schéma quasi-compact et quasi-séparé. Soit $\mathcal{E}$ un $\mathcal{O}_Y$-module localement libre de rang fini $n + 1$, de déterminant $\mathcal{L} = \wedge^{n + 1}(\mathcal{E})$. Soit $f : X = \mathbf{P}(\mathcal{E}) \to Y$ la projection. Soit $a$ l'adjoint à droite de $Rf_* : D_\QCoh(\mathcal{O}_X) \to D_\QCoh(\mathcal{O}_Y)$ du lemme \ref{lemma-twisted-inverse-image}. Il existe alors un isomorphisme
$$
c : f^*\mathcal{L}(-n - 1)[n] \longrightarrow a(\mathcal{O}_Y)
$$
En particulier, si $\mathcal{E} = \mathcal{O}_Y^{\oplus n + 1}$, alors $X = \mathbf{P}^n_Y$ et l'on obtient $a(\mathcal{O}_Y) = \mathcal{O}_X(-n - 1)[n]$.
\end{lemma}

\begin{proof}
Dans (la démonstration de) Cohomologie des schémas, lemme \ref{coherent-lemma-cohomology-projective-bundle}, nous avons construit un isomorphisme canonique
$$
R^nf_*(f^*\mathcal{L}(-n - 1)) \longrightarrow \mathcal{O}_Y
$$
De plus, $Rf_*(f^*\mathcal{L}(-n - 1))[n] = R^nf_*(f^*\mathcal{L}(-n - 1))$, c'est-à-dire que les autres images directes supérieures sont nulles. On obtient donc un isomorphisme
$$
Rf_*(f^*\mathcal{L}(-n - 1)[n]) \longrightarrow \mathcal{O}_Y
$$
Cet isomorphisme détermine $c$ comme dans l'énoncé, car $a$ est l'adjoint à droite de $Rf_*$. D'après le lemme \ref{lemma-proper-noetherian}, la construction des $a$ est locale sur la base. En particulier, pour vérifier que $c$ est un isomorphisme, on peut travailler localement sur $Y$. Autrement dit, on peut supposer $Y$ affine et $\mathcal{E} = \mathcal{O}_Y^{\oplus n + 1}$. Dans ce cas, les faisceaux $\mathcal{O}_X, \mathcal{O}_X(-1), \ldots,
\mathcal{O}_X(-n)$ engendrent $D_\QCoh(X)$, voir Catégories dérivées des schémas, lemme \ref{perfect-lemma-generator-P1}. Il suffit donc de montrer que
$c : \mathcal{O}_X(-n - 1)[n] \to a(\mathcal{O}_Y)$
est transformée en un isomorphisme par les foncteurs
$$
F_{i, p}(-) = \Hom_{D(\mathcal{O}_X)}(\mathcal{O}_X(i), (-)[p])
$$
pour $i \in \{-n, \ldots, 0\}$ et $p \in \mathbf{Z}$. Pour $F_{0, p}$, c'est vrai par construction de la flèche $c$ ! Pour $i \in \{-n, \ldots, -1\}$, on a
$$
\Hom_{D(\mathcal{O}_X)}(\mathcal{O}_X(i), \mathcal{O}_X(-n - 1)[n + p]) =
H^p(X, \mathcal{O}_X(-n - 1 - i)) = 0
$$
par le calcul de la cohomologie de l'espace projectif (Cohomologie des schémas, lemme \ref{coherent-lemma-cohomology-projective-space-over-ring}), et l'on a
$$
\Hom_{D(\mathcal{O}_X)}(\mathcal{O}_X(i), a(\mathcal{O}_Y)[p]) =
\Hom_{D(\mathcal{O}_Y)}(Rf_*\mathcal{O}_X(i), \mathcal{O}_Y[p]) = 0
$$
car $Rf_*\mathcal{O}_X(i) = 0$ d'après le même lemme. La source et le but de $F_{i, p}(c)$ sont donc nuls et $F_{i, p}(c)$ est nécessairement un isomorphisme. Cela achève la démonstration.
\end{proof}

\begin{example}
\label{example-base-change-wrong}
L'application de changement de base (\ref{equation-base-change-map}) n'est pas nécessairement un isomorphisme lorsque $f$ est parfait et propre et que $g$ est parfait. Soit $k$ un corps. Soit $Y = \mathbf{A}^2_k$ et soit $f : X \to Y$ l'éclatement de $Y$ à l'origine. Notons $E \subset X$ le diviseur exceptionnel. On peut factoriser $f$ sous la forme
$$
X \xrightarrow{i} \mathbf{P}^1_Y \xrightarrow{p} Y
$$
On obtient ainsi une factorisation $a = c \circ b$, où $a$, $b$ et $c$ sont les adjoints à droite du lemme \ref{lemma-twisted-inverse-image} de $Rf_*$, $Rp_*$ et $Ri_*$. Notons $\mathcal{O}(n)$ le faisceau structural tordu au sens de Serre sur $\mathbf{P}^1_Y$, et notons $\mathcal{O}_X(n)$ sa restriction à $X$. Remarquons que $X \subset \mathbf{P}^1_Y$ est défini par une équation de degré un, donc $\mathcal{O}(X) = \mathcal{O}(1)$. Le lemme \ref{lemma-upper-shriek-P1} donne $b(\mathcal{O}_Y) = \mathcal{O}(-2)[1]$. D'après le lemme \ref{lemma-twisted-inverse-image-closed}, on a
$$
a(\mathcal{O}_Y) = c(b(\mathcal{O}_Y)) =
c(\mathcal{O}(-2)[1]) =
R\SheafHom(\mathcal{O}_X, \mathcal{O}(-2)[1]) =
\mathcal{O}_X(-1)
$$
La dernière égalité résulte du lemme \ref{lemma-sheaf-with-exact-support-effective-Cartier}. Soit $Y' = \Spec(k)$ l'origine de $Y$. La restriction de $a(\mathcal{O}_Y)$ à $X' = E = \mathbf{P}^1_k$ est un faisceau inversible de degré $-1$ placé en degré cohomologique $0$. Mais, d'autre part,
$a'(\mathcal{O}_{\Spec(k)}) = \mathcal{O}_E(-2)[1]$
est un faisceau inversible de degré $-2$ placé en degré cohomologique $-1$, donc différent. Dans cet exemple, l'hypothèse de Tor-indépendance du lemme \ref{lemma-more-base-change} n'est pas satisfaite.
\end{example}

\begin{lemma}
\label{lemma-ext}
Soit $Y$ un espace annelé. Soit $\mathcal{I} \subset \mathcal{O}_Y$ un faisceau d'idéaux. Posons $\mathcal{O}_X = \mathcal{O}_Y/\mathcal{I}$ et $\mathcal{N} =
\SheafHom_{\mathcal{O}_Y}(\mathcal{I}/\mathcal{I}^2, \mathcal{O}_X)$. Il existe un isomorphisme canonique $c : \mathcal{N} \to
\SheafExt^1_{\mathcal{O}_Y}(\mathcal{O}_X, \mathcal{O}_X)
$.
\end{lemma}

\begin{proof}
Considérons la suite exacte courte canonique
\begin{equation}
\label{equation-second-order-thickening}
0 \to \mathcal{I}/\mathcal{I}^2 \to \mathcal{O}_Y/\mathcal{I}^2 \to
\mathcal{O}_X \to 0
\end{equation}
Soit $U \subset X$ un ouvert et soit $s \in \mathcal{N}(U)$. On peut former la somme amalgamée de (\ref{equation-second-order-thickening}) suivant $s$ pour obtenir une extension $E_s$ de $\mathcal{O}_X|_U$ par $\mathcal{O}_X|_U$. Celle-ci définit à son tour une section $c(s)$ de
$\SheafExt^1_{\mathcal{O}_Y}(\mathcal{O}_X, \mathcal{O}_X)$
sur $U$. Voir Cohomologie, lemme \ref{cohomology-lemma-section-RHom-over-U}, et Catégories dérivées, lemme \ref{derived-lemma-ext-1}. Réciproquement, étant donnée une extension
$$
0 \to \mathcal{O}_X|_U \to \mathcal{E} \to \mathcal{O}_X|_U \to 0
$$
de $\mathcal{O}_U$-modules, on peut trouver un recouvrement ouvert $U = \bigcup U_i$ et des sections $e_i \in \mathcal{E}(U_i)$ d'image $1 \in \mathcal{O}_X(U_i)$. Alors $e_i$ définit une application $\mathcal{O}_Y|_{U_i} \to \mathcal{E}|_{U_i}$ dont le noyau contient $\mathcal{I}^2$. On voit ainsi que $\mathcal{E}|_{U_i}$ provient d'une somme amalgamée comme ci-dessus. Cela montre que $c$ est surjective. Nous omettons la démonstration de l'injectivité.
\end{proof}

\begin{lemma}
\label{lemma-regular-ideal-ext}
Soit $Y$ un espace annelé. Soit $\mathcal{I} \subset \mathcal{O}_Y$ un faisceau d'idéaux. Posons $\mathcal{O}_X = \mathcal{O}_Y/\mathcal{I}$. Si $\mathcal{I}$ est Koszul-régulier (Diviseurs, définition \ref{divisors-definition-regular-ideal-sheaf}), la composition sur $R\SheafHom_{\mathcal{O}_Y}(\mathcal{O}_X, \mathcal{O}_X)$ définit des isomorphismes
$$
\wedge^i(\SheafExt^1_{\mathcal{O}_Y}(\mathcal{O}_X, \mathcal{O}_X))
\longrightarrow
\SheafExt^i_{\mathcal{O}_Y}(\mathcal{O}_X, \mathcal{O}_X)
$$
pour tout $i$.
\end{lemma}

\begin{proof}
Par composition, nous entendons l'application
$$
R\SheafHom_{\mathcal{O}_Y}(\mathcal{O}_X, \mathcal{O}_X)
\otimes_{\mathcal{O}_Y}^\mathbf{L}
R\SheafHom_{\mathcal{O}_Y}(\mathcal{O}_X, \mathcal{O}_X)
\longrightarrow
R\SheafHom_{\mathcal{O}_Y}(\mathcal{O}_X, \mathcal{O}_X)
$$
de Cohomologie, lemme \ref{cohomology-lemma-internal-hom-composition}. Elle induit des applications de multiplication
$$
\SheafExt^a_{\mathcal{O}_Y}(\mathcal{O}_X, \mathcal{O}_X)
\otimes_{\mathcal{O}_Y}
\SheafExt^b_{\mathcal{O}_Y}(\mathcal{O}_X, \mathcal{O}_X)
\longrightarrow
\SheafExt^{a + b}_{\mathcal{O}_Y}(\mathcal{O}_X, \mathcal{O}_X)
$$
À comparer avec Compléments d'algèbre, équation (\ref{more-algebra-equation-simple-tor-product}). Le lemme signifie que l'application induite
$$
\SheafExt^1_{\mathcal{O}_Y}(\mathcal{O}_X, \mathcal{O}_X)
\otimes \ldots \otimes
\SheafExt^1_{\mathcal{O}_Y}(\mathcal{O}_X, \mathcal{O}_X)
\longrightarrow
\SheafExt^i_{\mathcal{O}_Y}(\mathcal{O}_X, \mathcal{O}_X)
$$
se factorise par le produit extérieur et induit alors un isomorphisme. Pour le vérifier, on peut travailler localement sur $Y$. On peut donc supposer qu'il existe des sections globales $f_1, \ldots, f_r$ de $\mathcal{O}_Y$ qui engendrent $\mathcal{I}$ et forment une suite Koszul-régulière. Notons
$$
\mathcal{A} = \mathcal{O}_Y\langle \xi_1, \ldots, \xi_r\rangle
$$
le faisceau de $\mathcal{O}_Y$-algèbres différentielles graduées strictement commutatives qui est une algèbre de polynômes (à puissances divisées) sur les $\xi_1, \ldots, \xi_r$ en degré $-1$ sur $\mathcal{O}_Y$, de différentielle $\text{d}$ définie par $\text{d}\xi_i = f_i$. Notons $\mathcal{A}^\bullet$ le complexe sous-jacent de $\mathcal{O}_Y$-modules, qui est le complexe de Koszul mentionné ci-dessus. L'application canonique
$\mathcal{A}^\bullet \to \mathcal{O}_X$
est donc un quasi-isomorphisme. On obtient des quasi-isomorphismes
$$
R\SheafHom_{\mathcal{O}_Y}(\mathcal{O}_X, \mathcal{O}_X) \to
\SheafHom^\bullet(\mathcal{A}^\bullet, \mathcal{A}^\bullet) \to
\SheafHom^\bullet(\mathcal{A}^\bullet, \mathcal{O}_X)
$$
d'après Cohomologie, lemme \ref{cohomology-lemma-Rhom-strictly-perfect}. Les différentielles du dernier complexe sont nulles, d'où
$$
\SheafExt^i_{\mathcal{O}_Y}(\mathcal{O}_X, \mathcal{O}_X)
\cong \SheafHom_{\mathcal{O}_Y}(\mathcal{A}^{-i}, \mathcal{O}_X)
$$
Pour $j \in \{1, \ldots, r\}$, soit $\delta_j : \mathcal{A} \to \mathcal{A}$ la dérivation de degré $1$ telle que $\delta_j(\xi_i) = \delta_{ij}$ (symbole de Kronecker). Un calcul donne $\delta_j \circ \text{d} = - \text{d} \circ \delta_j$, ce qui fournit un morphisme de complexes
$$
\delta_j : \mathcal{A}^\bullet \to \mathcal{A}^\bullet[1].
$$
Ainsi, $\delta_j$ définit une section du faisceau $\SheafExt$ correspondant. Un autre calcul montre que les images de $\delta_1, \ldots, \delta_r$ forment une base de $\SheafHom_{\mathcal{O}_Y}(\mathcal{A}^{-1}, \mathcal{O}_X)$ sur $\mathcal{O}_X$. Puisqu'il est clair que $\delta_j \circ \delta_j = 0$ et $\delta_j \circ \delta_{j'} = - \delta_{j'} \circ \delta_j$ comme endomorphismes de $\mathcal{A}$, donc dans les faisceaux $\SheafExt$, on obtient la factorisation de notre application par la puissance extérieure. Pour vérifier que l'on obtient l'isomorphisme voulu, le lecteur vérifiera que les images des éléments
$$
\delta_{j_1} \circ \ldots \circ \delta_{j_i}
$$
pour $j_1 < \ldots < j_i$ forment une base du faisceau
$\SheafHom_{\mathcal{O}_Y}(\mathcal{A}^{-i}, \mathcal{O}_X)$
sur $\mathcal{O}_X$.
\end{proof}

\begin{lemma}
\label{lemma-regular-immersion-ext}
Soit $Y$ un espace annelé. Soit $\mathcal{I} \subset \mathcal{O}_Y$ un faisceau d'idéaux. Posons $\mathcal{O}_X = \mathcal{O}_Y/\mathcal{I}$ et $\mathcal{N} =
\SheafHom_{\mathcal{O}_Y}(\mathcal{I}/\mathcal{I}^2, \mathcal{O}_X)$. Si $\mathcal{I}$ est Koszul-régulier (Diviseurs, définition \ref{divisors-definition-regular-ideal-sheaf}), alors
$$
R\SheafHom_{\mathcal{O}_Y}(\mathcal{O}_X, \mathcal{O}_Y) =
\wedge^r \mathcal{N}[-r]
$$
où $r : Y \to \{1, 2, 3, \ldots \}$ associe à $y$ le nombre minimal de générateurs de $\mathcal{I}$ nécessaires dans un voisinage de $y$.
\end{lemma}

\begin{proof}
Les lemmes \ref{lemma-ext} et \ref{lemma-regular-ideal-ext} donnent des isomorphismes
$\wedge^i\mathcal{N} \to
\SheafExt^i_{\mathcal{O}_Y}(\mathcal{O}_X, \mathcal{O}_X)$
pour $i \geq 0$. Il suffit donc de montrer que l'application $\mathcal{O}_Y \to \mathcal{O}_X$ induit un isomorphisme
$$
\SheafExt^r_{\mathcal{O}_Y}(\mathcal{O}_X, \mathcal{O}_Y)
\longrightarrow
\SheafExt^r_{\mathcal{O}_Y}(\mathcal{O}_X, \mathcal{O}_X)
$$
et que
$\SheafExt^i_{\mathcal{O}_Y}(\mathcal{O}_X, \mathcal{O}_Y)$
est nul pour $i \not = r$. Ces énoncés sont locaux sur $Y$. On peut donc supposer qu'il existe des sections globales $f_1, \ldots, f_r$ de $\mathcal{O}_Y$ qui engendrent $\mathcal{I}$ et forment une suite Koszul-régulière. Soit $\mathcal{A}^\bullet$ le complexe de Koszul sur $f_1, \ldots, f_r$ introduit dans la démonstration du lemme \ref{lemma-regular-ideal-ext}. Alors
$$
R\SheafHom_{\mathcal{O}_Y}(\mathcal{O}_X, \mathcal{O}_Y) =
\SheafHom^\bullet(\mathcal{A}^\bullet, \mathcal{O}_Y)
$$
d'après Cohomologie, lemme \ref{cohomology-lemma-Rhom-strictly-perfect}. Notons $1 : \mathcal{A}^\bullet \to \mathcal{O}_Y$ l'application de $\mathcal{O}_Y$-algèbres différentielles graduées donnée par l'identité de $\mathcal{A}^0 = \mathcal{O}_Y \to \mathcal{O}_Y$ en degré $0$. Avec $\delta_j$ comme dans la démonstration du lemme \ref{lemma-regular-ideal-ext}, on obtient un isomorphisme de $\mathcal{O}_Y$-modules gradués
$$
\mathcal{O}_Y\langle \delta_1, \ldots, \delta_r\rangle
\longrightarrow
\SheafHom^\bullet(\mathcal{A}^\bullet, \mathcal{O}_Y)
$$
en envoyant $\delta_{j_1} \ldots \delta_{j_i}$ sur $1 \circ \delta_{j_1} \circ \ldots \circ \delta_{j_i}$ en degré $i$. Par cet isomorphisme, la différentielle du membre de droite induit une différentielle $\text{d}$ sur celui de gauche. Nos règles de signes donnent $\text{d}(1) = - \sum f_j \delta_j$. Puisque $\delta_j : \mathcal{A}^\bullet \to \mathcal{A}^\bullet[1]$ est un morphisme de complexes, on en déduit que
$$
\text{d}(\delta_{j_1} \ldots \delta_{j_i}) =
(- \sum f_j \delta_j )\delta_{j_1} \ldots \delta_{j_i}
$$
Observons que $\text{d} = \sum f_j \delta_j$ sur l'algèbre différentielle graduée $\mathcal{A}$. L'application définie par la règle
$$
1 \circ \delta_{j_1} \ldots \delta_{j_i} \longmapsto
(\delta_{j_1} \circ \ldots \circ \delta_{j_i})(\xi_1 \ldots \xi_r)
$$
définit donc un isomorphisme de complexes
$$
\SheafHom^\bullet(\mathcal{A}^\bullet, \mathcal{O}_Y)
\longrightarrow \mathcal{A}^\bullet[-r]
$$
si $r$ est impair, et commute aux différentielles au signe près si $r$ est pair. Dans tous les cas, ces complexes ont des cohomologies isomorphes, ce qui donne l'annulation cherchée. L'isomorphisme sur la cohomologie en degré $r$ induit par l'application
$$
\SheafHom^\bullet(\mathcal{A}^\bullet, \mathcal{O}_Y)
\longrightarrow
\SheafHom^\bullet(\mathcal{A}^\bullet, \mathcal{O}_X)
$$
s'en déduit aussi directement. (On observe que notre choix de conventions pour les complexes de Koszul intervient effectivement dans la définition de l'isomorphisme $R\SheafHom_{\mathcal{O}_X}(\mathcal{O}_X, \mathcal{O}_Y) =
\wedge^r \mathcal{N}[-r]$.)
\end{proof}

\begin{lemma}
\label{lemma-regular-immersion}
Soit $Y$ un schéma quasi-compact et quasi-séparé. Soit $i : X \to Y$ une immersion fermée Koszul-régulière. Soit $a$ l'adjoint à droite de $Ri_* : D_\QCoh(\mathcal{O}_X) \to D_\QCoh(\mathcal{O}_Y)$ du lemme \ref{lemma-twisted-inverse-image}. Il existe alors un isomorphisme
$$
\wedge^r\mathcal{N}[-r] \longrightarrow a(\mathcal{O}_Y)
$$
où
$\mathcal{N} = \SheafHom_{\mathcal{O}_X}(\mathcal{C}_{X/Y}, \mathcal{O}_X)$
est le faisceau normal de $i$ (Morphismes, section \ref{morphisms-section-conormal-sheaf}), et où $r$ est son rang, considéré comme fonction localement constante sur $X$.
\end{lemma}

\begin{proof}
Rappelons, d'après les lemmes \ref{lemma-twisted-inverse-image-closed} et \ref{lemma-sheaf-with-exact-support-ext}, que $a(\mathcal{O}_Y)$ est un objet de $D_\QCoh(\mathcal{O}_X)$ dont l'image directe sur $Y$ est $R\SheafHom_{\mathcal{O}_Y}(i_*\mathcal{O}_X, \mathcal{O}_Y)$. Le résultat découle donc du lemme \ref{lemma-regular-immersion-ext}.
\end{proof}

\begin{lemma}
\label{lemma-smooth-proper}
Soit $S$ un schéma noethérien. Soit $f : X \to S$ un morphisme lisse et propre de dimension relative $d$. Soit $a$ l'adjoint à droite de $Rf_* : D_\QCoh(\mathcal{O}_X) \to D_\QCoh(\mathcal{O}_S)$ du lemme \ref{lemma-twisted-inverse-image}. Il existe alors un isomorphisme
$$
\wedge^d \Omega_{X/S}[d] \longrightarrow a(\mathcal{O}_S)
$$
dans $D(\mathcal{O}_X)$.
\end{lemma}

\begin{proof}
Posons $\omega_{X/S}^\bullet = a(\mathcal{O}_S)$ comme dans la remarque \ref{remark-relative-dualizing-complex}. Soit $c$ l'adjoint à droite du lemme \ref{lemma-twisted-inverse-image} pour $\Delta : X \to X \times_S X$. Puisque $\Delta$ est la diagonale d'un morphisme lisse, c'est une immersion Koszul-régulière, voir Diviseurs, lemme \ref{divisors-lemma-immersion-smooth-into-smooth-regular-immersion}. En particulier, $\Delta$ est un morphisme parfait et propre (Compléments sur les morphismes, lemme \ref{more-morphisms-lemma-regular-immersion-perfect}), et l'on obtient
\begin{align*}
\mathcal{O}_X
& =
c(L\text{pr}_1^*\omega_{X/S}^\bullet) \\
& =
L\Delta^*(L\text{pr}_1^*\omega_{X/S}^\bullet)
\otimes_{\mathcal{O}_X}^\mathbf{L}
c(\mathcal{O}_{X \times_S X}) \\
& =
\omega_{X/S}^\bullet \otimes_{\mathcal{O}_X}^\mathbf{L}
c(\mathcal{O}_{X \times_S X}) \\
& =
\omega_{X/S}^\bullet
\otimes_{\mathcal{O}_X}^\mathbf{L}
\wedge^d(\mathcal{N}_\Delta)[-d]
\end{align*}
La première égalité est (\ref{equation-pre-rigid}), car $\omega_{X/S}^\bullet = a(\mathcal{O}_S)$. La deuxième résulte du lemme \ref{lemma-compare-with-pullback-flat-proper-noetherian}. La troisième résulte de $\text{pr}_1 \circ \Delta = \text{id}_X$. La quatrième résulte du lemme \ref{lemma-regular-immersion}. Observons que $\wedge^d(\mathcal{N}_\Delta)$ est un $\mathcal{O}_X$-module inversible. Ainsi, $\wedge^d(\mathcal{N}_\Delta)[-d]$ est un objet inversible de $D(\mathcal{O}_X)$, et l'on en conclut que $a(\mathcal{O}_S) = \omega_{X/S}^\bullet = \wedge^d(\mathcal{C}_\Delta)[d]$. Puisque le faisceau conormal $\mathcal{C}_\Delta$ de $\Delta$ est $\Omega_{X/S}$ d'après Morphismes, lemme \ref{morphisms-lemma-differentials-diagonal}, la démonstration est achevée.
\end{proof}







\section{Foncteurs image inverse extraordinaire}
\label{section-upper-shriek}

\noindent
Dans cette section, nous construisons, à l'aide de compactifications, les foncteurs $f^!$ pour les morphismes entre schémas séparés de type fini sur une base noethérienne fixée. Comme il est habituel en dualité cohérente, il faut vérifier la commutativité d'un certain nombre de diagrammes. Après avoir lu la construction, on pourra omettre les démonstrations des lemmes et passer à la section suivante, consacrée aux propriétés des foncteurs image inverse extraordinaire.

\begin{situation}
\label{situation-shriek}
Ici, $S$ est un schéma noethérien et $\textit{FTS}_S$ est la catégorie dont
\begin{enumerate}
\item les objets sont les schémas $X$ sur $S$ dont le morphisme structural $X \to S$ est séparé et de type fini ;
\item les morphismes $f : X \to Y$ entre objets sont les morphismes de schémas sur $S$.
\end{enumerate}
\end{situation}

\noindent
Dans la situation \ref{situation-shriek}, étant donné un morphisme $f : X \to Y$ de $\textit{FTS}_S$, nous définirons un foncteur exact
$$
f^! : D_\QCoh^+(\mathcal{O}_Y) \to D_\QCoh^+(\mathcal{O}_X)
$$
entre catégories triangulées. Choisissons pour cela une compactification $X \to \overline{X}$ sur $Y$, ce qui est possible d'après Compléments sur la platitude, théorème \ref{flat-theorem-nagata} et lemme \ref{flat-lemma-compactifyable}. Notons $\overline{f} : \overline{X} \to Y$ le morphisme structural. Soit
$\overline{a} : D_\QCoh(\mathcal{O}_Y) \to D_\QCoh(\mathcal{O}_{\overline{X}})$
l'adjoint à droite de $R\overline{f}_*$ construit au lemme \ref{lemma-twisted-inverse-image}. On pose alors
$$
f^!K  = \overline{a}(K)|_X
$$
pour $K \in D_\QCoh^+(\mathcal{O}_Y)$. Le résultat est un objet de $D_\QCoh^+(\mathcal{O}_X)$ d'après le lemme \ref{lemma-twisted-inverse-image-bounded-below}.

\begin{lemma}
\label{lemma-shriek-well-defined}
Dans la situation \ref{situation-shriek}, soit $f : X \to Y$ un morphisme de $\textit{FTS}_S$. Le foncteur $f^!$ est indépendant, à isomorphisme canonique près, du choix de la compactification.
\end{lemma}

\begin{proof}
La catégorie des compactifications de $X$ sur $Y$ est définie dans Compléments sur la platitude, section \ref{flat-section-compactify}. Elle est non vide d'après Compléments sur la platitude, théorème \ref{flat-theorem-nagata} et lemme \ref{flat-lemma-compactifyable}. À tout choix d'une compactification
$$
j : X \to \overline{X},\quad \overline{f} : \overline{X} \to Y
$$
la construction précédente associe le foncteur $j^* \circ \overline{a} :
D_\QCoh^+(\mathcal{O}_Y) \to D_\QCoh^+(\mathcal{O}_X)$, où $\overline{a}$ est l'adjoint à droite de $R\overline{f}_*$ construit au lemme \ref{lemma-twisted-inverse-image}.

\medskip\noindent
Supposons donné un morphisme $g : \overline{X}_1 \to \overline{X}_2$ entre compactifications $j_i : X \to \overline{X}_i$ sur $Y$ tel que $g^{-1}(j_2(X)) = j_1(X)$\footnote{Cela peut être faux avec notre définition de compactification. Voir Compléments sur la platitude, section \ref{flat-section-compactify}.}. Soit $\overline{c}$ l'adjoint à droite du lemme \ref{lemma-twisted-inverse-image} pour $g$. Alors $\overline{c} \circ \overline{a}_2 = \overline{a}_1$, car ces foncteurs sont adjoints de $R\overline{f}_{2, *} \circ Rg_* = R(\overline{f}_2 \circ g)_*$. Par (\ref{equation-sheafy}), on dispose d'une transformation canonique
$$
j_1^* \circ \overline{c} \longrightarrow j_2^*
$$
de foncteurs
$D^+_\QCoh(\mathcal{O}_{\overline{X}_2}) \to D^+_\QCoh(\mathcal{O}_X)$
qui est un isomorphisme d'après le lemme \ref{lemma-proper-noetherian}. La composée
$$
j_1^* \circ \overline{a}_1 \longrightarrow
j_1^* \circ \overline{c} \circ \overline{a}_2 \longrightarrow
j_2^* \circ \overline{a}_2
$$
est un isomorphisme de foncteurs que nous noterons $\alpha_g$.

\medskip\noindent
Considérons deux compactifications $j_i : X \to \overline{X}_i$, $i = 1, 2$ de $X$ sur $Y$. D'après Compléments sur la platitude, lemme \ref{flat-lemma-compactifications-cofiltered}, (b), on peut trouver une compactification $j : X \to \overline{X}$ d'image dense et des morphismes $g_i : \overline{X} \to \overline{X}_i$ de compactifications. D'après Compléments sur la platitude, lemme \ref{flat-lemma-compactifications-cofiltered}, (c), on a $g_i^{-1}(j_i(X)) = j(X)$. Le paragraphe précédent donne donc des isomorphismes
$$
\alpha_{g_i} :
j^* \circ \overline{a}
\longrightarrow
j_i^* \circ \overline{a}_i
$$
On obtient un isomorphisme
$$
\alpha_{g_2} \circ \alpha_{g_1}^{-1} :
j_1^* \circ \overline{a}_1 \to j_2^* \circ \overline{a}_2
$$
Pour achever la démonstration, il faut montrer que ces isomorphismes sont bien définis. Nous affirmons qu'il suffit de montrer que la composée d'isomorphismes construits au paragraphe précédent en est encore un (voir le paragraphe suivant pour l'énoncé précis). Le lecteur pourra le vérifier rapidement, mais nous détaillons aussi entièrement l'argument dans la suite de ce paragraphe. Considérons en effet un second choix de compactification $j' : X \to \overline{X}'$ d'image dense, avec des morphismes de compactifications $g'_i : \overline{X}' \to \overline{X}_i$. D'après Compléments sur la platitude, lemme \ref{flat-lemma-compactifications-cofiltered}, on peut trouver une compactification $j'' : X \to \overline{X}''$ d'image dense et des morphismes de compactifications $h : \overline{X}'' \to \overline{X}$ et $h' : \overline{X}'' \to \overline{X}'$. On peut même supposer $g_1 \circ h = g'_1 \circ h'$ et $g_2 \circ h = g'_2 \circ h'$. Le résultat du paragraphe suivant donne
$$
\alpha_{g_i} \circ \alpha_h = \alpha_{g_i \circ h} =
\alpha_{g'_i \circ h'} = \alpha_{g'_i} \circ \alpha_{h'}
$$
pour $i = 1, 2$. Puisque ce sont tous des isomorphismes de foncteurs, on en conclut que $\alpha_{g_2} \circ \alpha_{g_1}^{-1} =
\alpha_{g'_2} \circ \alpha_{g'_1}^{-1}$, comme voulu.

\medskip\noindent
Supposons données des compactifications $j_i : X \to \overline{X}_i$ pour $i = 1, 2, 3$, ainsi que des morphismes $g : \overline{X}_1 \to \overline{X}_2$ et $h : \overline{X}_2 \to \overline{X}_3$ de compactifications tels que $g^{-1}(j_2(X)) = j_1(X)$ et $h^{-1}(j_2(X)) = j_3(X)$. Soit $\overline{a}_i$ comme ci-dessus. L'affirmation précédente signifie que
$$
\alpha_g \circ \alpha_h = \alpha_{g \circ h} :
j_1^* \circ \overline{a}_1 \to j_3^* \circ \overline{a}_3
$$
Soit $\overline{c}$, resp.\ $\overline{d}$, l'adjoint à droite du lemme \ref{lemma-twisted-inverse-image} pour $g$, resp.\ $h$. Alors $\overline{c} \circ \overline{a}_2 = \overline{a}_1$ et $\overline{d} \circ \overline{a}_3 = \overline{a}_2$, et l'on dispose de transformations canoniques
$$
j_1^* \circ \overline{c} \longrightarrow j_2^*
\quad\text{et}\quad
j_2^* \circ \overline{d} \longrightarrow j_3^*
$$
de foncteurs
$D^+_\QCoh(\mathcal{O}_{\overline{X}_2}) \to D^+_\QCoh(\mathcal{O}_X)$
et
$D^+_\QCoh(\mathcal{O}_{\overline{X}_3}) \to D^+_\QCoh(\mathcal{O}_X)$
pour les mêmes raisons que précédemment. Notons $\overline{e}$ l'adjoint à droite du lemme \ref{lemma-twisted-inverse-image} pour $h \circ g$. On dispose d'une transformation canonique
$$
j_1^* \circ \overline{e} \longrightarrow j_3^*
$$
de foncteurs
$D^+_\QCoh(\mathcal{O}_{\overline{X}_3}) \to D^+_\QCoh(\mathcal{O}_X)$
donnée par (\ref{equation-sheafy}). Explicitement, il faut montrer que la composée
$$
\alpha_h \circ \alpha_g :
j_1^* \circ \overline{a}_1 \to
j_1^* \circ \overline{c} \circ \overline{a}_2 \to
j_2^* \circ \overline{a}_2 \to
j_2^* \circ \overline{d} \circ \overline{a}_3 \to
j_3^* \circ \overline{a}_3
$$
est égale à la composée
$$
\alpha_{h \circ g} :
j_1^* \circ \overline{a}_1 \to
j_1^* \circ \overline{e} \circ \overline{a}_3 \to
j_3^* \circ \overline{a}_3
$$
On procède en deux étapes. Il faut d'abord montrer que le diagramme
$$
\xymatrix{
\overline{a}_1 \ar[r] \ar[d] & \overline{c} \circ \overline{a}_2 \ar[d] \\
\overline{e} \circ \overline{a}_3 \ar[r] &
\overline{c} \circ \overline{d} \circ \overline{a}_3
}
$$
est commutatif, où la flèche horizontale inférieure provient de l'identification $\overline{e} = \overline{c} \circ \overline{d}$. Cela résulte de la commutativité du diagramme correspondant de foncteurs image directe totale
$$
\xymatrix{
R\overline{f}_{1, *} \ar[r] \ar[d] & Rg_* \circ R\overline{f}_{2, *} \ar[d] \\
R(h \circ g)_* \circ R\overline{f}_{3, *} \ar[r] &
Rg_* \circ Rh_* \circ R\overline{f}_{3, *}
}
$$
(insérer ici une référence ultérieure). Il faut ensuite montrer que la composée
$$
j_1^* \circ \overline{c} \circ \overline{d} \to
j_2^* \circ \overline{d} \to j_3^*
$$
est égale à l'application
$$
j_1^* \circ \overline{e} \to j_3^*
$$
via l'identification $\overline{e} = \overline{c} \circ \overline{d}$. Cela a été démontré au lemme \ref{lemma-compose-base-change-maps} (notons que, dans le cas présent, les morphismes $f', g'$ de ce lemme sont égaux à $\text{id}_X$).
\end{proof}

\begin{lemma}
\label{lemma-upper-shriek-composition}
Dans la situation \ref{situation-shriek}, soient $f : X \to Y$ et $g : Y \to Z$ deux morphismes composables de $\textit{FTS}_S$. Il existe alors un isomorphisme canonique $(g \circ f)^! \to f^! \circ g^!$.
\end{lemma}

\begin{proof}
Choisissons une compactification $i : Y \to \overline{Y}$ de $Y$ sur $Z$, puis une compactification $X \to \overline{X}$ de $X$ sur $\overline{Y}$. On utilise ici deux fois Compléments sur la platitude, théorème \ref{flat-theorem-nagata} et lemme \ref{flat-lemma-compactifyable}. Soit $\overline{a}$ l'adjoint à droite du lemme \ref{lemma-twisted-inverse-image} pour $\overline{X} \to \overline{Y}$, et soit $\overline{b}$ l'adjoint à droite du lemme \ref{lemma-twisted-inverse-image} pour $\overline{Y} \to Z$. Alors $\overline{a} \circ \overline{b}$ est l'adjoint à droite du lemme \ref{lemma-twisted-inverse-image} pour la composée $\overline{X} \to Z$. Ainsi, $g^! = i^* \circ \overline{b}$ et $(g \circ f)^! = (X \to \overline{X})^* \circ \overline{a} \circ \overline{b}$. Soit $U$ l'image inverse de $Y$ dans $\overline{X}$, de sorte que l'on obtient le diagramme commutatif
$$
\xymatrix{
X \ar[r]_j \ar[d] & U \ar[dl] \ar[r]_{j'} & \overline{X} \ar[dl] \\
Y \ar[r]_i \ar[d] & \overline{Y} \ar[dl] \\
Z
}
$$
Soit $\overline{a}'$ l'adjoint à droite du lemme \ref{lemma-twisted-inverse-image} pour $U \to Y$. Alors $f^! = j^* \circ \overline{a}'$. On obtient
$$
\gamma : (j')^* \circ \overline{a} \to \overline{a}' \circ i^*
$$
par (\ref{equation-sheafy}), ce qui permet de définir
$$
(g \circ f)^! =
(j' \circ j)^* \circ \overline{a} \circ \overline{b} =
j^* \circ (j')^* \circ \overline{a} \circ \overline{b}
\to
j^* \circ \overline{a}' \circ i^* \circ \overline{b} =
f^! \circ g^!
$$
qui est un isomorphisme sur les objets de $D_\QCoh^+(\mathcal{O}_Z)$ d'après le lemme \ref{lemma-proper-noetherian}. Pour achever la démonstration, montrons que cet isomorphisme ne dépend pas des choix effectués.

\medskip\noindent
Supposons donnés deux diagrammes
$$
\vcenter{
\xymatrix{
X \ar[r]_{j_1} \ar[d] & U_1 \ar[dl] \ar[r]_{j'_1} & \overline{X}_1 \ar[dl] \\
Y \ar[r]_{i_1} \ar[d] & \overline{Y}_1 \ar[dl] \\
Z
}
}
\quad\text{et}\quad
\vcenter{
\xymatrix{
X \ar[r]_{j_2} \ar[d] & U_2 \ar[dl] \ar[r]_{j'_2} & \overline{X}_2 \ar[dl] \\
Y \ar[r]_{i_2} \ar[d] & \overline{Y}_2 \ar[dl] \\
Z
}
}
$$
On peut d'abord choisir une compactification $i : Y \to \overline{Y}$ d'image dense de $Y$ sur $Z$ dominant à la fois $\overline{Y}_1$ et $\overline{Y}_2$, voir Compléments sur la platitude, lemme \ref{flat-lemma-compactifications-cofiltered}. D'après Compléments sur la platitude, lemme \ref{flat-lemma-right-multiplicative-system}, et Catégories, lemmes \ref{categories-lemma-morphisms-right-localization} et \ref{categories-lemma-equality-morphisms-right-localization}, on peut choisir une compactification $X \to \overline{X}$ d'image dense de $X$ sur $\overline{Y}$, munie de morphismes $\overline{X} \to \overline{X}_1$ et $\overline{X} \to \overline{X}_2$, telle que la composée $\overline{X} \to \overline{Y} \to \overline{Y}_1$ soit égale à la composée $\overline{X} \to \overline{X}_1 \to \overline{Y}_1$ et que la composée $\overline{X} \to \overline{Y} \to \overline{Y}_2$ soit égale à la composée $\overline{X} \to \overline{X}_2 \to \overline{Y}_2$. Il suffit donc de comparer les applications déterminées par nos diagrammes lorsqu'on dispose d'un diagramme commutatif de la forme
$$
\xymatrix{
X \ar[rr]_{j_1} \ar@{=}[d] & &
U_1 \ar[d] \ar[ddll] \ar[rr]_{j'_1} & &
\overline{X}_1 \ar[d] \ar[ddll] \\
X \ar'[r][rr]^-{j_2} \ar[d] & &
U_2 \ar'[dl][ddll] \ar'[r][rr]^-{j'_2} & &
\overline{X}_2 \ar[ddll] \\
Y \ar[rr]^{i_1} \ar@{=}[d] & & \overline{Y}_1 \ar[d] \\
Y \ar[rr]^{i_2} \ar[d] & & \overline{Y}_2 \ar[dll] \\
Z
}
$$
et que les compactifications $X \to \overline{X}_1$ et $Y \to \overline{Y}_2$ sont en outre d'image dense. On note $\overline{a}_i$, $\overline{a}'_i$, $\overline{c}$ et $\overline{c}'$ les adjoints à droite du lemme \ref{lemma-twisted-inverse-image} pour $\overline{X}_i \to \overline{Y}_i$, $U_i \to Y$, $\overline{X}_1 \to \overline{X}_2$ et $U_1 \to U_2$. Chacun des carrés
$$
\xymatrix{
X \ar[r] \ar[d] \ar@{}[dr]|A & U_1 \ar[d] \\
X \ar[r] & U_2
}
\quad
\xymatrix{
U_2 \ar[r] \ar[d] \ar@{}[dr]|B & \overline{X}_2 \ar[d] \\
Y \ar[r] & \overline{Y}_2
}
\quad
\xymatrix{
U_1 \ar[r] \ar[d] \ar@{}[dr]|C & \overline{X}_1 \ar[d] \\
Y \ar[r] & \overline{Y}_1
}
\quad
\xymatrix{
Y \ar[r] \ar[d] \ar@{}[dr]|D & \overline{Y}_1 \ar[d] \\
Y \ar[r] & \overline{Y}_2
}
\quad
\xymatrix{
X \ar[r] \ar[d] \ar@{}[dr]|E & \overline{X}_1 \ar[d] \\
X \ar[r] & \overline{X}_2
}
$$
est cartésien (voir Compléments sur la platitude, lemme \ref{flat-lemma-compactifications-cofiltered}, (c), pour A, D, E ; pour B, C, rappeler que $U_i$ est l'image inverse de $Y$ par $\overline{X}_i \to \overline{Y}_i$), et donne donc une application de changement de base (\ref{equation-sheafy}) comme suit :
$$
\begin{matrix}
\gamma_A : j_1^* \circ \overline{c}' \to j_2^* &
\gamma_B : (j_2')^* \circ \overline{a}_2 \to \overline{a}'_2 \circ i_2^* &
\gamma_C : (j_1')^* \circ \overline{a}_1 \to \overline{a}'_1 \circ i_1^* \\
\gamma_D : i_1^* \circ \overline{d} \to i_2^* &
\gamma_E : (j'_1 \circ j_1)^* \circ \overline{c} \to (j'_2 \circ j_2)^*
\end{matrix}
$$
Notons $f_1^! = j_1^* \circ \overline{a}'_1$, $f_2^! = j_2^* \circ \overline{a}'_2$, $g_1^! = i_1^* \circ \overline{b}_1$, $g_2^! = i_2^* \circ \overline{b}_2$, $(g \circ f)_1^! =
(j_1' \circ j_1)^* \circ \overline{a}_1 \circ \overline{b}_1$ et $(g \circ f)^!_2 =
(j_2' \circ j_2)^* \circ \overline{a}_2 \circ \overline{b}_2$. La construction du premier paragraphe de la démonstration et du lemme \ref{lemma-shriek-well-defined} utilise
\begin{enumerate}
\item $\gamma_C$ pour l'application $(g \circ f)^!_1 \to f_1^! \circ g_1^!$ ;
\item $\gamma_B$ pour l'application $(g \circ f)^!_2 \to f_2^! \circ g_2^!$ ;
\item $\gamma_A$ pour l'application $f_1^! \to f_2^!$ ;
\item $\gamma_D$ pour l'application $g_1^! \to g_2^!$ ;
\item $\gamma_E$ pour l'application $(g \circ f)^!_1 \to (g \circ f)^!_2$.
\end{enumerate}
Il faut montrer que le diagramme
$$
\xymatrix{
(g \circ f)^!_1 \ar[r]_{\gamma_E} \ar[d]_{\gamma_C} &
(g \circ f)^!_2 \ar[d]_{\gamma_B} \\
f_1^! \circ g_1^! \ar[r]^{\gamma_A \circ \gamma_D} & f_2^! \circ g_2^!
}
$$
est commutatif. Nous utiliserons les lemmes \ref{lemma-compose-base-change-maps} et \ref{lemma-compose-base-change-maps-horizontal}, avec les notations (et abus de notation) de la remarque \ref{remark-going-around} (en omettant notamment les produits $\star$ avec les transformations identités). On peut écrire $\gamma_E = \gamma_A \circ \gamma_F$, où
$$
\xymatrix{
U_1 \ar[r] \ar[d] \ar@{}[rd]|F & \overline{X}_1 \ar[d] \\
U_2 \ar[r] & \overline{X}_2
}
$$
On voit donc que
$$
\gamma_B \circ \gamma_E = \gamma_B \circ \gamma_A  \circ \gamma_F
= \gamma_A \circ \gamma_B \circ \gamma_F
$$
la dernière égalité résultant du fait que les deux carrés $A$ et $B$ ne se rencontrent qu'en un sommet (comme dans le dernier argument de la remarque \ref{remark-going-around}). Il suffit donc de montrer que $\gamma_D \circ \gamma_C = \gamma_B \circ \gamma_F$. Puisque ces deux expressions sont égales à l'application (\ref{equation-sheafy}) pour le carré
$$
\xymatrix{
U_1 \ar[r] \ar[d] & \overline{X}_1 \ar[d] \\
Y \ar[r] & \overline{Y}_2
}
$$
on conclut.
\end{proof}

\begin{lemma}
\label{lemma-pseudo-functor}
Dans la situation \ref{situation-shriek}, les constructions des lemmes \ref{lemma-shriek-well-defined} et \ref{lemma-upper-shriek-composition} définissent un pseudo-foncteur de la catégorie $\textit{FTS}_S$ vers la $2$-catégorie des catégories (voir Catégories, définition \ref{categories-definition-functor-into-2-category}).
\end{lemma}

\begin{proof}
Il faut montrer que, pour des morphismes $f : X \to Y$, $g : Y \to Z$, $h : Z \to T$, le diagramme
$$
\xymatrix{
(h \circ g \circ f)^! \ar[r]_{\gamma_{A + B}} \ar[d]_{\gamma_{B + C}} &
f^! \circ (h \circ g)^! \ar[d]^{\gamma_C} \\
(g \circ f)^! \circ h^! \ar[r]^{\gamma_A} & f^! \circ g^! \circ h^!
}
$$
est commutatif (voir ci-dessous pour la signification des $\gamma$). Pour cela, choisissons une compactification $\overline{Z}$ de $Z$ sur $T$, puis une compactification $\overline{Y}$ de $Y$ sur $\overline{Z}$, et enfin une compactification $\overline{X}$ de $X$ sur $\overline{Y}$. On utilise Compléments sur la platitude, théorème \ref{flat-theorem-nagata} et lemme \ref{flat-lemma-compactifyable}. Soit $W \subset \overline{Y}$ l'image inverse de $Z$ par $\overline{Y} \to \overline{Z}$, et soient $U \subset V \subset \overline{X}$ les images inverses de $Y \subset W$ par $\overline{X} \to \overline{Y}$. On obtient le diagramme suivant :
$$
\xymatrix{
X \ar[d]_f \ar[r] & U \ar[r] \ar[d] \ar@{}[dr]|A &
V \ar[d] \ar[r] \ar@{}[rd]|B & \overline{X} \ar[d] \\
Y \ar[d]_g \ar[r] & Y \ar[r] \ar[d] & W \ar[r] \ar[d] \ar@{}[rd]|C &
\overline{Y} \ar[d] \\
Z \ar[d]_h \ar[r] & Z \ar[d] \ar[r] & Z \ar[d] \ar[r] & \overline{Z} \ar[d] \\
T \ar[r] & T \ar[r] & T \ar[r] & T
}
$$
Sans introduire une profusion de notations, mais en raisonnant exactement comme dans la démonstration du lemme \ref{lemma-upper-shriek-composition}, on voit que les applications du premier diagramme affiché utilisent les applications (\ref{equation-sheafy}) associées aux rectangles $A + B$, $B + C$, $A$ et $C$, comme indiqué. D'après les lemmes \ref{lemma-compose-base-change-maps} et \ref{lemma-compose-base-change-maps-horizontal}, on a $\gamma_{A + B} = \gamma_A \circ \gamma_B$ et $\gamma_{B + C} = \gamma_C \circ \gamma_B$ ; l'égalité cherchée est donc vraie pourvu que $\gamma_A \circ \gamma_C = \gamma_C \circ \gamma_A$. Cela résulte du fait que les deux carrés $A$ et $C$ ne se rencontrent qu'en un sommet (comme dans le dernier argument de la remarque \ref{remark-going-around}).
\end{proof}

\begin{lemma}
\label{lemma-map-pullback-to-shriek-well-defined}
Dans la situation \ref{situation-shriek}, soit $f : X \to Y$ un morphisme de $\textit{FTS}_S$. Il existe des applications canoniques
$$
\mu_{f, K} :
Lf^*K \otimes_{\mathcal{O}_X}^\mathbf{L} f^!\mathcal{O}_Y
\longrightarrow
f^!K
$$
fonctorielles en $K$ dans $D^+_\QCoh(\mathcal{O}_Y)$. Si $g : Y \to Z$ est un autre morphisme de $\textit{FTS}_S$, le diagramme
$$
\xymatrix{
Lf^*(Lg^*K \otimes_{\mathcal{O}_Y}^\mathbf{L} g^!\mathcal{O}_Z)
\otimes_{\mathcal{O}_X}^\mathbf{L} f^!\mathcal{O}_Y
\ar@{=}[d] \ar[r]_-{\mu_f} &
f^!(Lg^*K \otimes_{\mathcal{O}_Y}^\mathbf{L} g^!\mathcal{O}_Z)
\ar[r]_-{f^!\mu_g} &
f^!g^!K \ar@{=}[d] \\
Lf^*Lg^*K \otimes_{\mathcal{O}_X}^\mathbf{L} Lf^* g^!\mathcal{O}_Z
\otimes_{\mathcal{O}_X}^\mathbf{L} f^!\mathcal{O}_Y \ar[r]^-{\mu_f} &
Lf^*Lg^*K \otimes_{\mathcal{O}_X}^\mathbf{L} f^!g^!\mathcal{O}_Z
\ar[r]^-{\mu_{g \circ f}} & f^!g^!K
}
$$
est commutatif pour tout $K \in D^+_\QCoh(\mathcal{O}_Z)$.
\end{lemma}

\begin{proof}
Si $f$ est propre, alors $f^! = a$ et l'on peut utiliser (\ref{equation-compare-with-pullback}) ; si $g$ est également propre, le lemme \ref{lemma-transitivity-compare-with-pullback} démontre la commutativité du diagramme (dans un cadre plus général).

\medskip\noindent
Définissons l'application $\mu_{f, K}$. Choisissons une compactification $j : X \to \overline{X}$ de $X$ sur $Y$. Puisque $f^!$ est défini comme $j^* \circ \overline{a}$, on obtient $\mu_{f, K}$ en restreignant l'application (\ref{equation-compare-with-pullback})
$$
L\overline{f}^*K \otimes_{\mathcal{O}_{\overline{X}}}^\mathbf{L}
\overline{a}(\mathcal{O}_Y)
\longrightarrow
\overline{a}(K)
$$
à $X$. Pour voir que cela ne dépend pas du choix de la compactification, on raisonne comme dans la démonstration du lemme \ref{lemma-shriek-well-defined}. Nous recommandons au lecteur de lire d'abord cette démonstration.

\medskip\noindent
Supposons donné un morphisme $g : \overline{X}_1 \to \overline{X}_2$ entre compactifications $j_i : X \to \overline{X}_i$ sur $Y$ tel que $g^{-1}(j_2(X)) = j_1(X)$. Notons $\overline{c}$ l'adjoint à droite de l'image directe du lemme \ref{lemma-twisted-inverse-image} pour le morphisme $g$. Les applications
$$
L\overline{f}_1^*K \otimes_{\mathcal{O}_{\overline{X}}}^\mathbf{L}
\overline{a}_1(\mathcal{O}_Y)
\longrightarrow
\overline{a}_1(K)
\quad\text{et}\quad
L\overline{f}_2^*K \otimes_{\mathcal{O}_{\overline{X}}}^\mathbf{L}
\overline{a}_2(\mathcal{O}_Y)
\longrightarrow
\overline{a}_2(K)
$$
s'insèrent dans le diagramme commutatif
$$
\xymatrix{
Lg^*(L\overline{f}_2^*K \otimes^\mathbf{L}
\overline{a}_2(\mathcal{O}_Y))
\otimes^\mathbf{L} \overline{c}(\mathcal{O}_{\overline{X}_2})
\ar@{=}[d] \ar[r]_-\sigma &
\overline{c}(L\overline{f}_2^*K \otimes^\mathbf{L}
\overline{a}_2(\mathcal{O}_Y)) \ar[r] &
\overline{c}(\overline{a}_2(K)) \ar@{=}[d] \\
L\overline{f}_1^*K \otimes^\mathbf{L} Lg^*\overline{a}_2(\mathcal{O}_Y)
\otimes^\mathbf{L} \overline{c}(\mathcal{O}_{\overline{X}_2})
\ar[r]^-{1 \otimes \tau} &
L\overline{f}_1^*K \otimes^\mathbf{L} \overline{a}_1(\mathcal{O}_Y) \ar[r] &
\overline{a}_1(K)
}
$$
d'après le lemme \ref{lemma-transitivity-compare-with-pullback}. D'après le lemme \ref{lemma-compare-on-open}, les applications $\sigma$ et $\tau$ se restreignent en des isomorphismes au-dessus de $X$. On peut en fait préciser davantage. Rappelons que la démonstration du lemme \ref{lemma-shriek-well-defined} utilisait l'application (\ref{equation-sheafy}) $\gamma : j_1^* \circ \overline{c} \to j_2^*$ pour construire l'isomorphisme $\alpha_g : j_1^* \circ \overline{a}_1 \to j_2^* \circ \overline{a}_2$. L'image inverse de l'application $\sigma$ par $j_1$ est l'identité de $j_2^*\left(L\overline{f}_2^*K \otimes^\mathbf{L}
\overline{a}_2(\mathcal{O}_Y)\right)$ si l'on identifie $j_1^*\overline{c}(\mathcal{O}_{\overline{X}_2})$ à $\mathcal{O}_X$ par $j_1^* \circ \overline{c} \to j_2^*$, voir le lemme \ref{lemma-restriction-compare-with-pullback}. De même, l'image inverse de l'application $\tau : Lg^*\overline{a}_2(\mathcal{O}_Y)
\otimes^\mathbf{L} \overline{c}(\mathcal{O}_{\overline{X}_2}) \to
\overline{a}_1(\mathcal{O}_Y) = \overline{c}(\overline{a}_2(\mathcal{O}_Y))$ est l'identité de $j_2^*\overline{a}_2(\mathcal{O}_Y)$. En prenant l'image inverse par $j_1$ et en appliquant $\gamma$ partout où cela est possible, on obtient donc un diagramme commutatif
$$
\xymatrix{
j_2^*\left(L\overline{f}_2^*K \otimes^\mathbf{L}
\overline{a}_2(\mathcal{O}_Y)\right) \ar[r] \ar[d] &
j_2^*\overline{a}_2(K) \\
j_1^*L\overline{f}_1^*K \otimes^\mathbf{L} j_2^*\overline{a}_2(\mathcal{O}_Y) &
j_1^*(L\overline{f}_1^*K \otimes^\mathbf{L} \overline{a}_1(\mathcal{O}_Y))
\ar[r] \ar[l]_{1 \otimes \alpha_g} &
j_1^* \overline{a}_1(K) \ar[lu]_{\alpha_g}
}
$$
Sa commutativité signifie exactement que l'application $\mu_{f, K}$ construite à l'aide de la compactification $\overline{X}_1$ est égale à l'application $\mu_{f, K}$ construite à l'aide de la compactification $\overline{X}_2$, via l'identification $\alpha_g$ utilisée dans la démonstration du lemme \ref{lemma-shriek-well-defined}. Des arguments catégoriques identiques à ceux de la démonstration du lemme \ref{lemma-shriek-well-defined} montrent alors que $\mu_{f, K}$ est bien définie (un petit détail est omis).

\medskip\noindent
Cela étant, la commutativité du diagramme de notre lemme résulte de la construction de l'isomorphisme $(g \circ f)^! \to f^! \circ g^!$ (première partie de la démonstration du lemme \ref{lemma-upper-shriek-composition}, utilisant $\overline{X} \to \overline{Y} \to Z$) et du résultat du lemme \ref{lemma-transitivity-compare-with-pullback} pour $\overline{X} \to \overline{Y} \to Z$.
\end{proof}



\section{Propriétés des foncteurs image inverse extraordinaire}
\label{section-upper-shriek-properties}

\noindent
Voici quelques propriétés des foncteurs image inverse extraordinaire.

\begin{lemma}
\label{lemma-shriek-open-immersion}
Dans la situation \ref{situation-shriek}, soit $Y$ un objet de $\textit{FTS}_S$ et soit $j : X \to Y$ une immersion ouverte. Il existe alors un isomorphisme canonique de foncteurs $j^! = j^*$.
\end{lemma}

\noindent
Pour un morphisme \'etale $f : X \to Y$ de $\textit{FTS}_S$, on a aussi $f^* \cong f^!$, voir le lemme \ref{lemma-shriek-etale}.

\begin{proof}
Dans ce cas, on peut choisir $\overline{X} = Y$ comme compactification. L'adjoint à droite du lemme \ref{lemma-twisted-inverse-image} pour $\text{id} : Y \to Y$ est alors le foncteur identité, d'où $j^! = j^*$ par définition.
\end{proof}

\begin{lemma}
\label{lemma-restrict-before-or-after}
Dans la situation \ref{situation-shriek}, soit
$$
\xymatrix{
U \ar[r]_j \ar[d]_g & X \ar[d]^f \\
V \ar[r]^{j'} & Y
}
$$
un diagramme commutatif de $\textit{FTS}_S$ où $j$ et $j'$ sont des immersions ouvertes. Alors $j^* \circ f^! = g^! \circ (j')^*$ comme foncteurs $D^+_\QCoh(\mathcal{O}_Y) \to D^+(\mathcal{O}_U)$.
\end{lemma}

\begin{proof}
Soit $h = f \circ j = j' \circ g$. D'après le lemme \ref{lemma-upper-shriek-composition}, on a $h^! = j^! \circ f^! = g^! \circ (j')^!$. D'après le lemme \ref{lemma-shriek-open-immersion}, on a $j^! = j^*$ et $(j')^! = (j')^*$.
\end{proof}

\begin{lemma}
\label{lemma-shriek-affine-line}
Dans la situation \ref{situation-shriek}, soit $Y$ un objet de $\textit{FTS}_S$ et soit $f : X = \mathbf{A}^1_Y \to Y$ la projection. Il existe alors un isomorphisme (non canonique) de foncteurs $f^!(-) \cong Lf^*(-) [1]$.
\end{lemma}

\begin{proof}
Puisque $X = \mathbf{A}^1_Y \subset \mathbf{P}^1_Y$ et puisque $\mathcal{O}_{\mathbf{P}^1_Y}(-2)|_X \cong \mathcal{O}_X$, cela résulte des lemmes \ref{lemma-upper-shriek-P1} et \ref{lemma-compare-with-pullback-flat-proper-noetherian}.
\end{proof}

\begin{lemma}
\label{lemma-shriek-closed-immersion}
Dans la situation \ref{situation-shriek}, soit $Y$ un objet de $\textit{FTS}_S$ et soit $i : X \to Y$ une immersion fermée. Il existe alors un isomorphisme canonique de foncteurs $i^!(-) = R\SheafHom(\mathcal{O}_X, -)$.
\end{lemma}

\begin{proof}
C'est une reformulation du lemme \ref{lemma-twisted-inverse-image-closed}.
\end{proof}

\begin{remark}[Description locale de l'image inverse extraordinaire]
\label{remark-local-calculation-shriek}
Dans la situation \ref{situation-shriek}, soit $f : X \to Y$ un morphisme de $\textit{FTS}_S$. Les lemmes précédents permettent de calculer $f^!$ localement comme suit. Supposons donnés des ouverts affines
$$
\xymatrix{
U \ar[r]_j \ar[d]_g & X \ar[d]^f \\
V \ar[r]^i & Y
}
$$
Puisque $j^! \circ f^! = g^! \circ i^!$ (lemme \ref{lemma-upper-shriek-composition}) et que $j^!$ et $i^!$ sont donnés par restriction (lemme \ref{lemma-shriek-open-immersion}), on voit que
$$
(f^!E)|_U = g^!(E|_V)
$$
pour tout $E \in D^+_\QCoh(\mathcal{O}_X)$. Écrivons $U = \Spec(A)$ et $V = \Spec(R)$, et soit $\varphi : R \to A$ l'homomorphisme d'anneaux de type fini correspondant à $g$. Choisissons une présentation $A = P/I$ où $P = R[x_1, \ldots, x_n]$ est une algèbre de polynômes en $n$ variables sur $R$. Choisissons un objet $K \in D^+(R)$ correspondant à $E|_V$ (Catégories dérivées des schémas, lemme \ref{perfect-lemma-affine-compare-bounded}). Nous affirmons alors que $f^!E|_U$ correspond à
$$
\varphi^!(K) = R\Hom(A, K \otimes_R^\mathbf{L} P)[n]
$$
où $R\Hom(A, -) : D(P) \to D(A)$ est le foncteur de Complexes dualisants, section \ref{dualizing-section-trivial}, et où $\varphi^! : D(R) \to D(A)$ est celui de Complexes dualisants, section \ref{dualizing-section-relative-dualizing-complex-algebraic}. En effet, le choix de la présentation donne une factorisation
$$
U \rightarrow \mathbf{A}^n_V \to \mathbf{A}^{n - 1}_V \to \ldots \to
\mathbf{A}^1_V \to V
$$
En appliquant le lemme \ref{lemma-shriek-affine-line} exactement $n$ fois, on voit que $(\mathbf{A}^n_V \to V)^!(E|_V)$ correspond à $K \otimes_R^\mathbf{L} P[n]$. D'après les lemmes \ref{lemma-sheaf-with-exact-support-quasi-coherent} et \ref{lemma-shriek-closed-immersion}, la dernière étape correspond à l'application de $R\Hom(A, -)$.
\end{remark}

\begin{lemma}
\label{lemma-shriek-coherent}
Dans la situation \ref{situation-shriek}, soit $f : X \to Y$ un morphisme de $\textit{FTS}_S$. Alors $f^!$ envoie $D_{\textit{Coh}}^+(\mathcal{O}_Y)$ dans $D_{\textit{Coh}}^+(\mathcal{O}_X)$.
\end{lemma}

\begin{proof}
La question est locale sur $X$ ; on peut donc supposer que $X$ et $Y$ sont affines. On peut alors factoriser $f : X \to Y$ sous la forme
$$
X \xrightarrow{i} \mathbf{A}^n_Y \to \mathbf{A}^{n - 1}_Y \to \ldots \to
\mathbf{A}^1_Y \to Y
$$
où $i$ est une immersion fermée. Le lemme résulte des lemmes \ref{lemma-shriek-affine-line} et \ref{lemma-sheaf-with-exact-support-coherent}, de Complexes dualisants, lemme \ref{dualizing-lemma-dualizing-polynomial-ring}, et d'une récurrence.
\end{proof}

\begin{lemma}
\label{lemma-shriek-dualizing}
Dans la situation \ref{situation-shriek}, soit $f : X \to Y$ un morphisme de $\textit{FTS}_S$. Si $K$ est un complexe dualisant sur $Y$, alors $f^!K$ est un complexe dualisant sur $X$.
\end{lemma}

\begin{proof}
La question est locale sur $X$ ; on peut donc supposer que $X$ et $Y$ sont affines. On peut alors factoriser $f : X \to Y$ sous la forme
$$
X \xrightarrow{i} \mathbf{A}^n_Y \to \mathbf{A}^{n - 1}_Y \to \ldots \to
\mathbf{A}^1_Y \to Y
$$
où $i$ est une immersion fermée. Le lemme \ref{lemma-shriek-affine-line} et Complexes dualisants, lemme \ref{dualizing-lemma-dualizing-polynomial-ring}, montrent par récurrence que $p^!K$ est un complexe dualisant sur $\mathbf{A}^n_Y$, où $p : \mathbf{A}^n_Y \to Y$ est la projection. De même, Complexes dualisants, lemme \ref{dualizing-lemma-dualizing-quotient}, et les lemmes \ref{lemma-sheaf-with-exact-support-quasi-coherent} et \ref{lemma-shriek-closed-immersion} montrent que $i^!$ transforme les complexes dualisants en complexes dualisants.
\end{proof}

\begin{lemma}
\label{lemma-shriek-via-duality}
Dans la situation \ref{situation-shriek}, soit $f : X \to Y$ un morphisme de $\textit{FTS}_S$. Soit $K$ un complexe dualisant sur $Y$. Posons $D_Y(M) = R\SheafHom_{\mathcal{O}_Y}(M, K)$ pour $M \in D_{\textit{Coh}}(\mathcal{O}_Y)$ et $D_X(E) = R\SheafHom_{\mathcal{O}_X}(E, f^!K)$ pour $E \in D_{\textit{Coh}}(\mathcal{O}_X)$. Il existe alors un isomorphisme canonique
$$
f^!M \longrightarrow D_X(Lf^*D_Y(M))
$$
pour $M \in D_{\textit{Coh}}^+(\mathcal{O}_Y)$.
\end{lemma}

\begin{proof}
Choisissons une compactification $j : X \subset \overline{X}$ de $X$ sur $Y$ (Compléments sur la platitude, théorème \ref{flat-theorem-nagata} et lemme \ref{flat-lemma-compactifyable}). Soit $a$ l'adjoint à droite du lemme \ref{lemma-twisted-inverse-image} pour $\overline{X} \to Y$. Posons
$D_{\overline{X}}(E) = R\SheafHom_{\mathcal{O}_{\overline{X}}}(E, a(K))$
pour $E \in D_{\textit{Coh}}(\mathcal{O}_{\overline{X}})$. Comme la formation de $R\SheafHom$ commute à la restriction aux ouverts et que $f^! = j^* \circ a$, il suffit de montrer qu'il existe un isomorphisme canonique
$$
a(M) \longrightarrow D_{\overline{X}}(L\overline{f}^*D_Y(M))
$$
pour $M \in D_{\textit{Coh}}(\mathcal{O}_Y)$. Pour $F \in D_\QCoh(\mathcal{O}_X)$, on a
\begin{align*}
\Hom_{\overline{X}}(
F, D_{\overline{X}}(L\overline{f}^*D_Y(M)))
& =
\Hom_{\overline{X}}(
F \otimes_{\mathcal{O}_X}^\mathbf{L} L\overline{f}^*D_Y(M), a(K)) \\
& =
\Hom_Y(
R\overline{f}_*(F \otimes_{\mathcal{O}_X}^\mathbf{L} L\overline{f}^*D_Y(M)),
K) \\
& =
\Hom_Y(
R\overline{f}_*(F) \otimes_{\mathcal{O}_Y}^\mathbf{L} D_Y(M),
K) \\
& =
\Hom_Y(
R\overline{f}_*(F), D_Y(D_Y(M))) \\
& =
\Hom_Y(R\overline{f}_*(F), M) \\
& = \Hom_{\overline{X}}(F, a(M))
\end{align*}
La première égalité résulte de Cohomologie, lemme \ref{cohomology-lemma-internal-hom}. La deuxième résulte de la définition de $a$. La troisième résulte de Catégories dérivées des schémas, lemme \ref{perfect-lemma-cohomology-base-change}. La quatrième résulte de Cohomologie, lemme \ref{cohomology-lemma-internal-hom}, et de la définition de $D_Y$. La cinquième résulte du lemme \ref{lemma-dualizing-schemes}. La dernière résulte de la définition de $a$. Le lemme de Yoneda donne donc $a(M) = D_{\overline{X}}(L\overline{f}^*D_Y(M))$.
\end{proof}

\begin{lemma}
\label{lemma-perfect-comparison-shriek}
Dans la situation \ref{situation-shriek}, soit $f : X \to Y$ un morphisme de $\textit{FTS}_S$. Supposons $f$ parfait (par exemple plat). Alors
\begin{enumerate}
\item[(a)] $f^!\mathcal{O}_Y$ appartient à $D_{\textit{Coh}}^b(\mathcal{O}_X)$ ;
\item[(b)] $f^!\mathcal{O}_Y$ est de Tor-dimension finie dans $D(f^{-1}\mathcal{O}_Y)$ ;
\item[(c)] $\mathcal{O}_X \to
R\SheafHom_{\mathcal{O}_X}(f^!\mathcal{O}_Y, f^!\mathcal{O}_Y)$ est un isomorphisme ;
\item[(d)] $f^!$ envoie $D_{\textit{Coh}}^b(\mathcal{O}_Y)$ dans $D_{\textit{Coh}}^b(\mathcal{O}_X)$ ;
\item[(e)] l'application
$\mu_{f,  K} :
Lf^*K \otimes_{\mathcal{O}_X}^\mathbf{L} f^!\mathcal{O}_Y
\to
f^!K$
du lemme \ref{lemma-map-pullback-to-shriek-well-defined} est un isomorphisme pour tout $K \in D_\QCoh^+(\mathcal{O}_Y)$.
\end{enumerate}
\end{lemma}

\begin{proof}
(Un morphisme plat de présentation finie est parfait, voir Compléments sur les morphismes, lemme \ref{more-morphisms-lemma-flat-finite-presentation-perfect}.) Les assertions (a), (b) et (c) sont locales sur $X$. On peut donc supposer $X$ et $Y$ affines. La remarque \ref{remark-local-calculation-shriek} transforme alors (a), (b) et (c) en (1)(a), (1)(b) et (1)(c) de Complexes dualisants, lemme \ref{dualizing-lemma-relative-dualizing-algebraic}. (Utiliser Catégories dérivées des schémas, lemmes \ref{perfect-lemma-tor-dimension-rel-affine}, \ref{perfect-lemma-pseudo-coherent}, \ref{perfect-lemma-identify-pseudo-coherent-noetherian} \ref{perfect-lemma-quasi-coherence-internal-hom}, pour faire correspondre les énoncés.)

\medskip\noindent
Pour démontrer (d) et (e), commençons par quelques remarques préliminaires.
\begin{enumerate}
\item Nous savons déjà que $f^!$ envoie $D_{\textit{Coh}}^+(\mathcal{O}_Y)$ dans $D_{\textit{Coh}}^+(\mathcal{O}_X)$, voir le lemme \ref{lemma-shriek-coherent}.
\item Si $f$ est une immersion ouverte, alors $f^! = f^*$, et (d) et (e) sont vraies, car on peut prendre $\overline{X} = Y$ dans la construction de $f^!$ et $\mu_f$, voir le lemme \ref{lemma-shriek-open-immersion}.
\item Si $f$ est un morphisme parfait et propre, (e) est vraie d'après le lemme \ref{lemma-compare-with-pullback-flat-proper-noetherian}.
\item S'il existe un recouvrement ouvert $X = \bigcup U_i$ tel que (d) soit vraie pour $U_i \to Y$, alors (d) est vraie pour $X \to Y$. Il en va de même pour (e). Cela tient au fait que la construction de $f^!$ et $\mu_f$ commute au passage aux sous-schémas ouverts.
\item Si $g : Y \to Z$ est un second morphisme parfait dans $\textit{FTS}_S$ et si (e) vaut pour $f$ et $g$, alors
$f^!g^!\mathcal{O}_Z =
Lf^*g^!\mathcal{O}_Z \otimes_{\mathcal{O}_X}^\mathbf{L} f^!\mathcal{O}_Y$
et (e) vaut pour $g \circ f$ par le diagramme commutatif du lemme \ref{lemma-map-pullback-to-shriek-well-defined}.
\item Si (d) et (e) valent pour $f$ et $g$, elles valent pour $g \circ f$. En effet, $f^!g^!\mathcal{O}_Z$ est alors borné supérieurement (par le point précédent), $L(g \circ f)^*$ est de dimension cohomologique finie, et (d) résulte de (e), déjà démontrée ci-dessus.
\end{enumerate}
Ces remarques montrent qu'il suffit de démontrer (d) et (e) lorsque $X$ est affine. Choisissons une immersion $X \to \mathbf{A}^n_Y$ (Morphismes, lemme \ref{morphisms-lemma-quasi-affine-finite-type-over-S}), que l'on factorise en $X \to U \to \mathbf{A}^n_Y \to Y$, où $X \to U$ est une immersion fermée et $U \subset \mathbf{A}^n_Y$ est ouverte. Remarquons que $X \to U$ est une immersion fermée parfaite d'après Compléments sur les morphismes, lemme \ref{more-morphisms-lemma-perfect-permanence}. Il suffit donc de démontrer le lemme pour une immersion fermée parfaite et pour la projection $\mathbf{A}^n_Y \to Y$.

\medskip\noindent
Soit $f : X \to Y$ une immersion fermée parfaite. Nous savons déjà que (d) est vraie. Soit $K \in D^b_{\textit{Coh}}(\mathcal{O}_Y)$. Alors $f^!K = R\SheafHom(\mathcal{O}_X, K)$ (lemme \ref{lemma-shriek-closed-immersion}) et $f_*f^!K = R\SheafHom_{\mathcal{O}_Y}(f_*\mathcal{O}_X, K)$. Puisque $f$ est parfait, le complexe $f_*\mathcal{O}_X$ est parfait ; ainsi, $R\SheafHom_{\mathcal{O}_Y}(f_*\mathcal{O}_X, K)$ est borné supérieurement. Cela démontre (d). Certains détails sont omis.

\medskip\noindent
Soit $f : \mathbf{A}^n_Y \to Y$ la projection. Alors (d) résulte d'applications successives du lemme \ref{lemma-shriek-affine-line}. Enfin, (e) est vraie, car elle vaut pour $\mathbf{P}^n_Y \to Y$ (plat et propre) et parce que $\mathbf{A}^n_Y \subset \mathbf{P}^n_Y$ est un ouvert.
\end{proof}

\begin{lemma}
\label{lemma-flat-shriek-relatively-perfect}
Dans la situation \ref{situation-shriek}, soit $f : X \to Y$ un morphisme de $\textit{FTS}_S$. Si $f$ est plat, alors $f^!\mathcal{O}_Y$ est un objet $Y$-parfait de $D(\mathcal{O}_X)$, et
\par\noindent\hfil
$\mathcal{O}_X \to
R\SheafHom_{\mathcal{O}_X}(f^!\mathcal{O}_Y, f^!\mathcal{O}_Y)$
\hfil\par\noindent
est un isomorphisme.
\end{lemma}

\begin{proof}
Les deux assertions sont locales sur $X$. On peut donc supposer $X$ et $Y$ affines. La remarque \ref{remark-local-calculation-shriek} transforme alors le lemme en un lemme algébrique, à savoir Complexes dualisants, lemme \ref{dualizing-lemma-relative-dualizing-algebraic}. (Utiliser Catégories dérivées des schémas, lemme \ref{perfect-lemma-affine-locally-rel-perfect}, pour passer d'un langage à l'autre.)
\end{proof}

\begin{lemma}
\label{lemma-lci-shriek}
Dans la situation \ref{situation-shriek}, soit $f : X \to Y$ un morphisme de $\textit{FTS}_S$. Supposons que $f : X \to Y$ soit un morphisme localement d'intersection complète. Alors
\begin{enumerate}
\item $f^!\mathcal{O}_Y$ est un objet inversible de $D(\mathcal{O}_X)$ ;
\item $f^!$ envoie les complexes parfaits sur des complexes parfaits.
\end{enumerate}
\end{lemma}

\begin{proof}
Rappelons qu'un morphisme localement d'intersection complète est parfait, voir Compléments sur les morphismes, lemme \ref{more-morphisms-lemma-lci-properties}. D'après le lemme \ref{lemma-perfect-comparison-shriek}, il suffit de montrer que $f^!\mathcal{O}_Y$ est un objet inversible de $D(\mathcal{O}_X)$. Cette question est locale sur $X$ et $Y$. On peut donc supposer que $X \to Y$ se factorise en $X \to \mathbf{A}^n_Y \to Y$, où la première flèche est une immersion Koszul-régulière. Voir Compléments sur les morphismes, section \ref{more-morphisms-section-lci}. Le résultat vaut pour $\mathbf{A}^n_Y \to Y$ d'après le lemme \ref{lemma-shriek-affine-line}. Il suffit donc de démontrer le lemme lorsque $f$ est une immersion Koszul-régulière. En travaillant encore localement, on se ramène au cas $X = \Spec(A)$ et $Y = \Spec(B)$, où $A = B/(f_1, \ldots, f_r)$ pour une suite régulière $f_1, \ldots, f_r \in B$ (utiliser le fait que, sur les anneaux locaux noethériens, les notions de suite Koszul-régulière et de suite régulière coïncident, voir Compléments d'algèbre, lemme \ref{more-algebra-lemma-noetherian-finite-all-equivalent}). Ainsi, $X \to Y$ est une composée
$$
X = X_r \to X_{r - 1} \to \ldots \to X_1 \to X_0 = Y
$$
où chaque flèche est l'inclusion d'un diviseur de Cartier effectif. On se ramène ainsi à l'inclusion d'un diviseur de Cartier effectif $i : D \to X$. Dans ce cas, $i^!\mathcal{O}_X = \mathcal{N}[1]$ d'après le lemme \ref{lemma-compute-for-effective-Cartier}, ce qui achève la démonstration.
\end{proof}







\section{Changement de base pour l'image inverse extraordinaire}
\label{section-base-change-shriek}

\noindent
Dans la situation \ref{situation-shriek}, soit
$$
\xymatrix{
X' \ar[r]_{g'} \ar[d]_{f'} & X \ar[d]^f \\
Y' \ar[r]^g & Y
}
$$
un diagramme cartésien de $\textit{FTS}_S$ tel que $X$ et $Y'$ soient Tor-indépendants sur $Y$. Notre cadre actuel ne suffit pas à construire une application de changement de base $L(g')^* \circ f^! \to (f')^! \circ Lg^*$ dans cette généralité. En effet, on ne peut pas en général choisir une compactification $j : X \to \overline{X}$ sur $Y$ telle que $\overline{X}$ et $Y'$ soient Tor-indépendants sur $Y$ ; la construction de l'application de changement de base donnée à la section \ref{section-base-change-map} ne s'applique donc pas\footnote{Le lecteur familier de la géométrie algébrique dérivée reconnaîtra qu'il ne s'agit pas d'une difficulté « réelle ». En prenant $\overline{X}'$ égal au produit fibré dérivé de $\overline{X}$ et $Y'$ sur $Y$, on peut raisonner exactement comme dans la démonstration du lemme \ref{lemma-base-change-shriek-flat} pour définir cette application. Après tout, la Tor-indépendance de $X$ et $Y'$ garantit que $X'$ sera un sous-schéma ouvert du schéma dérivé $\overline{X}'$.}.

\medskip\noindent
On trouvera une solution partielle à la section \ref{section-relative-dualizing-complexes}. En effet, si le morphisme $f$ est plat, on dispose d'une bonne notion de complexe dualisant relatif ; les lemmes \ref{lemma-compactifyable-relative-dualizing} \ref{lemma-base-change-relative-dualizing} et \ref{lemma-perfect-comparison-shriek} permettent alors de construire un isomorphisme canonique de changement de base. Si nous devons l'utiliser, nous ajouterons des énoncés précis et leurs démonstrations plus loin dans ce chapitre.

\begin{lemma}
\label{lemma-base-change-shriek-flat}
Dans la situation \ref{situation-shriek}, soit
$$
\xymatrix{
X' \ar[r]_{g'} \ar[d]_{f'} & X \ar[d]^f \\
Y' \ar[r]^g & Y
}
$$
un diagramme cartésien de $\textit{FTS}_S$, avec $g$ plat. Il existe alors un isomorphisme $L(g')^* \circ f^! \to (f')^! \circ Lg^*$ sur $D_\QCoh^+(\mathcal{O}_Y)$.
\end{lemma}

\begin{proof}
Puisque $g$ est plat, pour tout choix de compactification $j : X \to \overline{X}$ de $X$ sur $Y$, le schéma $\overline{X}$ est Tor-indépendant de $Y'$. Notons $j' : X' \to \overline{X}'$ le changement de base de $j$ et $\overline{g}' : \overline{X}' \to \overline{X}$ la projection. On définit l'application de changement de base comme la composée
$$
L(g')^* \circ f^! = L(g')^* \circ j^* \circ a =
(j')^* \circ L(\overline{g}')^* \circ a \longrightarrow
(j')^* \circ a' \circ Lg^* = (f')^! \circ Lg^*
$$
où la flèche centrale est l'application de changement de base (\ref{equation-base-change-map}), et où $a$ et $a'$ sont les adjoints à droite de l'image directe du lemme \ref{lemma-twisted-inverse-image} pour $\overline{X} \to Y$ et $\overline{X}' \to Y'$. Cette construction ne dépend pas du choix de la compactification (nous formulerons et démontrerons un lemme précis si nous avons besoin de ce résultat).

\medskip\noindent
Pour achever la démonstration, il suffit de montrer que l'application de changement de base $L(g')^* \circ a \to a' \circ Lg^*$ est un isomorphisme sur $D_\QCoh^+(\mathcal{O}_Y)$. D'après le lemme \ref{lemma-proper-noetherian}, la formation de $a$ et $a'$ commute à la restriction aux ouverts affines de $Y$ et $Y'$. D'après la remarque \ref{remark-check-over-affines}, on peut donc supposer $Y$ et $Y'$ affines. Le résultat découle alors du lemme \ref{lemma-more-base-change}.
\end{proof}

\begin{lemma}
\label{lemma-shriek-etale}
Dans la situation \ref{situation-shriek}, soit $f : X \to Y$ un morphisme \'etale de $\textit{FTS}_S$. Alors $f^! \cong f^*$ comme foncteurs sur $D^+_\QCoh(\mathcal{O}_Y)$.
\end{lemma}

\begin{proof}
Nous allons utiliser le fait qu'un morphisme \'etale est plat, syntomique et localement d'intersection complète (Morphismes, lemmes \ref{morphisms-lemma-etale-syntomic} et \ref{morphisms-lemma-etale-flat}, et Compléments sur les morphismes, lemme \ref{more-morphisms-lemma-flat-lci}). D'après le lemme \ref{lemma-perfect-comparison-shriek}, il suffit de montrer que $f^!\mathcal{O}_Y = \mathcal{O}_X$. Le lemme \ref{lemma-lci-shriek} affirme que $f^!\mathcal{O}_Y$ est un module inversible. Considérons le diagramme commutatif
$$
\xymatrix{
X \times_Y X \ar[r]_{p_2} \ar[d]_{p_1} & X \ar[d]^f \\
X \ar[r]^f & Y
}
$$
et la diagonale $\Delta : X \to X \times_Y X$. Puisque $\Delta$ est une immersion ouverte (Morphismes, lemmes \ref{morphisms-lemma-diagonal-unramified-morphism} et \ref{morphisms-lemma-etale-smooth-unramified}), le lemme \ref{lemma-shriek-open-immersion} donne $\Delta^! = \Delta^*$. Le lemme \ref{lemma-upper-shriek-composition} donne $\Delta^! \circ p_1^! \circ f^! = f^!$. En appliquant le lemme \ref{lemma-base-change-shriek-flat} au diagramme, on obtient $p_1^!\mathcal{O}_X = p_2^*f^!\mathcal{O}_Y$. On en conclut que
$$
f^!\mathcal{O}_Y = \Delta^!p_1^!f^!\mathcal{O}_Y =
\Delta^*(p_1^*f^!\mathcal{O}_Y \otimes p_1^!\mathcal{O}_X) =
\Delta^*(p_2^*f^!\mathcal{O}_Y \otimes p_1^*f^!\mathcal{O}_Y) =
(f^!\mathcal{O}_Y)^{\otimes 2}
$$
où l'on utilise encore le lemme \ref{lemma-perfect-comparison-shriek} à la deuxième étape. Ainsi, $f^!\mathcal{O}_Y = \mathcal{O}_X$, comme voulu.
\end{proof}


\noindent
Dans la suite de cette section, nous énonçons quelques résultats faciles à démontrer qui découleraient d'une bonne théorie de l'application de changement de base.

\begin{lemma}[Changement de base de remplacement]
\label{lemma-base-change-locally}
Dans la situation \ref{situation-shriek}, soit
$$
\xymatrix{
X' \ar[r]_{g'} \ar[d]_{f'} & X \ar[d]^f \\
Y' \ar[r]^g & Y
}
$$
un diagramme cartésien de $\textit{FTS}_S$. Soit $E \in D^+_\QCoh(\mathcal{O}_Y)$ un objet tel que $Lg^*E$ appartienne à $D^+(\mathcal{O}_Y)$. Si $f$ est plat, alors $L(g')^*f^!E$ et $(f')^!Lg^*E$ se restreignent en des objets isomorphes de $D(\mathcal{O}_{U'})$ pour tout ouvert affine $U' \subset X'$ dont les images sont contenues dans des ouverts affines de $Y$, $Y'$ et $X$.
\end{lemma}

\begin{proof}
Nos hypothèses permettent de se ramener immédiatement au cas où $X$, $Y$, $Y'$ et $X'$ sont affines. Écrivons $Y = \Spec(R)$, $Y' = \Spec(R')$, $X = \Spec(A)$ et $X' = \Spec(A')$. Alors $A' = A \otimes_R R'$. Soit $E$ correspondant à $K \in D^+(R)$. En notant $\varphi : R \to A$ et $\varphi' : R' \to A'$ les applications données, la remarque \ref{remark-local-calculation-shriek} montre que $L(g')^*f^!E$ et $(f')^!Lg^*E$ correspondent à $\varphi^!(K) \otimes_A^\mathbf{L} A'$ et $(\varphi')^!(K \otimes_R^\mathbf{L} R')$, où $\varphi^!$ et $(\varphi')^!$ sont les foncteurs de Complexes dualisants, section \ref{dualizing-section-relative-dualizing-complex-algebraic}. Le résultat découle de Complexes dualisants, lemme \ref{dualizing-lemma-bc-flat}.
\end{proof}

\begin{lemma}
\label{lemma-relative-dualizing-fibres}
Dans la situation \ref{situation-shriek}, soit $f : X \to Y$ un morphisme de $\textit{FTS}_S$. Supposons $f$ plat. Posons $\omega_{X/Y}^\bullet = f^!\mathcal{O}_Y$ dans $D^b_{\textit{Coh}}(X)$. Soit $y \in Y$ et soit $h : X_y \to X$ la projection. Alors $Lh^*\omega_{X/Y}^\bullet$ est un complexe dualisant sur $X_y$.
\end{lemma}

\begin{proof}
Le complexe $\omega_{X/Y}^\bullet$ appartient à $D^b_{\textit{Coh}}$ d'après le lemme \ref{lemma-perfect-comparison-shriek}. Être dualisant est une propriété locale. D'après le lemme \ref{lemma-base-change-locally}, il suffit donc de montrer que $(X_y \to y)^!\mathcal{O}_y$ est un complexe dualisant sur $X_y$. Cela résulte du lemme \ref{lemma-shriek-dualizing}.
\end{proof}








\section{Une théorie de la dualité}
\label{section-duality}

\noindent
Nous explicitons dans cette section la théorie de la dualité que donnent les résultats très généraux précédents pour les schémas séparés de type fini sur un schéma de base noethérien fixé.

\medskip\noindent
Rappelons qu'un complexe dualisant sur un schéma noethérien $X$ est un objet de $D(\mathcal{O}_X)$ qui donne localement, sur les ouverts affines, un complexe dualisant sur les anneaux correspondants, voir la définition \ref{definition-dualizing-scheme}.

\medskip\noindent
Étant donné un schéma noethérien $S$, notons $\textit{FTS}_S$ la catégorie des schémas séparés de type fini sur $S$. Alors :
\begin{enumerate}
\item les foncteurs $f^!$ font de $D_\QCoh^+$ un pseudo-foncteur sur $\textit{FTS}_S$ ;
\item si $f : X \to Y$ est un morphisme propre de $\textit{FTS}_S$, alors $f^!$ est la restriction de l'adjoint à droite de $Rf_* : D_\QCoh(\mathcal{O}_X) \to D_\QCoh(\mathcal{O}_Y)$ à $D_\QCoh^+(\mathcal{O}_Y)$, et il existe un isomorphisme canonique
$$
Rf_*R\SheafHom_{\mathcal{O}_X}(K, f^!M)
\to
R\SheafHom_{\mathcal{O}_Y}(Rf_*K, M)
$$
pour tous $K \in D_{\textit{Coh}}^-(\mathcal{O}_X)$ et $M \in D_\QCoh^+(\mathcal{O}_Y)$ ;
\item si un objet $X$ de $\textit{FTS}_S$ admet un complexe dualisant $\omega_X^\bullet$, le foncteur
$D_X = R\SheafHom_{\mathcal{O}_X}(-, \omega_X^\bullet)$
définit une involution de $D_{\textit{Coh}}(\mathcal{O}_X)$ échangeant $D_{\textit{Coh}}^+(\mathcal{O}_X)$ et $D_{\textit{Coh}}^-(\mathcal{O}_X)$ et préservant $D_{\textit{Coh}}^b(\mathcal{O}_X)$ ;
\item si $f : X \to Y$ est un morphisme de $\textit{FTS}_S$ et si $\omega_Y^\bullet$ est un complexe dualisant sur $Y$, alors
\begin{enumerate}
\item $\omega_X^\bullet = f^!\omega_Y^\bullet$ est un complexe dualisant sur $X$ ;
\item $f^!M = D_X(Lf^*D_Y(M))$ canoniquement pour $M \in D_{\textit{Coh}}^+(\mathcal{O}_Y)$ ;
\item si de plus $f$ est propre, alors
$$
Rf_*R\SheafHom_{\mathcal{O}_X}(K, \omega_X^\bullet) =
R\SheafHom_{\mathcal{O}_Y}(Rf_*K, \omega_Y^\bullet)
$$
pour $K$ dans $D^-_{\textit{Coh}}(\mathcal{O}_X)$ ;
\end{enumerate}
\item si $f : X \to Y$ est une immersion fermée dans $\textit{FTS}_S$, alors
\par\noindent\hfil
$f^!(-) = R\SheafHom(\mathcal{O}_X, -)$ ;
\hfil\par
\item si $f : Y \to X$ est un morphisme fini dans $\textit{FTS}_S$, alors
\par\noindent\hfil
$f_*f^!(-) = R\SheafHom_{\mathcal{O}_X}(f_*\mathcal{O}_Y, -)$ ;
\hfil\par
\item si $f : X \to Y$ est l'inclusion d'un diviseur de Cartier effectif dans un objet de $\textit{FTS}_S$, alors $f^!(-) = Lf^*(-) \otimes_{\mathcal{O}_X} f^*\mathcal{O}_Y(X)[-1]$ ;
\item si $f : X \to Y$ est une immersion Koszul-régulière de codimension $c$ dans un objet de $\textit{FTS}_S$, alors $f^!(-) \cong Lf^*(-) \otimes_{\mathcal{O}_X} \wedge^c\mathcal{N}[-c]$ ;
\item si $f : X \to Y$ est un morphisme lisse et propre de dimension relative $d$ dans $\textit{FTS}_S$, alors $f^!(-) \cong Lf^*(-)  \otimes_{\mathcal{O}_X} \Omega^d_{X/Y}[d]$.
\end{enumerate}
Cela résulte des lemmes \ref{lemma-dualizing-schemes}, \ref{lemma-iso-on-RSheafHom}, \ref{lemma-twisted-inverse-image-closed}, \ref{lemma-finite-twisted}, \ref{lemma-sheaf-with-exact-support-effective-Cartier}, \ref{lemma-regular-immersion}, \ref{lemma-smooth-proper}, \ref{lemma-upper-shriek-composition}, \ref{lemma-pseudo-functor}, \ref{lemma-shriek-closed-immersion}, \ref{lemma-shriek-dualizing}, \ref{lemma-shriek-via-duality} et \ref{lemma-perfect-comparison-shriek}, ainsi que de l'exemple \ref{example-iso-on-RSheafHom-noetherian}. Nous avons obtenu nos foncteurs par un procédé très abstrait qui repose en dernier ressort sur un théorème d'existence (Catégories dérivées, proposition \ref{derived-proposition-brown}). Cela signifie que nous n'avons, en général, aucune description explicite des foncteurs $f^!$. Cela peut parfois poser problème. En fait, il suffit pourtant souvent de connaître l'existence d'un complexe dualisant et l'isomorphisme de dualité pour déterminer $f^!$.






\section{Recollement des complexes dualisants}
\label{section-glue}

\noindent
Nous allons maintenant utiliser le recollement des complexes dualisants pour obtenir une théorie valable pour tous les schémas de type fini sur $S$, étant donné un couple $(S, \omega_S^\bullet)$ comme dans la situation \ref{situation-dualizing}. Cette construction est analogue à \cite[remarque on page 310]{RD}.

\begin{situation}
\label{situation-dualizing}
Ici, $S$ est un schéma noethérien et $\omega_S^\bullet$ est un complexe dualisant.
\end{situation}

\noindent
Dans la situation \ref{situation-dualizing}, soit $X$ un schéma de type fini sur $S$. Soit $\mathcal{U} : X = \bigcup_{i = 1, \ldots, n} U_i$ un recouvrement ouvert fini de $X$ par des objets de $\textit{FTS}_S$, voir la situation \ref{situation-shriek}. Cela signifie simplement que les morphismes $U_i \to S$ sont séparés (ils sont déjà de type fini). Tout schéma affine de type fini sur $S$ est un objet de $\textit{FTS}_S$ d'après Schémas, lemme \ref{schemes-lemma-compose-after-separated} ; de tels recouvrements ouverts existent donc bien. Pour tout $i, j, k \in \{1, \ldots, n\}$, les morphismes $p_i : U_i \to S$, $p_{ij} : U_i \cap U_j \to S$ et $p_{ijk} : U_i \cap U_j \cap U_k \to S$ sont alors séparés, et chacun de ces schémas est un objet de $\textit{FTS}_S$. Un tel recouvrement fournit
\begin{enumerate}
\item $\omega_i^\bullet = p_i^!\omega_S^\bullet$, complexe dualisant sur $U_i$, voir la section \ref{section-duality} ;
\item pour tout $i, j$, un isomorphisme canonique $\varphi_{ij} :
\omega_i^\bullet|_{U_i \cap U_j} \to \omega_j^\bullet|_{U_i \cap U_j}$ ;
\item
\label{item-cocycle-glueing}
pour tout $i, j, k$, on a
$$
\varphi_{ik}|_{U_i \cap U_j \cap U_k} =
\varphi_{jk}|_{U_i \cap U_j \cap U_k} \circ
\varphi_{ij}|_{U_i \cap U_j \cap U_k}
$$
dans $D(\mathcal{O}_{U_i \cap U_j \cap U_k})$.
\end{enumerate}
Dans (2), on utilise le fait que $(U_i \cap U_j \to U_i)^!$ est donné par restriction (lemme \ref{lemma-shriek-open-immersion}), ainsi que les isomorphismes canoniques
$$
(U_i \cap U_j \to U_i)^! \circ p_i^! = p_{ij}^! =
(U_i \cap U_j \to U_j)^! \circ p_j^!
$$
du lemme \ref{lemma-upper-shriek-composition} ; pour obtenir (3), on utilise le fait que les foncteurs image inverse extraordinaire forment un pseudo-foncteur d'après le lemme \ref{lemma-pseudo-functor}.

\medskip\noindent
Dans la situation qui vient d'être décrite, un {\it complexe dualisant normalisé relativement à $\omega_S^\bullet$ et $\mathcal{U}$} est un couple $(K, \alpha_i)$ où $K \in D(\mathcal{O}_X)$ et où les $\alpha_i : K|_{U_i} \to \omega_i^\bullet$ sont des isomorphismes tels que $\varphi_{ij}$ soit donné par $\alpha_j|_{U_i \cap U_j} \circ \alpha_i^{-1}|_{U_i \cap U_j}$. Puisqu'être un complexe dualisant sur un schéma est une propriété locale, les complexes dualisants normalisés relativement à $\omega_S^\bullet$ et $\mathcal{U}$ sont bien des complexes dualisants.

\begin{lemma}
\label{lemma-good-dualizing-unique}
Dans la situation \ref{situation-dualizing}, soit $X$ un schéma de type fini sur $S$ et soit $\mathcal{U}$ un recouvrement ouvert fini de $X$ par des schémas séparés sur $S$. S'il existe un complexe dualisant normalisé relativement à $\omega_S^\bullet$ et $\mathcal{U}$, il est unique à unique isomorphisme près.
\end{lemma}

\begin{proof}
Si $(K, \alpha_i)$ et $(K', \alpha_i')$ sont deux tels complexes, considérons $L = R\SheafHom_{\mathcal{O}_X}(K, K')$. D'après le lemme \ref{lemma-dualizing-unique-schemes} et sa démonstration, c'est un objet inversible de $D(\mathcal{O}_X)$. À l'aide de $\alpha_i$ et $\alpha'_i$, on obtient un isomorphisme
$$
\alpha_i^t \otimes \alpha'_i :
L|_{U_i} \longrightarrow
R\SheafHom_{\mathcal{O}_X}(\omega_i^\bullet, \omega_i^\bullet) =
\mathcal{O}_{U_i}[0]
$$
Cela implique déjà que $L = H^0(L)[0]$ dans $D(\mathcal{O}_X)$. De plus, $H^0(L)$ est un faisceau inversible muni de trivialisations sur les ouverts $U_i$ de $X$. Enfin, le fait que $\alpha_j|_{U_i \cap U_j} \circ \alpha_i^{-1}|_{U_i \cap U_j}$ et $\alpha'_j|_{U_i \cap U_j} \circ (\alpha'_i)^{-1}|_{U_i \cap U_j}$ donnent tous deux $\varphi_{ij}$ implique que les applications de transition sont $1$, et l'on obtient un isomorphisme $H^0(L) = \mathcal{O}_X$.
\end{proof}

\begin{lemma}
\label{lemma-good-dualizing-independence-covering}
Dans la situation \ref{situation-dualizing}, soit $X$ un schéma de type fini sur $S$ et soient $\mathcal{U}$, $\mathcal{V}$ deux recouvrements ouverts finis de $X$ par des schémas séparés sur $S$. S'il existe un complexe dualisant normalisé relativement à $\omega_S^\bullet$ et $\mathcal{U}$, il en existe un normalisé relativement à $\omega_S^\bullet$ et $\mathcal{V}$, et ces complexes sont canoniquement isomorphes.
\end{lemma}

\begin{proof}
Il suffit de le démontrer lorsque $\mathcal{U}$ est donné par les ouverts $U_1, \ldots, U_n$ et $\mathcal{V}$ par les ouverts $U_1, \ldots, U_{n + m}$. On peut même supposer, et l'on suppose, $m = 1$. Pour passer d'un complexe dualisant $(K, \alpha_i)$ normalisé relativement à $\omega_S^\bullet$ et $\mathcal{V}$ à un complexe dualisant normalisé relativement à $\omega_S^\bullet$ et $\mathcal{U}$, il suffit d'oublier $\alpha_i$ pour $i = n + 1$. Réciproquement, soit $(K, \alpha_i)$ un complexe dualisant normalisé relativement à $\omega_S^\bullet$ et $\mathcal{U}$. Pour achever la démonstration, il faut construire une application $\alpha_{n + 1} : K|_{U_{n + 1}} \to \omega_{n + 1}^\bullet$ satisfaisant les conditions requises. Observons pour cela que $U_{n + 1} = \bigcup U_i \cap U_{n + 1}$ est un recouvrement ouvert. Il est clair que $(K|_{U_{n + 1}}, \alpha_i|_{U_i \cap U_{n + 1}})$ est un complexe dualisant normalisé relativement à $\omega_S^\bullet$ et au recouvrement $U_{n + 1} = \bigcup U_i \cap U_{n + 1}$. D'autre part, par la condition (\ref{item-cocycle-glueing}), le couple
$(\omega_{n + 1}^\bullet|_{U_{n + 1}}, \varphi_{n + 1i})$
est un autre complexe dualisant normalisé relativement à $\omega_S^\bullet$ et au recouvrement $U_{n + 1} = \bigcup U_i \cap U_{n + 1}$. Le lemme \ref{lemma-good-dualizing-unique} donne un unique isomorphisme
$$
\alpha_{n + 1} : K|_{U_{n + 1}} \longrightarrow \omega_{n + 1}^\bullet
$$
compatible avec les isomorphismes locaux donnés. Le lecteur pourra vérifier, à titre d'exercice, que cela signifie qu'il possède la propriété requise.
\end{proof}

\begin{lemma}
\label{lemma-existence-good-dualizing}
Dans la situation \ref{situation-dualizing}, soit $X$ un schéma de type fini sur $S$ et soit $\mathcal{U}$ un recouvrement ouvert fini de $X$ par des schémas séparés sur $S$. Il existe alors un complexe dualisant normalisé relativement à $\omega_S^\bullet$ et $\mathcal{U}$.
\end{lemma}

\begin{proof}
Écrivons $\mathcal{U} : X = \bigcup_{i = 1, \ldots, n} U_i$. Démontrons le lemme par récurrence sur $n$. Le cas initial $n = 1$ est immédiat. Supposons $n > 1$. Posons $X' = U_1 \cup \ldots \cup U_{n - 1}$ et soit $(K', \{\alpha'_i\}_{i = 1, \ldots, n - 1})$ un complexe dualisant normalisé relativement à $\omega_S^\bullet$ et $\mathcal{U}' : X' = \bigcup_{i = 1, \ldots, n - 1} U_i$. Il est clair que $(K'|_{X' \cap U_n}, \alpha'_i|_{U_i \cap U_n})$ est un complexe dualisant normalisé relativement à $\omega_S^\bullet$ et au recouvrement $X' \cap U_n = \bigcup_{i = 1, \ldots, n - 1} U_i \cap U_n$. D'autre part, par la condition (\ref{item-cocycle-glueing}), le couple
$(\omega_n^\bullet|_{X' \cap U_n}, \varphi_{ni})$
est un autre complexe dualisant normalisé relativement à $\omega_S^\bullet$ et au recouvrement $X' \cap U_n = \bigcup_{i = 1, \ldots, n - 1} U_i \cap U_n$. Le lemme \ref{lemma-good-dualizing-unique} donne un unique isomorphisme
$$
\epsilon : K'|_{X' \cap U_n} \longrightarrow \omega_i^\bullet|_{X' \cap U_n}
$$
compatible avec les isomorphismes locaux donnés. D'après Cohomologie, lemme \ref{cohomology-lemma-glue}, on obtient $K \in D(\mathcal{O}_X)$ muni d'isomorphismes $\beta : K|_{X'} \to K'$ et $\gamma : K|_{U_n} \to \omega_n^\bullet$ tels que $\epsilon = \gamma|_{X'\cap U_n} \circ \beta|_{X' \cap U_n}^{-1}$. On définit alors
$$
\alpha_i = \alpha'_i \circ \beta|_{U_i}, i = 1, \ldots, n - 1,
\text{ et }
\alpha_n = \gamma
$$
Il reste à vérifier que $\varphi_{ij}$ est donné par $\alpha_j|_{U_i \cap U_j} \circ \alpha_i^{-1}|_{U_i \cap U_j}$. Pour $i, j \leq n - 1$, cela résulte de la condition correspondante pour $\alpha_i'$. Pour $i = j = n$, c'est également clair. Si $i < j = n$, on obtient
$$
\alpha_n|_{U_i \cap U_n} \circ \alpha_i^{-1}|_{U_i \cap U_n} =
\gamma|_{U_i \cap U_n} \circ \beta^{-1}|_{U_i \cap U_n}
\circ (\alpha'_i)^{-1}|_{U_i \cap U_n} =
\epsilon|_{U_i \cap U_n} \circ (\alpha'_i)^{-1}|_{U_i \cap U_n}
$$
Cela est égal à $\alpha_{in}$ précisément parce que $\epsilon$ est l'unique application compatible avec $\alpha_i'$ et $\alpha_{ni}$.
\end{proof}

\noindent
Soit $(S, \omega_S^\bullet)$ comme dans la situation \ref{situation-dualizing}. Les lemmes précédents montrent que, pour tout schéma $X$ de type fini sur $S$, il existe un couple $(K, \alpha_U)$, donné à unique isomorphisme près, constitué d'un objet $K \in D(\mathcal{O}_X)$ et d'isomorphismes $\alpha_U : K|_U \to \omega_U^\bullet$ pour tout sous-schéma ouvert $U \subset X$ séparé sur $S$. Ici, $\omega_U^\bullet = (U \to S)^!\omega_S^\bullet$ est un complexe dualisant sur $U$, voir la section \ref{section-duality}. De plus, si $\mathcal{U} : X = \bigcup U_i$ est un recouvrement ouvert fini par des ouverts séparés sur $S$, alors $(K, \alpha_{U_i})$ est un complexe dualisant normalisé relativement à $\omega_S^\bullet$ et $\mathcal{U}$. En effet, l'unicité à unique isomorphisme près résulte du lemme \ref{lemma-good-dualizing-unique}, l'existence pour un recouvrement ouvert du lemme \ref{lemma-existence-good-dualizing}, et le fait que $K$ convienne alors pour tous les recouvrements ouverts du lemme \ref{lemma-good-dualizing-independence-covering}.

\begin{definition}
\label{definition-good-dualizing}
Soit $S$ un schéma noethérien et soit $\omega_S^\bullet$ un complexe dualisant sur $S$. Soit $X$ un schéma de type fini sur $S$. Le complexe $K$ construit ci-dessus est appelé {\it complexe dualisant normalisé relativement à $\omega_S^\bullet$} et est noté $\omega_X^\bullet$.
\end{definition}

\noindent
Comme le suggère la terminologie, un complexe dualisant normalisé relativement à $\omega_S^\bullet$ n'est pas seulement un objet de la catégorie dérivée de $X$ : il est muni des isomorphismes locaux décrits ci-dessus. Cela reste compatible avec la définition $\omega_X^\bullet = p^!\omega_S^\bullet$, où $p : X \to S$ est le morphisme structural, si $X$ est séparé sur $S$. Plus généralement, on dispose de la vérification de cohérence suivante.

\begin{lemma}
\label{lemma-good-over-both}
Soit $(S, \omega_S^\bullet)$ comme dans la situation \ref{situation-dualizing}. Soit $f : X \to Y$ un morphisme de schémas de type fini sur $S$. Soient $\omega_X^\bullet$ et $\omega_Y^\bullet$ des complexes dualisants normalisés relativement à $\omega_S^\bullet$. Alors $\omega_X^\bullet$ est un complexe dualisant normalisé relativement à $\omega_Y^\bullet$.
\end{lemma}

\begin{proof}
Il s'agit simplement de vérifier les compatibilités. Choisissons un recouvrement ouvert affine fini $\mathcal{V} : Y = \bigcup V_j$. Pour tout $j$, choisissons un recouvrement ouvert affine fini $f^{-1}(V_j) = U_{ji}$. Posons $\mathcal{U} : X = \bigcup U_{ji}$. Les schémas $V_j$ et $U_{ji}$ sont séparés sur $S$ ; on dispose donc des foncteurs image inverse extraordinaire pour $q_j : V_j \to S$, $p_{ji} : U_{ji} \to S$, $f_{ji} : U_{ji} \to V_j$ et $f_{ji}' : U_{ji} \to Y$. Soit $(L, \beta_j)$ un complexe dualisant normalisé relativement à $\omega_S^\bullet$ et $\mathcal{V}$. Soit $(K, \gamma_{ji})$ un complexe dualisant normalisé relativement à $\omega_S^\bullet$ et $\mathcal{U}$. (Autrement dit, $L = \omega_Y^\bullet$ et $K = \omega_X^\bullet$.) On peut définir
$$
\alpha_{ji} :
K|_{U_{ji}} \xrightarrow{\gamma_{ji}}
p_{ji}^!\omega_S^\bullet = f_{ji}^!q_j^!\omega_S^\bullet
\xrightarrow{f_{ji}^!\beta_j^{-1}} f_{ji}^!(L|_{V_j}) =
(f_{ji}')^!(L)
$$
Pour achever la démonstration, il faut montrer que
$\alpha_{ji}|_{U_{ji} \cap U_{j'i'}}
\circ \alpha_{j'i'}^{-1}|_{U_{ji} \cap U_{j'i'}}$
est l'isomorphisme canonique $(f_{ji}')^!(L)|_{U_{ji} \cap U_{j'i'}} \to
(f_{j'i'}')^!(L)|_{U_{ji} \cap U_{j'i'}}$. C'est formel ; nous omettons les détails.
\end{proof}

\begin{lemma}
\label{lemma-open-immersion-good-dualizing-complex}
Soit $(S, \omega_S^\bullet)$ comme dans la situation \ref{situation-dualizing}. Soit $j : X \to Y$ une immersion ouverte de schémas de type fini sur $S$. Soient $\omega_X^\bullet$ et $\omega_Y^\bullet$ des complexes dualisants normalisés relativement à $\omega_S^\bullet$. Il existe alors un isomorphisme canonique $\omega_X^\bullet = \omega_Y^\bullet|_X$.
\end{lemma}

\begin{proof}
Cela résulte immédiatement de la construction des complexes dualisants normalisés donnée juste avant la définition \ref{definition-good-dualizing}.
\end{proof}

\begin{lemma}
\label{lemma-proper-map-good-dualizing-complex}
Soit $(S, \omega_S^\bullet)$ comme dans la situation \ref{situation-dualizing}. Soit $f : X \to Y$ un morphisme propre de schémas de type fini sur $S$. Soient $\omega_X^\bullet$ et $\omega_Y^\bullet$ des complexes dualisants normalisés relativement à $\omega_S^\bullet$. Soit $a$ l'adjoint à droite du lemme \ref{lemma-twisted-inverse-image} pour $f$. Il existe alors un isomorphisme canonique $a(\omega_Y^\bullet) = \omega_X^\bullet$.
\end{lemma}

\begin{proof}
Soient $p : X \to S$ et $q : Y \to S$ les morphismes structuraux. Si $X$ et $Y$ sont séparés sur $S$, cela résulte de $\omega_X^\bullet = p^!\omega_S^\bullet$, $\omega_Y^\bullet = q^!\omega_S^\bullet$, $f^! = a$ et $f^! \circ q^! = p^!$ (lemme \ref{lemma-upper-shriek-composition}). Dans le cas général, on utilise d'abord le lemme \ref{lemma-good-over-both} pour se ramener au cas $Y = S$. Alors $X$ et $Y$ sont séparés sur $S$, et le résultat vient d'être établi.
\end{proof}

\noindent
Soit $(S, \omega_S^\bullet)$ comme dans la situation \ref{situation-dualizing}. Pour un schéma $X$ de type fini sur $S$, notons $\omega_X^\bullet$ le complexe dualisant sur $X$ normalisé relativement à $\omega_S^\bullet$. Définissons $D_X(-) = R\SheafHom_{\mathcal{O}_X}(-, \omega_X^\bullet)$ comme au lemme \ref{lemma-dualizing-schemes}. Soit $f : X \to Y$ un morphisme de schémas de type fini sur $S$. Définissons
$$
f_{new}^! = D_X \circ Lf^* \circ D_Y :
D_{\textit{Coh}}^+(\mathcal{O}_Y)
\to
D_{\textit{Coh}}^+(\mathcal{O}_X)
$$
Si $f : X \to Y$ et $g : Y \to Z$ sont des morphismes composables entre schémas de type fini sur $S$, définissons
\begin{align*}
(g \circ f)^!_{new} & = D_X \circ L(g \circ f)^* \circ D_Z \\
& = D_X \circ Lf^* \circ Lg^* \circ D_Z \\
& \to D_X \circ Lf^* \circ D_Y \circ D_Y \circ Lg^* \circ D_Z \\
& = f^!_{new} \circ g^!_{new}
\end{align*}
où la flèche est celle du lemme \ref{lemma-dualizing-schemes}. Le lemme suivant rassemble les résultats obtenus.

\begin{lemma}
\label{lemma-duality-bootstrap}
Soit $(S, \omega_S^\bullet)$ comme dans la situation \ref{situation-dualizing}. Avec $f^!_{new}$ et $\omega_X^\bullet$ définis comme ci-dessus pour tous les schémas (et morphismes de schémas) de type fini sur $S$ :
\begin{enumerate}
\item les foncteurs $f^!_{new}$ et les flèches $(g \circ f)^!_{new} \to f^!_{new} \circ g^!_{new}$ font de $D_{\textit{Coh}}^+$ un pseudo-foncteur de la catégorie des schémas de type fini sur $S$ vers la $2$-catégorie des catégories ;
\item $\omega_X^\bullet = (X \to S)^!_{new} \omega_S^\bullet$ ;
\item le foncteur $D_X$ définit une involution de $D_{\textit{Coh}}(\mathcal{O}_X)$ échangeant $D_{\textit{Coh}}^+(\mathcal{O}_X)$ et $D_{\textit{Coh}}^-(\mathcal{O}_X)$ et préservant $D_{\textit{Coh}}^b(\mathcal{O}_X)$ ;
\item $\omega_X^\bullet = f^!_{new}\omega_Y^\bullet$ pour $f : X \to Y$ un morphisme de schémas de type fini sur $S$ ;
\item $f^!_{new}M = D_X(Lf^*D_Y(M))$ pour $M \in D_{\textit{Coh}}^+(\mathcal{O}_Y)$ ;
\item si de plus $f$ est propre, alors $f^!_{new}$ est isomorphe à la restriction de l'adjoint à droite de $Rf_* : D_\QCoh(\mathcal{O}_X) \to D_\QCoh(\mathcal{O}_Y)$ à $D_{\textit{Coh}}^+(\mathcal{O}_Y)$, et il existe un isomorphisme canonique
$$
Rf_*R\SheafHom_{\mathcal{O}_X}(K, f^!_{new}M)
\to
R\SheafHom_{\mathcal{O}_Y}(Rf_*K, M)
$$
pour $K \in D^-_{\textit{Coh}}(\mathcal{O}_X)$ et $M \in D_{\textit{Coh}}^+(\mathcal{O}_Y)$, et
$$
Rf_*R\SheafHom_{\mathcal{O}_X}(K, \omega_X^\bullet) =
R\SheafHom_{\mathcal{O}_Y}(Rf_*K, \omega_Y^\bullet)
$$
pour $K \in D^-_{\textit{Coh}}(\mathcal{O}_X)$, et
\end{enumerate}
\enlargethispage\smallskipamount
Si $X$ est séparé sur $S$, alors $\omega_X^\bullet$ est canoniquement isomorphe à $(X \to S)^!\omega_S^\bullet$ ; si $f$ est un morphisme entre schémas séparés sur $S$, il existe un isomorphisme canonique\footnote{Nous n'avons pas vérifié leur compatibilité avec les isomorphismes $(g \circ f)^! \to f^! \circ g^!$ et $(g \circ f)^!_{new} \to f^!_{new} \circ g^!_{new}$. Nous le ferons ici si nous en avons besoin ultérieurement.} $f_{new}^!K = f^!K$ pour $K$ dans $D_{\textit{Coh}}^+$.
\end{lemma}

\begin{proof}
Soient $f : X \to Y$, $g : Y \to Z$, $h : Z \to T$ des morphismes de schémas de type fini sur $S$. Il faut montrer que
$$
\xymatrix{
(h \circ g \circ f)^!_{new} \ar[r] \ar[d] &
f^!_{new} \circ (h \circ g)^!_{new} \ar[d] \\
(g \circ f)^!_{new} \circ h^!_{new} \ar[r] &
f^!_{new} \circ g^!_{new} \circ h^!_{new}
}
$$
est commutatif. Soient $\eta_Y : \text{id} \to D_Y^2$ et $\eta_Z : \text{id} \to D_Z^2$ les isomorphismes canoniques du lemme \ref{lemma-dualizing-schemes}. En utilisant Catégories, lemme \ref{categories-lemma-properties-2-cat-cats}, un calcul (omis) montre que les deux flèches
$(h \circ g \circ f)^!_{new} \to f^!_{new} \circ g^!_{new} \circ h^!_{new}$
sont données par
$$
1 \star \eta_Y \star 1 \star \eta_Z \star 1 :
D_X \circ Lf^* \circ Lg^* \circ Lh^* \circ D_T
\longrightarrow
D_X \circ Lf^* \circ D_Y^2 \circ Lg^* \circ D_Z^2 \circ Lh^* \circ D_T
$$
Cela démontre (1). L'assertion (2) résulte immédiatement de la définition de $(X \to S)^!_{new}$ et du fait que $D_S(\omega_S^\bullet) = \mathcal{O}_S$. L'assertion (3) est le lemme \ref{lemma-dualizing-schemes}. L'assertion (4) se démontre comme (2). L'assertion (5) est la définition de $f^!_{new}$.

\medskip\noindent
Démonstration de (6). Soit $a$ l'adjoint à droite du lemme \ref{lemma-twisted-inverse-image} pour le morphisme propre $f : X \to Y$ de schémas de type fini sur $S$. La difficulté est que l'on ne sait pas si $X$ ou $Y$ est séparé sur $S$ (et ce sera faux en général) ; on ne peut donc pas appliquer directement le lemme \ref{lemma-shriek-via-duality} à $f$ sur $S$. Pour contourner cette difficulté, on utilise l'identification canonique $\omega_X^\bullet = a(\omega_Y^\bullet)$ du lemme \ref{lemma-proper-map-good-dualizing-complex}. Ainsi, $f^!_{new}$ est la restriction de $a$ à $D_{\textit{Coh}}^+(\mathcal{O}_Y)$ d'après le lemme \ref{lemma-shriek-via-duality} appliqué à $f : X \to Y$ sur le schéma de base $Y$ ! Les égalités affichées résultent de l'exemple \ref{example-iso-on-RSheafHom-noetherian}.

\medskip\noindent
Les dernières assertions résultent de la construction des complexes dualisants normalisés et du lemme \ref{lemma-shriek-via-duality}, déjà utilisé.
\end{proof}

\begin{remark}
\label{remark-independent-omega-S}
Soit $S$ un schéma noethérien admettant un complexe dualisant. Soit $f : X \to Y$ un morphisme de schémas de type fini sur $S$. Alors le foncteur
$$
f_{new}^! : D^+_{Coh}(\mathcal{O}_Y) \to D^+_{Coh}(\mathcal{O}_X)
$$
est indépendant, à isomorphisme canonique près, du choix du complexe dualisant $\omega_S^\bullet$. Esquissons la démonstration. Tout second complexe dualisant est de la forme $\omega_S^\bullet \otimes_{\mathcal{O}_S}^\mathbf{L} \mathcal{L}$, où $\mathcal{L}$ est un objet inversible de $D(\mathcal{O}_S)$, voir le lemme \ref{lemma-dualizing-unique-schemes}. Pour tout morphisme séparé $p : U \to S$ de type fini, on a
$p^!(\omega_S^\bullet \otimes^\mathbf{L}_{\mathcal{O}_S} \mathcal{L}) =
p^!(\omega_S^\bullet) \otimes^\mathbf{L}_{\mathcal{O}_U} Lp^*\mathcal{L}$
d'après le lemme \ref{lemma-compare-with-pullback-perfect}. Si $\omega_X^\bullet$ et $\omega_Y^\bullet$ sont les complexes dualisants normalisés relativement à $\omega_S^\bullet$, alors $\omega_X^\bullet \otimes_{\mathcal{O}_X}^\mathbf{L} La^*\mathcal{L}$ et $\omega_Y^\bullet \otimes_{\mathcal{O}_Y}^\mathbf{L} Lb^*\mathcal{L}$ sont donc les complexes dualisants normalisés relativement à $\omega_S^\bullet \otimes_{\mathcal{O}_S}^\mathbf{L} \mathcal{L}$ (où $a : X \to S$ et $b : Y \to S$ sont les morphismes structuraux). Le résultat s'en déduit, car
\begin{align*}
& R\SheafHom_{\mathcal{O}_X}(Lf^*R\SheafHom_{\mathcal{O}_Y}(K,
\omega_Y^\bullet \otimes_{\mathcal{O}_Y}^\mathbf{L} Lb^*\mathcal{L}),
\omega_X^\bullet \otimes_{\mathcal{O}_X}^\mathbf{L} La^*\mathcal{L}) \\
& = R\SheafHom_{\mathcal{O}_X}(Lf^*R(\SheafHom_{\mathcal{O}_Y}(K,
\omega_Y^\bullet) \otimes_{\mathcal{O}_Y}^\mathbf{L} Lb^*\mathcal{L}),
\omega_X^\bullet \otimes_{\mathcal{O}_X}^\mathbf{L} La^*\mathcal{L}) \\
& = R\SheafHom_{\mathcal{O}_X}(Lf^*R\SheafHom_{\mathcal{O}_Y}(K,
\omega_Y^\bullet) \otimes_{\mathcal{O}_X}^\mathbf{L} La^*\mathcal{L},
\omega_X^\bullet \otimes_{\mathcal{O}_X}^\mathbf{L} La^*\mathcal{L}) \\
& = R\SheafHom_{\mathcal{O}_X}(Lf^*R\SheafHom_{\mathcal{O}_Y}(K,
\omega_Y^\bullet), \omega_X^\bullet)
\end{align*}
pour $K \in D^+_{Coh}(\mathcal{O}_Y)$. La dernière égalité résulte de l'inversibilité de $La^*\mathcal{L}$ dans $D(\mathcal{O}_X)$.
\end{remark}


\begin{example}
\label{example-trace-proper}
Soit $S$ un schéma noethérien et soit $\omega_S^\bullet$ un complexe dualisant. Soit $f : X \to Y$ un morphisme propre de schémas de type fini sur $S$. Soient $\omega_X^\bullet$ et $\omega_Y^\bullet$ des complexes dualisants normalisés relativement à $\omega_S^\bullet$. Dans cette situation, on a $a(\omega_Y^\bullet) = \omega_X^\bullet$ (lemme \ref{lemma-proper-map-good-dualizing-complex}) ; le morphisme trace (section \ref{section-trace}) est donc une flèche canonique
$$
\text{Tr}_f : Rf_*\omega_X^\bullet \longrightarrow \omega_Y^\bullet
$$
qui donne les isomorphismes (lemme \ref{lemma-duality-bootstrap})
$$
\Hom_X(L, \omega_X^\bullet) = \Hom_Y(Rf_*L, \omega_Y^\bullet)
$$
et
$$
Rf_*R\SheafHom_{\mathcal{O}_X}(L, \omega_X^\bullet) =
R\SheafHom_{\mathcal{O}_Y}(Rf_*L, \omega_Y^\bullet)
$$
pour $L$ dans $D_\QCoh(\mathcal{O}_X)$.
\end{example}

\begin{remark}
\label{remark-dualizing-finite}
Soit $S$ un schéma noethérien et soit $\omega_S^\bullet$ un complexe dualisant. Soit $f : X \to Y$ un morphisme fini entre schémas de type fini sur $S$. Soient $\omega_X^\bullet$ et $\omega_Y^\bullet$ des complexes dualisants normalisés relativement à $\omega_S^\bullet$. On a alors
$$
f_*\omega_X^\bullet = R\SheafHom(f_*\mathcal{O}_X, \omega_Y^\bullet)
$$
dans $D_\QCoh^+(f_*\mathcal{O}_X)$ d'après les lemmes \ref{lemma-finite-twisted} et \ref{lemma-proper-map-good-dualizing-complex}, et le morphisme trace de l'exemple \ref{example-trace-proper} est le morphisme
$$
\text{Tr}_f : Rf_*\omega_X^\bullet = f_*\omega_X^\bullet =
R\SheafHom(f_*\mathcal{O}_X, \omega_Y^\bullet) \longrightarrow
\omega_Y^\bullet
$$
souvent appelée « évaluation en $1$ ».
\end{remark}

\begin{remark}
\label{remark-relative-dualizing-complex-shriek}
Soit $f : X \to Y$ un morphisme plat et propre de schémas de type fini sur un couple $(S, \omega_S^\bullet)$ comme dans la situation \ref{situation-dualizing}. Le complexe dualisant relatif (remarque \ref{remark-relative-dualizing-complex}) est $\omega_{X/Y}^\bullet = a(\mathcal{O}_Y)$. Le lemme \ref{lemma-proper-map-good-dualizing-complex} donne le premier isomorphisme canonique de
$$
\omega_X^\bullet = a(\omega_Y^\bullet) =
Lf^*\omega_Y^\bullet \otimes_{\mathcal{O}_X}^\mathbf{L} \omega_{X/Y}^\bullet
$$
dans $D(\mathcal{O}_X)$. Le second isomorphisme canonique résulte de la discussion de la remarque \ref{remark-relative-dualizing-complex}.
\end{remark}




\section{Fonctions de dimension}
\label{section-dimension-functions}

\noindent
Nous avons besoin de quelques précisions sur la variation des fonctions de dimension lorsqu'on passe à un schéma de type fini sur un autre.

\begin{lemma}
\label{lemma-good-dualizing-normalized}
Soit $S$ un schéma noethérien et soit $\omega_S^\bullet$ un complexe dualisant. Soit $X$ un schéma de type fini sur $S$ et soit $\omega_X^\bullet$ le complexe dualisant normalisé relativement à $\omega_S^\bullet$. Si $x \in X$ est un point fermé au-dessus d'un point fermé $s$ de $S$, alors $\omega_{X, x}^\bullet$ est un complexe dualisant normalisé sur $\mathcal{O}_{X, x}$, pourvu que $\omega_{S, s}^\bullet$ soit un complexe dualisant normalisé sur $\mathcal{O}_{S, s}$.
\end{lemma}

\begin{proof}
On peut remplacer $X$ par un voisinage affine de $x$ ; on peut donc supposer, et l'on suppose, $f : X \to S$ séparé. Alors $\omega_X^\bullet = f^!\omega_S^\bullet$. Il faut montrer que
$R\Hom_{\mathcal{O}_{X, x}}(\kappa(x), \omega_{X, x}^\bullet)$
est concentré en degré $0$. Notons $i_x : x \to X$ le morphisme d'inclusion, qui est une immersion fermée puisque $x$ est un point fermé. Ainsi, $R\Hom_{\mathcal{O}_{X, x}}(\kappa(x), \omega_{X, x}^\bullet)$ représente $i_x^!\omega_X^\bullet$ d'après le lemme \ref{lemma-shriek-closed-immersion}. Considérons le diagramme commutatif
$$
\xymatrix{
x \ar[r]_{i_x} \ar[d]_\pi & X \ar[d]^f \\
s \ar[r]^{i_s} & S
}
$$
D'après Morphismes, lemme \ref{morphisms-lemma-closed-point-fibre-locally-finite-type}, l'extension $\kappa(x)/\kappa(s)$ est finie ; ainsi, $\pi$ est un morphisme fini. On en conclut que
$$
i_x^!\omega_X^\bullet = i_x^! f^! \omega_S^\bullet =
\pi^! i_s^! \omega_S^\bullet
$$
Si $\omega_{S, s}^\bullet$ est un complexe dualisant normalisé sur $\mathcal{O}_{S, s}$, le même raisonnement donne donc $i_s^!\omega_S^\bullet = \kappa(s)[0]$. On a
$$
R\pi_*(\pi^!(\kappa(s)[0])) =
R\SheafHom_{\mathcal{O}_s}(R\pi_*(\kappa(x)[0]), \kappa(s)[0]) =
\widetilde{\Hom_{\kappa(s)}(\kappa(x), \kappa(s))}
$$
La première égalité résulte de l'exemple \ref{example-iso-on-RSheafHom-noetherian} appliqué avec $L = \kappa(x)[0]$. La seconde résulte de l'exactitude de $\pi_*$. Ainsi, $\pi^!(\kappa(s)[0])$ est concentré en degré $0$, ce qui conclut.
\end{proof}

\begin{lemma}
\label{lemma-good-dualizing-dimension-function}
Soit $S$ un schéma noethérien et soit $\omega_S^\bullet$ un complexe dualisant. Soit $f : X \to S$ de type fini et soit $\omega_X^\bullet$ le complexe dualisant normalisé relativement à $\omega_S^\bullet$. Pour tout $x \in X$, on a
$$
\delta_X(x) - \delta_S(f(x)) = \text{trdeg}_{\kappa(f(x))}(\kappa(x))
$$
où $\delta_S$, resp.\ $\delta_X$, est la fonction de dimension de $\omega_S^\bullet$, resp.\ $\omega_X^\bullet$, voir le lemme \ref{lemma-dimension-function-scheme}.
\end{lemma}

\begin{proof}
On peut remplacer $X$ par un voisinage affine de $x$. On peut donc supposer, et l'on suppose, qu'il existe une compactification $X \subset \overline{X}$ sur $S$. On peut alors remplacer $X$ par $\overline{X}$ et supposer $X$ propre sur $S$. On peut aussi supposer $X$ connexe en remplaçant $X$ par la composante connexe de $X$ contenant $x$. Rappelons ensuite que $\delta_X$ et la fonction
$x \mapsto \delta_S(f(x)) + \text{trdeg}_{\kappa(f(x))}(\kappa(x))$
sont toutes deux des fonctions de dimension sur $X$, voir Morphismes, lemme \ref{morphisms-lemma-dimension-function-propagates} (et le fait que $S$ est universellement caténaire d'après le lemme \ref{lemma-dimension-function-scheme}). D'après Topologie, lemme \ref{topology-lemma-dimension-function-unique}, leur différence est localement constante, donc constante puisque $X$ est connexe. Il suffit donc de montrer l'égalité en un point quelconque de $X$. D'après Propriétés, lemme \ref{properties-lemma-locally-Noetherian-closed-point}, le schéma $X$ possède un point fermé $x$. Puisque $X \to S$ est propre, l'image $s$ de $x$ est fermée dans $S$. On peut donc appliquer le lemme \ref{lemma-good-dualizing-normalized} pour conclure.
\end{proof}

\begin{lemma}
\label{lemma-shriek}
Dans la situation \ref{situation-shriek}, soit $f : X \to Y$ un morphisme de $\textit{FTS}_S$. Soit $x \in X$ d'image $y \in Y$. Alors
$$
H^i(f^!\mathcal{O}_Y)_x \not = 0
\Rightarrow - \dim_x(X_y) \leq i.
$$
\end{lemma}

\begin{proof}
L'énoncé étant local sur $X$, on peut supposer que $X$ et $Y$ sont affines. Écrivons $X = \Spec(A)$ et $Y = \Spec(R)$. Alors $f^!\mathcal{O}_Y$ correspond au complexe dualisant relatif $\omega_{A/R}^\bullet$ de Complexes dualisants, section \ref{dualizing-section-relative-dualizing-complexes-Noetherian}, d'après la remarque \ref{remark-local-calculation-shriek}. Le lemme résulte donc de Complexes dualisants, lemme \ref{dualizing-lemma-relative-dualizing-trivial-vanishing}.
\end{proof}

\begin{lemma}
\label{lemma-flat-shriek}
Dans la situation \ref{situation-shriek}, soit $f : X \to Y$ un morphisme de $\textit{FTS}_S$. Soit $x \in X$ d'image $y \in Y$. Si $f$ est plat, alors
$$
H^i(f^!\mathcal{O}_Y)_x \not = 0
\Rightarrow - \dim_x(X_y) \leq i \leq 0.
$$
En fait, si toutes les fibres de $f$ sont de dimension $\leq d$, alors $f^!\mathcal{O}_Y$ est de Tor-amplitude dans $[-d, 0]$ comme objet de $D(X, f^{-1}\mathcal{O}_Y)$.
\end{lemma}

\begin{proof}
En raisonnant exactement comme dans la démonstration du lemme \ref{lemma-shriek}, cela résulte de Complexes dualisants, lemme \ref{dualizing-lemma-relative-dualizing-flat-vanishing}.
\end{proof}

\begin{lemma}
\label{lemma-shriek-over-CM}
Dans la situation \ref{situation-shriek}, soit $f : X \to Y$ un morphisme de $\textit{FTS}_S$. Soit $x \in X$ d'image $y \in Y$. Supposons que
\begin{enumerate}
\item $\mathcal{O}_{Y, y}$ soit de Cohen–Macaulay ;
\item $\text{trdeg}_{\kappa(f(\xi))}(\kappa(\xi)) \leq r$ pour tout point générique $\xi$ d'une composante irréductible de $X$ contenant $x$.
\end{enumerate}
Alors
$$
H^i(f^!\mathcal{O}_Y)_x \not = 0
\Rightarrow - r \leq i
$$
et la fibre $H^{-r}(f^!\mathcal{O}_Y)_x$ est $(S_2)$ comme $\mathcal{O}_{X, x}$-module.
\end{lemma}

\begin{proof}
Après avoir remplacé $X$ par un voisinage ouvert de $x$, on peut supposer que toute composante irréductible de $X$ passe par $x$. En raisonnant exactement comme dans la démonstration du lemme \ref{lemma-shriek}, cela résulte alors de Complexes dualisants, lemme \ref{dualizing-lemma-relative-dualizing-CM-vanishing}.
\end{proof}

\begin{lemma}
\label{lemma-flat-quasi-finite-shriek}
Dans la situation \ref{situation-shriek}, soit $f : X \to Y$ un morphisme de $\textit{FTS}_S$. Si $f$ est plat et quasi-fini, alors
$$
f^!\mathcal{O}_Y = \omega_{X/Y}[0]
$$
pour un $\mathcal{O}_X$-module cohérent $\omega_{X/Y}$ plat sur $Y$.
\end{lemma}

\begin{proof}
Cela résulte du lemme \ref{lemma-flat-shriek} et du fait que les faisceaux de cohomologie de $f^!\mathcal{O}_Y$ sont cohérents d'après le lemme \ref{lemma-shriek-coherent}.
\end{proof}

\begin{lemma}
\label{lemma-CM-shriek}
Dans la situation \ref{situation-shriek}, soit $f : X \to Y$ un morphisme de $\textit{FTS}_S$. Si $f$ est de Cohen–Macaulay (Compléments sur les morphismes, définition \ref{more-morphisms-definition-CM}), alors
$$
f^!\mathcal{O}_Y = \omega_{X/Y}[d]
$$
pour un $\mathcal{O}_X$-module cohérent $\omega_{X/Y}$ plat sur $Y$, où $d$ est la fonction localement constante sur $X$ donnant la dimension relative de $X$ sur $Y$.
\end{lemma}

\begin{proof}
La dimension relative $d$ est bien définie et localement constante d'après Morphismes, lemme \ref{morphisms-lemma-flat-finite-presentation-CM-fibres-relative-dimension}. Les faisceaux de cohomologie de $f^!\mathcal{O}_Y$ sont cohérents d'après le lemme \ref{lemma-shriek-coherent}. On obtiendra la platitude de $\omega_{X/Y}$ par le lemme \ref{lemma-flat-shriek} si l'on montre que les autres faisceaux de cohomologie de $f^!\mathcal{O}_Y$ sont nuls.

\medskip\noindent
La question est locale sur $X$ ; on peut donc supposer $X$ et $Y$ affines et le morphisme de dimension relative $d$. Si $d = 0$, le résultat découle directement du lemme \ref{lemma-flat-quasi-finite-shriek}. Si $d > 0$, on peut supposer qu'il existe une factorisation
$$
X \xrightarrow{g} \mathbf{A}^d_Y \xrightarrow{p} Y
$$
avec $g$ quasi-fini et plat, voir Compléments sur les morphismes, lemme \ref{more-morphisms-lemma-flat-finite-presentation-characterize-CM}. Alors $f^! = g^! \circ p^!$. Le lemme \ref{lemma-shriek-affine-line} montre que $p^!\mathcal{O}_Y \cong \mathcal{O}_{\mathbf{A}^d_Y}[-d]$. On conclut par le cas $d = 0$.
\end{proof}

\begin{remark}
\label{remark-the-same-is-true}
Soit $S$ un schéma noethérien muni d'un complexe dualisant $\omega_S^\bullet$. Dans ce cas, les lemmes \ref{lemma-shriek}, \ref{lemma-flat-shriek}, \ref{lemma-flat-quasi-finite-shriek} et \ref{lemma-CM-shriek} sont vrais pour tout morphisme $f : X \to Y$ de schémas de type fini sur $S$, en remplaçant toutefois $f^!$ par $f_{new}^!$. C'est clair, car chaque démonstration se ramène immédiatement au cas affine, où $f^! = f_{new}^!$ d'après le lemme \ref{lemma-duality-bootstrap}.
\end{remark}




\section{Modules dualisants}
\label{section-dualizing-module}


\noindent
Cette section prolonge Complexes dualisants, section \ref{dualizing-section-dualizing-module}.

\medskip\noindent
Soit $X$ un schéma noethérien et soit $\omega_X^\bullet$ un complexe dualisant. Soit $n \in \mathbf{Z}$ le plus petit entier tel que $H^n(\omega_X^\bullet)$ soit non nul. Autrement dit, $-n$ est la valeur maximale de la fonction de dimension associée à $\omega_X^\bullet$ (lemme \ref{lemma-dimension-function-scheme}). On appelle parfois $H^n(\omega_X^\bullet)$ un {\it module dualisant} ou un {\it faisceau dualisant} sur $X$, et on le note alors souvent $\omega_X$. Nous écrirons « soit $\omega_X$ un module dualisant » pour désigner cette situation.

\medskip\noindent
\begingroup\emergencystretch=2em
L'utilisation des modules dualisants $\omega_X$ sur les schémas noethériens $X$ exige quelques précautions :\par\endgroup
\begin{enumerate}
\item l'entier $n$ peut changer lorsqu'on passe de $X$ à un ouvert $U$ de $X$, et l'on n'aura alors pas nécessairement $\omega_X|_U = \omega_U$ ;
\item le complexe dualisant n'est pas unique ; le module dualisant n'est unique qu'à tensorisation par un module inversible près.
\end{enumerate}
La seconde difficulté sera souvent sans conséquence, car nous travaillerons avec $X$ de type fini sur une base $S$ munie d'un complexe dualisant fixé $\omega_S^\bullet$, et $\omega_X^\bullet$ sera le complexe dualisant normalisé relativement à $\omega_S^\bullet$. La première difficulté ne se présente pas si $X$ est équidimensionnel ; plus précisément, si la fonction de dimension associée à $\omega_X^\bullet$ (lemme \ref{lemma-dimension-function-scheme}) prend la même valeur entière en tout point générique de $X$.

\begin{example}
\label{example-proper-over-local}
Supposons $S = \Spec(A)$, où $(A, \mathfrak m, \kappa)$ est un anneau local noethérien, et où $\omega_S^\bullet$ correspond à un complexe dualisant normalisé $\omega_A^\bullet$. Alors, si $f : X \to S$ est propre sur $S$ et si $\omega_X^\bullet = f^!\omega_S^\bullet$, le faisceau cohérent
$$
\omega_X = H^{-\dim(X)}(\omega_X^\bullet)
$$
est un module dualisant, souvent appelé le module dualisant de $X$ ($S$ et $\omega_S^\bullet$ étant sous-entendus). Nous verrons qu'il possède de bonnes propriétés.
\end{example}

\begin{example}
\label{example-equidimensional-over-field}
Supposons que $X$ soit un schéma équidimensionnel de type fini sur un corps $k$. Il est alors usuel de prendre pour $\omega_X^\bullet$ le complexe dualisant normalisé relativement à $k[0]$, et d'appeler
$$
\omega_X = H^{-\dim(X)}(\omega_X^\bullet)
$$
le module dualisant de $X$. Si $X$ est séparé sur $k$, alors $\omega_X^\bullet = f^!\mathcal{O}_{\Spec(k)}$, où $f : X \to \Spec(k)$ est le morphisme structural, d'après le lemme \ref{lemma-duality-bootstrap}. Si $X$ est propre sur $k$, c'est un cas particulier de l'exemple \ref{example-proper-over-local}.
\end{example}

\begin{lemma}
\label{lemma-dualizing-module}
Soit $X$ un schéma noethérien connexe et soit $\omega_X$ un module dualisant sur $X$. Le support de $\omega_X$ est la réunion des composantes irréductibles de dimension maximale pour toute fonction de dimension, et $\omega_X$ est un $\mathcal{O}_X$-module cohérent satisfaisant la propriété $(S_2)$.
\end{lemma}

\begin{proof}
Selon les conventions précédentes, il existe un complexe dualisant $\omega_X^\bullet$ dont $\omega_X$ est le premier faisceau de cohomologie non nul à partir de la gauche. Puisque $X$ est connexe, deux fonctions de dimension diffèrent d'une constante (Topologie, lemme \ref{topology-lemma-dimension-function-unique}). On peut donc utiliser la fonction de dimension associée à $\omega_X^\bullet$ (lemme \ref{lemma-dimension-function-scheme}). Ces remarques faites, le lemme résulte de Complexes dualisants, lemme \ref{dualizing-lemma-depth-dualizing-module}, et des définitions (en particulier Cohomologie des schémas, définition \ref{coherent-definition-depth}).
\end{proof}

\begin{lemma}
\label{lemma-vanishing-good-dualizing}
Soient $X/A$, avec $\omega_X^\bullet$, et $\omega_X$ comme dans l'exemple \ref{example-proper-over-local}. Alors
\begin{enumerate}
\item $H^i(\omega_X^\bullet) \not = 0 \Rightarrow
i \in \{-\dim(X), \ldots, 0\}$ ;
\item la dimension du support de $H^i(\omega_X^\bullet)$ est au plus $-i$ ;
\item $\text{Supp}(\omega_X)$ est la réunion des composantes de dimension $\dim(X)$ ;
\item $\omega_X$ satisfait la propriété $(S_2)$.
\end{enumerate}
\end{lemma}

\begin{proof}
Soient $\delta_X$ et $\delta_S$ les fonctions de dimension associées à $\omega_X^\bullet$ et $\omega_S^\bullet$ comme au lemme \ref{lemma-good-dualizing-dimension-function}. Puisque $X$ est propre sur $A$, tout sous-schéma fermé de $X$ contient un point fermé $x$ dont l'image est le point fermé $s \in S$, et $\delta_X(x) = \delta_S(s) = 0$. Ainsi, $\delta_X(\xi) = \dim(\overline{\{\xi\}})$ pour tout point $\xi \in X$. On peut donc vérifier chacune des assertions du lemme en considérant ce qui se passe sur $\Spec(\mathcal{O}_{X, x})$ ; le résultat découle alors de Complexes dualisants, lemmes \ref{dualizing-lemma-sitting-in-degrees} et \ref{dualizing-lemma-depth-dualizing-module}. Certains détails sont omis. Les deux dernières assertions peuvent aussi se déduire du lemme \ref{lemma-dualizing-module}.
\end{proof}

\begin{lemma}
\label{lemma-dualizing-module-proper-over-A}
Soit $X/A$, de module dualisant $\omega_X$, comme dans l'exemple \ref{example-proper-over-local}. Soit $d = \dim(X_s)$ la dimension de la fibre fermée. Si $\dim(X) = d + \dim(A)$, alors le module dualisant $\omega_X$ représente le foncteur
$$
\mathcal{F} \longmapsto \Hom_A(H^d(X, \mathcal{F}), \omega_A)
$$
sur la catégorie des $\mathcal{O}_X$-modules cohérents.
\end{lemma}

\begin{proof}
On a
\begin{align*}
\Hom_X(\mathcal{F}, \omega_X)
& =
\Ext^{-\dim(X)}_X(\mathcal{F}, \omega_X^\bullet) \\
& =
\Hom_X(\mathcal{F}[\dim(X)], \omega_X^\bullet) \\
& =
\Hom_X(\mathcal{F}[\dim(X)], f^!(\omega_A^\bullet)) \\
& =
\Hom_S(Rf_*\mathcal{F}[\dim(X)], \omega_A^\bullet) \\
& =
\Hom_A(H^d(X, \mathcal{F}), \omega_A)
\end{align*}
La première égalité résulte de $H^i(\omega_X^\bullet) = 0$ pour $i < -\dim(X)$, voir le lemme \ref{lemma-vanishing-good-dualizing} et Catégories dérivées, lemme \ref{derived-lemma-negative-exts}. La deuxième résulte de la définition des groupes Ext. La troisième est notre choix de $\omega_X^\bullet$. La quatrième résulte du fait que $f^!$ est l'adjoint à droite du lemme \ref{lemma-twisted-inverse-image} pour $f$, voir la section \ref{section-duality}. La dernière égalité résulte du fait que $R^if_*\mathcal{F}$ est nul pour $i > d$ (Cohomologie des schémas, lemme \ref{coherent-lemma-higher-direct-images-zero-above-dimension-fibre}), et que $H^j(\omega_A^\bullet)$ est nul pour $j < -\dim(A)$.
\end{proof}








\section{Schémas de Cohen--Macaulay}
\label{section-CM}

\noindent
Cette section prolonge Complexes dualisants, section \ref{dualizing-section-CM}. La dualité prend une forme particulièrement simple pour les schémas de Cohen–Macaulay.

\begin{lemma}
\label{lemma-dualizing-module-CM-scheme}
Soit $X$ un schéma localement noethérien muni d'un complexe dualisant $\omega_X^\bullet$.
\begin{enumerate}
\item $X$ est de Cohen–Macaulay $\Leftrightarrow$ $\omega_X^\bullet$ possède localement un unique faisceau de cohomologie non nul ;
\item $\mathcal{O}_{X, x}$ est de Cohen–Macaulay $\Leftrightarrow$ $\omega_{X, x}^\bullet$ possède une unique cohomologie non nulle ;
\item $U = \{x \in X \mid \mathcal{O}_{X, x}\text{ est de Cohen--Macaulay}\}$ est ouvert et de Cohen–Macaulay.
\end{enumerate}
Si $X$ est connexe et de Cohen–Macaulay, il existe un entier $n$ et un $\mathcal{O}_X$-module cohérent de Cohen–Macaulay $\omega_X$ tels que $\omega_X^\bullet = \omega_X[-n]$.
\end{lemma}

\begin{proof}
Par définition et d'après Complexes dualisants, lemme \ref{dualizing-lemma-dualizing-localize}, pour tout $x \in X$, le complexe $\omega_{X, x}^\bullet$ est un complexe dualisant sur $\mathcal{O}_{X, x}$. Complexes dualisants, lemme \ref{dualizing-lemma-apply-CM}, donne donc (2).

\medskip\noindent
Pour démontrer (3), supposons $\mathcal{O}_{X, x}$ de Cohen–Macaulay. Soit $n_x$ l'unique entier tel que $H^{n_{x}}(\omega_{X, x}^\bullet)$ soit non nul. Pour un voisinage affine $V \subset X$ de $x$, on a $\omega_X^\bullet|_V$ dans $D^b_{\textit{Coh}}(\mathcal{O}_V)$ ; il n'y a donc qu'un nombre fini de modules cohérents non nuls $H^i(\omega_X^\bullet)|_V$. Après avoir rétréci $V$, on peut ainsi supposer que seul $H^{n_x}$ est non nul, voir Modules, lemme \ref{modules-lemma-finite-type-stalk-zero}. On voit de cette façon que $\mathcal{O}_{X, v}$ est de Cohen–Macaulay pour tout $v \in V$. Cela montre que $U$ est ouvert et est un schéma de Cohen–Macaulay.

\medskip\noindent
Démonstration de (1). L'implication $\Leftarrow$ résulte de (2). L'implication $\Rightarrow$ résulte du paragraphe précédent, où l'on a montré que, si $\mathcal{O}_{X, x}$ est de Cohen–Macaulay, alors, dans un voisinage de $x$, le complexe $\omega_X^\bullet$ ne possède qu'un seul faisceau de cohomologie non nul.

\medskip\noindent
Supposons $X$ connexe et de Cohen–Macaulay. Ce qui précède montre que l'application $x \mapsto n_x$ est localement constante. Puisque $X$ est connexe, elle est constante, de valeur $n$, disons. En posant $\omega_X = H^n(\omega_X^\bullet)$, on obtient le lemme, car $\omega_X$ est de Cohen–Macaulay d'après Complexes dualisants, lemme \ref{dualizing-lemma-apply-CM} (et Cohomologie des schémas, définition \ref{coherent-definition-Cohen-Macaulay}).
\end{proof}

\begin{lemma}
\label{lemma-has-dualizing-module-CM-scheme}
Soit $X$ un schéma localement noethérien. S'il existe un faisceau cohérent $\omega_X$ tel que $\omega_X[0]$ soit un complexe dualisant sur $X$, alors $X$ est un schéma de Cohen–Macaulay.
\end{lemma}

\begin{proof}
Cela résulte immédiatement de Complexes dualisants, lemme \ref{dualizing-lemma-has-dualizing-module-CM}, et de nos définitions.
\end{proof}

\begin{lemma}
\label{lemma-affine-flat-Noetherian-CM}
Dans la situation \ref{situation-shriek}, soit $f : X \to Y$ un morphisme de $\textit{FTS}_S$. Soit $x \in X$. Si $f$ est plat, les conditions suivantes sont équivalentes :
\begin{enumerate}
\item $f$ est de Cohen–Macaulay en $x$ ;
\item $f^!\mathcal{O}_Y$ possède un unique faisceau de cohomologie non nul dans un voisinage de $x$.
\end{enumerate}
\end{lemma}

\begin{proof}
Un sens du lemme résulte du lemme \ref{lemma-CM-shriek}. Pour démontrer la réciproque, on peut supposer que $f^!\mathcal{O}_Y$ possède un unique faisceau de cohomologie non nul. Soit $y = f(x)$. Soient $\xi_1, \ldots, \xi_n \in X_y$ les points génériques de la fibre $X_y$ se spécialisant en $x$. Soient $d_1, \ldots, d_n$ les dimensions des composantes irréductibles correspondantes de $X_y$. Le morphisme $f : X \to Y$ est de Cohen–Macaulay en $\eta_i$ d'après Compléments sur les morphismes, lemme \ref{more-morphisms-lemma-flat-finite-presentation-CM-open}. Le lemme \ref{lemma-CM-shriek} donne donc $d_1 = \ldots = d_n$. Si $d$ désigne la valeur commune, alors $d = \dim_x(X_y)$. Après avoir rétréci $X$, on peut supposer toutes les fibres de dimension au plus $d$ (Morphismes, lemme \ref{morphisms-lemma-openness-bounded-dimension-fibres}). L'unique faisceau de cohomologie non nul $\omega = H^{-d}(f^!\mathcal{O}_Y)$ est alors plat sur $Y$ d'après le lemme \ref{lemma-flat-shriek}. Si $h : X_y \to X$ désigne le morphisme canonique, on a donc
$Lh^*(f^!\mathcal{O}_Y) = Lh^*(\omega[d]) = (h^*\omega)[d]$
d'après Catégories dérivées des schémas, lemme \ref{perfect-lemma-tor-independence-and-tor-amplitude}. Ainsi, $h^*\omega[d]$ est le complexe dualisant de $X_y$ d'après le lemme \ref{lemma-relative-dualizing-fibres}. Par conséquent, $X_y$ est de Cohen–Macaulay d'après le lemme \ref{lemma-dualizing-module-CM-scheme}. Cela montre que $f$ est de Cohen–Macaulay en $x$, comme voulu.
\end{proof}

\begin{remark}
\label{remark-CM-morphism-compare-dualizing}
Dans la situation \ref{situation-shriek}, soit $f : X \to Y$ un morphisme de $\textit{FTS}_S$. Supposons que $f$ soit un morphisme de Cohen–Macaulay de dimension relative $d$. Soit $\omega_{X/Y} = H^{-d}(f^!\mathcal{O}_Y)$ l'unique faisceau de cohomologie non nul de $f^!\mathcal{O}_Y$, voir le lemme \ref{lemma-CM-shriek}. Il existe alors un isomorphisme canonique
$$
f^!K = Lf^*K \otimes_{\mathcal{O}_X}^\mathbf{L} \omega_{X/Y}[d]
$$
pour $K \in D^+_\QCoh(\mathcal{O}_Y)$, voir le lemme \ref{lemma-perfect-comparison-shriek}. En particulier, si $S$ admet un complexe dualisant $\omega_S^\bullet$, si $\omega_Y^\bullet = (Y \to S)^!\omega_S^\bullet$ et si $\omega_X^\bullet = (X \to S)^!\omega_S^\bullet$, on a
$$
\omega_X^\bullet =
Lf^*\omega_Y^\bullet \otimes_{\mathcal{O}_X}^\mathbf{L} \omega_{X/Y}[d]
$$
Si de plus $X$ et $Y$ sont connexes et de Cohen–Macaulay, et si $\omega_Y$ et $\omega_X$ désignent les uniques faisceaux de cohomologie non nuls de $\omega_Y^\bullet$ et $\omega_X^\bullet$, on a donc
$$
\omega_X = f^*\omega_Y \otimes_{\mathcal{O}_X} \omega_{X/Y}.
$$
Des résultats analogues valent pour $X$ et $Y$ schémas quelconques de type fini sur $S$ (donc pas nécessairement séparés sur $S$), avec des complexes dualisants normalisés relativement à $\omega_S^\bullet$ comme à la section \ref{section-glue}.
\end{remark}





\section{Schémas de Gorenstein}
\label{section-gorenstein}

\noindent
Cette section prolonge Complexes dualisants, section \ref{dualizing-section-gorenstein}.

\begin{definition}
\label{definition-gorenstein}
Soit $X$ un schéma. On dit que $X$ est {\it de Gorenstein} si $X$ est localement noethérien et si $\mathcal{O}_{X, x}$ est de Gorenstein pour tout $x \in X$.
\end{definition}

\noindent
Cette définition a un sens, car un anneau noethérien est dit de Gorenstein si et seulement si tous ses anneaux locaux sont de Gorenstein, voir Complexes dualisants, définition \ref{dualizing-definition-gorenstein}.

\begin{lemma}
\label{lemma-gorenstein-CM}
Un schéma de Gorenstein est de Cohen–Macaulay.
\end{lemma}

\begin{proof}
En travaillant localement sur les ouverts affines, cela résulte du résultat algébrique correspondant, à savoir Complexes dualisants, lemme \ref{dualizing-lemma-gorenstein-CM}.
\end{proof}

\begin{lemma}
\label{lemma-regular-gorenstein}
Un schéma régulier est de Gorenstein.
\end{lemma}

\begin{proof}
En travaillant localement sur les ouverts affines, cela résulte du résultat algébrique correspondant, à savoir Complexes dualisants, lemme \ref{dualizing-lemma-regular-gorenstein}.
\end{proof}

\begin{lemma}
\label{lemma-gorenstein}
Soit $X$ un schéma localement noethérien.
\begin{enumerate}
\item Si $X$ admet un complexe dualisant $\omega_X^\bullet$, alors
\begin{enumerate}
\item $X$ est de Gorenstein $\Leftrightarrow$ $\omega_X^\bullet$ est un objet inversible de $D(\mathcal{O}_X)$ ;
\item $\mathcal{O}_{X, x}$ est de Gorenstein $\Leftrightarrow$ $\omega_{X, x}^\bullet$ est un objet inversible de $D(\mathcal{O}_{X, x})$ ;
\item $U = \{x \in X \mid \mathcal{O}_{X, x}\text{ est de Gorenstein}\}$ est un sous-schéma ouvert de Gorenstein.
\end{enumerate}
\item Si $X$ est de Gorenstein, alors $X$ admet un complexe dualisant si et seulement si $\mathcal{O}_X[0]$ est un complexe dualisant.
\end{enumerate}
\end{lemma}

\begin{proof}
En travaillant localement sur les ouverts affines, cela résulte du résultat algébrique correspondant, à savoir Complexes dualisants, lemme \ref{dualizing-lemma-gorenstein}.
\end{proof}

\begin{lemma}
\label{lemma-gorenstein-lci}
Si $f : Y \to X$ est un morphisme localement d'intersection complète et si $X$ est un schéma de Gorenstein, alors $Y$ est de Gorenstein.
\end{lemma}

\begin{proof}
D'après Compléments sur les morphismes, lemme \ref{more-morphisms-lemma-affine-lci}, il suffit de démontrer l'énoncé correspondant pour les homomorphismes d'anneaux. C'est Complexes dualisants, lemme \ref{dualizing-lemma-gorenstein-lci}.
\end{proof}

\begin{lemma}
\label{lemma-gorenstein-local-syntomic}
La propriété $\mathcal{P}(S) =$« $S$ est de Gorenstein » est locale pour la topologie syntomique.
\end{lemma}

\begin{proof}
Soit $\{S_i \to S\}$ un recouvrement syntomique. Le schéma $S$ est localement noethérien si et seulement si chaque $S_i$ est noethérien, voir Descente, lemme \ref{descent-lemma-Noetherian-local-fppf}. On peut donc supposer $S$ et $S_i$ localement noethériens. Si $S$ est de Gorenstein, chaque $S_i$ est de Gorenstein d'après le lemme \ref{lemma-gorenstein-lci}. Réciproquement, si chaque $S_i$ est de Gorenstein, on peut, pour tout point $s \in S$, choisir $i$ et $t \in S_i$ d'image $s$. Alors $\mathcal{O}_{S, s} \to \mathcal{O}_{S_i, t}$ est un homomorphisme local plat d'anneaux, avec $\mathcal{O}_{S_i, t}$ de Gorenstein. Ainsi, $\mathcal{O}_{S, s}$ est de Gorenstein d'après Complexes dualisants, lemme \ref{dualizing-lemma-flat-under-gorenstein}.
\end{proof}






\section{Morphismes de Gorenstein}
\label{section-gorenstein-morphisms}

\noindent
Cette section s'inscrit dans une série. Les sections correspondantes pour les morphismes normaux, réguliers et de Cohen–Macaulay se trouvent dans Compléments sur les morphismes, sections \ref{more-morphisms-section-normal}, \ref{more-morphisms-section-regular} et \ref{more-morphisms-section-CM}.

\medskip\noindent
Le lemme suivant montre qu'il n'y a pas lieu de définir les schémas géométriquement de Gorenstein, car ils coïncideraient avec les schémas de Gorenstein.

\begin{lemma}
\label{lemma-gorenstein-base-change}
Soit $X$ un schéma localement noethérien sur le corps $k$. Soit $k'/k$ une extension de corps de type fini. Soit $x \in X$ un point, et soit $x' \in X_{k'}$ un point au-dessus de $x$. Alors
$$
\mathcal{O}_{X, x}\text{ est de Gorenstein}
\Leftrightarrow
\mathcal{O}_{X_{k'}, x'}\text{ est de Gorenstein}
$$
Si $X$ est localement de type fini sur $k$, le même énoncé vaut pour toute extension de corps $k'/k$.
\end{lemma}

\begin{proof}
Dans les deux cas, l'homomorphisme d'anneaux $\mathcal{O}_{X, x} \to \mathcal{O}_{X_{k'}, x'}$ est un homomorphisme local fidèlement plat d'anneaux locaux noethériens. Si $\mathcal{O}_{X_{k'}, x'}$ est de Gorenstein, il en est donc de même de $\mathcal{O}_{X, x}$ d'après Complexes dualisants, lemme \ref{dualizing-lemma-flat-under-gorenstein}. Pour remonter, on utilise également Complexes dualisants, lemme \ref{dualizing-lemma-flat-under-gorenstein}. Il faut donc montrer que
$$
\mathcal{O}_{X_{k'}, x'}/\mathfrak m_x \mathcal{O}_{X_{k'}, x'} =
\kappa(x) \otimes_k k'
$$
est de Gorenstein. Remarquons que, dans le premier cas, $k \to k'$ est de type fini, et que, dans le second, $k \to \kappa(x)$ est de type fini. Le résultat découle donc de la propriété (A) pour les anneaux de Gorenstein, voir Complexes dualisants, lemme \ref{dualizing-lemma-formal-fibres-gorenstein}.
\end{proof}

\noindent
Le lemme précédent garantit que la définition suivante des morphismes de Gorenstein est la bonne.

\begin{definition}
\label{definition-gorenstein-morphism}
Soit $f : X \to Y$ un morphisme de schémas. Supposons que toutes les fibres $X_y$ soient des schémas localement noethériens.
\begin{enumerate}
\item Soient $x \in X$ et $y = f(x)$. On dit que $f$ est {\it de Gorenstein en $x$} si $f$ est plat en $x$ et si l'anneau local du schéma $X_y$ en $x$ est de Gorenstein.
\item On dit que $f$ est un {\it morphisme de Gorenstein} si $f$ est de Gorenstein en tout point de $X$.
\end{enumerate}
\end{definition}

\noindent
Voici une reformulation.

\begin{lemma}
\label{lemma-gorenstein-morphism}
Soit $f : X \to Y$ un morphisme de schémas. Supposons toutes les fibres de $f$ localement noethériennes. Les conditions suivantes sont équivalentes :
\begin{enumerate}
\item $f$ est de Gorenstein ;
\item $f$ est plat et ses fibres sont des schémas de Gorenstein.
\end{enumerate}
\end{lemma}

\begin{proof}
Cela résulte directement des définitions.
\end{proof}

\begin{lemma}
\label{lemma-gorenstein-CM-morphism}
Un morphisme de Gorenstein est de Cohen–Macaulay.
\end{lemma}

\begin{proof}
Cela résulte du lemme \ref{lemma-gorenstein-CM} et des définitions.
\end{proof}

\begin{lemma}
\label{lemma-lci-gorenstein}
Un morphisme syntomique est de Gorenstein. De manière équivalente, un morphisme plat localement d'intersection complète est de Gorenstein.
\end{lemma}

\begin{proof}
Rappelons qu'un morphisme syntomique est plat et que ses fibres sont localement des intersections complètes sur des corps, voir Morphismes, lemme \ref{morphisms-lemma-syntomic-flat-fibres}. Puisqu'une intersection complète locale sur un corps est un schéma de Gorenstein d'après le lemme \ref{lemma-gorenstein-lci}, on conclut. Les propriétés « syntomique » et « plat et localement d'intersection complète » sont équivalentes d'après Compléments sur les morphismes, lemme \ref{more-morphisms-lemma-flat-lci}.
\end{proof}

\begin{lemma}
\label{lemma-composition-gorenstein}
Soient $f : X \to Y$ et $g : Y \to Z$ des morphismes. Supposons que les fibres $X_y$, $Y_z$ et $X_z$ de $f$, $g$ et $g \circ f$ soient localement noethériennes.
\begin{enumerate}
\item Si $f$ est de Gorenstein en $x$ et si $g$ est de Gorenstein en $f(x)$, alors $g \circ f$ est de Gorenstein en $x$.
\item Si $f$ et $g$ sont de Gorenstein, alors $g \circ f$ est de Gorenstein.
\item Si $g \circ f$ est de Gorenstein en $x$ et si $f$ est plat en $x$, alors $f$ est de Gorenstein en $x$ et $g$ est de Gorenstein en $f(x)$.
\item Si $g \circ f$ est de Gorenstein et si $f$ est plat, alors $f$ est de Gorenstein et $g$ est de Gorenstein en tout point de l'image de $f$.
\end{enumerate}
\end{lemma}

\begin{proof}
Après traduction en algèbre, cela résulte de Complexes dualisants, lemme \ref{dualizing-lemma-flat-under-gorenstein}.
\end{proof}

\begin{lemma}
\label{lemma-flat-morphism-from-gorenstein-scheme}
\begin{slogan}
Le caractère de Gorenstein de l'espace total d'une fibration plate entraîne le même caractère pour la base et les fibres
\end{slogan}
Soit $f : X \to Y$ un morphisme plat de schémas localement noethériens. Si $X$ est de Gorenstein, alors $f$ est de Gorenstein et $\mathcal{O}_{Y, f(x)}$ est de Gorenstein pour tout $x \in X$.
\end{lemma}

\begin{proof}
Après traduction en algèbre, cela résulte de Complexes dualisants, lemme \ref{dualizing-lemma-flat-under-gorenstein}.
\end{proof}

\begin{lemma}
\label{lemma-base-change-gorenstein}
Soit $f : X \to Y$ un morphisme de schémas. Supposons que toutes les fibres $X_y$ soient des schémas localement noethériens. Soit $Y' \to Y$ localement de type fini. Soit $f' : X' \to Y'$ le changement de base de $f$. Soit $x' \in X'$ un point d'image $x \in X$.
\begin{enumerate}
\item Si $f$ est de Gorenstein en $x$, alors $f' : X' \to Y'$ est de Gorenstein en $x'$.
\item Si $f$ est plat en $x$ et si $f'$ est de Gorenstein en $x'$, alors $f$ est de Gorenstein en $x$.
\item Si $Y' \to Y$ est plat en $f'(x')$ et si $f'$ est de Gorenstein en $x'$, alors $f$ est de Gorenstein en $x$.
\end{enumerate}
\end{lemma}

\begin{proof}
Remarquons que l'hypothèse sur $Y' \to Y$ implique que, pour $y' \in Y'$ d'image $y \in Y$, l'extension de corps $\kappa(y')/\kappa(y)$ est de type fini. Par conséquent, toutes les fibres $X'_{y'} = (X_y)_{\kappa(y')}$ sont elles aussi localement noethériennes, voir Variétés, lemme \ref{varieties-lemma-locally-Noetherian-base-change}. Le lemme a donc un sens. Posons $y' = f'(x')$ et $y = f(x)$. On obtient ainsi le diagramme commutatif suivant d'anneaux locaux
$$
\xymatrix{
\mathcal{O}_{X', x'} & \mathcal{O}_{X, x} \ar[l] \\
\mathcal{O}_{Y', y'} \ar[u] & \mathcal{O}_{Y, y} \ar[l] \ar[u]
}
$$
dans lequel le coin supérieur gauche est un localisé du produit tensoriel des coins supérieur droit et inférieur gauche sur le coin inférieur droit.

\medskip\noindent
Supposons que $f$ soit de Gorenstein en $x$. La platitude de $\mathcal{O}_{Y, y} \to \mathcal{O}_{X, x}$ entraîne celle de $\mathcal{O}_{Y', y'} \to \mathcal{O}_{X', x'}$, voir Algèbre, lemme \ref{algebra-lemma-base-change-flat-up-down}. Le fait que $\mathcal{O}_{X, x}/\mathfrak m_y\mathcal{O}_{X, x}$ soit de Gorenstein implique que $\mathcal{O}_{X', x'}/\mathfrak m_{y'}\mathcal{O}_{X', x'}$ est de Gorenstein, voir le lemme \ref{lemma-gorenstein-base-change}. On voit donc que $f'$ est de Gorenstein en $x'$.

\medskip\noindent
Supposons que $f$ soit plat en $x$ et que $f'$ soit de Gorenstein en $x'$. Le fait que $\mathcal{O}_{X', x'}/\mathfrak m_{y'}\mathcal{O}_{X', x'}$ soit de Gorenstein implique que $\mathcal{O}_{X, x}/\mathfrak m_y\mathcal{O}_{X, x}$ est de Gorenstein, voir le lemme \ref{lemma-gorenstein-base-change}. On voit donc que $f$ est de Gorenstein en $x$.

\medskip\noindent
Supposons que $Y' \to Y$ soit plat en $y'$ et que $f'$ soit de Gorenstein en $x'$. La platitude de $\mathcal{O}_{Y', y'} \to \mathcal{O}_{X', x'}$ et de $\mathcal{O}_{Y, y} \to \mathcal{O}_{Y', y'}$ entraîne celle de $\mathcal{O}_{Y, y} \to \mathcal{O}_{X, x}$, voir Algèbre, lemme \ref{algebra-lemma-base-change-flat-up-down}. Le fait que $\mathcal{O}_{X', x'}/\mathfrak m_{y'}\mathcal{O}_{X', x'}$ soit de Gorenstein implique que $\mathcal{O}_{X, x}/\mathfrak m_y\mathcal{O}_{X, x}$ est de Gorenstein, voir le lemme \ref{lemma-gorenstein-base-change}. On voit donc que $f$ est de Gorenstein en $x$.
\end{proof}

\begin{lemma}
\label{lemma-flat-lft-base-change-gorenstein}
Soit $f : X \to Y$ un morphisme de schémas plat et localement de type fini. Alors la formation de l'ensemble
$\{x \in X \mid f\text{ est de Gorenstein en }x\}$
commute à tout changement de base.
\end{lemma}

\begin{proof}
L'hypothèse implique que toute fibre de $f$ est localement de type fini sur un corps, donc localement noethérienne, et il en est de même après tout changement de base. L'énoncé a donc un sens. En examinant les fibres, on se ramène au problème suivant : soit $X$ un schéma localement de type fini sur un corps $k$, soit $K/k$ une extension de corps, et soit $x_K \in X_K$ un point d'image $x \in X$. Problème : montrer que $\mathcal{O}_{X_K, x_K}$ est de Gorenstein si et seulement si $\mathcal{O}_{X, x}$ est de Gorenstein.

\medskip\noindent
On peut résoudre le problème avec un peu d'algèbre comme suit. Choisissons un ouvert affine $\Spec(A) \subset X$ contenant $x$. Écrivons que $x$ correspond à $\mathfrak p \subset A$. Avec $A_K = A \otimes_k K$, on voit que $\Spec(A_K) \subset X_K$ contient $x_K$. Écrivons que $x_K$ correspond à $\mathfrak p_K \subset A_K$. Soit $\omega_A^\bullet$ un complexe dualisant pour $A$. D'après Complexes dualisants, lemme \ref{dualizing-lemma-base-change-dualizing-over-field}, $\omega_A^\bullet \otimes_A A_K$ est un complexe dualisant pour $A_K$. On a terminé, car $A_\mathfrak p \to (A_K)_{\mathfrak p_K}$ est un homomorphisme local plat d'anneaux noethériens et, par conséquent, $(\omega_A^\bullet)_\mathfrak p$ est un objet inversible de $D(A_\mathfrak p)$ si et seulement si $(\omega_A^\bullet)_\mathfrak p \otimes_{A_\mathfrak p} (A_K)_{\mathfrak p_K}$ est un objet inversible de $D((A_K)_{\mathfrak p_K})$. Quelques détails sont omis ; indication : examiner les modules de cohomologie.
\end{proof}

\begin{lemma}
\label{lemma-affine-flat-Noetherian-gorenstein}
Dans la situation \ref{situation-shriek}, soit $f : X \to Y$ un morphisme de $\textit{FTS}_S$. Soit $x \in X$. Si $f$ est plat, les conditions suivantes sont équivalentes :
\begin{enumerate}
\item $f$ est de Gorenstein en $x$ ;
\item $f^!\mathcal{O}_Y$ est isomorphe à un objet inversible dans un voisinage de $x$.
\end{enumerate}
En particulier, l'ensemble des points où $f$ est de Gorenstein est ouvert dans $X$.
\end{lemma}

\begin{proof}
Posons $\omega^\bullet = f^!\mathcal{O}_Y$. D'après le lemme \ref{lemma-relative-dualizing-fibres}, c'est un complexe borné à faisceaux de cohomologie cohérents dont la restriction dérivée $Lh^*\omega^\bullet$ à la fibre $X_y$ est un complexe dualisant sur $X_y$. Notons $i : x \to X_y$ l'inclusion d'un point. Les conditions suivantes sont alors équivalentes :
\begin{enumerate}
\item $f$ est de Gorenstein en $x$ ;
\item $\mathcal{O}_{X_y, x}$ est de Gorenstein ;
\item $Lh^*\omega^\bullet$ est inversible dans un voisinage de $x$ ;
\item $Li^* Lh^* \omega^\bullet$ possède exactement une cohomologie non nulle, de dimension $1$ sur $\kappa(x)$ ;
\item $L(h \circ i)^* \omega^\bullet$ possède exactement une cohomologie non nulle, de dimension $1$ sur $\kappa(x)$ ;
\item $\omega^\bullet$ est inversible dans un voisinage de $x$.
\end{enumerate}
L'équivalence de (1) et (2) résulte de la définition (puisque $f$ est plat). L'équivalence de (2) et (3) résulte du lemme \ref{lemma-gorenstein}. L'équivalence de (3) et (4) résulte de Compléments sur l'algèbre, lemme \ref{more-algebra-lemma-lift-bounded-pseudo-coherent-to-perfect}. L'équivalence de (4) et (5) vaut parce que $Li^* Lh^* = L(h \circ i)^*$. L'équivalence de (5) et (6) résulte de Compléments sur l'algèbre, lemme \ref{more-algebra-lemma-lift-bounded-pseudo-coherent-to-perfect}. Le lemme est donc clair.
\end{proof}

\begin{lemma}
\label{lemma-flat-finite-presentation-characterize-gorenstein}
Soit $f : X \to S$ un morphisme de schémas plat et localement de présentation finie. Soit $x \in X$ d'image $s \in S$. Posons $d = \dim_x(X_s)$. Les conditions suivantes sont équivalentes :
\begin{enumerate}
\item $f$ est de Gorenstein en $x$ ;
\item il existe un voisinage ouvert $U \subset X$ de $x$ et un morphisme localement quasi-fini $U \to \mathbf{A}^d_S$ sur $S$ qui est de Gorenstein en $x$ ;
\item il existe un voisinage ouvert $U \subset X$ de $x$ et un morphisme de Gorenstein localement quasi-fini $U \to \mathbf{A}^d_S$ sur $S$ ;
\item pour tout $S$-morphisme $g : U \to \mathbf{A}^d_S$ d'un voisinage ouvert $U \subset X$ de $x$, on a : $g$ est quasi-fini en $x$ $\Rightarrow$ $g$ est de Gorenstein en $x$.
\end{enumerate}
En particulier, l'ensemble des points où $f$ est de Gorenstein est ouvert dans $X$.
\end{lemma}

\begin{proof}
Choisissons des ouverts affines $U = \Spec(A) \subset X$ avec $x \in U$, et $V = \Spec(R) \subset S$ avec $f(U) \subset V$. Alors $R \to A$ est un homomorphisme d'anneaux plat de présentation finie. Soit $\mathfrak p \subset A$ l'idéal premier correspondant à $x$. Après avoir remplacé $A$ par une localisation principale, on peut supposer qu'il existe un homomorphisme quasi-fini $R[x_1, \ldots, x_d]  \to A$, voir Algèbre, lemme \ref{algebra-lemma-quasi-finite-over-polynomial-algebra}. Il existe donc au moins une paire $(U, g)$ formée d'un voisinage ouvert $U \subset X$ de $x$ et d'un morphisme localement\footnote{Si $S$ est quasi-séparé, alors $g$ sera quasi-fini.} quasi-fini $g : U \to \mathbf{A}^d_S$.

\medskip\noindent
Cela étant dit, le lemme se traduit par le problème algébrique suivant (traduction omise). Étant donnés $R \to A$ plat et de présentation finie, un idéal premier $\mathfrak p \subset A$ et $\varphi : R[x_1, \ldots, x_d] \to A$ quasi-fini en $\mathfrak p$, les conditions suivantes sont équivalentes :
\begin{enumerate}
\item[(a)] $\Spec(\varphi)$ est de Gorenstein en $\mathfrak p$ ;
\item[(b)] $\Spec(A) \to \Spec(R)$ est de Gorenstein en $\mathfrak p$ ;
\item[(c)] $\Spec(A) \to \Spec(R)$ est de Gorenstein dans un voisinage ouvert de $\mathfrak p$.
\end{enumerate}
Dans chaque cas, $R[x_1, \ldots, x_n] \to A$ est plat en $\mathfrak p$ ; l'ouverture de la platitude (Algèbre, théorème \ref{algebra-theorem-openness-flatness}) permet donc de supposer que $R[x_1, \ldots, x_n] \to A$ est plat (en remplaçant $A$ par une localisation principale convenable). D'après Algèbre, lemme \ref{algebra-lemma-flat-finite-presentation-limit-flat}, il existe $R_0 \subset R$ et $R_0[x_1, \ldots, x_n] \to A_0$ tels que $R_0$ soit de type fini sur $\mathbf{Z}$, que $R_0 \to A_0$ soit de type fini et que $R_0[x_1, \ldots, x_n] \to A_0$ soit plat. Remarquons que l'ensemble des points où un morphisme plat de type fini est de Gorenstein commute au changement de base d'après le lemme \ref{lemma-base-change-gorenstein}. On se ramène ainsi au cas où $R$ est noethérien.

\medskip\noindent
On peut donc supposer $X$ et $S$ affines et disposer d'une factorisation de $f$ de la forme
$$
X \xrightarrow{g} \mathbf{A}^n_S \xrightarrow{p} S
$$
où $g$ est plat et quasi-fini et où $S$ est noethérien. Alors $X$ et $\mathbf{A}^n_S$ sont séparés sur $S$, et l'on a
$$
f^!\mathcal{O}_S = g^!p^!\mathcal{O}_S = g^!\mathcal{O}_{\mathbf{A}^n_S}[n]
$$
d'après les propriétés connues des foncteurs image inverse extraordinaire (lemmes \ref{lemma-upper-shriek-composition} et \ref{lemma-shriek-affine-line}). L'équivalence de (a), (b) et (c) résulte donc du lemme \ref{lemma-affine-flat-Noetherian-gorenstein}.
\end{proof}

\begin{lemma}
\label{lemma-gorenstein-local-source-and-target}
La propriété
$\mathcal{P}(f)=$« les fibres de $f$ sont localement noethériennes et $f$ est de Gorenstein »
est locale pour la topologie fppf sur le but et locale pour la topologie syntomique sur la source.
\end{lemma}

\begin{proof}
On a
$\mathcal{P}(f) =
\mathcal{P}_1(f) \wedge \mathcal{P}_2(f)$
où
$\mathcal{P}_1(f)=$« $f$ est plat » et
$\mathcal{P}_2(f)=$« les fibres de $f$ sont localement noethériennes et de Gorenstein ».
On sait que $\mathcal{P}_1$ est locale pour la topologie fppf sur la source et sur le but, voir Descente, lemmes \ref{descent-lemma-descending-property-flat} et \ref{descent-lemma-flat-fpqc-local-source}. Il reste donc à traiter $\mathcal{P}_2$.

\medskip\noindent
Soit $f : X \to Y$ un morphisme de schémas. Soit $\{\varphi_i : Y_i \to Y\}_{i \in I}$ un recouvrement fppf de $Y$. Notons $f_i : X_i \to Y_i$ le changement de base de $f$ par $\varphi_i$. Soit $i \in I$, et soit $y_i \in Y_i$ un point. Posons $y = \varphi_i(y_i)$. Remarquons que
$$
X_{i, y_i} = \Spec(\kappa(y_i)) \times_{\Spec(\kappa(y))} X_y.
$$
et que $\kappa(y_i)/\kappa(y)$ est une extension de corps de type fini. Ainsi, si $X_y$ est localement noethérien, alors $X_{i, y_i}$ est localement noethérien, voir Variétés, lemme \ref{varieties-lemma-locally-Noetherian-base-change}. Et si, de plus, $X_y$ est de Gorenstein, alors $X_{i, y_i}$ est de Gorenstein, voir le lemme \ref{lemma-gorenstein-base-change}. Par conséquent, $\mathcal{P}_2$ est locale pour la topologie fppf sur le but.

\medskip\noindent
Soit $\{X_i \to X\}$ un recouvrement syntomique de $X$. Soit $y \in Y$. Dans ce cas, $\{X_{i, y} \to X_y\}$ est un recouvrement syntomique de la fibre. La localité de $\mathcal{P}_2$ pour la topologie syntomique sur la source résulte donc du lemme \ref{lemma-gorenstein-local-syntomic}.
\end{proof}




\section{Compléments sur les complexes dualisants}
\label{section-more-dualizing}

\noindent
Quelques lemmes qui ne trouvent pas très bien leur place ailleurs.

\begin{lemma}
\label{lemma-descent-ascent}
Soit $f : X \to Y$ un morphisme de schémas localement noethériens. Supposons que
\begin{enumerate}
\item $f$ soit syntomique et surjectif ; ou que
\item $f$ soit un morphisme surjectif, plat et localement d'intersection complète ; ou que
\item $f$ soit un morphisme de Gorenstein surjectif de type fini.
\end{enumerate}
Alors $K \in D_\QCoh(\mathcal{O}_Y)$ est un complexe dualisant sur $Y$ si et seulement si $Lf^*K$ est un complexe dualisant sur $X$.
\end{lemma}

\begin{proof}
En prenant des ouverts affines et en utilisant Catégories dérivées de schémas, lemme \ref{perfect-lemma-affine-compare-bounded}, cela se traduit par Complexes dualisants, lemme \ref{dualizing-lemma-descent-ascent}.
\end{proof}








\section{Dualité pour les schémas propres sur des corps}
\label{section-duality-proper-over-field}

\noindent
Dans cette section, nous développons les conséquences des résultats très généraux ci-dessus sur les complexes dualisants et la dualité pour les schémas propres sur des corps.

\begin{lemma}
\label{lemma-duality-proper-over-field}
Soit $X$ un schéma propre sur un corps $k$. Il existe un complexe dualisant $\omega_X^\bullet$ possédant les propriétés suivantes :
\begin{enumerate}
\item $H^i(\omega_X^\bullet)$ n'est non nul que pour $i \in [-\dim(X), 0]$ ;
\item $\omega_X = H^{-\dim(X)}(\omega_X^\bullet)$ est un $(S_2)$-module cohérent dont le support est la réunion des composantes irréductibles de dimension $\dim(X)$ ;
\item la dimension du support de $H^i(\omega_X^\bullet)$ est au plus $-i$ ;
\item pour $x \in X$ fermé, le module
$H^i(\omega_{X, x}^\bullet) \oplus \ldots \oplus H^0(\omega_{X, x}^\bullet)$
est non nul si et seulement si $\text{depth}(\mathcal{O}_{X, x}) \leq -i$ ;
\item pour $K \in D_\QCoh(\mathcal{O}_X)$, il existe des isomorphismes fonctoriels\footnote{Cette propriété caractérise $\omega_X^\bullet$ dans $D_\QCoh(\mathcal{O}_X)$ à isomorphisme unique près par le lemme de Yoneda. Puisque $\omega_X^\bullet$ appartient à $D^b_{\textit{Coh}}(\mathcal{O}_X)$, il suffit en fait de considérer $K \in D^b_{\textit{Coh}}(\mathcal{O}_X)$.}
\begingroup\small
$$
\Ext^i_X(K, \omega_X^\bullet) = \Hom_k(H^{-i}(X, K), k)
$$
\endgroup
compatibles aux décalages et aux triangles distingués ;
\item il existe des isomorphismes fonctoriels
\par\noindent\hfil
$\Hom(\mathcal{F}, \omega_X) = \Hom_k(H^{\dim(X)}(X, \mathcal{F}), k)$
\hfil\par\noindent
pour $\mathcal{F}$ quasi-cohérent sur $X$ ;
\item si $X \to \Spec(k)$ est lisse de dimension relative $d$, alors $\omega_X^\bullet \cong \wedge^d\Omega_{X/k}[d]$ et $\omega_X \cong \wedge^d\Omega_{X/k}$.
\end{enumerate}
\end{lemma}

\begin{proof}
Notons $f : X \to \Spec(k)$ le morphisme structural. Soit $a$ l'adjoint à droite de l'image directe par ce morphisme, voir le lemme \ref{lemma-twisted-inverse-image}. Considérons le complexe dualisant relatif
$$
\omega_X^\bullet = a(\mathcal{O}_{\Spec(k)})
$$
Comparer avec la remarque \ref{remark-relative-dualizing-complex}. Puisque $f$ est propre, on a
\par\noindent\hfil
$f^!(\mathcal{O}_{\Spec(k)}) = a(\mathcal{O}_{\Spec(k)})$
\hfil\par\noindent
par définition, voir la section \ref{section-upper-shriek}. En appliquant le lemme \ref{lemma-shriek-dualizing}, on trouve que $\omega_X^\bullet$ est un complexe dualisant. En outre, on voit que $\omega_X^\bullet$ et $\omega_X$ sont comme dans les exemples \ref{example-proper-over-local} et \ref{example-equidimensional-over-field}.

\medskip\noindent
Les assertions (1), (2) et (3) résultent du lemme \ref{lemma-vanishing-good-dualizing}.

\medskip\noindent
Pour un point fermé $x \in X$, on voit que $\omega_{X, x}^\bullet$ est un complexe dualisant normalisé sur $\mathcal{O}_{X, x}$, voir le lemme \ref{lemma-good-dualizing-normalized}. L'assertion (4) résulte alors de Complexes dualisants, lemme \ref{dualizing-lemma-depth-in-terms-dualizing-complex}.

\medskip\noindent
L'assertion (5) vaut par construction, puisque $a$ est l'adjoint à droite de $Rf_* : D_\QCoh(\mathcal{O}_X) \to D(\mathcal{O}_{\Spec(k)}) = D(k)$, que l'on peut identifier à $K \mapsto R\Gamma(X, K)$. On utilise également que la catégorie dérivée $D(k)$ des $k$-modules est la même que la catégorie des $k$-espaces vectoriels gradués.

\medskip\noindent
L'assertion (6) résulte du lemme \ref{lemma-dualizing-module-proper-over-A} pour $\mathcal{F}$ cohérent, et, en général, en explicitant (5) pour $K = \mathcal{F}[0]$ et $i = -\dim(X)$.

\medskip\noindent
L'assertion (7) résulte du lemme \ref{lemma-smooth-proper}.
\end{proof}

\begin{remark}
\label{remark-duality-proper-over-field}
Soient $k$, $X$ et $\omega_X^\bullet$ comme dans le lemme \ref{lemma-duality-proper-over-field}. L'identité du complexe $\omega_X^\bullet$ correspond, via l'isomorphisme fonctoriel de l'assertion (5), à un morphisme
$$
t : H^0(X, \omega_X^\bullet) \longrightarrow k
$$
Pour un $K$ arbitraire dans $D_\QCoh(\mathcal{O}_X)$, l'identification
\par\noindent\hfil
$\Hom(K, \omega_X^\bullet)$ à $H^0(X, K)^\vee$
\hfil\par\noindent
dans l'assertion (5) correspond à l'accouplement
$$
\Hom_X(K, \omega_X^\bullet) \times H^0(X, K) \longrightarrow k,\quad
(\alpha, \beta) \longmapsto t(\alpha(\beta))
$$
Cela résulte de la fonctorialité des isomorphismes de (5). De même, pour tout $i \in \mathbf{Z}$, on obtient l'accouplement
$$
\Ext^i_X(K, \omega_X^\bullet) \times H^{-i}(X, K) \longrightarrow k,\quad
(\alpha, \beta) \longmapsto t(\alpha(\beta))
$$
Ici, on considère $\alpha$ comme un morphisme $K[-i] \to \omega_X^\bullet$ et $\beta$ comme un élément de $H^0(X, K[-i])$ afin de définir $\alpha(\beta)$. Remarquons que si $K$ est général, on sait seulement que cet accouplement est non dégénéré d'un côté : l'accouplement induit un isomorphisme de $\Hom_X(K, \omega_X^\bullet)$, resp.\ $\Ext^i_X(K, \omega_X^\bullet)$, sur le dual $k$-linéaire de $H^0(X, K)$, resp.\ $H^{-i}(X, K)$, mais en général pas réciproquement. Si $K$ appartient à $D^b_{\textit{Coh}}(\mathcal{O}_X)$, alors $\Hom_X(K, \omega_X^\bullet)$, $\Ext_X(K, \omega_X^\bullet)$, $H^0(X, K)$ et $H^i(X, K)$ sont des $k$-espaces vectoriels de dimension finie (d'après Catégories dérivées de schémas, lemmes \ref{perfect-lemma-coherent-internal-hom} et \ref{perfect-lemma-direct-image-coherent-bdd-below}), et les accouplements sont parfaits au sens usuel.
\end{remark}

\begin{remark}
\label{remark-coherent-duality-proper-over-field}
Nous poursuivons la discussion de la remarque \ref{remark-duality-proper-over-field} et conservons les mêmes notations $k$, $X$, $\omega_X^\bullet$ et $t$. Si $\mathcal{F}$ est un $\mathcal{O}_X$-module cohérent, on obtient des accouplements parfaits
$$
\langle -, - \rangle :
\Ext^i_X(\mathcal{F}, \omega_X^\bullet) \times H^{-i}(X,\mathcal{F})
\longrightarrow k,\quad
(\alpha, \beta) \longmapsto t(\alpha(\beta))
$$
d'espaces vectoriels de dimension finie sur $k$. Ces accouplements satisfont la fonctorialité (évidente) suivante : si $\varphi : \mathcal{F} \to \mathcal{G}$ est un homomorphisme de $\mathcal{O}_X$-modules cohérents, alors on a
$$
\langle \alpha \circ \varphi, \beta \rangle =
\langle \alpha, \varphi(\beta) \rangle
$$
pour $\alpha \in \Ext^i_X(\mathcal{G}, \omega_X^\bullet)$ et $\beta \in H^{-i}(X, \mathcal{F})$. Autrement dit, l'application $k$-linéaire $\Ext^i_X(\mathcal{G}, \omega_X^\bullet) \to
\Ext^i_X(\mathcal{F}, \omega_X^\bullet)$ induite par $\varphi$ est, via les accouplements, le dual $k$-linéaire de l'application $k$-linéaire $H^{-i}(X, \mathcal{F}) \to H^{-i}(X, \mathcal{G})$ induite par $\varphi$. Formulé ainsi, cela reste valable si $\varphi$ est un homomorphisme de $\mathcal{O}_X$-modules quasi-cohérents.
\end{remark}

\begin{lemma}
\label{lemma-duality-proper-over-field-perfect}
Soient $k$, $X$ et $\omega_X^\bullet$ comme dans le lemme \ref{lemma-duality-proper-over-field}. Soit $t : H^0(X, \omega_X^\bullet) \to k$ comme dans la remarque \ref{remark-duality-proper-over-field}. Supposons $E \in D(\mathcal{O}_X)$ parfait. Alors les accouplements
$$
H^i(X, \omega_X^\bullet \otimes_{\mathcal{O}_X}^\mathbf{L} E^\vee)
\times
H^{-i}(X, E) \longrightarrow k, \quad
(\xi, \eta) \longmapsto
t((1_{\omega_X^\bullet} \otimes \epsilon)(\xi \cup \eta))
$$
sont parfaits pour tout $i$. Ici, $\cup$ désigne le cup-produit de Cohomologie, section \ref{cohomology-section-cup-product}, et $\epsilon : E^\vee \otimes_{\mathcal{O}_X}^\mathbf{L} E \to \mathcal{O}_X$ est comme dans Cohomologie, exemple \ref{cohomology-example-dual-derived}.
\end{lemma}

\begin{proof}
En remplaçant $E$ par $E[-i]$, on se ramène au cas $i = 0$. D'après Cohomologie, lemme \ref{cohomology-lemma-ext-composition-is-cup}, l'accouplement est le même que celui étudié dans la remarque \ref{remark-duality-proper-over-field}, d'où le résultat par la discussion de cette remarque.
\end{proof}

\begin{lemma}
\label{lemma-duality-proper-over-field-CM}
Soit $X$ un schéma propre sur un corps $k$, de Cohen–Macaulay et équidimensionnel de dimension $d$. Le module $\omega_X$ du lemme \ref{lemma-duality-proper-over-field} possède les propriétés suivantes :
\begin{enumerate}
\item $\omega_X$ est un module dualisant sur $X$ (section \ref{section-dualizing-module}) ;
\item $\omega_X$ est un module cohérent de Cohen–Macaulay dont le support est $X$ ;
\item il existe des isomorphismes fonctoriels
\par\noindent\hfil
$\Ext^i_X(K, \omega_X[d]) = \Hom_k(H^{-i}(X, K), k)$
\hfil\par\noindent
compatibles aux décalages et aux triangles distingués, pour $K \in D_\QCoh(X)$ ;
\item il existe des isomorphismes fonctoriels
$\Ext^{d - i}(\mathcal{F}, \omega_X) = \Hom_k(H^i(X, \mathcal{F}), k)$
pour $\mathcal{F}$ quasi-cohérent sur $X$.
\end{enumerate}
\end{lemma}

\begin{proof}
Il résulte clairement du lemme \ref{lemma-duality-proper-over-field} que $\omega_X$ est un module dualisant (puisqu'il s'agit du faisceau de cohomologie non nul le plus à gauche d'un complexe dualisant). On a $\omega_X^\bullet = \omega_X[d]$, et $\omega_X$ est de Cohen–Macaulay puisque $X$ est de Cohen–Macaulay, voir le lemme \ref{lemma-dualizing-module-CM-scheme}. Les autres assertions en résultent, combinées aux assertions correspondantes du lemme \ref{lemma-duality-proper-over-field}.
\end{proof}

\begin{remark}
\label{remark-rework-duality-locally-free-CM}
Soit $X$ un schéma propre de Cohen–Macaulay sur un corps $k$, équidimensionnel de dimension $d$. Soient $\omega_X^\bullet$ et $\omega_X$ comme dans le lemme \ref{lemma-duality-proper-over-field}. D'après le lemme \ref{lemma-duality-proper-over-field-CM}, on a $\omega_X^\bullet = \omega_X[d]$. Soit $t : H^d(X, \omega_X) \to k$ le morphisme de la remarque \ref{remark-duality-proper-over-field}. Soit $\mathcal{E}$ un $\mathcal{O}_X$-module localement libre de type fini, de dual $\mathcal{E}^\vee$. On a alors des accouplements parfaits
$$
H^i(X, \omega_X \otimes_{\mathcal{O}_X} \mathcal{E}^\vee)
\times
H^{d - i}(X, \mathcal{E})
\longrightarrow
k,\quad
(\xi, \eta) \longmapsto t(1 \otimes \epsilon)(\xi \cup \eta))
$$
où $\cup$ est le cup-produit et $\epsilon : \mathcal{E}^\vee \otimes_{\mathcal{O}_X} \mathcal{E}
\to \mathcal{O}_X$ le morphisme d'évaluation. C'est un cas particulier du lemme \ref{lemma-duality-proper-over-field-perfect}.
\end{remark}

\noindent
Voici une vérification de cohérence pour le complexe dualisant.

\begin{lemma}
\label{lemma-sanity-check-duality}
Soit $X$ un schéma propre sur un corps $k$. Soient $\omega_X^\bullet$ et $\omega_X$ comme dans le lemme \ref{lemma-duality-proper-over-field}.
\begin{enumerate}
\item Si $X \to \Spec(k)$ se factorise sous la forme $X \to \Spec(k') \to \Spec(k)$ pour un corps $k'$, alors $\omega_X^\bullet$ et $\omega_X$ sont comme dans le lemme \ref{lemma-duality-proper-over-field} pour le morphisme $X \to \Spec(k')$.
\item Si $K/k$ est une extension de corps, alors les images inverses de $\omega_X^\bullet$ et $\omega_X$ sur le changement de base $X_K$ sont comme dans le lemme \ref{lemma-duality-proper-over-field} pour le morphisme $X_K \to \Spec(K)$.
\end{enumerate}
\end{lemma}

\begin{proof}
Notons $f : X \to \Spec(k)$ le morphisme structural et $f' : X \to \Spec(k')$ la factorisation donnée. Dans la démonstration du lemme \ref{lemma-duality-proper-over-field}, on a pris $\omega_X^\bullet = a(\mathcal{O}_{\Spec(k)})$, où $a$ est l'adjoint à droite du lemme \ref{lemma-twisted-inverse-image} pour $f$. Il faut donc montrer
$a(\mathcal{O}_{\Spec(k)}) \cong a'(\mathcal{O}_{\Spec(k)})$
où $a'$ est l'adjoint à droite du lemme \ref{lemma-twisted-inverse-image} pour $f'$. Puisque $k' \subset H^0(X, \mathcal{O}_X)$, on voit que $k'/k$ est une extension finie (Cohomologie des schémas, lemme \ref{coherent-lemma-proper-over-affine-cohomology-finite}). Par unicité des adjoints, on a $a = a' \circ b$, où $b$ est l'adjoint à droite du lemme \ref{lemma-twisted-inverse-image} pour $g : \Spec(k') \to \Spec(k)$. On peut aussi l'exprimer ainsi : on a $f^! = (f')^! \circ g^!$. Il suffit donc de montrer que $\Hom_k(k', k) \cong k'$ comme $k'$-modules, voir l'exemple \ref{example-affine-twisted-inverse-image}. Cela vaut parce que ce sont des $k'$-espaces vectoriels de même dimension (à savoir de dimension $1$).

\medskip\noindent
Démonstration de (2). Cela vaut parce que l'on dispose du changement de base pour $a$ d'après le lemme \ref{lemma-more-base-change}. Voir la discussion de la remarque \ref{remark-relative-dualizing-complex}.
\end{proof}








\section{Complexes dualisants relatifs}
\label{section-relative-dualizing-complexes}

\noindent
Pour un morphisme propre, plat et de présentation finie, on dispose d'un complexe dualisant relatif rigide, voir la remarque \ref{remark-relative-dualizing-complex} et le lemme \ref{lemma-van-den-bergh}. Pour un morphisme séparé et de type fini $f : X \to Y$ de schémas noethériens, on peut considérer $f^!\mathcal{O}_Y$. Dans cette section, nous définissons les complexes dualisants relatifs pour les morphismes plats et localement de présentation finie (mais pas nécessairement quasi-séparés ni quasi-compacts) entre schémas (pas nécessairement localement noethériens). Nous montrons que ces complexes existent, sont uniques à isomorphisme unique près et coïncident avec ceux des cas mentionnés ci-dessus. Avant de lire cette section, prière de lire Complexes dualisants, section \ref{dualizing-section-relative-dualizing-complexes}.

\begin{definition}
\label{definition-relative-dualizing-complex}
Soit $X \to S$ un morphisme de schémas plat et localement de présentation finie. Soit $W \subset X \times_S X$ un ouvert quelconque tel que la diagonale $\Delta_{X/S} : X \to X \times_S X$ se factorise par une immersion fermée $\Delta : X \to W$. Un {\it complexe dualisant relatif} est un couple $(K, \xi)$ formé d'un objet $K \in D(\mathcal{O}_X)$ et d'un morphisme
$$
\xi : \Delta_*\mathcal{O}_X \longrightarrow L\text{pr}_1^*K|_W
$$
dans $D(\mathcal{O}_W)$ tels que
\begin{enumerate}
\item $K$ soit $S$-parfait (Catégories dérivées de schémas, définition \ref{perfect-definition-relatively-perfect}), et que
\item $\xi$ définisse un isomorphisme de $\Delta_*\mathcal{O}_X$ sur
$R\SheafHom_{\mathcal{O}_W}(
\Delta_*\mathcal{O}_X, L\text{pr}_1^*K|_W)$.
\end{enumerate}
\end{definition}

\noindent
\begingroup\emergencystretch=4em
D'après le lemme \ref{lemma-sheaf-with-exact-support-ext}, la condition (2) équivaut à l'existence d'un isomorphisme
$$
\mathcal{O}_X \longrightarrow
R\SheafHom(\mathcal{O}_X, L\text{pr}_1^*K|_W)
$$
dans $D(\mathcal{O}_X)$ dont l'image directe par $\Delta$ est égale à $\xi$. Puisque
\par\noindent\hfil
$R\SheafHom(\mathcal{O}_X, L\text{pr}_1^*K|_W)$
\hfil\par\noindent
est indépendant du choix de l'ouvert $W$, il en est de même de la catégorie des couples $(K, \xi)$. Si $X \to S$ est séparé, on peut choisir $W = X \times_S X$. Nous ramènerons nombre des arguments au cas des anneaux au moyen du lemme suivant.\par\endgroup

\begin{lemma}
\label{lemma-relative-dualizing-complex-algebra}
Soit $X \to S$ un morphisme de schémas plat et localement de présentation finie. Soit $(K, \xi)$ un complexe dualisant relatif. Alors, pour tout diagramme commutatif
$$
\xymatrix{
\Spec(A) \ar[d] \ar[r] & X \ar[d] \\
\Spec(R) \ar[r] & S
}
$$
dont les flèches horizontales sont des immersions ouvertes, la restriction de $K$ à $\Spec(A)$ correspond, via Catégories dérivées de schémas, lemme \ref{perfect-lemma-affine-compare-bounded}, à un complexe dualisant relatif pour $R \to A$ au sens de Complexes dualisants, définition \ref{dualizing-definition-relative-dualizing-complex}.
\end{lemma}

\begin{proof}
Puisque la formation de $R\SheafHom$ commute aux restrictions aux ouverts, on peut aussi bien supposer $X = \Spec(A)$ et $S = \Spec(R)$. Remarquons que les objets relativement parfaits de $D(\mathcal{O}_X)$ sont pseudo-cohérents et appartiennent donc à $D_\QCoh(\mathcal{O}_X)$ (Catégories dérivées de schémas, lemme \ref{perfect-lemma-pseudo-coherent}). L'énoncé a donc un sens. Remarquons que les constructions de $\Delta_*$, $L\text{pr}_1^*$ et $R\SheafHom$ sont compatibles à ce qui se passe du côté algébrique d'après Catégories dérivées de schémas, lemmes \ref{perfect-lemma-quasi-coherence-pushforward}, \ref{perfect-lemma-quasi-coherence-pullback} et \ref{perfect-lemma-quasi-coherence-internal-hom}. Pour le dernier, on remarque que $L\text{pr}_1^*K$ est $S$-parfait (donc borné inférieurement) et que $\Delta_*\mathcal{O}_X$ est un objet pseudo-cohérent de $D(\mathcal{O}_W)$ ; traduit en algèbre, cela signifie que $A$ est pseudo-cohérent comme $A \otimes_R A$-module, ce qui résulte de Compléments sur l'algèbre, lemme \ref{more-algebra-lemma-more-relative-pseudo-coherent-is-moot}, appliqué à $R \to A \otimes_R A \to A$. On retrouve ainsi exactement les conditions de Complexes dualisants, définition \ref{dualizing-definition-relative-dualizing-complex}.
\end{proof}

\begin{lemma}
\label{lemma-relative-dualizing-RHom}
Soit $X \to S$ un morphisme de schémas plat et localement de présentation finie. Soit $(K, \xi)$ un complexe dualisant relatif. Alors
$\mathcal{O}_X \to R\SheafHom_{\mathcal{O}_X}(K, K)$
est un isomorphisme.
\end{lemma}

\begin{proof}
En raisonnant localement sur les ouverts affines, le lemme \ref{lemma-relative-dualizing-complex-algebra} ramène l'énoncé au cas algébrique, qui est Complexes dualisants, lemme \ref{dualizing-lemma-relative-dualizing-RHom}.
\end{proof}

\begin{lemma}
\label{lemma-uniqueness-relative-dualizing}
Soit $X \to S$ un morphisme de schémas plat et localement de présentation finie. Si $(K, \xi)$ et $(L, \eta)$ sont deux complexes dualisants relatifs sur $X/S$, il existe un unique isomorphisme $K \to L$ qui envoie $\xi$ sur $\eta$.
\end{lemma}

\begin{proof}
Soit $U \subset X$ un ouvert affine dont l'image est contenue dans un ouvert affine de $S$. Il existe alors un isomorphisme $K|_U \to L|_U$ d'après le lemme \ref{lemma-relative-dualizing-complex-algebra} et Complexes dualisants, lemme \ref{dualizing-lemma-uniqueness-relative-dualizing}. Le lecteur peut reprendre l'argument de ce lemme dans le cas des schémas pour obtenir une démonstration dans ce cas-ci. Nous utiliserons plutôt un argument de recollement.

\medskip\noindent
Supposons donné un isomorphisme $\alpha : K \to L$. Alors $\alpha(\xi) = u \eta$ pour une section inversible $u \in H^0(W, \Delta_*\mathcal{O}_X) = H^0(X, \mathcal{O}_X)$. (En effet, $\eta$ et $\alpha(\xi)$ sont tous deux des générateurs d'un $\Delta_*\mathcal{O}_X$-module inversible par hypothèse.) Après avoir remplacé $\alpha$ par $u^{-1}\alpha$, on voit donc que $\alpha(\xi) = \eta$. Puisque le groupe des automorphismes de $K$ est $H^0(X, \mathcal{O}_X^*)$ d'après le lemme \ref{lemma-relative-dualizing-RHom}, il existe au plus un tel $\alpha$.

\medskip\noindent
Soit $\mathcal{B}$ la collection des ouverts affines de $X$ dont l'image est contenue dans un ouvert affine de $S$. Pour chaque $U \in \mathcal{B}$, les deux paragraphes précédents fournissent un unique isomorphisme $\alpha_U : K|_U \to L|_U$ envoyant $\xi$ sur $\eta$. Remarquons que $\text{Ext}^i(K|_U, K|_U) = 0$ pour $i < 0$ et pour tout ouvert $U$ de $X$ d'après le lemme \ref{lemma-relative-dualizing-RHom}. En appliquant Cohomologie, lemme \ref{cohomology-lemma-vanishing-and-glueing}, à $\text{id} : X \to X$, on obtient un unique morphisme $\alpha : K \to L$ qui coïncide avec $\alpha_U$ pour tout $U \in \mathcal{B}$. Alors $\alpha$ envoie $\xi$ sur $\eta$ puisque c'est vrai localement.
\end{proof}

\begin{lemma}
\label{lemma-existence-relative-dualizing}
Soit $X \to S$ un morphisme de schémas plat et localement de présentation finie. Il existe un complexe dualisant relatif $(K, \xi)$.
\end{lemma}

\begin{proof}
Soit $\mathcal{B}$ la collection des ouverts affines de $X$ dont l'image est contenue dans un ouvert affine de $S$. Pour chaque $U$, on dispose d'un complexe dualisant relatif $(K_U, \xi_U)$ pour $U$ sur $S$. En effet, choisissons un ouvert affine $V \subset S$ tel que $U \to X \to S$ se factorise par $V$. Écrivons $U = \Spec(A)$ et $V = \Spec(R)$. D'après Complexes dualisants, lemme \ref{dualizing-lemma-base-change-relative-dualizing}, il existe un complexe dualisant relatif $K_A \in D(A)$ pour $R \to A$. En remontant l'argument de la démonstration du lemme \ref{lemma-relative-dualizing-complex-algebra}, celui-ci détermine un objet $V$-parfait $K_U \in D(\mathcal{O}_U)$ et un morphisme
$$
\xi : \Delta_*\mathcal{O}_U \to L\text{pr}_1^*K_U
$$
dans $D(\mathcal{O}_{U \times_V U})$. Comme être $V$-parfait équivaut à être $S$-parfait et comme $U \times_V U = U \times_S U$, on trouve que $(K_U, \xi_U)$ convient.

\medskip\noindent
Si $U' \subset U \subset X$ avec $U', U \in \mathcal{B}$, alors il existe un unique isomorphisme $\rho_{U'}^U : K_U|_{U'} \to K_{U'}$ dans $D(\mathcal{O}_{U'})$ envoyant $\xi_U|_{U' \times_S U'}$ sur $\xi_{U'}$ d'après le lemme \ref{lemma-uniqueness-relative-dualizing} (remarquons que la restriction d'un complexe dualisant relatif à un ouvert est trivialement un complexe dualisant relatif). L'unicité garantit que
$\rho^U_{U''} = \rho^V_{U''} \circ \rho ^U_{U'}|_{U''}$
pour $U'' \subset U' \subset U$ dans $\mathcal{B}$. Remarquons que $\text{Ext}^i(K_U, K_U) = 0$ pour $i < 0$, pour $U \in \mathcal{B}$, d'après le lemme \ref{lemma-relative-dualizing-RHom} appliqué à $U/S$ et $K_U$. Le lemme de recollement BBD (Cohomologie, théorème \ref{cohomology-theorem-glueing-bbd-general}) affirme donc qu'il existe une unique solution, à savoir un objet $K \in D(\mathcal{O}_X)$ et des isomorphismes $\rho_U : K|_U \to K_U$ tels que l'on ait $\rho^U_{U'} \circ \rho_U|_{U'} = \rho_{U'}$ pour tout $U' \subset U$, $U, U' \in \mathcal{B}$.

\medskip\noindent
Pour achever la démonstration, il faut construire le morphisme
$$
\xi : \Delta_*\mathcal{O}_X \longrightarrow L\text{pr}_1^*K|_W
$$
dans $D(\mathcal{O}_W)$ induisant un isomorphisme de $\Delta_*\mathcal{O}_X$ sur
\par\noindent\hfil
$R\SheafHom_{\mathcal{O}_W}(\Delta_*\mathcal{O}_X, L\text{pr}_1^*K|_W)$.
\hfil\par\noindent
Puisque l'on peut changer $W$, choisissons $W = \bigcup_{U \in \mathcal{B}} U \times_S U$. On peut utiliser $\rho_U$ pour obtenir des isomorphismes
$$
R\SheafHom_{\mathcal{O}_W}(
\Delta_*\mathcal{O}_X, L\text{pr}_1^*K|_W)|_{U \times_S U}
\xrightarrow{\rho_U}
R\SheafHom_{\mathcal{O}_{U \times_S U}}(
\Delta_*\mathcal{O}_U, L\text{pr}_1^*K_U)
$$
Comme $W$ est recouvert par les ouverts $U \times_S U$, on conclut que les faisceaux de cohomologie de
$R\SheafHom_{\mathcal{O}_W}(\Delta_*\mathcal{O}_X, L\text{pr}_1^*K|_W)$
sont nuls sauf en degré $0$. En outre, on obtient des isomorphismes
\begingroup\small
$$
H^0\left(U \times_S U, R\SheafHom_{\mathcal{O}_W}(\Delta_*\mathcal{O}_X,
L\text{pr}_1^*K|_W)\right)
\xrightarrow{\rho_U}
H^0\left((R\SheafHom_{\mathcal{O}_{U \times_S U}}(
\Delta_*\mathcal{O}_U, L\text{pr}_1^*K_U)\right)
$$
\endgroup
Soit $\tau_U$ dans le membre de gauche un élément envoyé sur $\xi_U$ par ce morphisme. Les compatibilités entre $\rho^U_{U'}$, $\xi_U$, $\xi_{U'}$, $\rho_U$ et $\rho_{U'}$ pour $U' \subset U \subset X$ ouvert de $U', U \in \mathcal{B}$ impliquent que $\tau_U|_{U' \times_S U'} = \tau_{U'}$. On obtient ainsi une section globale $\tau$ du $0$-ième faisceau de cohomologie $H^0(R\SheafHom_{\mathcal{O}_W}(\Delta_*\mathcal{O}_X, L\text{pr}_1^*K|_W))$. Puisque les autres faisceaux de cohomologie de
$R\SheafHom_{\mathcal{O}_W}(\Delta_*\mathcal{O}_X, L\text{pr}_1^*K|_W)$
sont nuls, cette section globale $\tau$ détermine comme voulu un morphisme $\xi$. Puisque la restriction de $\xi$ à $U \times_S U$ donne $\xi_U$, on voit qu'il satisfait la condition finale de la définition \ref{definition-relative-dualizing-complex}.
\end{proof}

\begin{lemma}
\label{lemma-base-change-relative-dualizing}
Considérons un carré cartésien
$$
\xymatrix{
X' \ar[d]_{f'} \ar[r]_{g'} & X \ar[d]^f \\
S' \ar[r]^g & S
}
$$
de schémas. Supposons $X \to S$ plat et localement de présentation finie. Soit $(K, \xi)$ un complexe dualisant relatif pour $f$. Posons $K' = L(g')^*K$. Soit $\xi'$ le changement de base dérivé de $\xi$ (voir la démonstration). Alors $(K', \xi')$ est un complexe dualisant relatif pour $f'$.
\end{lemma}

\begin{proof}
Considérons le carré cartésien
$$
\xymatrix{
X' \ar[d]_{\Delta_{X'/S'}} \ar[r] & X \ar[d]^{\Delta_{X/S}} \\
X' \times_{S'} X' \ar[r]^{g' \times g'} & X \times_S X
}
$$
Choisissons un ouvert $W \subset X \times_S X$ tel que $\Delta_{X/S}$ se factorise par une immersion fermée $\Delta : X \to W$. Choisissons un ouvert $W' \subset X' \times_{S'} X'$ tel que $\Delta_{X'/S'}$ se factorise par une immersion fermée $\Delta' : X \to W'$ et tel que $(g' \times g')(W') \subset W$. Notons encore $g' \times g' : W' \to W$ le morphisme induit. On a
$$
L(g' \times g')^*\Delta_*\mathcal{O}_X =
\Delta'_*\mathcal{O}_{X'}
\quad\text{et}\quad
L(g' \times g')^*L\text{pr}_1^*K|_W =
L\text{pr}_1^*K'|_{W'}
$$
La première égalité vaut parce que $X$ et $X' \times_{S'} X'$ sont Tor-indépendants sur $X \times_S X$ (voir par exemple Compléments sur les morphismes, lemme \ref{more-morphisms-lemma-case-of-tor-independence}). La seconde résulte de la transitivité de l'image inverse dérivée (Cohomologie, lemme \ref{cohomology-lemma-derived-pullback-composition}). Ainsi, $\xi' = L(g' \times g')^*\xi$ peut être considéré comme un morphisme
$$
\xi' : \Delta'_*\mathcal{O}_{X'} \longrightarrow L\text{pr}_1^*K'|_{W'}
$$
Cela étant dit, la démonstration du lemme est immédiate. Premièrement, $K'$ est $S'$-parfait d'après Catégories dérivées de schémas, lemme \ref{perfect-lemma-base-change-relatively-perfect}. Pour vérifier que $\xi'$ induit un isomorphisme de $\Delta'_*\mathcal{O}_{X'}$ sur
$R\SheafHom_{\mathcal{O}_{W'}}(
\Delta'_*\mathcal{O}_{X'}, L\text{pr}_1^*K'|_{W'})$
on peut travailler localement sur les ouverts affines. Le lemme \ref{lemma-relative-dualizing-complex-algebra} ramène au résultat correspondant en algèbre, démontré dans Complexes dualisants, lemme \ref{dualizing-lemma-base-change-relative-dualizing}.
\end{proof}

\begin{lemma}
\label{lemma-flat-proper-relative-dualizing}
Soit $S$ un schéma quasi-compact et quasi-séparé. Soit $f : X \to S$ un morphisme propre et plat de présentation finie. Le complexe dualisant relatif $\omega_{X/S}^\bullet$ de la remarque \ref{remark-relative-dualizing-complex}, muni de (\ref{equation-pre-rigid}), est un complexe dualisant relatif au sens de la définition \ref{definition-relative-dualizing-complex}.
\end{lemma}

\begin{proof}
Dans le lemme \ref{lemma-properties-relative-dualizing}, on a démontré que $\omega_{X/S}^\bullet$ est $S$-parfait. Soit $c$ l'adjoint à droite du lemme \ref{lemma-twisted-inverse-image} pour la diagonale $\Delta : X \to X \times_S X$. On peut alors appliquer $\Delta_*$ à (\ref{equation-pre-rigid}) pour obtenir un isomorphisme
$$
\Delta_*\mathcal{O}_X \to
\Delta_*(c(L\text{pr}_1^*\omega_{X/S}^\bullet)) =
R\SheafHom_{\mathcal{O}_{X \times_S X}}(
\Delta_*\mathcal{O}_X, L\text{pr}_1^*\omega_{X/S}^\bullet)
$$
L'égalité résulte des lemmes \ref{lemma-twisted-inverse-image-closed} et \ref{lemma-sheaf-with-exact-support-ext}. Cela achève la démonstration.
\end{proof}

\begin{remark}
\label{remark-relative-dualizing-complex-bis}
Soit $X \to S$ un morphisme de schémas plat, propre et de présentation finie. D'après le lemme \ref{lemma-existence-relative-dualizing}, il existe un complexe dualisant relatif $(\omega_{X/S}^\bullet, \xi)$ au sens de la définition \ref{definition-relative-dualizing-complex}. Considérons un morphisme quelconque $g : S' \to S$, où $S'$ est quasi-compact et quasi-séparé (par exemple, un ouvert affine de $S$). Le lemme \ref{lemma-base-change-relative-dualizing} montre que $(L(g')^*\omega_{X/S}^\bullet, L(g')^*\xi)$ est un complexe dualisant relatif pour le changement de base $f' : X' \to S'$ au sens de la définition \ref{definition-relative-dualizing-complex}. Soit $\omega_{X'/S'}^\bullet$ le complexe dualisant relatif pour $X' \to S'$ au sens de la remarque \ref{remark-relative-dualizing-complex}. En combinant les lemmes \ref{lemma-flat-proper-relative-dualizing} et \ref{lemma-uniqueness-relative-dualizing}, on voit qu'il existe un unique isomorphisme
$$
\omega_{X'/S'}^\bullet \longrightarrow L(g')^*\omega_{X/S}^\bullet
$$
compatible à (\ref{equation-pre-rigid}) et à $L(g')^*\xi$. Ces isomorphismes sont compatibles aux morphismes entre schémas quasi-compacts et quasi-séparés sur $S$ et aux isomorphismes de changement de base du lemme \ref{lemma-proper-flat-base-change} (si nous avons un jour besoin de cette compatibilité, nous l'énoncerons et la démontrerons soigneusement ici).
\end{remark}

\begin{lemma}
\label{lemma-compactifyable-relative-dualizing}
Dans la situation \ref{situation-shriek}, soit $f : X \to Y$ un morphisme de $\textit{FTS}_S$. Si $f$ est plat, alors $f^!\mathcal{O}_Y$ est (la première composante d')un complexe dualisant relatif pour $X$ sur $Y$ au sens de la définition \ref{definition-relative-dualizing-complex}.
\end{lemma}

\begin{proof}
D'après le lemme \ref{lemma-flat-shriek-relatively-perfect}, on sait que $f^!\mathcal{O}_Y$ est $Y$-parfait. Comme $f$ est séparé, la diagonale $\Delta : X \to X \times_Y X$ est une immersion fermée et $\Delta_*\Delta^!(-) =
R\SheafHom_{\mathcal{O}_{X \times_Y X}}(\mathcal{O}_X, -)$, voir les lemmes \ref{lemma-twisted-inverse-image-closed} et \ref{lemma-sheaf-with-exact-support-ext}. Pour achever la démonstration, il suffit donc de montrer
$\Delta^!(L\text{pr}_1^*f^!(\mathcal{O}_Y)) \cong \mathcal{O}_X$
où $\text{pr}_1 : X \times_Y X \to X$ est la première projection. On a
$$
\mathcal{O}_X = \Delta^! \text{pr}_1^!\mathcal{O}_X =
\Delta^! \text{pr}_1^! L\text{pr}_2^*\mathcal{O}_Y =
\Delta^!(L\text{pr}_1^* f^!\mathcal{O}_Y)
$$
où $\text{pr}_2 : X \times_Y X \to X$ est la seconde projection et où l'on a utilisé l'isomorphisme de changement de base $\text{pr}_1^! \circ L\text{pr}_2^* = L\text{pr}_1^* \circ f^!$ du lemme \ref{lemma-base-change-shriek-flat}.
\end{proof}

\begin{lemma}
\label{lemma-relative-dualizing-composition}
Soient $f : Y \to X$ et $X \to S$ des morphismes de schémas plats et de présentation finie. Soient $(K, \xi)$ et $(M, \eta)$ des complexes dualisants relatifs pour $X \to S$ et $Y \to X$. Posons $E = M \otimes_{\mathcal{O}_Y}^\mathbf{L} Lf^*K$. Alors $(E, \zeta)$ est un complexe dualisant relatif pour $Y \to S$ pour un $\zeta$ convenable.
\end{lemma}

\begin{proof}
En utilisant le lemme \ref{lemma-relative-dualizing-complex-algebra} et la version algébrique du présent lemme (Complexes dualisants, lemme \ref{dualizing-lemma-relative-dualizing-composition}), on voit que $E$ est, localement sur les ouverts affines, la première composante d'un complexe dualisant relatif. En particulier, on voit que $E$ est $S$-parfait, puisque cela se vérifie localement sur les ouverts affines, voir Catégories dérivées de schémas, lemme \ref{perfect-lemma-affine-locally-rel-perfect}.

\medskip\noindent
Démontrons d'abord l'existence de $\zeta$ lorsque les morphismes $X \to S$ et $Y \to X$ sont séparés, de sorte que $\Delta_{X/S}$, $\Delta_{Y/X}$ et $\Delta_{Y/S}$ sont des immersions fermées. Considérons le diagramme suivant
$$
\xymatrix{
& & Y \ar@{=}[r] & Y \ar[d]^f \\
Y \ar[r]_{\Delta_{Y/X}} &
Y \times_X Y \ar[d]_m \ar[r]_\delta \ar[ru]_q &
Y \times_S Y \ar[d]^{f \times f} \ar[ru]_p & X\\
& X \ar[r]^{\Delta_{X/S}} & X \times_S X \ar[ru]_r
}
$$
où $p$, $q$ et $r$ sont les premières projections. D'après le lemme \ref{lemma-sheaf-with-exact-support-internal-home}, on a
\begingroup\footnotesize
$$
R\SheafHom_{\mathcal{O}_{Y \times_S Y}}(
\Delta_{Y/S, *}\mathcal{O}_Y, Lp^*E) =
R\delta_*\left(R\SheafHom_{\mathcal{O}_{Y \times_X Y}}(
\Delta_{Y/X, *}\mathcal{O}_Y,
R\SheafHom(\mathcal{O}_{Y \times_X Y}, Lp^*E))\right)
$$
\endgroup
D'après le lemme \ref{lemma-sheaf-with-exact-support-tensor}, on a
$$
R\SheafHom(\mathcal{O}_{Y \times_X Y}, Lp^*E) =
R\SheafHom(\mathcal{O}_{Y \times_X Y}, L(f \times f)^*Lr^*K)
\otimes_{\mathcal{O}_{Y \times_S Y}}^\mathbf{L} Lq^*M
$$
D'après le lemme \ref{lemma-flat-bc-sheaf-with-exact-support}, on a
$$
R\SheafHom(\mathcal{O}_{Y \times_X Y}, L(f \times f)^*Lr^*K) =
Lm^*R\SheafHom(\mathcal{O}_X, Lr^*K)
$$
La dernière expression est isomorphe (via $\xi$) à $Lm^*\mathcal{O}_X = \mathcal{O}_{Y \times_X Y}$. Par conséquent, l'expression qui la précède est isomorphe à $Lq^*M$. Ainsi,
$$
R\SheafHom_{\mathcal{O}_{Y \times_S Y}}(
\Delta_{Y/S, *}\mathcal{O}_Y, Lp^*E) =
R\delta_*\left(R\SheafHom_{\mathcal{O}_{Y \times_X Y}}(
\Delta_{Y/X, *}\mathcal{O}_Y, Lq^*M)\right)
$$
Le terme entre parenthèses est isomorphe à $\Delta_{Y/X, *}*\mathcal{O}_X$ via $\eta$. Cela achève la démonstration dans le cas séparé.

\medskip\noindent
Dans le cas général, on choisit un ouvert $W \subset X \times_S X$ tel que $\Delta_{X/S}$ se factorise par une immersion fermée $\Delta : X \to W$, et un ouvert $V \subset Y \times_X Y$ tel que $\Delta_{Y/X}$ se factorise par une immersion fermée $\Delta' : Y \to V$. Enfin, choisissons un ouvert $W' \subset Y \times_S Y$ dont l'intersection avec $Y \times_X Y$ donne $V$ et dont l'image est contenue dans $W$. On considère alors le diagramme
$$
\xymatrix{
& & Y \ar@{=}[r] & Y \ar[d]^f \\
Y \ar[r]_{\Delta'} &
V \ar[d]_m \ar[r]_\delta \ar[ru]_q &
W' \ar[d]^{f \times f} \ar[ru]_p & X\\
& X \ar[r]^\Delta & W \ar[ru]_r
}
$$
et l'on utilise exactement le même argument que précédemment.
\end{proof}






\section{La classe fondamentale d'un morphisme localement d'intersection complète}
\label{section-fundamental-class}

\noindent
Dans cette section, nous utiliserons les calculs effectués dans la section \ref{section-examples}. Notre résultat souffrira donc du même type de défaut d'unicité que celui rencontré dans cette section.

\begin{lemma}
\label{lemma-determinant}
Soit $X$ un espace localement annelé. Soit
$$
\mathcal{E}_1 \xrightarrow{\alpha} \mathcal{E}_0 \to \mathcal{F} \to 0
$$
une suite exacte courte de $\mathcal{O}_X$-modules. Supposons $\mathcal{E}_1$ et $\mathcal{E}_0$ localement libres de rangs $r_1, r_0$. Il existe alors un morphisme canonique
$$
\wedge^{r_0 - r_1}\mathcal{F}
\longrightarrow
\wedge^{r_1}(\mathcal{E}_1^\vee) \otimes \wedge^{r_0}\mathcal{E}_0
$$
qui est un isomorphisme sur la fibre en $x \in X$ si et seulement si $\mathcal{F}$ est localement libre de rang $r_0 - r_1$ dans un voisinage ouvert de $x$.
\end{lemma}

\begin{proof}
Si $r_1 > r_0$, alors $\wedge^{r_0 - r_1}\mathcal{F} = 0$ par convention, et l'unique morphisme ne peut être un isomorphisme. On peut donc supposer $r = r_0 - r_1 \geq 0$. Définissons le morphisme par la formule
$$
s_1 \wedge \ldots \wedge s_r \mapsto
t_1^\vee \wedge \ldots \wedge t_{r_1}^\vee \otimes
\alpha(t_1) \wedge \ldots \wedge \alpha(t_{r_1}) \wedge
\tilde s_1 \wedge \ldots \wedge \tilde s_r
$$
où $t_1, \ldots, t_{r_1}$ est une base locale de $\mathcal{E}_1$, où, en conséquence, $t_1^\vee, \ldots, t_{r_1}^\vee$ est la base duale de $\mathcal{E}_1^\vee$, et où $s'_i$ est un relèvement local de $s_i$ en une section de $\mathcal{E}_0$. Nous omettons la vérification que cette définition ne dépend pas des choix.

\medskip\noindent
Si $\mathcal{F}$ est localement libre de rang $r$, il est immédiat de vérifier que le morphisme est un isomorphisme. Réciproquement, supposons que le morphisme soit un isomorphisme sur les fibres en $x$. Alors $\wedge^r\mathcal{F}_x$ est inversible. Cela implique que $\mathcal{F}_x$ est engendré par au plus $r$ éléments. Cela n'est possible que si $\alpha$ est de rang $r$ modulo $\mathfrak m_x$, c'est-à-dire si $\alpha$ est de rang maximal modulo $\mathfrak m_x$. Il s'ensuit que $\alpha$ est de rang maximal dans un voisinage de $x$ et donc que $\mathcal{F}$ est localement libre de rang $r$ dans un voisinage, comme souhaité.
\end{proof}

\begin{lemma}
\label{lemma-fundamental-class-lci}
Soit $Y$ un schéma noethérien. Soit $f : X \to Y$ un morphisme localement d'intersection complète qui se factorise en une immersion $X \to P$ suivie d'un morphisme propre et lisse $P \to Y$. Soit $r$ la fonction localement constante sur $X$ telle que $\omega_{X/Y} = H^{-r}(f^!\mathcal{O}_Y)$ soit l'unique faisceau de cohomologie non nul de $f^!\mathcal{O}_Y$, voir le lemme \ref{lemma-lci-shriek}. Il existe alors un morphisme
$$
\wedge^r\Omega_{X/Y} \longrightarrow \omega_{X/Y}
$$
qui est un isomorphisme sur la fibre en un point $x$ si et seulement si $f$ est lisse en $x$.
\end{lemma}

\begin{proof}
L'hypothèse implique que $X$ est compactifiable sur $Y$ ; ainsi $f^!$ est défini, voir la section \ref{section-upper-shriek}. Soit $j : W \to P$ un sous-schéma ouvert tel que $X \to P$ se factorise par une immersion fermée $i : X \to W$. En outre, on a $f^! = i^! \circ j^! \circ g^!$, où $g : P \to Y$ est le morphisme donné. On a $g^!\mathcal{O}_Y = \wedge^d\Omega_{P/Y}[d]$ d'après le lemme \ref{lemma-smooth-proper}, où $d$ est la fonction localement constante donnant la dimension relative de $P$ sur $Y$. On a $j^! = j^*$. On a $i^!\mathcal{O}_W = \wedge^c\mathcal{N}[-c]$, où $c$ est la codimension de $X$ dans $W$ (une fonction localement constante sur $X$) et où $\mathcal{N}$ est le faisceau normal de l'immersion régulière au sens de Koszul $i$, voir le lemme \ref{lemma-regular-immersion}. En combinant ce qui précède, on trouve
$$
f^!\mathcal{O}_Y =
\left(\wedge^c\mathcal{N} \otimes_{\mathcal{O}_X}
\wedge^d\Omega_{P/Y}|_X\right)[d - c]
$$
où l'on a également utilisé le lemme \ref{lemma-perfect-comparison-shriek}. Ainsi, $r = d|_X - c$ comme fonctions localement constantes sur $X$. Le faisceau conormal de $X \to P$ est le module $\mathcal{I}/\mathcal{I}^2$, où $\mathcal{I} \subset \mathcal{O}_W$ est le faisceau d'idéaux de $i$, voir Morphismes, section \ref{morphisms-section-conormal-sheaf}. Considérons la suite exacte canonique
$$
\mathcal{I}/\mathcal{I}^2 \to
\Omega_{P/Y}|_X \to \Omega_{X/Y} \to 0
$$
de Morphismes, lemme \ref{morphisms-lemma-differentials-relative-immersion}. Notre morphisme s'obtient en appliquant le lemme \ref{lemma-determinant}.

\medskip\noindent
Si $f$ est lisse en $x$, alors le morphisme est un isomorphisme par application du lemme \ref{lemma-determinant} et parce que $\Omega_{X/Y}$ est localement libre en $x$ de rang $r$. Réciproquement, supposons que notre morphisme soit un isomorphisme sur les fibres en $x$. Le lemme montre alors que $\Omega_{X/Y}$ est libre de rang $r$ après avoir remplacé $X$ par un voisinage ouvert de $x$. D'autre part, on peut également supposer que $X = \Spec(A)$ et $Y = \Spec(R)$, où $A = R[x_1, \ldots, x_n]/(f_1, \ldots, f_m)$ et où $f_1, \ldots, f_m$ est une suite régulière au sens de Koszul (cela résulte de la définition des morphismes localement d'intersection complète). Cela implique clairement $r = n - m$. On conclut que le rang de la matrice des dérivées partielles $\partial f_j/\partial x_i$ dans le corps résiduel en $x$ est $m$. Après avoir réordonné les variables, on peut donc supposer que le déterminant de $(\partial f_j/\partial x_i)_{1 \leq i, j \leq m}$ est inversible dans un voisinage ouvert de $x$. Il s'ensuit que $R \to A$ est lisse en ce point, voir par exemple Algèbre, exemple \ref{algebra-example-make-standard-smooth}.
\end{proof}

\begin{lemma}
\label{lemma-fundamental-class-almost-lci}
Soit $f : X \to Y$ un morphisme de schémas. Soit $r \geq 0$. Supposons que
\begin{enumerate}
\item $Y$ soit de Cohen–Macaulay (Propriétés, définition \ref{properties-definition-Cohen-Macaulay}) ;
\item $f$ se factorise sous la forme $X \to P \to Y$, où le premier morphisme est une immersion et le second est lisse et propre ;
\item si $x \in X$ et $\dim(\mathcal{O}_{X, x}) \leq 1$, alors $f$ soit de Koszul en $x$ (Compléments sur les morphismes, définition \ref{more-morphisms-definition-lci}) ; et
\item si $\xi$ est un point générique d'une composante irréductible de $X$, alors on ait
$\text{trdeg}_{\kappa(f(\xi))} \kappa(\xi) = r$.
\end{enumerate}
Alors, avec $\omega_{X/Y} = H^{-r}(f^!\mathcal{O}_Y)$, il existe un morphisme
$$
\wedge^r\Omega_{X/Y} \longrightarrow \omega_{X/Y}
$$
qui est un isomorphisme sur le lieu où $f$ est lisse.
\end{lemma}

\begin{proof}
Soit $U \subset X$ le sous-schéma ouvert au-dessus duquel $f$ est un morphisme localement d'intersection complète. Puisque $f$ est de dimension relative $r$ en tout point générique par l'hypothèse (4), on voit que la fonction localement constante du lemme \ref{lemma-fundamental-class-lci} est constante de valeur $r$, et l'on obtient un morphisme
$$
\wedge^r\Omega_{X/Y}|_U = \wedge^r \Omega_{U/Y}
\longrightarrow
\omega_{U/Y} = \omega_{X/Y}|_U
$$
qui est un isomorphisme aux points lisses de $f$ (ce lieu est contenu dans $U$, puisqu'un morphisme lisse est localement d'intersection complète). D'après le lemme \ref{lemma-shriek-over-CM} et l'hypothèse que $Y$ est de Cohen–Macaulay, le module $\omega_{X/Y}$ est $(S_2)$. Puisque $U$ contient tous les points de codimension $1$ d'après la condition (3), et en utilisant Diviseurs, lemme \ref{divisors-lemma-depth-2-hartog}, on voit que $j_*\omega_{U/Y} = \omega_{X/Y}$. Par conséquent, le morphisme sur $U$ s'étend à $X$, ce qui achève la démonstration.
\end{proof}






\section{Prolongement par zéro des modules cohérents}
\label{section-extension-by-zero}

\noindent
Le contenu de cette section et des quelques sections suivantes se trouve dans l'appendice de Deligne de \cite{RD}.

\medskip\noindent
Dans cette section, $j : U \to X$ désignera une immersion ouverte de schémas noethériens. Nous allons considérer des systèmes projectifs $(K_n)$ dans $D^b_{\textit{Coh}}(\mathcal{O}_X)$ construits comme suit. Soit $\mathcal{F}^\bullet$ un complexe borné de $\mathcal{O}_X$-modules cohérents. Soit $\mathcal{I} \subset \mathcal{O}_X$ un faisceau quasi-cohérent d'idéaux tel que $V(\mathcal{I}) = X \setminus U$. On peut alors poser
$$
K_n = \mathcal{I}^n\mathcal{F}^\bullet
$$
Plus précisément, $K_n$ est l'objet de $D^b_{\textit{Coh}}(\mathcal{O}_X)$ représenté par le complexe dont le terme en degré $q$ est le sous-module cohérent $\mathcal{I}^n\mathcal{F}^q$ de $\mathcal{F}^q$. Remarquons que les morphismes $\ldots \to K_3 \to K_2 \to K_1$ induisent des isomorphismes après restriction à $U$. Appelons un tel système un {\it système de Deligne}.

\begin{lemma}
\label{lemma-lift-map}
Soit $j : U \to X$ une immersion ouverte de schémas noethériens. Soit $(K_n)$ un système de Deligne, et notons $K \in D^b_{\textit{Coh}}(\mathcal{O}_U)$ la valeur du système constant $(K_n|_U)$. Soit $L$ un objet de $D^b_{\textit{Coh}}(\mathcal{O}_X)$. Alors $\colim \Hom_X(K_n, L) = \Hom_U(K, L|_U)$.
\end{lemma}

\begin{proof}
Soit $L \to M \to N \to L[1]$ un triangle distingué dans $D^b_{\textit{Coh}}(\mathcal{O}_X)$. On obtient alors un diagramme commutatif
\begingroup\scriptsize
$$
\xymatrix{
\ldots \ar[r] &
\colim \Hom_X(K_n, L) \ar[r] \ar[d] &
\colim \Hom_X(K_n, M) \ar[r] \ar[d] &
\colim \Hom_X(K_n, N) \ar[r] \ar[d] &
\ldots \\
\ldots \ar[r] &
\Hom_U(K, L|_U) \ar[r] &
\Hom_U(K, M|_U) \ar[r] &
\Hom_U(K, N|_U) \ar[r] &
\ldots
}
$$
\endgroup
dont les lignes sont exactes d'après Catégories dérivées, lemme \ref{derived-lemma-representable-homological}, et Algèbre, lemme \ref{algebra-lemma-directed-colimit-exact}. Ainsi, si l'énoncé du lemme vaut pour $N[-1]$, $L$, $N$ et $L[1]$, alors il vaut pour $M$ d'après le lemme des cinq. En utilisant les triangles distingués des troncatures canoniques de $L$ (voir Catégories dérivées, remarque \ref{derived-remark-truncation-distinguished-triangle}), on se ramène donc au cas où $L$ ne possède qu'un seul faisceau de cohomologie non nul.

\medskip\noindent
Choisissons un complexe borné $\mathcal{F}^\bullet$ de $\mathcal{O}_X$-modules cohérents et un idéal quasi-cohérent $\mathcal{I} \subset \mathcal{O}_X$ définissant $X \setminus U$, tels que $K_n$ soit représenté par $\mathcal{I}^n\mathcal{F}^\bullet$. Les troncatures « brutales » donnent des suites exactes courtes compatibles de complexes, scindées terme à terme,
$$
0 \to \sigma_{\geq a + 1} \mathcal{I}^n\mathcal{F}^\bullet \to
\mathcal{I}^n\mathcal{F}^\bullet \to
\sigma_{\leq a} \mathcal{I}^n\mathcal{F}^\bullet \to 0
$$
qui correspondent à leur tour à des systèmes compatibles de triangles distingués dans $D^b_{\textit{Coh}}(\mathcal{O}_X)$. En raisonnant comme ci-dessus, on se ramène au cas où $\mathcal{F}^\bullet$ ne possède qu'un seul terme non nul. On est ainsi ramené au cas examiné dans le paragraphe suivant.

\medskip\noindent
Étant donnés un $\mathcal{O}_X$-module cohérent $\mathcal{F}$ et un $\mathcal{O}_X$-module cohérent $\mathcal{G}$, il faut montrer que le morphisme canonique
$$
\colim \Ext^i_X(\mathcal{I}^n\mathcal{F}, \mathcal{G})
\longrightarrow
\Ext^i_U(\mathcal{F}|_U, \mathcal{G}|_U)
$$
est un isomorphisme pour tout $i \geq 0$. Pour $i = 0$, c'est Cohomologie des schémas, lemme \ref{coherent-lemma-homs-over-open}. Supposons $i > 0$.

\medskip\noindent
Injectivité. Soit $\xi \in \Ext^i_X(\mathcal{I}^n\mathcal{F}, \mathcal{G})$ un élément dont la restriction à $U$ est nulle. Il faut montrer qu'il existe un $m \geq n$ tel que la restriction de $\xi$ à $\mathcal{I}^m\mathcal{F} = \mathcal{I}^{m - n}\mathcal{I}^n\mathcal{F}$ soit nulle. Après avoir remplacé $\mathcal{F}$ par $\mathcal{I}^n\mathcal{F}$, on peut supposer $n = 0$, c'est-à-dire que l'on dispose de
$\xi \in \Ext^i_X(\mathcal{F}, \mathcal{G})$
dont la restriction à $U$ est nulle. D'après Catégories dérivées de schémas, proposition \ref{perfect-proposition-DCoh}, on a
$D^b_{\textit{Coh}}(\mathcal{O}_X) = D^b(\textit{Coh}(\mathcal{O}_X))$.
On peut donc calculer le groupe $\Ext$ dans la catégorie abélienne des $\mathcal{O}_X$-modules cohérents. Cela implique qu'il existe une surjection $\alpha : \mathcal{F}'' \to \mathcal{F}$ telle que $\xi \circ \alpha = 0$ (c'est ici que l'on utilise $i > 0$). Posons $\mathcal{F}' = \Ker(\alpha)$, de sorte que l'on ait une suite exacte courte
$$
0 \to \mathcal{F}' \to \mathcal{F}'' \to \mathcal{F} \to 0
$$
Il s'ensuit que $\xi$ est l'image d'un élément $\xi' \in \Ext^{i - 1}_X(\mathcal{F}', \mathcal{G})$ dont la restriction à $U$ appartient à l'image de $\Ext^{i - 1}_U(\mathcal{F}''|_U, \mathcal{G}|_U) \to
\Ext^{i - 1}_U(\mathcal{F}'|_U, \mathcal{G}|_U)$. Par Artin–Rees, les systèmes projectifs $(\mathcal{I}^n\mathcal{F}')$ et $(\mathcal{I}^n \mathcal{F}'' \cap \mathcal{F}')$ sont pro-isomorphes, voir Cohomologie des schémas, lemme \ref{coherent-lemma-Artin-Rees}. Puisque l'on dispose du système compatible de suites exactes courtes
$$
0 \to
\mathcal{F}' \cap \mathcal{I}^n\mathcal{F}'' \to
\mathcal{I}^n\mathcal{F}'' \to
\mathcal{I}^n\mathcal{F} \to 0
$$
on obtient un diagramme commutatif
\begingroup\small
$$
\xymatrix{
\colim \Ext^{i - 1}_X(\mathcal{I}^n\mathcal{F}'', \mathcal{G})
\ar[r] \ar[d] &
\colim \Ext^{i - 1}_X(\mathcal{F}' \cap \mathcal{I}^n\mathcal{F}'', \mathcal{G})
\ar[r] \ar[d] &
\colim \Ext^i_X(\mathcal{I}^n\mathcal{F}, \mathcal{G})
\ar[d] \\
\Ext^{i - 1}_U(\mathcal{F}''|_U, \mathcal{G}|_U) \ar[r] &
\Ext^{i - 1}_U(\mathcal{F}'|_U, \mathcal{G}|_U) \ar[r] &
\Ext^{i - 1}_U(\mathcal{F}|_U, \mathcal{G}|_U)
}
$$
\endgroup
à lignes exactes. Par récurrence sur $i$ et d'après l'observation ci-dessus sur les systèmes projectifs, on trouve que les deux flèches verticales de gauche sont des isomorphismes. Maintenant, $\xi$ fournit un élément du groupe supérieur droit qui est l'image de $\xi'$ dans le groupe supérieur médian, lequel s'envoie à son tour sur un élément du groupe inférieur médian provenant d'un élément du groupe inférieur gauche. On conclut que $\xi$ s'envoie sur zéro dans
$\Ext^i_X(\mathcal{I}^n\mathcal{F}, \mathcal{G})$
pour un certain $n$, comme souhaité.

\medskip\noindent
Surjectivité. Soit $\xi \in \Ext^i_U(\mathcal{F}|_U, \mathcal{G}|_U)$. En raisonnant comme ci-dessus et en utilisant $i > 0$, on peut trouver une surjection $\mathcal{H} \to \mathcal{F}|_U$ de $\mathcal{O}_U$-modules cohérents telle que $\xi$ s'envoie sur zéro dans $\Ext^i_U(\mathcal{H}, \mathcal{G}|_U)$. On peut alors trouver un morphisme $\varphi : \mathcal{F}'' \to \mathcal{F}$ de $\mathcal{O}_X$-modules cohérents dont la restriction à $U$ est $\mathcal{H} \to \mathcal{F}|_U$, voir Propriétés, lemme \ref{properties-lemma-extend-finite-presentation}. Remarquons que le lemme ne garantit pas que $\varphi$ soit surjectif, mais cela n'aura pas d'importance (on peut choisir un $\varphi$ surjectif au prix d'un peu de travail supplémentaire). Notons $\mathcal{F}' = \Ker(\varphi)$. La suite exacte courte
$$
0 \to \mathcal{F}'|_U \to \mathcal{F}''|_U \to \mathcal{F}|_U \to 0
$$
montre que $\xi$ est l'image de $\xi'$ dans $\Ext^{i - 1}_U(\mathcal{F}'|_U, \mathcal{G}|_U)$. Par récurrence sur $i$, on peut trouver un $n$ tel que $\xi'$ soit l'image d'un certain $\xi'_n$ dans $\Ext^{i - 1}_X(\mathcal{I}^n\mathcal{F}', \mathcal{G})$. Par Artin–Rees, on peut trouver un $m \geq n$ tel que $\mathcal{F}' \cap \mathcal{I}^m\mathcal{F}'' \subset
\mathcal{I}^n\mathcal{F}'$. En utilisant la suite exacte courte
$$
0 \to \mathcal{F}' \cap \mathcal{I}^m\mathcal{F}'' \to
\mathcal{I}^m\mathcal{F}'' \to \mathcal{I}^m\Im(\varphi) \to 0
$$
l'image de $\xi'_n$ dans
$\Ext^{i - 1}_X(\mathcal{F}' \cap \mathcal{I}^m\mathcal{F}'', \mathcal{G})$
est envoyée par le morphisme de bord sur un élément $\xi_m$ de $\Ext^i_X(\mathcal{I}^m\Im(\varphi), \mathcal{G})$ qui s'envoie sur $\xi$. Puisque $\Im(\varphi)$ et $\mathcal{F}$ coïncident sur $U$, on voit que $\mathcal{F}/\mathcal{I}^m\Im(\varphi)$ est à support dans $X \setminus U$. Il existe donc un $l \geq m$ tel que $\mathcal{I}^l\mathcal{F} \subset \mathcal{I}^m\Im(\varphi)$, voir Cohomologie des schémas, lemme \ref{coherent-lemma-power-ideal-kills-sheaf}. En prenant l'image de $\xi_m$ dans $\Ext^i_X(\mathcal{I}^l\mathcal{F}, \mathcal{G})$, on conclut.
\end{proof}

\begin{lemma}
\label{lemma-lift-map-plus}
Le résultat du lemme \ref{lemma-lift-map} vaut même pour $L \in D^+_{\textit{Coh}}(\mathcal{O}_X)$.
\end{lemma}

\begin{proof}
En effet, si $(K_n)$ est un système de Deligne, il existe un $b \in \mathbf{Z}$ tel que $H^i(K_n) = 0$ pour $i > b$. Alors $\Hom(K_n, L) = \Hom(K_n, \tau_{\leq b}L)$ et $\Hom(K, L) = \Hom(K, \tau_{\leq b}L)$. En utilisant le résultat du lemme pour $\tau_{\leq b}L$, on conclut.
\end{proof}

\begin{lemma}
\label{lemma-extension-by-zero}
Soit $j : U \to X$ une immersion ouverte de schémas noethériens.
\begin{enumerate}
\item Soient $(K_n)$ et $(L_n)$ des systèmes de Deligne. Soient $K$ et $L$ les valeurs des systèmes constants $(K_n|_U)$ et $(L_n|_U)$. Étant donné un morphisme $\alpha : K \to L$ de $D(\mathcal{O}_U)$, il existe un unique morphisme de pro-systèmes $(K_n) \to (L_n)$ de $D^b_{\textit{Coh}}(\mathcal{O}_X)$ dont la restriction à $U$ est $\alpha$.
\item Étant donné $K \in D^b_{\textit{Coh}}(\mathcal{O}_U)$, il existe un système de Deligne $(K_n)$ tel que $(K_n|_U)$ soit constant de valeur $K$.
\item Le pro-objet $(K_n)$ de $D^b_{\textit{Coh}}(\mathcal{O}_X)$ de (2) est unique à isomorphisme unique près (comme pro-objet).
\end{enumerate}
\end{lemma}

\begin{proof}
L'assertion (1) résulte immédiatement du lemme \ref{lemma-lift-map} et du fait que les morphismes entre pro-systèmes sont les mêmes que les morphismes entre les foncteurs qu'ils corereprésentent, voir Catégories, remarque \ref{categories-remark-pro-category-copresheaves}.

\medskip\noindent
Soit $K$ comme dans (2). On peut choisir $K' \in D^b_{\textit{Coh}}(\mathcal{O}_X)$ dont la restriction à $U$ est isomorphe à $K$, voir Catégories dérivées de schémas, lemme \ref{perfect-lemma-lift-coherent}. D'après Catégories dérivées de schémas, proposition \ref{perfect-proposition-DCoh}, on peut représenter $K'$ par un complexe borné $\mathcal{F}^\bullet$ de $\mathcal{O}_X$-modules cohérents. Choisissons un faisceau quasi-cohérent d'idéaux $\mathcal{I} \subset \mathcal{O}_X$ dont le lieu d'annulation est $X \setminus U$ (par exemple, choisissons $\mathcal{I}$ correspondant à la structure de sous-schéma induit réduit sur $X \setminus U$). On peut alors prendre pour $K_n$ l'objet représenté par le complexe $\mathcal{I}^n\mathcal{F}^\bullet$ comme dans l'introduction de cette section.

\medskip\noindent
L'assertion (3) résulte immédiatement des assertions (1) et (2).
\end{proof}

\begin{lemma}
\label{lemma-extension-by-zero-triangle}
Soit $j : U \to X$ une immersion ouverte de schémas noethériens. Soit
$$
K \to L \to M \to K[1]
$$
un triangle distingué de $D^b_{\textit{Coh}}(\mathcal{O}_U)$. Il existe alors un système projectif de triangles distingués
$$
K_n \to L_n \to M_n \to K_n[1]
$$
dans $D^b_{\textit{Coh}}(\mathcal{O}_X)$ tel que $(K_n)$, $(L_n)$ et $(M_n)$ soient des systèmes de Deligne et que la restriction de ces triangles distingués à $U$ soit isomorphe au triangle distingué dont nous sommes partis.
\end{lemma}

\begin{proof}
Soit $(K_n)$ comme dans le lemme \ref{lemma-extension-by-zero}, assertion (2). Choisissons un objet $L'$ de $D^b_{\textit{Coh}}(\mathcal{O}_X)$ dont la restriction à $U$ est $L$ (cela est possible d'après le lemme). D'après le lemme \ref{lemma-lift-map}, on peut trouver un $n$ et un morphisme $K_n \to L'$ sur $X$ dont la restriction à $U$ est la flèche donnée $K \to L$. On conclut à l'existence d'un morphisme $K' \to L'$ de $D^b_{\textit{Coh}}(\mathcal{O}_X)$ dont la restriction à $U$ est la flèche donnée $K \to L$.

\medskip\noindent
D'après Catégories dérivées de schémas, proposition \ref{perfect-proposition-DCoh}, on peut trouver un morphisme $\alpha^\bullet : \mathcal{F}^\bullet \to
\mathcal{G}^\bullet$ de complexes bornés de $\mathcal{O}_X$-modules cohérents représentant $K' \to L'$. Choisissons un faisceau quasi-cohérent d'idéaux $\mathcal{I} \subset \mathcal{O}_X$ dont le lieu d'annulation est $X \setminus U$. Posons alors
$K_n = \mathcal{I}^n\mathcal{F}^\bullet$
et
$L_n = \mathcal{I}^n\mathcal{G}^\bullet$.
Remarquons que $\alpha^\bullet$ induit un morphisme de complexes $\alpha_n^\bullet : \mathcal{I}^n\mathcal{F}^\bullet \to
\mathcal{I}^n\mathcal{G}^\bullet$. La construction des cônes dans Catégories dérivées, section \ref{derived-section-cones}, montre clairement que
$$
C(\alpha_n)^\bullet = \mathcal{I}^nC(\alpha^\bullet)
$$
et l'on peut donc poser $M_n = C(\alpha_n)^\bullet$. On dispose en effet d'un système compatible de triangles distingués (voir la discussion de Catégories dérivées, section \ref{derived-section-canonical-delta-functor})
$$
K_n \to L_n \to M_n \to K_n[1]
$$
dont la restriction à $U$ est isomorphe au triangle distingué dont nous sommes partis, d'après l'axiome TR3 et Catégories dérivées, lemme \ref{derived-lemma-third-isomorphism-triangle}.
\end{proof}

\begin{remark}
\label{remark-extension-by-zero}
Soit $j : U \to X$ une immersion ouverte de schémas noethériens. Envoyer $K \in D^b_{\textit{Coh}}(\mathcal{O}_U)$ sur un système de Deligne dont la restriction à $U$ est $K$ détermine un foncteur
$$
Rj_! :
D^b_{\textit{Coh}}(\mathcal{O}_U)
\longrightarrow
\text{Pro-}D^b_{\textit{Coh}}(\mathcal{O}_X)
$$
qui est « exact » d'après le lemme \ref{lemma-extension-by-zero-triangle} et qui est « adjoint à gauche » du foncteur $j^* : D^b_{\textit{Coh}}(\mathcal{O}_X) \to
D^b_{\textit{Coh}}(\mathcal{O}_U)$ d'après le lemme \ref{lemma-lift-map}.
\end{remark}

\begin{remark}
\label{remark-extension-by-zero-linear-pro-system}
Soient $(A_n)$ et $(B_n)$ des systèmes projectifs d'une catégorie $\mathcal{C}$. Appelons pro-morphisme linéaire de $(A_n)$ vers $(B_n)$ la donnée d'une famille compatible de morphismes $\varphi_n : A_{cn + d} \to B_n$ pour tout $n \geq 1$, pour des entiers fixes $c, d \geq 1$. On dira que $(\varphi_n : A_{cn + d} \to B_n)$ et $(\psi_n : A_{c'n + d'} \to B_n)$ déterminent le même morphisme s'il existe $c'' \geq \max(c, c')$ et $d'' \geq \max(d, d')$ tels que les deux morphismes induits $A_{c'' n + d''} \to B_n$ soient égaux pour tout $n$. Il semble probable que les systèmes de Deligne $(K_n)$ de valeur donnée sur $U$ soient bien définis à pro-isomorphisme linéaire près. Si nous en avons un jour besoin, nous le formulerons et le démontrerons soigneusement ici.
\end{remark}

\begin{lemma}
\label{lemma-deligne-system-2-out-of-3}
Soit $j : U \to X$ une immersion ouverte de schémas noethériens. Soit
$$
K_n \to L_n \to M_n \to K_n[1]
$$
un système projectif de triangles distingués dans $D^b_{\textit{Coh}}(\mathcal{O}_X)$. Si $(K_n)$ et $(M_n)$ sont pro-isomorphes à des systèmes de Deligne, alors il en est de même de $(L_n)$.
\end{lemma}

\begin{proof}
Remarquons que les systèmes $(K_n|_U)$ et $(M_n|_U)$ sont essentiellement constants puisqu'ils sont pro-isomorphes à des systèmes constants. Notons $K$ et $M$ leurs valeurs. D'après Catégories dérivées, lemme \ref{derived-lemma-essentially-constant-2-out-of-3}, le système projectif $L_n|_U$ est lui aussi essentiellement constant. Notons $L$ sa valeur. Soit $N \in D^b_{\textit{Coh}}(\mathcal{O}_X)$. Considérons le diagramme commutatif
\begingroup\scriptsize
$$
\xymatrix{
\ldots \ar[r] &
\colim \Hom_X(M_n, N) \ar[r] \ar[d] &
\colim \Hom_X(L_n, N) \ar[r] \ar[d] &
\colim \Hom_X(K_n, N) \ar[r] \ar[d] &
\ldots \\
\ldots \ar[r] &
\Hom_U(M, N|_U) \ar[r] &
\Hom_U(L, N|_U) \ar[r] &
\Hom_U(K, N|_U) \ar[r] &
\ldots
}
$$
\endgroup
D'après le lemme \ref{lemma-lift-map} et le fait que des ind-systèmes isomorphes ont même colimite, les deux flèches verticales à droite et les deux à gauche de celle du milieu sont des isomorphismes. Le lemme des cinq montre que la flèche verticale du milieu est un isomorphisme. Maintenant, si $(L'_n)$ est un système de Deligne dont la restriction à $U$ a pour valeur constante $L$ (un tel système existe d'après le lemme \ref{lemma-extension-by-zero}), alors on a aussi $\colim \Hom_X(L'_n, N) = \Hom_U(L, N|_U)$. Les pro-systèmes $(L_n)$ et $(L'_n)$ sont donc pro-isomorphes d'après Catégories, remarque \ref{categories-remark-pro-category-copresheaves}.
\end{proof}

\begin{lemma}
\label{lemma-consequence-Artin-Rees-bis}
Soit $X$ un schéma noethérien. Soit $\mathcal{I} \subset \mathcal{O}_X$ un faisceau quasi-cohérent d'idéaux. Soit $\mathcal{F}^\bullet$ un complexe de $\mathcal{O}_X$-modules cohérents. Soit $p \in \mathbf{Z}$. Posons $\mathcal{H} = H^p(\mathcal{F}^\bullet)$ et $\mathcal{H}_n = H^p(\mathcal{I}^n\mathcal{F}^\bullet)$. Il existe alors des morphismes canoniques de $\mathcal{O}_X$-modules
$$
\ldots \to \mathcal{H}_3 \to \mathcal{H}_2 \to \mathcal{H}_1 \to \mathcal{H}
$$
Il existe un $c > 0$ tel que, pour $n \geq c$, l'image de $\mathcal{H}_n \to \mathcal{H}$ soit contenue dans $\mathcal{I}^{n - c}\mathcal{H}$, et il existe un morphisme canonique de $\mathcal{O}_X$-modules $\mathcal{I}^n\mathcal{H} \to \mathcal{H}_{n - c}$ tel que les composés
$$
\mathcal{I}^n \mathcal{H} \to \mathcal{H}_{n - c} \to
\mathcal{I}^{n - 2c}\mathcal{H}
\quad\text{et}\quad
\mathcal{H}_n \to \mathcal{I}^{n - c}\mathcal{H} \to \mathcal{H}_{n - 2c}
$$
soient les morphismes canoniques. En particulier, les systèmes projectifs $(\mathcal{H}_n)$ et $(\mathcal{I}^n\mathcal{H})$ sont isomorphes comme pro-objets de $\textit{Mod}(\mathcal{O}_X)$.
\end{lemma}

\begin{proof}
Si $X$ est affine, la traduction en algèbre donne Compléments sur l'algèbre, lemme \ref{more-algebra-lemma-consequence-Artin-Rees-bis}. Dans le cas général, on raisonne exactement comme dans la démonstration de ce lemme, en remplaçant la référence à Artin–Rees en algèbre par une référence à Cohomologie des schémas, lemme \ref{coherent-lemma-Artin-Rees}. Les détails sont omis.
\end{proof}

\begin{lemma}
\label{lemma-characterize-extension-by-zero-algebra}
Soit $j : U \to X$ une immersion ouverte de schémas noethériens. Soient $a \leq b$ des entiers. Soit $(K_n)$ un système projectif de $D^b_{\textit{Coh}}(\mathcal{O}_X)$ tel que $H^i(K_n) = 0$ pour $i \not \in [a, b]$. Les conditions suivantes sont équivalentes :
\begin{enumerate}
\item $(K_n)$ est pro-isomorphe à un système de Deligne ;
\item pour tout $p \in \mathbf{Z}$, il existe un $\mathcal{O}_X$-module cohérent $\mathcal{F}$ tel que les pro-systèmes $(H^p(K_n))$ et $(\mathcal{I}^n\mathcal{F})$ soient pro-isomorphes.
\end{enumerate}
\end{lemma}

\begin{proof}
Supposons (1). Pour démontrer (2), on peut supposer que $(K_n)$ est un système de Deligne. Par définition, on peut choisir un complexe borné $\mathcal{F}^\bullet$ de $\mathcal{O}_X$-modules cohérents et un faisceau quasi-cohérent d'idéaux définissant $X \setminus U$, tels que $K_n$ soit représenté par $\mathcal{I}^n\mathcal{F}^\bullet$. Le résultat résulte donc du lemme \ref{lemma-consequence-Artin-Rees-bis}.

\medskip\noindent
Supposons (2). Nous allons démontrer que $(K_n)$ est comme dans (1), par récurrence sur $b - a$. Si $a = b$, alors (1) vaut essentiellement par hypothèse. Si $a < b$, considérons le système compatible de triangles distingués
$$
\tau_{\leq a}K_n \to K_n \to \tau_{\geq a + 1}K_n \to (\tau_{\leq a}K_n)[1]
$$
Voir Catégories dérivées, remarque \ref{derived-remark-truncation-distinguished-triangle}. Par récurrence sur $b - a$, on sait que $\tau_{\leq a}K_n$ et $\tau_{\geq a + 1}K_n$ sont pro-isomorphes à des systèmes de Deligne. On conclut par le lemme \ref{lemma-deligne-system-2-out-of-3}.
\end{proof}

\begin{lemma}
\label{lemma-characterize-extension-by-zero}
Soit $j : U \to X$ une immersion ouverte de schémas noethériens. Soit $(K_n)$ un système projectif dans $D^b_{\textit{Coh}}(\mathcal{O}_X)$. Soit $X = W_1 \cup \ldots \cup W_r$ un recouvrement ouvert. Les conditions suivantes sont équivalentes :
\begin{enumerate}
\item $(K_n)$ est pro-isomorphe à un système de Deligne ;
\item pour chaque $i$, la restriction $(K_n|_{W_i})$ est pro-isomorphe à un système de Deligne relativement à l'immersion ouverte $U \cap W_i \to W_i$.
\end{enumerate}
\end{lemma}

\begin{proof}
On raisonne par récurrence sur $r$. Si $r = 1$, le résultat est clair. Supposons $r > 1$. Posons $V = W_1 \cup \ldots \cup W_{r - 1}$. Par récurrence, on voit que $(K_n|_V)$ est un système de Deligne. On est ainsi ramené à la discussion du paragraphe suivant.

\medskip\noindent
Supposons que $X = V \cup W$ soit un recouvrement ouvert et que $(K_n|_W)$ et $(K_n|_V)$ soient pro-isomorphes à des systèmes de Deligne. Il faut montrer que $(K_n)$ est pro-isomorphe à un système de Deligne. Remarquons que $(K_n|_{V \cap W})$ est pro-isomorphe à un système de Deligne (il résulte immédiatement de la construction des systèmes de Deligne que leur restriction aux ouverts les préserve). En particulier, les pro-systèmes
$(K_n|_{U \cap V})$,
$(K_n|_{U \cap W})$ et
$(K_n|_{U \cap V \cap W})$
sont essentiellement constants. Il résulte des triangles distingués de Cohomologie, lemme \ref{cohomology-lemma-exact-sequence-j-star}, et de Catégories dérivées, lemme \ref{derived-lemma-essentially-constant-2-out-of-3}, que $(K_n|_U)$ est essentiellement constant. Notons $K \in D^b_{\textit{Coh}}(\mathcal{O}_U)$ la valeur de ce système. Soit $L$ un objet de $D^b_{\textit{Coh}}(\mathcal{O}_X)$. Considérons le diagramme
\begingroup\tiny
$$
\xymatrix{
\colim \Ext^{-1}(K_n|_V, L|_V) \oplus
\colim \Ext^{-1}(K_n|_W, L|_W) \ar[r] \ar[d] &
\Ext^{-1}(K|_{U \cap V}, L|_{U \cap V}) \oplus
\Ext^{-1}(K|_{U \cap W}, L|_{U \cap W}) \ar[d] \\
\colim \Ext^{-1}(K_n|_{V \cap W}, L|_{V \cap W}) \ar[r] \ar[d] &
\Ext^{-1}(K|_{U \cap V \cap W}, L|_{U \cap V \cap W}) \ar[d] \\
\colim \Hom(K_n, L) \ar[d] \ar[r] &
\Hom(K|_U, L|_U) \ar[d] \\
\colim \Hom(K_n|_V, L|_V) \oplus \colim \Hom(K_n|_W, L|_W) \ar[r] \ar[d] &
\Hom(K|_{U \cap V}, L|_{U \cap V}) \oplus
\Hom(K|_{U \cap W}, L|_{U \cap W}) \ar[d] \\
\colim \Hom(K_n|_{V \cap W}, L|_{V \cap W}) \ar[r] &
\Hom(K|_{U \cap V \cap W}, L|_{U \cap V \cap W})
}
$$
\endgroup
Les suites verticales sont exactes d'après Cohomologie, lemme \ref{cohomology-lemma-mayer-vietoris-hom}, et parce que les colimites filtrantes sont exactes. Toutes les flèches horizontales, sauf celle du milieu, sont des isomorphismes d'après le lemme \ref{lemma-lift-map} et le fait que des systèmes pro-isomorphes ont mêmes colimites. Par le lemme des cinq, celle du milieu est donc aussi un isomorphisme. Il s'ensuit que $(K_n)$ est pro-isomorphe à un système de Deligne pour $K$. En effet, si $(K'_n)$ est un système de Deligne dont la restriction à $U$ a pour valeur constante $K$ (un tel système existe d'après le lemme \ref{lemma-extension-by-zero}), alors on a aussi $\colim \Hom_X(K'_n, L) = \Hom_U(K, L|_U)$. Les pro-systèmes $(K_n)$ et $(K'_n)$ sont donc pro-isomorphes d'après Catégories, remarque \ref{categories-remark-pro-category-copresheaves}.
\end{proof}

\begin{lemma}
\label{lemma-tensoring-Deligne-system}
Soit $j : U \to X$ une immersion ouverte de schémas noethériens. Soit $\mathcal{I} \subset \mathcal{O}_X$ un faisceau quasi-cohérent d'idéaux tel que $V(\mathcal{I}) = X \setminus U$. Soit $K$ dans $D^b_{\textit{Coh}}(\mathcal{O}_X)$. Alors
$$
K \otimes_{\mathcal{O}_X}^\mathbf{L} \mathcal{I}^n
$$
est pro-isomorphe à un système de Deligne de valeur constante $K|_U$ sur $U$.
\end{lemma}

\begin{proof}
D'après le lemme \ref{lemma-characterize-extension-by-zero}, la question est locale sur $X$. On peut donc supposer que $X$ est le spectre d'un anneau noethérien. Dans ce cas, l'énoncé résulte de la version algébrique, qui est Compléments sur l'algèbre, lemme \ref{more-algebra-lemma-tensoring-Deligne-system}.
\end{proof}







\section{Préliminaires à la cohomologie à support compact}
\label{section-preliminaries-compactly-supported}

\noindent
Dans la situation \ref{situation-shriek}, soit $f : X \to Y$ un morphisme de la catégorie $\textit{FTS}_S$. À l'aide des constructions de la section précédente, nous construirons un foncteur
$$
Rf_! :
D^b_{\textit{Coh}}(\mathcal{O}_X)
\longrightarrow
\text{Pro-}D^b_{\textit{Coh}}(\mathcal{O}_Y)
$$
qui se réduit au foncteur de la remarque \ref{remark-extension-by-zero} si $f$ est une immersion ouverte et qui, en général, se construit au moyen d'une compactification de $f$. Avant cela, nous avons besoin des lemmes suivants pour démontrer que notre construction est bien définie.

\begin{lemma}
\label{lemma-well-defined-pre}
Soit $f : X \to Y$ un morphisme propre de schémas noethériens. Soit $V \subset Y$ un sous-schéma ouvert et posons $U = f^{-1}(V)$. Diagramme
$$
\xymatrix{
U \ar[r]_j \ar[d]_g & X \ar[d]^f \\
V \ar[r]^{j'} & Y
}
$$
Il existe alors un isomorphisme canonique $Rj'_! \circ Rg_* \to Rf_* \circ Rj_!$ de foncteurs $D^b_{\textit{Coh}}(\mathcal{O}_U) \to
\text{Pro-}D^b_{\textit{Coh}}(\mathcal{O}_Y)$, où $Rj_!$ et $Rj'_!$ sont comme dans la remarque \ref{remark-extension-by-zero}.
\end{lemma}

\begin{proof}[Première démonstration]
Soit $K$ un objet de $D^b_{\textit{Coh}}(\mathcal{O}_U)$. Soit $(K_n)$ un système de Deligne pour $U \to X$ dont la restriction à $U$ est constante de valeur $K$. Cela signifie bien sûr que $(K_n)$ représente $Rj_!K$ dans $\text{Pro-}D^b_{\textit{Coh}}(\mathcal{O}_X)$. Remarquons que $Rj'_!Rg_*K$ et $Rf_*Rj_!K$ se restreignent tous deux au pro-objet constant de valeur $Rg_*K$ sur $V$. C'est immédiat pour le premier et, pour le second, cela résulte de $(Rf_*K_n)|_V = Rg_*(K_n|_U) = Rg_*K$. Par l'unicité des systèmes de Deligne du lemme \ref{lemma-extension-by-zero}, il suffit de montrer que $(Rf_*K_n)$ est pro-isomorphe à un système de Deligne. Le lemme cité montrera aussi que l'isomorphisme obtenu est fonctoriel.

\medskip\noindent
Démonstration que $(Rf_*K_n)$ est pro-isomorphe à un système de Deligne. Tout d'abord, remarquons que la question est indépendante du choix du système de Deligne $(K_n)$ correspondant à $K$ (par l'unicité précitée). D'après les lemmes \ref{lemma-extension-by-zero-triangle} et \ref{lemma-deligne-system-2-out-of-3}, si l'on dispose d'un triangle distingué
$$
K \to L \to M \to K[1]
$$
dans $D^b_{\textit{Coh}}(\mathcal{O}_U)$ et que le résultat vaut pour $K$ et $M$, alors il vaut pour $L$. En utilisant les triangles distingués des troncatures canoniques (Catégories dérivées, remarque \ref{derived-remark-truncation-distinguished-triangle}), on se ramène au problème étudié dans le paragraphe suivant.

\medskip\noindent
Soit $\mathcal{F}$ un $\mathcal{O}_X$-module cohérent. Soit $\mathcal{J} \subset \mathcal{O}_Y$ un faisceau quasi-cohérent d'idéaux définissant $Y \setminus V$. Notons $\mathcal{J}^n\mathcal{F}$ l'image de $f^*\mathcal{J}^n \otimes \mathcal{F} \to \mathcal{F}$. Il faut montrer que $(Rf_*(\mathcal{J}^n\mathcal{F}))$ est un système de Deligne. D'après le lemme \ref{lemma-characterize-extension-by-zero}, la question est locale sur $Y$. On peut donc supposer $Y = \Spec(A)$ affine et que $\mathcal{J}$ correspond à un idéal $I \subset A$. D'après le lemme \ref{lemma-characterize-extension-by-zero-algebra}, il suffit de montrer que le système projectif de modules de cohomologie $(H^p(X, I^n\mathcal{F}))$ est pro-isomorphe au système projectif $(I^n M)$ pour un certain $A$-module fini $M$. Cela est démontré dans Cohomologie des schémas, lemme \ref{coherent-lemma-cohomology-powers-ideal-application}.
\end{proof}

\begin{proof}[Seconde démonstration]
Soit $K$ un objet de $D^b_{\textit{Coh}}(\mathcal{O}_U)$. Soit $L$ un objet de $D^b_{\textit{Coh}}(\mathcal{O}_Y)$. Nous allons construire une bijection
$$
\Hom_{\text{Pro-}D^b_{\textit{Coh}}(\mathcal{O}_Y)}(Rj'_!Rg_*K, L)
\longrightarrow
\Hom_{\text{Pro-}D^b_{\textit{Coh}}(\mathcal{O}_Y)}(Rf_*Rj_!K, L)
$$
fonctorielle en $K$ et $L$. À $K$ fixé, cela déterminera un isomorphisme de pro-objets
$Rf_*Rj_!K \to Rj'_!Rg_*K$
d'après Catégories, remarque \ref{categories-remark-pro-category-copresheaves}, et, lorsque $K$ varie, cela déterminera un isomorphisme de foncteurs. Pour produire effectivement l'isomorphisme, on utilise la suite d'égalités fonctorielles
\begin{align*}
\Hom_{\text{Pro-}D^b_{\textit{Coh}}(\mathcal{O}_Y)}(Rj'_!Rg_*K, L)
& =
\Hom_V(Rg_*K, L|_V) \\
& =
\Hom_U(K, g^!(L|_V)) \\
& =
\Hom_U(K, f^!L|_U)) \\
& =
\Hom_{\text{Pro-}D^b_{\textit{Coh}}(\mathcal{O}_X)}(Rj_!K, f^!L) \\
& =
\Hom_{\text{Pro-}D^b_{\textit{Coh}}(\mathcal{O}_Y)}(Rf_*Rj_!K, L)
\end{align*}
La première égalité résulte du lemme \ref{lemma-lift-map}. La deuxième vaut parce que $g$ est propre (comme changement de base de $f$ à $V$) et que $g^!$ est donc, par construction, l'adjoint à droite de l'image directe, voir la section \ref{section-upper-shriek}. La troisième vaut puisque $g^!(L|_V) = f^!L|_U$ d'après le lemme \ref{lemma-restrict-before-or-after}. Comme $f^!L$ appartient à $D^+_{\textit{Coh}}(\mathcal{O}_X)$ d'après le lemme \ref{lemma-shriek-coherent}, la quatrième égalité résulte du lemme \ref{lemma-lift-map-plus}. La cinquième vaut à nouveau parce que $f^!$ est l'adjoint à droite de $Rf_*$ puisque $f$ est propre.
\end{proof}

\begin{lemma}
\label{lemma-well-defined}
Soit $j : U \to X$ une immersion ouverte de schémas noethériens. Soit $j' : U \to X'$ une compactification de $U$ sur $X$ (voir la démonstration), et notons $f : X' \to X$ le morphisme structural. Il existe alors un isomorphisme canonique $Rj_! \to Rf_* \circ R(j')_!$ de foncteurs $D^b_{\textit{Coh}}(\mathcal{O}_U) \to
\text{Pro-}D^b_{\textit{Coh}}(\mathcal{O}_X)$, où $Rj_!$ et $Rj'_!$ sont comme dans la remarque \ref{remark-extension-by-zero}.
\end{lemma}

\begin{proof}
Dire que $X'$ est une compactification de $U$ sur $X$ signifie précisément que $f : X' \to X$ est propre, que $j'$ est une immersion ouverte et que $j = f \circ j'$. Voir Compléments sur la platitude, section \ref{flat-section-compactify}. Si $j'(U) = f^{-1}(j(U))$, le lemme résulte immédiatement du lemme \ref{lemma-well-defined-pre}. Si $j'(U) \not = f^{-1}(j(U))$, notons $X'' \subset X'$ l'adhérence schématique de $j' : U \to X'$ et $j'' : U \to X''$ l'immersion ouverte correspondante. Diagramme
$$
\xymatrix{
& & X'' \ar[d]^{f'} \\
& & X' \ar[d]^f \\
U \ar[rr]^j \ar[rru]^{j'} \ar[rruu]^{j''} & & X
}
$$
D'après Compléments sur la platitude, lemme \ref{flat-lemma-compactifications-cofiltered}, assertion (c), et la discussion ci-dessus, on a des isomorphismes $Rf'_* \circ Rj''_! = Rj'_!$ et $R(f \circ f')_* \circ Rj''_! = Rj_!$. Puisque $R(f \circ f')_* = Rf_* \circ Rf'_*$, on conclut.
\end{proof}

\begin{remark}
\label{remark-covariance-open-j-lower-shriek}
Soient $X \supset U \supset U'$ des sous-schémas ouverts d'un schéma noethérien $X$. Notons $j : U \to X$ et $j' : U' \to X$ les morphismes d'inclusion. Nous affirmons qu'il existe un morphisme canonique
$$
Rj'_!(K|_{U'}) \longrightarrow Rj_!K
$$
fonctoriel en $K$ dans $D^b_{\textit{Coh}}(\mathcal{O}_U)$. En effet, d'après le lemme \ref{lemma-lift-map}, on dispose, pour tout $L$ dans $D^b_{\textit{Coh}}(\mathcal{O}_X)$, du morphisme
\begin{align*}
\Hom_{\text{Pro-}D^b_{\textit{Coh}}(\mathcal{O}_X)}(Rj_!K, L)
& =
\Hom_U(K, L|_U) \\
& \to
\Hom_{U'}(K|_{U'}, L|_{U'}) \\
& =
\Hom_{\text{Pro-}D^b_{\textit{Coh}}(\mathcal{O}_X)}(Rj'_!(K|_{U'}), L)
\end{align*}
fonctoriel en $L$ et $K'$. La fonctorialité en $L$ montre, d'après Catégories, remarque \ref{categories-remark-pro-category-copresheaves}, que l'on obtient un morphisme canonique $Rj'_!(K|_{U'}) \to Rj_!K$ qui est fonctoriel en $K$ par la fonctorialité en $K$ de la flèche ci-dessus.

\medskip\noindent
Voici une construction explicite de cette flèche. Supposons que $\mathcal{F}^\bullet$ soit un complexe borné de $\mathcal{O}_X$-modules cohérents dont la restriction à $U$ représente $K$ dans la catégorie dérivée. Nous avons vu dans la démonstration du lemme \ref{lemma-extension-by-zero} qu'un tel complexe existe toujours. Soient $\mathcal{I}$, resp.\ $\mathcal{I}'$, des faisceaux quasi-cohérents d'idéaux sur $X$ tels que $V(\mathcal{I}) = X \setminus U$, resp.\ $V(\mathcal{I}') = X \setminus U'$. Après avoir remplacé $\mathcal{I}$ par $\mathcal{I} + \mathcal{I}'$, on peut supposer $\mathcal{I}' \subset \mathcal{I}$. Par construction, $Rj_!K$, resp.\ $Rj'_!(K|_{U'})$, est représenté par le système projectif $(K_n)$, resp.\ $(K'_n)$, de $D^b_{\textit{Coh}}(\mathcal{O}_X)$, avec
$$
K_n = \mathcal{I}^n\mathcal{F}^\bullet
\quad\text{resp.}\quad
K'_n = (\mathcal{I}')^n\mathcal{F}^\bullet
$$
Il est clair que le morphisme construit ci-dessus est donné par les morphismes $K'_n \to K_n$ provenant des inclusions $(\mathcal{I}')^n \subset \mathcal{I}^n$.
\end{remark}






\section{Cohomologie à support compact pour les modules cohérents}
\label{section-compactly-supported}

\noindent
Dans la situation \ref{situation-shriek}, étant donné un morphisme $f : X \to Y$ de $\textit{FTS}_S$, nous allons définir un foncteur
$$
Rf_! : D^b_{\textit{Coh}}(\mathcal{O}_X)
\longrightarrow
\text{Pro-}D^b_{\textit{Coh}}(\mathcal{O}_Y)
$$
Choisissons une compactification $j : X \to \overline{X}$ sur $Y$, ce qui est possible d'après Compléments sur la platitude, théorème \ref{flat-theorem-nagata}, et lemme \ref{flat-lemma-compactifyable}. Notons $\overline{f} : \overline{X} \to Y$ le morphisme structural. Posons alors
$$
Rf_!K  = R\overline{f}_* Rj_! K
$$
pour $K \in D^b_{\textit{Coh}}(\mathcal{O}_X)$, où $Rj_!$ est comme dans la remarque \ref{remark-extension-by-zero}.

\begin{lemma}
\label{lemma-lower-shriek-well-defined}
Le foncteur $Rf_!$ est, à isomorphisme près, indépendant du choix de la compactification.
\end{lemma}

\noindent
En fait, le foncteur $Rf_!$ sera caractérisé comme un « adjoint à gauche » de $f^!$, ce qui le déterminera à isomorphisme unique près.

\begin{proof}
Considérons la catégorie des compactifications de $X$ sur $Y$, qui est cofiltrante d'après Compléments sur la platitude, théorème \ref{flat-theorem-nagata}, et lemmes \ref{flat-lemma-compactifications-cofiltered} et \ref{flat-lemma-compactifyable}. À chaque choix d'une compactification
$$
j : X \to \overline{X},\quad \overline{f} : \overline{X} \to Y
$$
la construction ci-dessus associe le foncteur $R\overline{f}_* \circ Rj_!$. Supposons donné un morphisme $g : \overline{X}_1 \to \overline{X}_2$ entre des compactifications $j_i : X \to \overline{X}_i$ sur $Y$. On obtient alors un isomorphisme
$$
R\overline{f}_{2, *} \circ Rj_{2, !} =
R\overline{f}_{2, *} \circ Rg_* \circ j_{1, !} =
R\overline{f}_{1, *} \circ Rj_{1, !}
$$
en utilisant le lemme \ref{lemma-well-defined} dans la première égalité. On voit ainsi que notre foncteur est indépendant, à isomorphisme près, du choix de la compactification.
\end{proof}

\begin{proposition}
\label{proposition-duality-compactly-supported}
Dans la situation \ref{situation-shriek}, soit $f : X \to Y$ un morphisme de $\textit{FTS}_S$. Les foncteurs $Rf_!$ et $f^!$ sont alors adjoints au sens suivant : pour tous $K \in D^b_{\textit{Coh}}(\mathcal{O}_X)$ et $L \in D^+_{\textit{Coh}}(\mathcal{O}_Y)$, on a
$$
\Hom_X(K, f^!L) =
\Hom_{\text{Pro-}D^+_{\textit{Coh}}(\mathcal{O}_Y)}(Rf_!K, L)
$$
bifonctoriellement en $K$ et $L$.
\end{proposition}

\begin{proof}
Choisissons une compactification $j : X \to \overline{X}$ sur $Y$ et notons $\overline{f} : \overline{X} \to Y$ le morphisme structural. On a alors
\begin{align*}
\Hom_X(K, f^!L)
& =
\Hom_X(K, j^*\overline{f}{}^!L) \\
& =
\Hom_{\text{Pro-}D^+_{\textit{Coh}}(\mathcal{O}_{\overline{X}})}
(Rj_!K, \overline{f}{}^!L) \\
& =
\Hom_{\text{Pro-}D^+_{\textit{Coh}}(\mathcal{O}_Y)}(Rf_*Rj_!K, L) \\
& =
\Hom_{\text{Pro-}D^+_{\textit{Coh}}(\mathcal{O}_Y)}(Rf_!K, L)
\end{align*}
La première égalité résulte immédiatement de la construction de $f^!$ dans la section \ref{section-upper-shriek}. D'après le lemme \ref{lemma-shriek-coherent}, on a $\overline{f}{}^!L$ dans $D^+_{\textit{Coh}}(\mathcal{O}_{\overline{X}})$ ; la deuxième égalité résulte donc du lemme \ref{lemma-lift-map-plus}. Puisque $\overline{f}$ est propre, le foncteur $\overline{f}{}^!$ est, par construction, l'adjoint à droite de l'image directe. C'est ce qui donne la troisième égalité. La quatrième vaut parce que $Rf_! = Rf_* Rj_!$.
\end{proof}

\begin{lemma}
\label{lemma-compactly-supported-triangle}
Dans la situation \ref{situation-shriek}, soit $f : X \to Y$ un morphisme de $\textit{FTS}_S$. Soit
$$
K \to L \to M \to K[1]
$$
un triangle distingué de $D^b_{\textit{Coh}}(\mathcal{O}_X)$. Il existe alors un système projectif de triangles distingués
$$
K_n \to L_n \to M_n \to K_n[1]
$$
dans $D^b_{\textit{Coh}}(\mathcal{O}_Y)$ tel que les pro-systèmes $(K_n)$, $(L_n)$ et $(M_n)$ donnent $Rf_!K$, $Rf_!L$ et $Rf_!M$.
\end{lemma}

\begin{proof}
Choisissons une compactification $j : X \to \overline{X}$ sur $Y$ et notons $\overline{f} : \overline{X} \to Y$ le morphisme structural. Choisissons un système projectif de triangles distingués
$$
\overline{K}_n \to \overline{L}_n \to \overline{M}_n \to \overline{K}_n[1]
$$
dans $D^b_{\textit{Coh}}(\mathcal{O}_{\overline{X}})$ comme dans le lemme \ref{lemma-extension-by-zero-triangle}, correspondant à l'immersion ouverte $j$ et au triangle distingué donné. Prenons $K_n = R\overline{f}_*\overline{K}_n$, et de même pour $L_n$ et $M_n$. Cela convient par la définition même de $Rf_!$.
\end{proof}

\begin{remark}
\label{remark-compose-inverse-systems}
Soit $\mathcal{C}$ une catégorie. Supposons donné un système projectif
$$
\ldots \xrightarrow{\alpha_4} (M_{3, n}) \xrightarrow{\alpha_3} (M_{2, n})
\xrightarrow{\alpha_2} (M_{1, n})
$$
de systèmes projectifs dans la catégorie des pro-objets de $\mathcal{C}$. Autrement dit, les flèches $\alpha_i$ sont des morphismes de pro-objets. D'après Catégories, exemple \ref{categories-example-pro-morphism-inverse-systems}, on peut représenter chaque $\alpha_i$ par un couple $(m_i, a_i)$, où $m_i : \mathbf{N} \to \mathbf{N}$ est une fonction croissante et $a_{i, n} : M_{i, m_i(n)} \to M_{i - 1, n}$ est un morphisme de $\mathcal{C}$ rendant commutatifs les diagrammes
$$
\xymatrix{
\ldots \ar[r] &
M_{i, m_i(3)} \ar[d]^{a_{i, 3}} \ar[r] &
M_{i, m_i(2)} \ar[d]^{a_{i, 2}} \ar[r] &
M_{i, m_i(1)} \ar[d]^{a_{i, 1}} \\
\ldots \ar[r] &
M_{i - 1, 3} \ar[r] &
M_{i - 1, 2} \ar[r] &
M_{i - 1, 1}
}
$$
En remplaçant $m_i(n)$ par $\max(n, m_i(n))$ et en ajustant les morphismes $a_i(n)$ en conséquence (comme dans l'exemple cité), on peut supposer $m_i(n) \geq n$. Dans cette situation, considérons le système projectif
$$
\ldots \to
M_{4, m_4(m_3(m_2(4)))} \to
M_{3, m_3(m_2(3))} \to
M_{2, m_2(2)} \to
M_{1, 1}
$$
de terme général
$$
M_k = M_{k, m_k(m_{k - 1}(\ldots (m_2(k))\ldots))}
$$
Pour tout objet $N$ de $\mathcal{C}$, on a
$$
\colim_i \colim_n \Mor_\mathcal{C}(M_{i, n}, N) =
\colim_k \Mor_\mathcal{C}(M_k, N) 
$$
Nous omettons les détails. Autrement dit, on voit que le système projectif $(M_k)$ possède la propriété
$$
\colim_i \Mor_{\text{Pro-}\mathcal{C}}((M_{i, n}), N) =
\Mor_{\text{Pro-}\mathcal{C}}((M_k), N)
$$
Cette propriété détermine le système projectif $(M_k)$ à pro-isomorphisme près d'après la discussion de Catégories, remarque \ref{categories-remark-pro-category-copresheaves}. On peut ainsi transformer certains systèmes projectifs de $\text{Pro-}\mathcal{C}$ en pro-objets dont les catégories d'indices sont dénombrables.
\end{remark}

\begin{remark}
\label{remark-composition-lower-shriek}
Dans la situation \ref{situation-shriek}, soient $f : X \to Y$ et $g : Y \to Z$ des morphismes composables de $\textit{FTS}_S$. Définissons la composition
$$
Rg_! \circ Rf_! :
D^b_{\textit{Coh}}(\mathcal{O}_X)
\longrightarrow
\text{Pro-}D^b_{\textit{Coh}}(\mathcal{O}_Z)
$$
Par la construction même de $Rf_!$ pour $K$ dans $D^b_{\textit{Coh}}(\mathcal{O}_X)$, l'objet $Rf_!K$ obtenu est la classe de pro-isomorphisme d'un système projectif $(M_n)$ dans $D^b_{\textit{Coh}}(\mathcal{O}_Y)$. Ensuite, puisque $Rg_!$ se construit de la même manière, on voit que
$$
\ldots \to Rg_!M_3 \to Rg_!M_2 \to Rg_!M_1
$$
est un système projectif de $\text{Pro-}D^b_{\textit{Coh}}(\mathcal{O}_Y)$. D'après la discussion de la remarque \ref{remark-compose-inverse-systems}, il existe une unique classe de pro-isomorphisme, que l'on notera $Rg_! Rf_! K$, représentée par des systèmes projectifs dans $D^b_{\textit{Coh}}(\mathcal{O}_Z)$ et caractérisée par l'égalité
$$
\Hom_{\text{Pro-}D^b_{\textit{Coh}}(\mathcal{O}_Z)}(Rg_!Rf_!K, L) =
\colim_n \Hom_{\text{Pro-}D^b_{\textit{Coh}}(\mathcal{O}_Z)}(Rg_!M_n, L)
$$
Nous omettons la discussion nécessaire pour voir que cette construction est fonctorielle en $K$, car cela résultera immédiatement du lemme suivant.
\end{remark}

\begin{lemma}
\label{lemma-composition-lower-shriek}
Dans la situation \ref{situation-shriek}, soient $f : X \to Y$ et $g : Y \to Z$ des morphismes composables de $\textit{FTS}_S$. Avec les notations de la remarque \ref{remark-composition-lower-shriek}, on a
$Rg_! \circ Rf_! = R(g \circ f)_!$.
\end{lemma}

\begin{proof}
D'après la discussion de Catégories, remarque \ref{categories-remark-pro-category-copresheaves}, il suffit de montrer que l'on obtient le même résultat si l'on calcule $\Hom$ à valeurs dans $L$ dans la catégorie $D^b_{\textit{Coh}}(\mathcal{O}_Z)$. Pour ce faire, avec les notations de la remarque \ref{remark-composition-lower-shriek}, on calcule comme suit :
\begin{align*}
\Hom_Z(Rg_!Rf_!K, L)
& =
\colim_n \Hom_Z(Rg_!M_n, L) \\
& =
\colim_n \Hom_Y(M_n, g^!L) \\
& =
\Hom_Y(Rf_!K, g^!L) \\
& =
\Hom_X(K, f^!g^!L) \\
& =
\Hom_X(K, (g \circ f)^!L) \\
& =
\Hom_Z(R(g \circ f)_!K, L)
\end{align*}
La première égalité est la définition de $Rg_!Rf_!K$. La deuxième est la proposition \ref{proposition-duality-compactly-supported} pour $g$. La troisième exprime que $Rf_!K$ est donné par $(M_n)$. La quatrième est la proposition \ref{proposition-duality-compactly-supported} pour $f$. La cinquième est le lemme \ref{lemma-upper-shriek-composition}. La sixième est la proposition \ref{proposition-duality-compactly-supported} pour $g \circ f$.
\end{proof}

\begin{remark}
\label{remark-covariance-open-lower-shriek}
Dans la situation \ref{situation-shriek}, soit $f : X \to Y$ un morphisme de $\textit{FTS}_S$, et soit $U \subset X$ un ouvert. Posons $g = f|_U : U \to Y$. Il existe alors un morphisme canonique
$$
Rg_!(K|_U) \longrightarrow Rf_!K
$$
fonctoriel en $K$ dans $D^b_{\textit{Coh}}(\mathcal{O}_X)$, que l'on peut définir d'au moins trois manières.
\begin{enumerate}
\item Notons $i : U \to X$ le morphisme d'inclusion. On a $Rg_! = Rf_! \circ Ri_!$ d'après le lemme \ref{lemma-composition-lower-shriek}, et l'on peut appliquer $Rf_!$ au morphisme $Ri_!(K|_U) \to K$, qui est un cas particulier de la remarque \ref{remark-covariance-open-j-lower-shriek}.
\item Choisissons une compactification $j : X \to \overline{X}$ de $X$ sur $Y$, de morphisme structural $\overline{f} : \overline{X} \to Y$. Posons $j' = j \circ i : U \to \overline{X}$. On peut utiliser $Rf_! = R\overline{f}_* \circ Rj_!$ et $Rg_! = R\overline{f}_* \circ Rj'_!$, puis appliquer $R\overline{f}_*$ au morphisme $Rj'_!(K|_U) \to Rj_!K$ de la remarque \ref{remark-covariance-open-j-lower-shriek}.
\item On peut utiliser
\begin{align*}
\Hom_{\text{Pro-}D^b_{\textit{Coh}}(\mathcal{O}_Y)}(Rf_!K, L)
& =
\Hom_X(K, f^!L) \\
& \to
\Hom_U(K|_U, f^!L|_U) \\
& =
\Hom_U(K|_U, g^!L) \\
& =
\Hom_{\text{Pro-}D^b_{\textit{Coh}}(\mathcal{O}_Y)}(Rg_!(K|_U), L)
\end{align*}
fonctoriel en $L$ et $K$. Ici, on a appliqué deux fois la proposition \ref{proposition-duality-compactly-supported}, ainsi que la construction des foncteurs image inverse extraordinaire, qui montre que $g^! = i^* \circ f^!$. La fonctorialité en $L$ montre, d'après Catégories, remarque \ref{categories-remark-pro-category-copresheaves}, que l'on obtient un morphisme canonique $Rg_!(K|_U) \to Rf_!K$ dans $\text{Pro-}D^b_{\textit{Coh}}(\mathcal{O}_Y)$, qui est fonctoriel en $K$ par la fonctorialité en $K$ de la flèche ci-dessus.
\end{enumerate}
Chacune de ces trois constructions donne la même flèche ; nous omettons les détails.
\end{remark}

\begin{remark}
\label{remark-covariance-etale-lower-shriek}
Généralisons la covariance de la cohomologie à support compact donnée dans la remarque \ref{remark-covariance-open-lower-shriek} aux morphismes \'etales. Dans la situation \ref{situation-shriek}, supposons en effet donné un diagramme commutatif
$$
\xymatrix{
U \ar[rr]_h \ar[rd]_g & & X \ar[ld]^f \\
& Y
}
$$
de $\textit{FTS}_S$, avec $h$ \'etale. Il existe alors un morphisme canonique
$$
Rg_!(h^*K) \longrightarrow Rf_!K
$$
fonctoriel en $K$ dans $D^b_{\textit{Coh}}(\mathcal{O}_X)$. Nous définissons cette transformation au moyen de la suite de morphismes
\begin{align*}
\Hom_{\text{Pro-}D^b_{\textit{Coh}}(\mathcal{O}_Y)}(Rf_!K, L)
& =
\Hom_X(K, f^!L) \\
& \to
\Hom_U(h^*K, h^*(f^!L)) \\
& =
\Hom_U(h^*K, h^!f^!L) \\
& =
\Hom_U(h^*K, g^!L) \\
& =
\Hom_{\text{Pro-}D^b_{\textit{Coh}}(\mathcal{O}_Y)}(Rg_!(h^*K), L)
\end{align*}
fonctorielle en $L$ et $K$. Ici, on a appliqué deux fois la proposition \ref{proposition-duality-compactly-supported}, utilisé l'égalité $h^* = h^!$ du lemme \ref{lemma-shriek-etale} et l'égalité $h^! \circ f^! = g^!$ du lemme \ref{lemma-upper-shriek-composition}. La fonctorialité en $L$ montre, d'après Catégories, remarque \ref{categories-remark-pro-category-copresheaves}, que l'on obtient un morphisme canonique $Rg_!(h^*K) \to Rf_!K$ dans $\text{Pro-}D^b_{\textit{Coh}}(\mathcal{O}_Y)$, qui est fonctoriel en $K$ par la fonctorialité en $K$ de la flèche ci-dessus.
\end{remark}

\begin{remark}
\label{remark-covariance-lower-shriek}
Dans les remarques \ref{remark-covariance-open-lower-shriek} et \ref{remark-covariance-etale-lower-shriek}, nous avons vu que la construction de la cohomologie à support compact est covariante relativement aux immersions ouvertes et aux morphismes \'etales. En fait, le bon degré de généralité est le suivant : étant donné un diagramme commutatif
$$
\xymatrix{
U \ar[rr]_h \ar[rd]_g & & X \ar[ld]^f \\
& Y
}
$$
de $\textit{FTS}_S$, avec $h$ plat et quasi-fini, il existe une transformation canonique
$$
Rg_! \circ h^* \longrightarrow Rf_!
$$
Comme dans la remarque \ref{remark-covariance-etale-lower-shriek}, ce morphisme peut se construire à l'aide d'une transformation de foncteurs $h^* \to h^!$ sur $D^+_{\textit{Coh}}(\mathcal{O}_X)$. Rappelons que $h^!K = h^*K \otimes \omega_{U/X}$, où $\omega_{U/X} = h^!\mathcal{O}_X$ est le faisceau dualisant relatif du morphisme plat quasi-fini $h$ (voir les lemmes \ref{lemma-perfect-comparison-shriek} et \ref{lemma-flat-quasi-finite-shriek}). Rappelons que $\omega_{U/X}$ est le même que le module dualisant relatif qui sera construit dans Discriminants, remarque \ref{discriminant-remark-relative-dualizing-for-quasi-finite}, d'après Discriminants, lemme \ref{discriminant-lemma-compare-dualizing}. On peut donc utiliser l'élément de trace
$\tau_{U/X} : \mathcal{O}_U \to \omega_{U/X}$
qui sera construit dans Discriminants, remarque \ref{discriminant-remark-relative-dualizing-for-flat-quasi-finite}, pour définir notre transformation. Si nous en avons un jour besoin, nous formulerons et démontrerons précisément le résultat ici.
\end{remark}







\section{Dualité pour la cohomologie à support compact}
\label{section-duality-compactly-supported}

\noindent
Soit $k$ un corps. Soit $U$ un schéma séparé de type fini sur $k$. Soit $K$ un objet de $D^b_{\textit{Coh}}(\mathcal{O}_U)$. Définissons la cohomologie à support compact $H^i_c(U, K)$ de $K$ comme suit. Choisissons une immersion ouverte $j : U \to X$ dans un schéma propre sur $k$ et un système de Deligne $(K_n)$ pour $j : U \to X$ dont la restriction à $U$ est constante de valeur $K$. Posons alors
$$
H^i_c(U, K) = \lim H^i(X, K_n)
$$
On le considère comme un $k$-espace vectoriel topologique muni de la topologie limite (voir Compléments sur l'algèbre, section \ref{more-algebra-section-topological-ring}). Plusieurs remarques s'imposent.

\medskip\noindent
Premièrement, cette définition est indépendante du choix de $X$ et de $(K_n)$. En effet, si $p : U \to \Spec(k)$ désigne le morphisme structural, on sait déjà que $Rp_!K = (R\Gamma(X, K_n))$ est bien défini à pro-isomorphisme près dans $D(k)$ ; il en est donc de même de la limite qui définit $H^i_c(U, K)$.

\medskip\noindent
Deuxièmement, il pourrait paraître plus naturel d'utiliser l'expression
$$
H^i(R\lim R\Gamma(X, K_n)) = R\Gamma(X, R\lim K_n)
$$
mais cela donnerait le même résultat : comme les $k$-espaces vectoriels $H^j(X, K_n)$ sont de dimension finie, ces systèmes projectifs satisfont la condition de Mittag-Leffler, et les termes $R^1\lim$ de Cohomologie, lemme \ref{cohomology-lemma-RGamma-commutes-with-Rlim}, sont donc nuls.

\medskip\noindent
Si $U' \subset U$ est un sous-schéma ouvert, il existe un morphisme canonique
$$
H^i_c(U', K|_{U'}) \longrightarrow H^i_c(U, K)
$$
fonctoriel en $K$ dans $D^b_{\textit{Coh}}(\mathcal{O}_U)$. Voir par exemple la remarque \ref{remark-covariance-open-lower-shriek}. En fait, en utilisant la remarque \ref{remark-covariance-etale-lower-shriek}, on voit que, plus généralement, un tel morphisme existe pour un morphisme \'etale $U' \to U$ de schémas séparés de type fini sur $k$.

\medskip\noindent
Si $V$ est un $k$-espace vectoriel, on munit $\Hom_k(V, k)$ d'une topologie comme suit : écrivons $V = \bigcup V_i$ comme la réunion filtrante de ses sous-$k$-espaces vectoriels de dimension finie et utilisons la topologie limite sur $\Hom_k(V, k) = \lim \Hom_k(V_i, k)$. Si $\dim_k V < \infty$, la topologie de $\Hom_k(V, k)$ est discrète. Plus généralement, si $V = \colim_n V_n$ est écrit comme une colimite filtrante d'espaces vectoriels de dimension finie, alors $\Hom_k(V, k) = \lim \Hom_k(V_n, k)$ comme espaces vectoriels topologiques.

\begin{lemma}
\label{lemma-duality-compact-support}
Soit $p : U \to \Spec(k)$ séparé et de type fini, où $k$ est un corps. Soit $\omega_{U/k}^\bullet = p^!\mathcal{O}_{\Spec(k)}$. Il existe des isomorphismes canoniques
$$
\Hom_k(H^i(U, K), k) =
H^{-i}_c(U, R\SheafHom_{\mathcal{O}_U}(K, \omega_{U/k}^\bullet))
$$
de $k$-espaces vectoriels topologiques, fonctoriels en $K$ dans $D^b_{\textit{Coh}}(\mathcal{O}_U)$.
\end{lemma}

\begin{proof}
Choisissons une compactification $j : U \to X$ sur $k$. Soit $\mathcal{I} \subset \mathcal{O}_X$ un faisceau quasi-cohérent d'idéaux tel que $V(\mathcal{I}) = X \setminus U$. D'après Catégories dérivées de schémas, proposition \ref{perfect-proposition-DCoh}, on peut choisir $M \in D^b_{\textit{Coh}}(\mathcal{O}_X)$ avec $K = M|_U$. On a
\begingroup\small
$$
H^i(U, K) =
\Ext^i_U(\mathcal{O}_U, M|_U) =
\colim \Ext^i_X(\mathcal{I}^n, M) =
\colim H^i(X, R\SheafHom_{\mathcal{O}_X}(\mathcal{I}^n, M))
$$
\endgroup
d'après le lemme \ref{lemma-lift-map}. Puisque $\mathcal{I}^n$ est un $\mathcal{O}_X$-module cohérent, on a $\mathcal{I}^n$ dans $D^-_{\textit{Coh}}(\mathcal{O}_X)$ ; ainsi $R\SheafHom_{\mathcal{O}_X}(\mathcal{I}^n, M)$ appartient à $D^+_{\textit{Coh}}(\mathcal{O}_X)$ d'après Catégories dérivées de schémas, lemme \ref{perfect-lemma-coherent-internal-hom}.

\medskip\noindent
Soit $\omega_{X/k}^\bullet = q^!\mathcal{O}_{\Spec(k)}$, où $q : X \to \Spec(k)$ est le morphisme structural, voir la section \ref{section-duality-proper-over-field}. On trouve que
\begin{align*}
\Hom_k(
&
H^i(X, R\SheafHom_{\mathcal{O}_X}(\mathcal{I}^n, M)), k) \\
& =
\Ext^{-i}_X(R\SheafHom_{\mathcal{O}_X}(\mathcal{I}^n, M),
\omega_{X/k}^\bullet) \\
& =
H^{-i}(X, R\SheafHom_{\mathcal{O}_X}(R\SheafHom_{\mathcal{O}_X}(
\mathcal{I}^n, M), \omega_{X/k}^\bullet))
\end{align*}
d'après le lemme \ref{lemma-duality-proper-over-field}. D'après l'assertion (1) du lemme \ref{lemma-internal-hom-evaluate-isom}, le morphisme canonique
$$
R\SheafHom_{\mathcal{O}_X}(M, \omega_{X/k}^\bullet)
\otimes_{\mathcal{O}_X}^\mathbf{L} \mathcal{I}^n
\longrightarrow
R\SheafHom_{\mathcal{O}_X}(R\SheafHom_{\mathcal{O}_X}(
\mathcal{I}^n, M), \omega_{X/k}^\bullet)
$$
est un isomorphisme. Remarquons que
$\omega^\bullet_{U/k} = \omega^\bullet_{X/k}|_U$
parce que $p^!$ se construit comme $q^!$ composé avec la restriction à $U$. Par conséquent, $R\SheafHom_{\mathcal{O}_X}(M, \omega_{X/k}^\bullet)$ est un objet de $D^b_{\textit{Coh}}(\mathcal{O}_X)$ dont la restriction est $R\SheafHom_{\mathcal{O}_U}(K, \omega_{U/k}^\bullet)$ sur $U$. D'après le lemme \ref{lemma-tensoring-Deligne-system}, on conclut donc que
$$
\lim H^{-i}(X, R\SheafHom_{\mathcal{O}_X}(M, \omega_{X/k}^\bullet)
\otimes_{\mathcal{O}_X}^\mathbf{L} \mathcal{I}^n)
$$
est un représentant du membre de droite de l'égalité du lemme. En combinant tous les isomorphismes ainsi obtenus, on obtient l'isomorphisme du lemme.
\end{proof}

\begin{lemma}
\label{lemma-duality-compact-support-restrict-open}
Avec les notations du lemme \ref{lemma-duality-compact-support}, supposons que $U' \subset U$ soit un sous-schéma ouvert. Alors le diagramme
$$
\xymatrix{
\Hom_k(H^i(U, K), k) \ar[rr] & &
H^{-i}_c(U, R\SheafHom_{\mathcal{O}_U}(K, \omega_{U/k}^\bullet)) \\
\Hom_k(H^i(U', K|_{U'}), k) \ar[rr] \ar[u] & &
H^{-i}_c(U', R\SheafHom_{\mathcal{O}_{U'}}(K, \omega_{U'/k}^\bullet)) \ar[u]
}
$$
est commutatif. Ici, les flèches horizontales sont les isomorphismes du lemme \ref{lemma-duality-compact-support}, la flèche verticale de gauche est la transposée du morphisme de restriction
\par\noindent\hfil
$H^i(U, K) \to H^i(U', K|_{U'})$,
\hfil\par\noindent
et la flèche verticale de droite est celle de la remarque \ref{remark-covariance-open-lower-shriek} (voir la discussion précédant le lemme).
\end{lemma}

\begin{proof}
Nous recommandons vivement au lecteur de sauter cette démonstration. Choisissons $X$ et $M$ comme dans la démonstration du lemme \ref{lemma-duality-compact-support}. Nous omettrons l'indice $\mathcal{O}_X$ de $R\SheafHom$ et $\otimes^\mathbf{L}$. Écrivons
$$
H^i(U, K) = \colim H^i(X, R\SheafHom(\mathcal{I}^n, M))
$$
et
$$
H^i(U', K|_{U'}) = \colim H^i(X, R\SheafHom((\mathcal{I}')^n, M))
$$
comme dans la démonstration du lemme \ref{lemma-duality-compact-support}, où l'on choisit $\mathcal{I}' \subset \mathcal{I}$ comme dans la discussion de la remarque \ref{remark-covariance-open-j-lower-shriek}, de sorte que le morphisme $H^i(U, K) \to H^i(U', K|_{U'})$ soit induit par les morphismes
$(\mathcal{I}')^n \to \mathcal{I}^n$.
De même, écrivons
$$
H^i_c(U, R\SheafHom(K, \omega_{U/k}^\bullet)) =
\lim H^i(X,  R\SheafHom(M, \omega_{X/k}^\bullet)
\otimes^\mathbf{L} \mathcal{I}^n)
$$
et
$$
H^i_c(U', R\SheafHom(K|_{U'}, \omega_{U'/k}^\bullet)) =
\lim H^i(X,  R\SheafHom(M, \omega_{X/k}^\bullet)
\otimes^\mathbf{L} (\mathcal{I}')^n)
$$
de sorte que la flèche $H^i_c(U', R\SheafHom(K|_{U'}, \omega_{U'/k}^\bullet))
\to H^i_c(U, R\SheafHom(K, \omega_{U/k}^\bullet))$ se déduise de façon analogue des morphismes $(\mathcal{I}')^n \to \mathcal{I}^n$. Les diagrammes
$$
\xymatrix{
R\SheafHom(M, \omega_{X/k}^\bullet)
\otimes^\mathbf{L} \mathcal{I}^n
\ar[rr] & &
R\SheafHom(R\SheafHom(\mathcal{I}^n, M), \omega_{X/k}^\bullet) \\
R\SheafHom(M, \omega_{X/k}^\bullet) \otimes^\mathbf{L} (\mathcal{I}')^n
\ar[rr] \ar[u] & &
R\SheafHom(R\SheafHom((\mathcal{I}')^n, M), \omega_{X/k}^\bullet) \ar[u]
}
$$
sont commutatifs parce que la construction des flèches horizontales dans Cohomologie, lemme \ref{cohomology-lemma-internal-hom-evaluate}, est fonctorielle dans les trois termes. On est ainsi finalement ramené à l'assertion que les diagrammes
\begingroup\small
$$
\xymatrix{
\Hom_k(H^i(X, R\SheafHom(\mathcal{I}^n, M)), k) \ar[r] &
H^{-i}(X, R\SheafHom(R\SheafHom(
\mathcal{I}^n, M), \omega_{X/k}^\bullet)) \\
\Hom_k(H^i(X, R\SheafHom((\mathcal{I}')^n, M)), k) \ar[r] \ar[u] &
H^{-i}(X, R\SheafHom(R\SheafHom(
(\mathcal{I}')^n, M), \omega_{X/k}^\bullet)) \ar[u]
}
$$
\endgroup
sont commutatifs. Cela vaut parce que l'isomorphisme de dualité
$$
\Hom_k(H^i(X, L), k) = \Ext^{-i}_X(L, \omega_{X/k}^\bullet) =
H^{-i}(X, R\SheafHom(L, \omega_{X/k}^\bullet))
$$
est fonctoriel en $L$ dans $D_\QCoh(\mathcal{O}_X)$.
\end{proof}

\begin{lemma}
\label{lemma-h0-compactly-supported}
Soit $X$ un schéma propre sur un corps $k$. Soit $K \in D^b_{\textit{Coh}}(\mathcal{O}_X)$ avec $H^i(K) = 0$ pour $i < 0$. Posons $\mathcal{F} = H^0(K)$. Soit $Z \subset X$ un fermé dont le complémentaire est $U = X \setminus U$. Alors le sous-espace
$$
H^0_c(U, K|_U) \subset H^0(X, \mathcal{F})
$$
est constitué des sections globales de $\mathcal{F}$ qui s'annulent dans un voisinage ouvert de $Z$.
\end{lemma}

\begin{proof}
Considérons le morphisme
\par\noindent\hfil
$H^0_c(U, K|_U) \to H^0_X(X, K) = H^0(X, K) = H^0(X, \mathcal{F})$
\hfil\par\noindent
de la remarque \ref{remark-covariance-open-lower-shriek}. Pour l'étudier, représentons $K$ par un complexe borné $\mathcal{F}^\bullet$ avec $\mathcal{F}^i = 0$ pour $i < 0$. On a alors, par définition,
$$
H^0_c(U, K|_U) = \lim H^0(X, \mathcal{I}^n\mathcal{F}^\bullet)
= \lim \Ker(
H^0(X, \mathcal{I}^n\mathcal{F}^0) \to
H^0(X, \mathcal{I}^n\mathcal{F}^1))
$$
Par Artin–Rees (Cohomologie des schémas, lemme \ref{coherent-lemma-Artin-Rees}), cela coïncide avec
\par\noindent\hfil
$\lim H^0(X, \mathcal{I}^n\mathcal{F})$.
\hfil\par\noindent
Ainsi, la flèche $H^0_c(U, K|_U) \to H^0(X, \mathcal{F})$ est injective, et son image est formée des sections globales de $\mathcal{F}$ qui sont contenues dans le sous-faisceau $\mathcal{I}^n\mathcal{F}$ pour tout $n$. Leur caractérisation comme les sections qui s'annulent dans un voisinage de $Z$ résulte du théorème d'intersection de Krull (Algèbre, lemme \ref{algebra-lemma-intersect-powers-ideal-module-zero}), en passant aux germes de $\mathcal{F}$. Voir la discussion d'Algèbre, remarque \ref{algebra-remark-intersection-powers-ideal}, pour le cas des fonctions.
\end{proof}




\section{Le théorème de Lichtenbaum}
\label{section-lichtenbaum}

\noindent
Le théorème ci-dessous a été conjecturé par Lichtenbaum et démontré par Grothendieck (voir \cite{Hartshorne-local-cohomology}). Kleiman en donne une très belle démonstration dans \cite{Kleiman-Lichtenbaum}. Une généralisation du théorème au cas de la cohomologie à supports se trouve dans \cite{Lyubeznik-Lichtenbaum}. La partie la plus intéressante de l'argument figure dans la démonstration du lemme suivant.

\begin{lemma}
\label{lemma-lichtenbaum}
Soit $U$ une variété. Soit $\mathcal{F}$ un $\mathcal{O}_U$-module cohérent. Si $H^d(U, \mathcal{F})$ est non nul, alors $\dim(U) \geq d$ et, en cas d'égalité, $U$ est propre.
\end{lemma}

\begin{proof}
Le théorème d'annulation de Grothendieck dans Cohomologie, proposition \ref{cohomology-proposition-vanishing-Noetherian}, permet de conclure que $\dim(U) \geq d$. Supposons $\dim(U) = d$. Choisissons une compactification $U \to X$ telle que $U$ soit dense dans $X$. (Cela est possible d'après Compléments sur la platitude, théorème \ref{flat-theorem-nagata}, et lemme \ref{flat-lemma-compactifyable}.) Après avoir remplacé $X$ par sa réduction, on trouve que $X$ est une variété propre de dimension $d$, et l'on voit que $U$ est propre si et seulement si $U = X$. Posons $Z = X \setminus U$. Nous montrerons que $H^d(U, \mathcal{F})$ est nul si $Z$ est non vide.

\medskip\noindent
Choisissons un $\mathcal{O}_X$-module cohérent $\mathcal{G}$ dont la restriction à $U$ est $\mathcal{F}$, voir Propriétés, lemme \ref{properties-lemma-lift-finite-presentation}. Notons $\omega_X^\bullet$ le complexe dualisant de $X$ comme dans la section \ref{section-duality-proper-over-field}. Posons $\omega_U^\bullet = \omega_X^\bullet|_U$. Alors $H^d(U, \mathcal{F})$ est dual de
$$
H^{-d}_c(U, R\SheafHom_{\mathcal{O}_U}(\mathcal{F}, \omega_U^\bullet))
$$
d'après le lemme \ref{lemma-duality-compact-support}. D'après le lemme \ref{lemma-duality-proper-over-field}, on voit que les faisceaux de cohomologie de $\omega_X^\bullet$ sont nuls en degrés $< -d$ et que $H^{-d}(\omega_X^\bullet) = \omega_X$ est un $\mathcal{O}_X$-module cohérent qui est $(S_2)$ et dont le support est $X$. En particulier, $\omega_X$ est sans torsion, voir Diviseurs, lemme \ref{divisors-lemma-torsion-free-finite-noetherian-domain}. On voit donc que le faisceau de cohomologie
$$
H^{-d}(R\SheafHom_{\mathcal{O}_X}(\mathcal{G}, \omega_X^\bullet)) =
\SheafHom(\mathcal{G}, \omega_X)
$$
est sans torsion, voir Diviseurs, lemme \ref{divisors-lemma-hom-into-torsion-free}. Par conséquent, ce faisceau ne possède aucune section non nulle qui s'annule sur un ouvert non vide de $X$ (une telle section serait de torsion). Le lemme \ref{lemma-h0-compactly-supported} montre donc que
$H^{-d}_c(U, R\SheafHom_{\mathcal{O}_U}(\mathcal{F}, \omega_U^\bullet))$
est nul, et ainsi $H^d(U, \mathcal{F})$ est nul, comme souhaité.
\end{proof}

\begin{theorem}
\label{theorem-lichtenbaum}
Soit $X$ un schéma non vide, séparé et de type fini sur un corps $k$. Soit $d = \dim(X)$. Les conditions suivantes sont équivalentes :
\begin{enumerate}
\item $H^d(X, \mathcal{F}) = 0$ pour tout $\mathcal{O}_X$-module cohérent $\mathcal{F}$ sur $X$ ;
\item $H^d(X, \mathcal{F}) = 0$ pour tout $\mathcal{O}_X$-module quasi-cohérent $\mathcal{F}$ sur $X$ ; et
\item aucune composante irréductible $X' \subset X$ de dimension $d$ n'est propre sur $k$.
\end{enumerate}
\end{theorem}

\begin{proof}
Supposons qu'il existe une composante irréductible $X' \subset X$ (que l'on considère comme un sous-schéma fermé intègre) qui soit propre et de dimension $d$. Soit $\omega_{X'}$ un module dualisant de $X'$ sur $k$, voir le lemme \ref{lemma-duality-proper-over-field}. Alors $H^d(X', \omega_{X'})$ est non nul puisqu'il est dual de $H^0(X', \mathcal{O}_{X'})$ d'après le lemme. On voit donc que $H^d(X, \omega_{X'}) = H^d(X', \omega_{X'})$ est non nul, et l'on conclut que (1) ne vaut pas. On voit ainsi que (1) implique (3).

\medskip\noindent
Démontrons que (3) implique (1). Soit $\mathcal{F}$ un $\mathcal{O}_X$-module cohérent tel que $H^d(X, \mathcal{F})$ soit non nul. Choisissons une filtration
$$
0 = \mathcal{F}_0 \subset \mathcal{F}_1 \subset
\ldots \subset \mathcal{F}_m = \mathcal{F}
$$
comme dans Cohomologie des schémas, lemme \ref{coherent-lemma-coherent-filter}. On obtient des suites exactes
$$
H^d(X, \mathcal{F}_i) \to H^d(X, \mathcal{F}_{i + 1}) \to
H^d(X, \mathcal{F}_{i + 1}/\mathcal{F}_i)
$$
Il existe donc un $i \in \{1, \ldots, m\}$ tel que $H^d(X, \mathcal{F}_{i + 1}/\mathcal{F}_i)$ soit non nul. Par le choix de la filtration, cela signifie qu'il existe un sous-schéma fermé intègre $Z \subset X$ et un faisceau cohérent non nul d'idéaux $\mathcal{I} \subset \mathcal{O}_Z$ tels que $H^d(Z, \mathcal{I})$ soit non nul. D'après le lemme \ref{lemma-lichtenbaum}, on conclut que $\dim(Z) = d$ et que $Z$ est propre sur $k$, en contradiction avec (3). Ainsi, (3) implique (1).

\medskip\noindent
Enfin, montrons que (1) et (2) sont équivalentes pour tout schéma noethérien $X$. En effet, (2) implique trivialement (1). Réciproquement, supposons (1) et soit $\mathcal{F}$ un $\mathcal{O}_X$-module quasi-cohérent. On peut alors écrire $\mathcal{F} = \colim \mathcal{F}_i$ comme la colimite filtrante de ses sous-modules cohérents, voir Propriétés, lemme \ref{properties-lemma-quasi-coherent-colimit-finite-type}. On a alors $H^d(X, \mathcal{F}) = \colim H^d(X, \mathcal{F}_i) = 0$ d'après Cohomologie, lemme \ref{cohomology-lemma-quasi-separated-cohomology-colimit}. Ainsi, (2) est vraie.
\end{proof}






\begin{multicols}{2}[\section{Autres chapitres}]
\noindent
Préliminaires
\begin{enumerate}
\item \hyperref[introduction-section-phantom]{Introduction}
\item \hyperref[conventions-section-phantom]{Conventions}
\item \hyperref[sets-section-phantom]{Théorie des ensembles}
\item \hyperref[categories-section-phantom]{Catégories}
\item \hyperref[topology-section-phantom]{Topologie}
\item \hyperref[sheaves-section-phantom]{Faisceaux sur les espaces}
\item \hyperref[sites-section-phantom]{Sites et faisceaux}
\item \hyperref[stacks-section-phantom]{Champs}
\item \hyperref[fields-section-phantom]{Corps}
\item \hyperref[algebra-section-phantom]{Algèbre commutative}
\item \hyperref[brauer-section-phantom]{Groupes de Brauer}
\item \hyperref[homology-section-phantom]{Algèbre homologique}
\item \hyperref[derived-section-phantom]{Catégories dérivées}
\item \hyperref[simplicial-section-phantom]{Méthodes simpliciales}
\item \hyperref[more-algebra-section-phantom]{Compléments d'algèbre}
\item \hyperref[smoothing-section-phantom]{Lissage des morphismes d'anneaux}
\item \hyperref[modules-section-phantom]{Faisceaux de modules}
\item \hyperref[sites-modules-section-phantom]{Modules sur les sites}
\item \hyperref[injectives-section-phantom]{Objets injectifs}
\item \hyperref[cohomology-section-phantom]{Cohomologie des faisceaux}
\item \hyperref[sites-cohomology-section-phantom]{Cohomologie sur les sites}
\item \hyperref[dga-section-phantom]{Algèbre différentielle graduée}
\item \hyperref[dpa-section-phantom]{Algèbre à puissances divisées}
\item \hyperref[sdga-section-phantom]{Faisceaux différentiels gradués}
\item \hyperref[hypercovering-section-phantom]{Hyperrecouvrements}
\end{enumerate}
Schémas
\begin{enumerate}
\setcounter{enumi}{25}
\item \hyperref[schemes-section-phantom]{Schémas}
\item \hyperref[constructions-section-phantom]{Constructions de schémas}
\item \hyperref[properties-section-phantom]{Propriétés des schémas}
\item \hyperref[morphisms-section-phantom]{Morphismes de schémas}
\item \hyperref[coherent-section-phantom]{Cohomologie des schémas}
\item \hyperref[divisors-section-phantom]{Diviseurs}
\item \hyperref[limits-section-phantom]{Limites de schémas}
\item \hyperref[varieties-section-phantom]{Variétés}
\item \hyperref[topologies-section-phantom]{Topologies sur les schémas}
\item \hyperref[descent-section-phantom]{Descente}
\item \hyperref[perfect-section-phantom]{Catégories dérivées de schémas}
\item \hyperref[more-morphisms-section-phantom]{Compléments sur les morphismes}
\item \hyperref[flat-section-phantom]{Compléments sur la platitude}
\item \hyperref[groupoids-section-phantom]{Schémas en groupoïdes}
\item \hyperref[more-groupoids-section-phantom]{Compléments sur les schémas en groupoïdes}
\item \hyperref[etale-section-phantom]{Morphismes étales de schémas}
\end{enumerate}
Thèmes de théorie des schémas
\begin{enumerate}
\setcounter{enumi}{41}
\item \hyperref[chow-section-phantom]{Homologie de Chow}
\item \hyperref[intersection-section-phantom]{Théorie de l'intersection}
\item \hyperref[pic-section-phantom]{Schémas de Picard des courbes}
\item \hyperref[weil-section-phantom]{Théories cohomologiques de Weil}
\item \hyperref[adequate-section-phantom]{Modules adéquats}
\item \hyperref[dualizing-section-phantom]{Complexes dualisants}
\item \hyperref[duality-section-phantom]{Dualité pour les schémas}
\item \hyperref[discriminant-section-phantom]{Discriminants et différentes}
\item \hyperref[derham-section-phantom]{Cohomologie de de Rham}
\item \hyperref[local-cohomology-section-phantom]{Cohomologie locale}
\item \hyperref[algebraization-section-phantom]{Géométrie algébrique et formelle}
\item \hyperref[curves-section-phantom]{Courbes algébriques}
\item \hyperref[resolve-section-phantom]{Résolution des surfaces}
\item \hyperref[models-section-phantom]{Réduction semi-stable}
\item \hyperref[functors-section-phantom]{Foncteurs et morphismes}
\item \hyperref[equiv-section-phantom]{Catégories dérivées de variétés}
\item \hyperref[pione-section-phantom]{Groupes fondamentaux de schémas}
\item \hyperref[etale-cohomology-section-phantom]{Cohomologie étale}
\item \hyperref[crystalline-section-phantom]{Cohomologie cristalline}
\item \hyperref[proetale-section-phantom]{Cohomologie pro-étale}
\item \hyperref[relative-cycles-section-phantom]{Cycles relatifs}
\item \hyperref[more-etale-section-phantom]{Compléments de cohomologie étale}
\item \hyperref[trace-section-phantom]{La formule des traces}
\end{enumerate}
Espaces algébriques
\begin{enumerate}
\setcounter{enumi}{64}
\item \hyperref[spaces-section-phantom]{Espaces algébriques}
\item \hyperref[spaces-properties-section-phantom]{Propriétés des espaces algébriques}
\item \hyperref[spaces-morphisms-section-phantom]{Morphismes d'espaces algébriques}
\item \hyperref[decent-spaces-section-phantom]{Espaces algébriques décents}
\item \hyperref[spaces-cohomology-section-phantom]{Cohomologie des espaces algébriques}
\item \hyperref[spaces-limits-section-phantom]{Limites d'espaces algébriques}
\item \hyperref[spaces-divisors-section-phantom]{Diviseurs sur les espaces algébriques}
\item \hyperref[spaces-over-fields-section-phantom]{Espaces algébriques sur des corps}
\item \hyperref[spaces-topologies-section-phantom]{Topologies sur les espaces algébriques}
\item \hyperref[spaces-descent-section-phantom]{Descente et espaces algébriques}
\item \hyperref[spaces-perfect-section-phantom]{Catégories dérivées d'espaces}
\item \hyperref[spaces-more-morphisms-section-phantom]{Compléments sur les morphismes d'espaces}
\item \hyperref[spaces-flat-section-phantom]{Platitude sur les espaces algébriques}
\item \hyperref[spaces-groupoids-section-phantom]{Groupoïdes dans les espaces algébriques}
\item \hyperref[spaces-more-groupoids-section-phantom]{Compléments sur les groupoïdes dans les espaces}
\item \hyperref[bootstrap-section-phantom]{Amorçage}
\item \hyperref[spaces-pushouts-section-phantom]{Sommes amalgamées d'espaces algébriques}
\end{enumerate}
Thèmes de géométrie
\begin{enumerate}
\setcounter{enumi}{81}
\item \hyperref[spaces-chow-section-phantom]{Groupes de Chow des espaces}
\item \hyperref[groupoids-quotients-section-phantom]{Quotients de groupoïdes}
\item \hyperref[spaces-more-cohomology-section-phantom]{Compléments sur la cohomologie des espaces}
\item \hyperref[spaces-simplicial-section-phantom]{Espaces simpliciaux}
\item \hyperref[spaces-duality-section-phantom]{Dualité pour les espaces}
\item \hyperref[formal-spaces-section-phantom]{Espaces algébriques formels}
\item \hyperref[restricted-section-phantom]{Algébrisation des espaces formels}
\item \hyperref[spaces-resolve-section-phantom]{Résolution des surfaces revisitée}
\end{enumerate}
Théorie des déformations
\begin{enumerate}
\setcounter{enumi}{89}
\item \hyperref[formal-defos-section-phantom]{Théorie formelle des déformations}
\item \hyperref[defos-section-phantom]{Théorie des déformations}
\item \hyperref[cotangent-section-phantom]{Le complexe cotangent}
\item \hyperref[examples-defos-section-phantom]{Problèmes de déformation}
\end{enumerate}
Champs algébriques
\begin{enumerate}
\setcounter{enumi}{93}
\item \hyperref[algebraic-section-phantom]{Champs algébriques}
\item \hyperref[examples-stacks-section-phantom]{Exemples de champs}
\item \hyperref[stacks-sheaves-section-phantom]{Faisceaux sur les champs algébriques}
\item \hyperref[criteria-section-phantom]{Critères de représentabilité}
\item \hyperref[artin-section-phantom]{Axiomes d'Artin}
\item \hyperref[quot-section-phantom]{Espaces de Quot et de Hilbert}
\item \hyperref[stacks-properties-section-phantom]{Propriétés des champs algébriques}
\item \hyperref[stacks-morphisms-section-phantom]{Morphismes de champs algébriques}
\item \hyperref[stacks-limits-section-phantom]{Limites de champs algébriques}
\item \hyperref[stacks-cohomology-section-phantom]{Cohomologie des champs algébriques}
\item \hyperref[stacks-perfect-section-phantom]{Catégories dérivées de champs}
\item \hyperref[stacks-introduction-section-phantom]{Introduction aux champs algébriques}
\item \hyperref[stacks-more-morphisms-section-phantom]{Compléments sur les morphismes de champs}
\item \hyperref[stacks-geometry-section-phantom]{La géométrie des champs}
\end{enumerate}
Thèmes de théorie des espaces de modules
\begin{enumerate}
\setcounter{enumi}{107}
\item \hyperref[moduli-section-phantom]{Champs de modules}
\item \hyperref[moduli-curves-section-phantom]{Espaces de modules de courbes}
\end{enumerate}
Divers
\begin{enumerate}
\setcounter{enumi}{109}
\item \hyperref[examples-section-phantom]{Exemples}
\item \hyperref[exercises-section-phantom]{Exercices}
\item \hyperref[guide-section-phantom]{Guide bibliographique}
\item \hyperref[desirables-section-phantom]{Desiderata}
\item \hyperref[coding-section-phantom]{Style de programmation}
\item \hyperref[obsolete-section-phantom]{Obsolète}
\item \hyperref[fdl-section-phantom]{Licence de documentation libre GNU}
\item \hyperref[index-section-phantom]{Index généré automatiquement}
\end{enumerate}
\end{multicols}


\bibliography{my}
\bibliographystyle{amsalpha}

\end{document}
